
%%%%%%%%%%%%%%%%%%%%%%%%%%%%%%%% SCRIPT FOR E-RARA AND MDZ FILES     %%%%%%%%%%%%%%%%%%%%%%%%%%%%%%%%%%%%%%%%%%%%%%%%
%%%%%%%%%%%%%%%%%%%%%%%%%%% fini le 30.04.2025 par F. GOY            %%%%%%%%%%%%%%%%%%%%%%%%%%%%%%%%%%%%%%%%%%%%%%%%
% !TeX TS-program = lualatex
\documentclass{article}
\usepackage[T1]{fontenc}
\usepackage{microtype}
\usepackage[pdfusetitle,hidelinks]{hyperref}

\usepackage{polyglossia}
\setmainlanguage{english}
\setotherlanguages{latin,greek}
\usepackage[series={},nocritical,noend,noeledsec,nofamiliar,noledgroup]{reledmac}
\usepackage{reledpar}

\usepackage{fontspec}
\setmainfont{TeX Gyre Termes}

\usepackage{sectsty}
\usepackage{xcolor}

\usepackage{fancyhdr}
\pagestyle{fancy}
\fancyhf{}
\fancyhead[LE,RO]{\nouppercase{\leftmark}}  
\cfoot{\thepage}
\renewcommand{\headrulewidth}{0.4pt}

% Redefine \section to remove numbering
\usepackage{titlesec}
\titleformat{\section}[block]{\normalfont\scshape\color{gray}}{}{0pt}{} % no number in heading
\titleformat{\subsection}[hang]{\normalfont}{}{0pt}{} % also remove subsection number
\titleformat{\subsubsection}[hang]{\normalfont\footnotesize\color{black}}{}{0pt}{}

% Modify how section marks are stored to exclude numbers
\makeatletter
\renewcommand{\sectionmark}[1]{%
	\markboth{#1}{}} % Only store the section title, without number
\renewcommand{\subsectionmark}[1]{%
	\markright{#1}} % Only store the subsection title, without number
\renewcommand{\numberline}[1]{} % Hide the section number in TOC
\makeatother

\begin{document}

\date{}
        \title{Commentarii in epistolas D. Pauli ad Timotheum: [Hyperius Andreae],[Ed. Mylius, Johannes], [1582]}
\maketitle
\tableofcontents
\clearpage
\begin{pages} 
\beginnumbering
        
\marginpar{[ p.1 ]}
\section{CAPVT I}
\phantomsection
\addcontentsline{toc}{subsection}{\textit{EPISTOLAE D. PAVLI PRIORIS AD TIMOTHEVM ARGVMENTVM. }}
\subsection*{\textit{EPISTOLAE D. PAVLI PRIORIS AD TIMOTHEVM ARGVMENTVM. }}\pstart \huge\textbf{T}\normalsize IMOTHEVS matre Eunica ludaea fideli, patre vero Graeco ni tus, a pueritia sacras literas didicerat, vt est Act. 16. et 2. ad Timoth.3. Paulus autem cum Lystram veniret, inuenit eum inter credentes egregie doctum, audiuitqueab omnibus fratribus, qui Lystris erant et Iconii, commendari vehementer. Itaque  permotus est, vt eum sibi adiungeret, (propter importuniratem eriam quorundam ludaeorum circumcidit) tum plenius instituendum, tum vsurus postea illius opera in negocio Euangelii ad diuersas ecclesias promouendo. Vnde etiam eum non vno loco modo in inscriptionibus, modon in aliis Epistolarum partibus, vt fidum collegam et ministrum Dei celebrat: ad Roman.6 1 ad Corinth.4.16.2. Corinth.1. Philip.1.2.Coloss.1. 1.Thessal.1.3. 2.Thessal.1. Et in Actis saepe praedicatur eiusdem dilsgentia. Hunc proinde Apostolus, cum relicta Epheso proficisceretur in Macedoniam, voluit ecclesiae Ephesinae suscipere gu bernationem, et bene a se constitutam in prima dignitate conseruare. Ibi itaque  agenti misit hanc de omni officio episcoporum eruditam epistolam. Ea proinde continentur hac epistola, vt de constituenda et ordinanda ecclesia inscribi iure optimo queat. Etenim omne officium corum, qui in ecclesiae versantur munere, hoc est, quae facere et do cere eos oporteat, qui ecclesiam sibi concreditam bene instituere cupiunt, mirabili ratione et accuratissime quidem tradit. Quem finem sibi fuisse propositum Apostolus suis ipse verbis cap.3. exprimit. Haec tibi scribo, inquit, vt noris, quomodo oporteat in domo Dei versari, quae est ecelesia Dei viuentis. Quemadmodum igitur multa sunt ad ecclessam recte ordinandam necessaria: ita in partes plures, quibus precipua capita vniuersae gubernationis ecclesiasticae comprehenduntur, praesentem epistolam Apostolus distribuit. Duo sunt primaria eorum, qui ecclesiae praesunt, officia: alterum et amplissimum in docendo consistit: alterum in Sacramentorum dispensatione, adeoquein procurandis iis omnibus, quae ad ritus, ceremonias, bonum ordinem, disciplinam eccle siasticam, ad conseruationem religionis pertinent. Atque  haec vt facerent, discipulis suis Christus mandauit, quando a morte sua ecclesiam suo iure constituens ait, Matth.28.  \pend\pstart Data est mihiomnis potestas in coelo et in terra. Euntes eroo docete omnes gentes, baptizantes eos in nomine patris, filii, et Spiritus Sancti, docentes eos se uare omnia, quae praecepi vobis. Apte ergo diuditur haec epistola in duas partes principales: priorem de doctrina, protrahitur haecad finem capitis quinti: posteriorem de ceremoniis, bono ordine, totaque  gubernatione, quae per capita reliqua in partes certas diuiditur. Itaque  prima parte compendium siue Epitomen vniuersae doctrinae Christianae exponit. Primo autem et postremo loco monet, vt in do ctrina semel accepta ministri ecclesiarum con stantes mancant. Dehinc cum doctrina sacra rite distribuitur in Legem et Euangelium, Apostolus ante omnia, et, vt apparent, propter quosdam, qui legem plus iusto vrgebant, et quasi ad iustitiam coram Deo foret neceslaria, commendabant: demonstrat breuiter, quis verus finis atque  vsus legis sit, nullo autem modo datam, vt per eam homi nes iustitiam spiritualem et aeternam apud Deum consequerentur. Finis pracepti est Charitas, inquit, ex puro corde: ett. et Scimus, quod bonasit lex, si quis ea legitime vitatur. Subiicit protsnus doctrinam de Euangelio, dicens, in iis, qui Euangelium amplectun tur, exuberare supra modum gratiam Domins nostri, cum fide et dilectione: vbi itidem verum vsum, effectum, et fructus Euangelii summatim demonstrat. Mox claudit hanc partem per exhortationem seu praeceptionem, vt ministri, qui pure et cum fructu docere volunt, ipsi in fide persistant, bonaque  cum conscientia negocium Christi agant: Exemplo addito de falsorum doctorum reiectione aut iusta excommunicatione.  \pend\pstart Hactenus de priore et excellentiore ministrorum officio, quod in docendo positum est In secunda parte epistolae, qua de reliquis ad imterium pertmenuibus: primus lo eus est capite secundo, de sacris coetibus, isquein duas partes diuiditur: priore ostendit q̃ uariae et pro quibarebus precationes ibi debeant fieri: Ante omnia, inquit, fiant depri catiomeob́ seorauoneset cAddit, quo assectu, quaque  modestia ta uir, s mulieres msacre  \pend
\vspace{0.5cm}\noindent
\vspace{0.2cm}\rule{1cm}{0.2pt}\\ 
\hspace{0.2cm}\textit{mg}
\footnotesize Act.16. 
\normalsize| 
\hspace{0.2cm}\textit{mg}
\footnotesize secunda partis distributio. 
\normalsize| 
\hspace{0.2cm}\textit{mg}
\footnotesize 1.Locus, de Sacris coetibus. 
\normalsize| 
\section*{IN I. EPIST. D. PAVLI. }
\marginpar{[ p.2 ]}\pstart illo coetu se component: Volo, inquit, orare viros in omni loco, etc. Posteriore parte agit de prophetia slue lectione atque  interpretatione Scripturarum in sacro coetu habenda.  \pend\pstart Cum autem hac in parte satis bene essent Ephesi omnia, et haud dubie ad praescriptum quod extat 1. Corinth. 14. Apostolus nihil aliud quam iubet mulieres in silentio discere, nec vllo modo ad publice docendum aspirare.  \pend\pstart Secundo loco, vbi distinguitur caput tertium, pergit explicare nonnulla, quae si be ne consideramus, pertinent ad senatum ecclesiasticum. Etenim, vt ad Tit.1. dicitur, et recte Ambrosius praesentis epistolae cap.5. docet: Vnaquaeque  ciuitas siue ecclesia habe re debet aliquot presbyteros seu seniores, qui constituunt senatum, consessum, collegium, consilium (non refert quomodo appellemus, modo de re ipsa constet) eceleliasticum: et cum hoc debet Episcopus de omnibus rebus ecclesiae conferre et consie? capere. Quae proinde sequuntur, aliquot locis explicata eiusmodisunt, vt in senatu ecclesiastico et per eundem expediri debeant. Est autem husus Senatus officium, ecclesiae proponere, vt idoneos ministros ante omnia eligant, Actor.6. Disserit ergo Apostolus primum de diuersis gradibus siue functionibus in ecclesia necessariis, de Episcopis videlicet et presbyteris. Scriptura namque  Episcopos et presbyteros subinde eosdem sine discrimine nominat, Act.20. et ad Tit.1. Et diaconos ostendit, quales ad haec munera eligentur et ordinabuntur. lta vero orationem instituit, vt intelligi pariter queat, quales esse tunc quoque  conueniat, cùm iam in dignitate sunt positi.  \pend\pstart Tertio loco praescribit, quomodo Episcopus, et quicunque  praeest ecclesiae, pruden ter admonebit fratres, inprimis autem eos qui populum docent, vt prouideant, ne qui aurem praebeant falsis doctoribus: cumque  hoc loco sublimis sit, initio reddit attentum, proposita tum ecclesiae ipsius, tum doctrinae sacrae dignitate: Haec tibi scribens, inquit, sperans fore, vt ad te veniam, etc. Arbitror autem hanc esse mentem Apostoli, debere episcopos non tantum hac de re publice ad populum disserere, verùm multo ma gis habere interdum Senatum siue conuentum ecclesiasticum, in quo ista fideliter proponet omnibus senioribus, quorum interest postea populum instituere ac hortari, vt   falsos doctores caute vitent. Etenim cùm dicit Apostolus in progressu cap.4. De his si   et commonefeceris fratres, bonus eris minister: intelligit inprimis presbyteros, qui sunt adiutores et collegae Episcoporum. Siquidem in hanc sententiam 2. ad Tim.2. ait, Haec commenda fidelibus hominibus, qui eruntidonei, ut alios quoque doceant. Idem Paulus prorsus simili admonitione vtitur apud presbyteros Actor.20. Ac re uera primum et excellentissimum munus senatus ecclesiastici est, accurate de doctrina conferre et prouidere, quae et qualsa populo expediant in ecclesiis salutariter proponi. Cuius rei exempla plura in Actis, ad Galath.12. Timoth.2. et c.habemus.  \pend\pstart Quartus locus continet quaedam de disciplina ecclesiastica, nempe de admonitionibus atque  obiurgationibus, quod interdum per Episcopum modó quidem publicem, modon vero coram senatu ecclesiastico sieri, ecclesiarum status atque  necessitas postulat. Quomodo igitur in hisce se gerat Episcopus, apostolus demonstrat: Exigitanteomnia, vt, qui alios sit reprehensurus vel moniturus, ipse irreprehensibilis et sine culpa esse modis omnibus studeat: Nemo inquit, tuam iuuentuse despiciat: postea addit cap.i.   Seniorem ne sapius obiurges, etc.  \pend\pstart Quintus locus est de aerario Ecclesiae, de eius recta administratione etiam senatus ecclesiastici interest, tum cognoscere, tum decernere. Diaconis interea incumbit, et colligere quae ad sumptus faciendos sunt necessaria, et collecta fideliter ex iudicio senatus ceclesiastici dispensare. Viduas honora, inquit, quae verae viduɇ sunt. Pluribus verbis mandat, ne grauetur aerarium eccleslae: Viduas quoad fieri potest, sustentari a liberis et nepotibus. Quin imo eadem de caussa vult vnumquemque  prospicere suis familiaribus. Praescribit insuper, quales eligi viduas deceat ad pauperum ministerium, de quibus Act.6. Dehinc addit, alendos ex eodem aerario presbyteros, maxime qui laborant in doctrina. Breuiter, duo hominum genera vult ecclesiarum beneficentia susten tari, pauperes omniope destitutos, et ministros verbi. Quemadmodum autem quarto loco agebatur de exhortationibus et obiurgationibus, in senatu ecclesiastico ad conseruandam disciplinam, fieri solitis: ita sexto loco disserit Apostolus de iudicio, quod ecclesia, hoc est, Episcopus cum senatu ecclesiastico exercet in eos, qui grauius deliquerunt,  \pend
\vspace{0.5cm}\noindent
\vspace{0.2cm}\rule{1cm}{0.2pt}\\ 
\hspace{0.2cm}\textit{mg}
\footnotesize 2.Locus de Senatu eccle siastico. Presbyteri necessarj ad Ecclesiae administrationem. Offcium Sena tus ecclesiastici. 
\normalsize| 
\hspace{0.2cm}\textit{mg}
\footnotesize 3. caput siue locus de fugi endis pseudodoctoribus. 
\normalsize| 
\hspace{0.2cm}\textit{mg}
\footnotesize 4. De disci plina ecdes astica. 
\normalsize| 
\hspace{0.2cm}\textit{mg}
\footnotesize 5. De dispenfaione facul tatam. Offieium diacouoram. 
\normalsize| 
\hspace{0.2cm}\textit{mg}
\footnotesize Quinam alen di sint ecclefiae facultati bus. 6. De iudicio ecclesiastico. 
\normalsize| 
\section*{AD TIMOTHEVM CAP. I. }
\marginpar{[ p.3 ]}\pstart deliquerunt, quos necesse est seuerius et quidem publice reprehendi, aut, si obtemperare noluerint, tandem tradi Satanae siue excommunicari. Breuiter vero de forma talis ecclesiastici iudicii, proposito exemplo de non damnando presbytero, nisi testibus aliquot conuicto: quaedam interserit, ex quibus pleniorem iudicii formam licet colligere. Aduersus presbyterum, inquit, accusationem ne admiseris, etc.  \pend\pstart Sentimo loco exponitur, non esse admittendos ordinandosue, nili quorum vita probe fuerit explorata. Qua inre etiam Senatus ecclesiastici iudicium episcopus adhibebit, quando capite quarto aperte dictum est, autoritate Τπρεσθυτερίς manus imponi. Debet examen vitae atque  doctrinae in hoc Senatu fieri: ait igitur, Manus citl * cui imponas.  \pend\pstart Octauo subiicit caput sextum, exhortationem nominatim faciendam ad Seruos, Iic interim paratam, vt consolationem illis adferre posset ingentem:cum Ephesi multi essent opulenti, haud dubie multi quoque  erant serui: quare postulat ecclesiae eius status huiusmodi exhortationem.  \pend\pstart Nono sequitur epistolae conclusio, grauis et affectibus referta, eademquedissecta in partes quatuor. Prima I’imotheo praecipit, vt doceat et exhortetur, quemadmodun epistola hac est praescriptum. Secunda, vt ab iis qui aliter docent vel exhortantur, se tempestiue seiungat. Tertia consolatur bonoque  iubet esse animo, etiam cum parum acci pit ab auditoribus, nullumquein ministerio quaestum facit, Est questus magnus, inquit, Pietas. Quarta hortatur et praecipit, vt strenue fideliterque  pergat boni pastoris praestare officium. Sectare autem iustitiam, inquit, Pietatem, fidem, charitatem, etc.  \pend\pstart Decimo appendix de exhortatione ad diuites, quales Ephesi erant plurimi. Vndecimo repetit ex superioribus admonitionem, quam Timothei animo altissime impressam cupit manere, de vitandis inutilibus disceptationibus, quae reuera sem per sunt ecclesiis exitiosae.  \pend\pstart Quod si iam rogat aliquis, ad quae genera concionum sacrarum partes huius episte lae referantur, facilis est responsio. Prima pars, in qua de lege atque  Euangelio agitur, ad διιδιασκαλιαν siue genus διδ ασκαλικόν pertinet. Quae vero sequuntur omnia, quoniam ad institutionem et rectam ordinationem ecelesiarum diriguntur, generis sunt παιδ οι τικς, id est, institutiui.  \pend\pstart Fui prolixior opinione: sed praesentis epistolae, quantumuis breuis sit, maxima et admiranda vtilitas me coëgit de singulis membris dicere aliquanto dilucidius. Scio mul torum conciliorum numerosos euoluissetis canones, nec potuissetis tam multa, tamque ́ accommoda ad ecclesiarum rectum ordinem colligere. Quamobrem aperte ae liberem dico, si quando de ecclesiarum conformatione serio agatur (agi autem omnino est neces sarium) non aliunde aptius, quam ex hac epistola, posse capi iudicium. Deinde et ex aliis noui testamenti scriptis sumentur quaecunque  ad ecclesiam sic instaurandam, vt et Deo placere, et piis mentibus tranquillitatem solatiumqueadferre queant, sunt necessaria. Hoc pro certo omnes habeamus, ex praescripto verbi Dei ecclesiam optime in omnibus reformari posse. Ne speremus, ab vllis hominibus nos meliora accepturos, quam Deus ipse in libris sacris extare uoluit. At vero nequaquam aerrat, qui scripturam, quam dictauit excellentissimus doctor Spiritus sanctus, ducem sequitur.  \pend
\phantomsection
\addcontentsline{toc}{subsection}{\textit{Paulus Apostolus Iesu Christi, iuxta delegationem Dei Seruatoris nostri, et Domini lesu Christi, qui est spes nostra, Timotheo germano filio in fide. Gratia, misericordia, pax, a Deo Patre nostro, et domino lesu Christo domino nostro. }}
\subsection*{\textit{Paulus Apostolus Iesu Christi, iuxta delegationem Dei Seruatoris nostri, et Domini lesu Christi, qui est spes nostra, Timotheo germano filio in fide. Gratia, misericordia, pax, a Deo Patre nostro, et domino lesu Christo domino nostro. }}\pstart Paulus Apostolus Iesu Christi.) Cùm asserit munus apostolatus sibi iuxta delegationem, hoc est, imperio, mandato peculiari ipsius Dei, iniunctum: scire debemus, quae in hac epistola proponuntur, Apostolo diuinitus tradita, atque  eam ob caussam omnibus ecclesuis serio proponenda et alacriter suscipienda. Fere consentiunt ea verba Apostoli cum iis, quae in aliis salutationibus habentur, cùm dicit, Per voluntatem Dei, duat τθελίματω † θεο͂. Sic in epistola ad Galatas negat se ab hominibus, aut per homines, sed per Christum Iesum et Deum Patrem ad munus apostolicum esse adoptatum, idquefacit propter pseudoapostolos, qui Euangelii doctrinam a Paulo praedicatam in odium inuidiamque  vocare nitebantur, contendebantes in operibus legis, quae per  \pend
\vspace{0.5cm}\noindent
\vspace{0.2cm}\rule{1cm}{0.2pt}\\ 
\hspace{0.2cm}\textit{mg}
\footnotesize Hdicimce desiae quale fit. 
\normalsize| 
\hspace{0.2cm}\textit{mg}
\footnotesize 7. N3 quosuisesseadmit tendos ad ec clesiae guber nationem e ministeriu. 8. De munere scruoran
\normalsize| 
\hspace{0.2cm}\textit{mg}
\footnotesize 9. Concuu sio quadrimeuibris. 
\normalsize| 
\hspace{0.2cm}\textit{mg}
\footnotesize 10. De diuitum studiis. 11. De uitaudis rixis iulutilibus. Genus cause. 
\normalsize| 
\hspace{0.2cm}\textit{mg}
\footnotesize Didescaies. παδιαιιηέμ 
\normalsize| 
\hspace{0.2cm}\textit{mg}
\footnotesize vfuspiftole. 
\normalsize| 
\hspace{0.2cm}\textit{mg}
\footnotesize Aduerbum Dei examinari omnes doctrinas et traditiones necesse est. Epigrapba
\normalsize| 
\hspace{0.2cm}\textit{mg}
\footnotesize Antorius te gationis Pau linae a Deo pendet: igitur Dogmatahu ius epistolae A spiritu San do sunt pro fecta. 
\normalsize| 
\section*{IN I. EPIST. D. PAVLI. }
\marginpar{[ p.4 ]}\pstart Mosen Dei digito lata ac scripta esset, plus fiduciae collocandum, quam Pauli apostatae do ctrinae fidendum. Vtuntur prophetae Dei eodem exordii genere, ad causae suae et doctrinae personaeque ́, quam sane gerunt diuinam, commendationem, vbi omnem vocationis, reuelationis, doctrinae, et factorum rationem, principium, et finem ad Deum summum suum principem regemque  referunt. Legatos itaque  Dei Prophetas et Apostolos esse meminerimus, qui, vt a Deo missi sunt et ad functionem tam praeclaram segregati, ita audiri hosce Dei loco oportet: sin quis repudiarit, filium Dei et Deum patrem is reiecerit. Lucae 1o. Sic non minus Apostolus peccasset, si vocem Dei et filii eius Iesu Christi vocantis ipsum audire recusasset: vel si etiam non vocatus cucurrisset: de vtroque  vide lerem.23. Amos3. Exod.3.4. Ierem. 1. Ionae1.  \pend\pstart Dei seruatoris nostri:) Accedit hinc dignitas muneri apostolico Diui Pauli, qu tam pe tris, que  filii mandato atque  imperio istud susceperit. Haecautem non propter Timotheum, qui sine dubio probe nouerat officium, dignitatem, fidem, diligentiam Apostoli: sed tum propter illos; apud quos ecclesiae erant ordinandae, in quibus semper aliqui inueniuntur vel sapientes, iuxta mundum, homines, vel maligni, qui resistunt et quibuscunque  modis conantur impedire rectam ordinationem, tum propter eos, qui cum Timotheo erant in gubernatione ec clesiastica in tempus posterum successuri, et similes labores pro ordinandis ecclesiis subire debent. Etenim pariter qui praesunt ecclesiis, et qui subsunt in ecclesiis, habent in hac epistola, quod amplectantur et tuto sequantur. Et si considerant, qualia Apostolus hic proponat, et qumo omnia ei mandata fuerint a Deo Patre et filio, his contenti erunt. Quod si aliqui praepostero z elo vel malitia excaecati confidunt, posse salubriora aliqua induci in ecclesias, quam quae noui testamenti libris sacris, maxime vero scriptis Apostoli Pauli sunt explicata, certe hi vehementer falluntur, et se non Paulo tantùm, sed Deo etiam patri ipsi et Dño lesu Christo temerario ausuopponunt. Caeterum Deum patrem seruatorem hoc loco vocat, quod nomem licet vsitato more scripturae Christo filio attribuatur, recte tamen et patri assignari potest, quia opera Trinitatis ad extra, vt vocant, sunt indi uisa, et personarum aequalis est potestas: Namque  et ipse pater seruat nos per filium et propter filium suum, et in filio. Sic Ephes.1. Elegit nos, inquit, in ipso, antequam iacerentur fundamenta mundi: praedestinauit nos, ut adoptaret in filios per lesum Christum. Et, Nemo potest venire ad Christum, nisi pater traxerit eum. lohan.6.  \pend\pstart Et Domini Iesu Christi.) Nam quod a patre constitutum est, idem constituitur et per filium. lohan.5.pater meus ad hoc vsque  tempus operatur, et ego operor. Ibid. Pater omne sudicium dedit filio, vt omnes honorent filium, sicut honorant patrem. Et Christus Actor.9. ad Ananiam, Organon electum est mibi iste, vt portet nomen meum coram gentibus et regibus acfiliis Israel. De potestate et dominio Christi lege Psal.2.8.72. ño.Matth.28.lohan.2..etc.  \pend\pstart Qui estspes nostra.) Additur hoc propter molem et difficultatem negociorum, itemque  pericula, quae, cum ordinantur ecclesiae, occurrunt. Infinita pericula et impedimenta experiuntur hi, qui laborant in ordinandis ecclesiis: sed quando nihil aliud quaerunt, quam gloriam Dei, et earum vtilitatem, atque  spem suam in Christum caput ecclesiae reiiciunt, sentiunt auxilia et mirabiles praeter omnium opinionem successus. Nituntur au tem hac spe siue fiducia, quod Deus per Christum faciet ipsos idoneos ministros noui testamenti, 2. Corinth.3. quod cum ipsi plantant, rigant, uel aliud quippiam faciunt, Deus dare uelit incrementum, 1. Corinth.3. Dicitur autem Christus spes nostra per Metalepsimquoniam in illo, qui post mortem suam et extremam deiectionem, resurrexit, omnem insuper accepit a Deo Patre potestatem in coelo et in terra, denique  ad coelos ascendit, ubi amplissima fruitur gloria: propositum est certum exemplum, cùm piis omnibus, tum his praecipue, qui fideliter laborant in ecclesia: quod quamlibet grauia pericula in hoc mundo patiantur, nunquam tamen extremum succumbent, sed a Christo habente omnem potestatem adiuuabuntur, et tandem cum Christo haereditatem filiorum Dei accipient. I.Petri 1. Deus pater regenuit nos in spem uiuam, perhoc, quod resurrexit lesus Christus ex mortuis, in haereditatem immortalem, ad Coloss.1.Christus spes gloriae, quem nos annunciamus.  \pend\pstart Proinde cum Christus sit spes nostra, id est, in quem et propter quem iubemur sperare, recte dictum est a leremia capite decimoseptimo, Maledictus homo, qui  \pend
\vspace{0.5cm}\noindent
\vspace{0.2cm}\rule{1cm}{0.2pt}\\ 
\hspace{0.2cm}\textit{mg}
\footnotesize Peudoapostolorum de gmata fase quia spiritum ueritatis non sequuntur di cem. 
\normalsize| 
\hspace{0.2cm}\textit{mg}
\footnotesize Pseudoapo stolorum d lus. Difgime u ror. et falsorum doctorum, Act.9. 
\normalsize| 
\hspace{0.2cm}\textit{mg}
\footnotesize Digmitas mi neris aposto ticia peie filii eius pe sona petenda est. 
\normalsize| 
\hspace{0.2cm}\textit{mg}
\footnotesize Deus pater, Seruator dicitur, quia seruat crede tes in per, et propter filin. 
\normalsize| 
\hspace{0.2cm}\textit{mg}
\footnotesize Scopus mint strorum De et ecdesie, est Dei gloria et eccle siae utilitas. 
\normalsize| 
\hspace{0.2cm}\textit{mg}
\footnotesize Christusspes fidelum. Matth.as. Act.. Marci 16. Christus fidis ministri? fubsidium pr. stat in uniuersis labobus atque  mo lestis. 
\normalsize| 
\section*{AD TIMOTHEVM CAP. I. }
\marginpar{[ p.5 ]}\pstart confidit in homine, et ponit carnem brachium suum, et a domino recedit cor eius. Sunt qui confidunt in virtute sua: alii in multitudine diuitiarum suarum gloriantur: alii suam iustitiam cum ludaeis volentes statuere, iustitiae Dei non subduntur: alii hypocritae fingunt se in Christum sperare, cum tamen Christum non habeant: verùm hi omnes in desperationem et laqueos diaboli tandem incidunt. Talium namque  spes, vt inquit sapl ens, est quasi lanugo, quae a vento tollitur, et tanquam spuma gracilis, quae a procella dispergitur, et tanquam fumus, qui a vento diffusus est, et tanquam memoria hospitis vnius diei praetereuntis.  \pend\pstart Timotheo germano filio.) Filium vocat hebraica consuetudine: quia seniores sic apnellare solent quoscunque  iuniores: item praeceptores suos discipulos: Vnde quidam filii prophetarum nominantur. Sed ideo maxime vocat filium, quia instituit illum in doctrina fidei ac religionis. Nam in huiusmodi generatio spiritualis spectatur:I. Corint.4. Etiamsi innumeros paedagogos habeatis in Christo, non multos tamen habetis patres: siquidem in Christo lesu per Euangelium ego vos genui. Et ibidem Timotheum vocat dilectum filium et fidelem in domino. Galat.4 Filioli mei, quos iterum parturio, donec formetur Christus in vobis. Sic etiam in veteri testamento (vt dixi) filii appellati sunt discipuli siue auditores a suis magistris: vnde Solomon in Prouerbiis saepe ait, Fili mi, cum tamen ad quemlibet doctrinae eapacem ei sermo sit. Et per generationem spiritu: Iem filii Abrahae vocantur, qui fidem atque  opera Abrahae faciunt, filii Sarae liberae, in spiritu non secundum carnem. Haec vera et firma generatio est, inquit Ambrosius, quae nescit oc casum, nescit morbos et pestilentiam, et quaeignorat famem et sitim. Ob has causas, et ma xime quia apostolus Timotheum spiritualiter Christo genuit, et quia Timotheus Paulum per omnia, praecipue autem in fide, referebat, ideo Pauli fillus in fide recte dicitur. Dignus tali patre Timotheus, et dignus tali filio Paulus vterque  ab inuicem gloriam accipiens, vterque  in propagando Euangelio iuxta diligens. Et cùm Apostolus filium in fide Timotheum appellat, tacite insinuat, dandam ei operam, vt in sana doctrina fidei perseueret, nihilque ́ sic requiri a quocumque  Episcopo, quam vt fidem incolumem tueatur et apud suos promoueat.  \pend\pstart Germanum autem vocat yνησιον, qu ingenio, fide, diligentia, ardore in promouendo Euangelio, iuuandisqueecclesiis Paulum per omnia referret exprimeretque ́. Vnde sic quoque  de eo apud Philippenses 2. Neminem habeo pari mecum animo praeditum, qui germa ne (ννοσίως) res vestras curaturus sit. Multum autem refert, qualibus hominibus committan tur ecclesiae gubernandae. Prouidere debent Episcopi, vt exploratos et satis notos ad hoc munus exemplo Pauli eligant. Necesse etiam est, eos, qui in munere ecclesiastico versantur, habere fidos et germanos collegas, qualem Apostolus Timotheum hic praedicat: quem alibi quoque  collegam et cooperarium sibi adiumgit: Vnde in quibusdam epistolis nomem Timothei cum suo in falutatione ponit, nempe in 2.ad Corinth ad Philippenses, ag Colossenses 1. et 2. Thessalonicenses, ad Philemonem: In quibusdam vero promittit, q velit eum mittere ad ecclesias visitandas, instaurandas, confirmandas. Laus haec excitare Timotheum ad singularem incredibilemquediligentiam potuit, facereque  ne quid se suaquepsona indignum committeret, nihilquepraetermitteret, qd'ad muneris dignitatem desideraretur.  \pend\pstart Gratia, misericordia pax) Gratiam precatur et pacem, quemadmodum alias in omnibus epistolis. Misericordiam autem tantum in duabus his ad Timoth. et illa ad Titum Etsi autem nomnullis placeat, adiici posse misericordiam, ad interpretandum, quid sibi velit gratia:nonnulli autem additam malunt vehementi quodam affectu, quo Apostolus de salute Timothei admodum erat solicitus. Attamem si respiciamus, quid tota hac epistola agatur, iudicabimus nomine misericordiae significari singulare Dei presidiq̃ et auxilium, itemqueeximia dona, quibus, qui ecclesiis praesunt, opus habent in tam sublimi atqu longe difficilimo mu nere. Reuera enim nisi Deus praeclaris donis illos ornet ac miserscordia singulari adit uet in laboribus constituendarum ecclesiarum, non potest sieri,ut ad gloriam Dei et utilita. tem hominum aliquid perficiant. Ac de singulari miserscordia, quae hie, qui ad ministerion ordinantur, est necessarsa:ita rursus apostolus2 ad Corinth.4. Cum ministeriu hocha beamus ut nostri misertus est Deus, non degeneramus, sed retecimus occultamenta dedecoris, non uersantes per astutiam, neque  dolo tractantes uerbu Dei, sed mansfestatione ueritatis commendamus nosipsos apud omnem conscientiam hominum in conspectu Del. Et sta  tim icapite presentis episstole, Gratiam habeo qui me potentem reddidit, Christo lesu domino te  \pend
\vspace{0.5cm}\noindent
\vspace{0.2cm}\rule{1cm}{0.2pt}\\ 
\hspace{0.2cm}\textit{mg}
\footnotesize Sapieuss. 
\normalsize| 
\hspace{0.2cm}\textit{mg}
\footnotesize Iobans. Rom.o. Gdats4 
\normalsize| 
\hspace{0.2cm}\textit{mg}
\footnotesize Multun re fert, delectum haberi eori, qui ad ecclesiae guberna tionem admittuntur. 
\normalsize| 
\hspace{0.2cm}\textit{mg}
\footnotesize Nifxiear? dia hic auxiium et cha rismarasigni ficat, quae in iis requirun tur, qui Ecclesiaeɇ mune racapessĩt. 
\normalsize| 
\section*{IN I. EPIST. D. PAVLI. }
\marginpar{[ p.6 ]}\pstart  noltro, qui fidelem me iudicauit ponendo in ministerium, qui prius eram blasphemus et n'penceutorvioientus, led meriimiertus eit, nain ignorans id facieba perincredulitate.  \pend\pstart A Deo patre nostro, et domino Iesu Christo Domino nostro.) Nam a Patre, qui est verus pa ter misericordiarum, et Christo, quem Esaias cap.9. vocat principem pacis, bona illa sunt impetranda. Bis autem Christus hic praedicatur dominus: Ac primùm quidem in genere, quia nimirum est verus Deus ac dominus dominantium, licet humiliatus sit apud homines vel p hominibus. Sic lohan.20. Thomas dicit, Dñs meus, et Deus meus. Si em̃ Deus, vtique  et dominus est. Deinde et dominus noster est, quia ipse nos a potestate Diaboli et mortis redemit, atque  ita plane suos fecit, quemadmodum docet Petrus 1.Epist.capa. Tertio autem loco vocatur dominus, quia caput ecclesiae est: vultqueApostolus significare, qu ab illo mandata, haec de ordinandis ecclesiis acceperit. Quamobcausam recipere eum merito omnes debesst. Deinde ab illo domino et summo Pontifice debent omnes fideles ministri singularem misericor. diam ac dona necessaria expectare. ls namque  est, qui ascendit in altum, et dedit dona hominibus  Caeterùm in quibusdam exemplaribus semel tantùm, posteriori videlicet loco legi tur vocabulum domini: quod admonendum duximus.  \pend
\phantomsection
\addcontentsline{toc}{subsection}{\textit{Quemadmodum rogaui te, vt remaneres Ephesi, cum proficiscerer in Macedoniam, permaneto: vt denuncies quibusdam, ne diuersam sequantur doctrinam, nec attendant fabulis et genealogiis nunquam finiendis, quae quaestiones prabent maquimncaquuanta denque atpa fuen. }}
\subsection*{\textit{Quemadmodum rogaui te, vt remaneres Ephesi, cum proficiscerer in Macedoniam, permaneto: vt denuncies quibusdam, ne diuersam sequantur doctrinam, nec attendant fabulis et genealogiis nunquam finiendis, quae quaestiones prabent maquimncaquuanta denque atpa fuen. }}\pstart Est haec epistola breuis, tot tamen insignes locos complectitur, quot nulla alia. In tanta breuitate autem nihil potest exordii loco sumi praeter haec pauca uerba, in quibus beneuolentia quidem captari videtur, dum reuocat in memoriam, qu antea amanter rogarit Timotheum vt Ephesi maneret, atque  vt idem adhuc faciat, rogat iterum. Attentum vero reddit, dum significat se cupere, vt ad ecclesiarum vtilitatem nomnulla ibi perficiat. Est autem oratio eclipti ca, siquidem graece haec tantùm dicuntur: Quemadmodum rogaui te, vt remaneres Ephesi cum proficiscerer in Atacedoniam: Addendum vero est: Ita nunc quoque  rogo, vt ibi maneas: vel, ita remaneto: siue, ut interpres suppleuit, ita facito. Neque  insolens est Apostolo sic concise lo qui, maxime quando vehementi quodam affectu ad aliqua valde utilia uel necessaria properat, quod 'eum hoc loco uoluisse admodum est probabile. Nam cum alias Apostolus in omnibus negociis, tum maxime in his, q̃ ad publicam spectant ecclesiarum utilitatem, uehemens erat et ardens. Eiusmodi autem ecclipseis passim in ipsius scriptis occurrunt, ad Ro.5.9. ad Gal.5.  \pend\pstart Prima parte huius epistolae disserit Apostolus erudite, breuiter tamen, de doctrina pure tradenda in ecclesia. Quandocunque  autem ecclesia aliqua est ab initio constituenda uel usquam collapsa fuerit restauranda ac reformanda, semper primo loco expediri ea oportet, quae ad doctrinam spectant, priusquam ad ceremonias uel quascumque  alias res fiat progres sus: rationes sunt in promptu. Primon Deus ipse eccleslam instituens formam hanc monstrauit, retinuit et seruauit persepe: Exod.2o. pmulgat Deus doctrinam legis, qua fides in Deum et Charitas in homines commendatur. Postea cap.23.proponit ceremonias. Et Christus Matth.28. Marci1o. mandat discipulis, primum ut peragrent mundum vniuersum, atque  Euangelium praedicent omni creaturae. Deinde ut credentes baptizent. Secundo, doctrina mul to est dignior atque  excellentior, quam quaecunque  res aliae: Siquidem propter se expetenda est: cum mentem plene informet et uoluntatem Dei tradat: et per se altius penetrat in corda, instituitquehominem internum: Dehine quibuscunque  hominibus tam doctis ac sapientibus, quam rudibus et infirmis congruit: Adhaec semper remansit eadem, nec mutatur. Denique  vera doctrina solius Dei propria est. Atqui ceremoniae ex altera parte propter aliquid aliud introducuntur, quippe quae signa tantum sunt, vel effecta, uel testimonia doctrinae: nec af ficiunt animos, nisi quatenus eorum ratio per doctrinam redditur, et seruiunt homini externo: Praeterea propter infirmos praecipue instituuntur, qui paedagogia quadam opus habent: Crebro insuper mutantur, et in plerisque  magna libertas relicta est: Denique  eaedem ceremoniae, quas obseruat ecclesia uera, subinde usurpantur ab his quoq, qui Deum non rite colunt. Tertio postulat hoc ipsa docendi ratio. Etenim non possunt ceremoniae institui uel approbari, nisi doctrina aliqua grauiore praeeunte, et per doctrinam, quae nimirum formam, causas, atque  usum ceremoniarum ostendit hominibusquecommendat. Porro ita comnectit Apostolus quae de doctrina primo loco adfert uerbis antecedenti  \pend
\vspace{0.5cm}\noindent
\vspace{0.2cm}\rule{1cm}{0.2pt}\\ 
\hspace{0.2cm}\textit{mg}
\footnotesize Christus De minus, 
\normalsize| 
\hspace{0.2cm}\textit{mg}
\footnotesize 1.Qa q uerus Deus. 2. Quiano redemit aptestate Sata ne, mortis et peccati. 2. Quiacaput est ecele siae. Ephes. 4.5. Ephes.4. Exordum. 
\normalsize| 
\hspace{0.2cm}\textit{mg}
\footnotesize Bencuolentiae captatio Attentio. 
\normalsize| 
\hspace{0.2cm}\textit{mg}
\footnotesize Prina pari de doctrins pure traden dain Eccle. sia. 
\normalsize| 
\hspace{0.2cm}\textit{mg}
\footnotesize Inecclesiinstauratione doctrina inprimis ha benda ratio est. 
\normalsize| 
\hspace{0.2cm}\textit{mg}
\footnotesize 2. Raioab exemplo et instituto dei 2. Ratio a do ctrinaeusuet necessitate 
\normalsize| 
\hspace{0.2cm}\textit{mg}
\footnotesize Doctrina se napropter seipsan ex petenda. 
\normalsize| 
\hspace{0.2cm}\textit{mg}
\footnotesize Discrimen doctrinae e uangelicae et ceremoniam̃ a fiueque  effectuabae ternitate it et tempore 3. Ratio. Ce remoniae nul Lae funt aut es se possunt, siteucbode 
\normalsize| 
\section*{AD TIMOTHEVM CAP. I. }
\marginpar{[ p.7 ]}\pstart bus, vt seiungi vix possint:qd'illi alias vsitatum est, maxime vbi breuitati studet, atque  ad docendum aliquid, vel praecipiendum magno adfectu rapitur. Vtitur ea sententiae forma, vt causam finalem notet, ob quam Timotheum Ephesi manere voluerit. Est autem, si quando pri mùm instituenda vel reformanda fuerit ecclesia aliqua, perutile obseruare et sequi methodum, qua Apostolus in tradenda doctrina vtitur. Primo euerti necesse est doctrinam sasam. Secundo loco explicabitur doctrina vera. Et rursus in explicatione doctrinae vera prior locus dabitur tractationilegis, posterior Euangelii: Sic namque  videmus Apostolum hoc loco atque  etiam alibi pgredi. Etenim qd'ad destructionem falsae doctrinae attinet, in graui illa disputatione ad Romanos, ante omnia euertit falsam psuasionem de lustitia co ram Deo obrinenda ex operibus, dum argumentis pspicuis et quasi experimentis euincit, omnes rum gentiles, que  ludaeos, si considerentur ipsorum opera, reos esse apud Deum et damnationi obnoxios. Inde vero pergit nouis argumentis demonstrare, qu homines iustificentur fide. Actor.17. in concione ad Athenienses statim reprehendit illorum superstitiones et idololatriam: postea annunciat vnum verum Deum ac lesum Christum. Quod ita fieri conuenit duabus potissimum de causis. Prima, non possunt simul consistere, multo minus misceri et conglutinari doctrina vera et falsa, sed necesse est, vt illa statuatur, hanc prorsus tolli: Negatic et affirmatio respectu eodem de eodem dicinequeunt. Secundo uidemus eademrarsone doceri in ordinandis quibuscunque  rebus, quae ante non fuerunt recte constitutae. Medici, si curare uolunt corpora aegrota, primo adhibent medicamenta, quae causas maliremo uent aut minuunt: Deinde porrigunt, quae confirmant ac uires augent. Qui colere acfer tilem reddere uult agrum siue hortum, eradicat et extirpatuepres, sentes, spinas: mox serit. ac plantat, quae profutura arbitratur. Qui emendare student eiuitatem, notant publica uitia, abrogant instituta iniqua, tùm coduntnouas et salutares leges. Eodem igitur pacto quibus propositum est ordinare uel reformare ecclesiam, instium facient a doctrina: sed prius. quidem destruent dogmata falsa:postea uero stabilient uera. Nunc uidemus, quando doctrinam falsam Apostolus euerti, uel prohiberi, antequam latius serpat, uelit. Denuncies,   inquit, quibusda. παραγγελ ο͂ν est cum autoritate et ex officio edicere, siue mandandum aliquid sit siue prohibendum: Vnde ea uoce Act.4. et s.utuntur principes populi et Seniores, cùm interdicunt Apostolis, ne loquantur de Christo. Apostolus idem2. ad Thes.. bis usurpat in praecepto de uitandis fratribus, qui inordinate se gerunt. Admonet proinde Apostolus, in his, quae ad Ecclesiam, maxime uero ad doctrinam seu dogma fidei pertinent, grauiter et cum autoritate debere per Episcopos cuncta fieri.  \pend\pstart Quibusdam:) Indefinita ac proinde generalis est Exhortatio: quo sciat Timotheus, sibi aduersus quoscunque  suspectos ac deprehensos, tam illo, quam alio quouis tempore, ta Ii denunciatione utendum. Quanquam ex sequentibus facile est animaduertere, quales prohibere maxime uelit. Erant sicut alibi, ita Ephesi quoque  ludaei, qui perturbabant ecclesias, qui multa et uaria spargebant uana et fabulosa: sed maximem urgebant, legem ad iustitiam esse necessariam: atque  cum Euangelio, cui suum concedebant locum simul debere obseruari: Verba sequentia et epistola ad Galatas, imo singulae prope epistolae id testantur.  \pend\pstart Ne diuersam sequantur doctrinam:) Maxime intelligi conuenit de Doctoribus ipsis, atque  his denunciandum, ne diuersa doceant: declarant id quae sequuntur, tam hoc capite, quam capite sexto, ubi ea lex recurrit. Non male autem interpretabimur:nihil autem refert, utrum ἐτεροδιιδασκαλέρ interpretemur aliud docere: quasi sermo sit de doctrinae substantia seu materia:an aliter docere:ut accipiatur de forma siue docendi modo. Nam pariter utrunque  Apostolus prohibet. Docent aliud, qui tradunt, quae in Scripturis nusquam sunt comprehensa, aut etiam aperte aduersantur Scripturis: cuius generis sunt, quae in hoc loco ab Apostolo nominatim perstringuntur, fabulae ineptae, uaniloquium. Et Christus Matt. 15. et Marci7. obiicit Pharisaeis, quod traditiones urgerent, per quas irrita fiebant Dei praecepta. Irenaeus liber 4. cap.25. qui fust Apostolorum uicinus temporibus, addit. ad suam usque  aetatem extitisse legem quandam appellatam Pharisaicam, legi Dei plane aduersantem. Docent aliter, qui proponunt quidem quae in Scripturis inueniuntur. sed non tractant ea synceriter: addunt aliquid, aut demunt, aut mutant, aut sententias peruertunt. De talibus est quod subnjcitur, quosdam praedicasse legem, sed legitimun eius usum atque  finem non reddidisse. Et capite 6́.carpuntur, qui logomachias mouent, dequeverbis magis quam rebus ipsis contendune Ad Galat.i. ait Apostolus, nonnut  \pend
\vspace{0.5cm}\noindent
\vspace{0.2cm}\rule{1cm}{0.2pt}\\ 
\hspace{0.2cm}\textit{mg}
\footnotesize Meibodis utilis et necessaria ad ecclesiae reformationem, Lex et Euangelium. Ab exemplo, Rom.2.5. 
\normalsize| 
\hspace{0.2cm}\textit{mg}
\footnotesize A natura con trariorum, uel cotradicterie pugnantium. A fimuli. Metephora a corpore ad anmum. Simile ab agrieultura 
\normalsize| 
\hspace{0.2cm}\textit{mg}
\footnotesize e̓τεροδιδα meoqd. 
\normalsize| 
\hspace{0.2cm}\textit{mg}
\footnotesize xnισ saica tempore trenci. 
\normalsize| 
\hspace{0.2cm}\textit{mg}
\footnotesize Rsecedi bist. liber .. cap
\normalsize| 
\section*{IN I. EPIST. D. PAVLI. }
\marginpar{[ p.8 ]}
\marginpar{[ p.9 ]}
\marginpar{[ p.10 ]}
\marginpar{[ p.11 ]}
\marginpar{[ p.12 ]}
\marginpar{[ p.13 ]}\pstart ,, los velle interuertere Euangelium. Actor: 2o. ad prebyteros Ephesios, Ex vobis ipfis , exorientur viri loquentes peruersa, ut abducant discipulos post se: Generatim autem Epiphanius haeres.15.refert, quadruplicem apud scribas et legisperitos expositionem in pretio habitam: primam in nomen Mosis praephetae fuisse collatam: secundam in praeceptorem ipsorum Azibam appellatum siue Barchochabum: Tertiam in Andam siue Annam, qui et ludas: Quartam in filios Assamonaei. Ex his quatuor repetitionibus siue expositionibus sunt, quae apud ipsos recepta sunt Sapientiae opinione: pleraque  vero insipientiae nomine iactantur, ac celebrantur: quaedam etiam loco electae doctrinae praedicantur ac diuulgantur. Hieronymo appellantur eae expositiones διουτερώτας: ldemqueenarrans cap.2.9. Esaiae, refert, Nazaraeos obiecisse scribis et pharisaeis, quod ipsorum διδεπερῶται defecissent, qui populo illuserant traditionibus pessimis, atque  ad decipiendos simplices die noctuque  vigilarant, qui peccare fecerant homines, vt Christum Dei filium negarent. Ex his proinde manifestum fit, quomodo multi praesertim ex ludaeis, aetate Apostolorum, ausi sunt partim aliud, partim aliter docere ac spargere perecclesias, atque  ita eas perturbare.  \pend\pstart Nec attendant fabulis et genealogiis, etc.) Per contemptum vocat fabulas, quaecunq ave ritate scripturarum sunt aliena, nullamque  vtilitatem aut aedificationem adferunt. Sic autem 2. ad » Timoth.4. Fabulas veritati opponit: A veritate, inquit, aures auertent, ad fabulas vero conuertentur. Ac notum omnibus est, in ludaeorum expositionibus, teste Epiphanio, et comprobantibus id Rabinorum commentariis, quas ludaei etiamnum hodie legunt, et suscipiunt, plurima fabulosa, stulta, ridicula, contineri. Sed exempli caussa addit Apostolus vnam speciem, genealogias inquam, ex qua de reliquis iudicium potest fieri.  \pend\pstart Caeterùm, quare genealogiis pertexendis tunc quidam fuerint curiosius, et vt reprehensionem incurrerent, dediti, operaepretium est consideremus.  \pend\pstart Herodes Antipatris filius, cum obscuris ortus esset parentibus, exussit libros diutilsime asseruatos in templo, in quibus familiarum genealogiae conscriptae fuerant: Vnde permoti ludaei, inprimis Nazaraei, multum operae posuerunt, in genealogiis partim memoriter, partim ex diariis et annalibus recensendis: atque  istud facientes ludaeam peragrabant vniuersam Refert id Africanus apud Eusebium Bcclesiast. histor.libus 1.cap.6.  \pend\pstart Deinde, colligebant genealogias suas ludaei, vt lingua et inani fastu iactarent eximios patriarchas, ac praerogatiuam quandam sibi prae gentilibus arrogarent, perinde ac si eo nomine etiam apud Deum praestantiores haberi deberent: ad Roman.3. 11. ad Philip.3. Tertio: Ex genealogiis conabantur patefacere veritatem ac certitudinem diuinarum promissionum, quae ad Abrahamum ac Dauidem de Christo ex semine istorum nascituro, tot ante seculis facte fuerant. Hac de caussa Euangelistae quoque  genealogiam Christi conscripserunt. Quarto: Dum ita gencalogiis digerendis plerique  erant intenti, saepenumeron dissensiones ac certamina acria exoriebantur. Quinetiam probabile est, etiam de Genealogia Christi, vt ab Euangelistis edita est, quosdam nimis que  odiose rixatos: Quod subindicat Epiphanius haeres.29. de Nazaraeis disserens: Apud ipsos, inquit, Euangelium Matthaei, quemadmodum ab initio scriptum est hebraicis literis, adhuc seruatur: nomnulli etiam genealogias ab Abraham vsq ad Christum amputauerunt.  \pend\pstart Et dum haec plerique  ludeorum ad Christum conuersi vana curiositate agerent, grauabant etiam gentiles recens ad ecclesiam accedentes, vt genealogiarum ordines perdiscerent, quemadmodum refert Theodoretus: lnterea autem aliarum rerum, quae vtiliores erant magisquenecessariae, tractationem negligebant, imo ueron et impediebant. Quae omnia cum sic se haberent, merito reprehendit Apostolus inane et noxium hoc studium, atque omnes ab eo detrahere desiderat.  \pend\pstart Quid? dicet aliquis: Christiani non curabunt genealogias? Aut peccant, qui ueteres illas ex scripturis inuestigent? uel, qui etiam hoc tempore illustrium familiarum, inprimis regiarum ac principaliq̃, maiorum, originem ordine digerunt? Respondeo aliquot propo sitionibus, Genealogias in libris sacris ueteris testamenti cognoscere studebimus, idque  easdem ob causas, ob quas Spiritus Sanctus mandari literis illas uoluit: pimô quidem, ut notemus rationem temporum atque  annorum mundi seriem. Secundo uero, ut deprehendamus, quomodo impletae sint promissiones, ad sanctos patres, de Christo ex ipsorum posteritate nascituro, factae. Atque  haec posterior causa cum primis obseruanda est. Ac dignum est, quod adnotetur, singulari Dei prouidentia factum ut paulatim et quasi per gra dus a promissionibus generalioribus ad speciales progressus fieret. Primum namque   \pend
\vspace{0.5cm}\noindent
\vspace{0.2cm}\rule{1cm}{0.2pt}\\ 
\hspace{0.2cm}\textit{mg}
\footnotesize rabile 44 
\normalsize| 
\hspace{0.2cm}\textit{mg}
\footnotesize Degencalo gis iudaeorum. 
\normalsize| 
\hspace{0.2cm}\textit{mg}
\footnotesize obiectio 
\normalsize| 
\hspace{0.2cm}\textit{mg}
\footnotesize Responsio. 1. Proposi tio. 
\normalsize| 
\hspace{0.2cm}\textit{mg}
\footnotesize Vtile est ueteris testamem ti obserugi genealogia 1.vtilitas. 2. Vtilitas. 
\normalsize| 
\section*{AD TIMOTHEVM CAP. I. }
\marginpar{[ p.14 ]}\pstart A damo denunciabatur in genere, ex muliere, sine ex semine mulieris proditurum seruz torem. Postea Abrahamo, qu’ex ipsius posteritate, ac certo populo, cireunciso videlicet Christus esset nasciturus: lacobo Patriarchae datum fuit praeuidere, ex qua tribu nomen trahente a quodam ipsius filio expectari Siloh deberet. Non recedet, inquit, sceptrum de Iudah, donec veniat Siloh, et ad illum gentes confluent. Dehinc Dauidi manifestatum ex qua familia eiusdem tribus nasci Christus vellet. I’andem etiam certa virgo est in eadem familia diuinitus monstrata, ex qua carnem Christus acciperet, misso nimirum Angelo ad Mariam, losepho, 9 ex domo Dauid erat, desponsatam. Ex his autem liquet, quam necessarium fuerit, genealogiam Chr sti ex libris veteris testamenti, quemadmo dum Euangelistas fecisse videmus, excerpi atque  Gmlbus cegnorceiuainl exhiberi.  \pend\pstart » Secundo: li quae praeterea causae siue fines adferuntur, ob quas gencalogiis digeren dis insistere conueniret, non tamen sunt magni ponderis, ac propterea illis minus mouebimur. Et cum supputatio annorum ex genealogiis obscure adhuc depromatur, nec satis haberi queat comprobata, quod testantur mirabiles eruditorum hac in parte dissensiones: attamen Christi populus, tribus, familia, ex genealogiis, vt Euangelistae Matthaeus et Lucas eas pertexuerunt, manifeste et sine errore deprehenduntur: quamobrem in Euangelistarum diligentia par est omnes acquiescamus. Si qui dissidere Euangelistas aliquo mo do suspicantur, at non est difficilis conciliandi ratio: quod doctissimem demonstrauit lulius Africanus liber  de hoc argumento edito: vt testatur Hieronymus in opere De scriptoribus Ecclesiasticis: et summam sententiae illius transcripsit Eusebius liber 1. Ecclesiast. histor. cap.6.dicens: Alterum Euangelistam sequi naturae ordinem, alterum niti instituto legis. Memi nit etia Africani probatqueeius sententiam Augustinus Retractat:libus 2. cap.7. Et consuli idem potest de consensu Euangelistarum lib2. cap.4. Contra faustum Manichaeum liber 3. et4. Cumañ Euangelistae suum munus in recensendis genealogiis cognitu necessariis, summa fide sint executi, consequitur, disceptationes omnes de genealogiis autoritate Euangelistarum dirimi: praeterea, finem esse debere omnis curiosae de iisdem inquisitionis, vbi precipuarum genealogiarum, q̃ habentur in sacris literis, verus vsus satis est explicatus, supuacaneus et poenitendus labor fuerit, plura de illis velle inquirere. Nec dubium, quin huc inprimis respexerit Apostolus, quando, ne ꝗ attendant fabulis atque  genealogiis, phibet. Vult tatum necessaria addisci. Quarto quod proinde Herodes sibros genealogiarum exuri curauerit, inter. pretari debemus iusto det iudicio esse factum. Etenim vaticinium lacob de ablato sceptro a ludah, aut de aduentu Siloh, impletum conspiciebatur, iamque  rempus appetebat, quo regnum sacerdotum, lex, ceremoniae, templum, et vrbs ipsa, debebant deleri. Quibus omnibus ita semel subuersis, necessario etiam perire oportebat genealogiarum sup putatlonem. Ac docet sane experientia, ludaeos, ex quo tempore Hierosolyma fuit po stremo per Titum Vespasiani filium solo aequata, non posse tribus vel familias suas ordinate recensere:idquetum, quia hinc illinc in remotissima loca fuêre per prouincias, nullo familiarum aut coloniarum seruato ordine, dispersi: tum quia temporis, quo varia subire exilia et sedes subinde mutare coacti sunt, longinquitas non permittit: Certum est hoc pa cto Deum velle gentis superbiam et stultam ob maiores gloriationem retundere.  \pend\pstart Sed neque  necessarium amplius est, pro annorum supputatione a genealogiis admi, niculum petere. Nam ex q̃ Christus generis humani seruator carnem induit, sapienter introducta receptaqueest ratio digerendi annos, non ab Adamo, Abrahamo, vel Dauide, sed ab ipsius Christi incarnatione: quo pacto anni sine disceptationibus  et absque  errore vllo hucusque  recensentur. Vbi namque  eum, ꝗ olim patribus  et familiis erat pmissus, praesentem recepi mus, nulla nos deinceps de illorum patrum familiis et genealogiis debet solicitudo premere. Imo si quis ad genealogias recurrere nunc vellet, suspectum sane se redderet, perin de ac si dubitaret de pmissionibus quondam factis, aut nondum crederet Christum in mun dum venisse. Ergo in Christo ac genealogiae eius descriptione vera per Euangelistas aedita finem omne genealogiarum studid̃ accepit. Tantum de genealogiis, q̃ in sacris literis recensen tur, pro vt de illis aetate diui Pauli disceptabatur, et quatenus negotio religionis sunt accommodatae. Quod porro ad genealogias regum, principum, aut quorumuis illustrium homi. num attinet, eas sine reprehensione licet Christianis hominibus digerere ac notare. Neque  em̃ ad religionem haec res spectat, sed plane politica est, ideoquesicut multae aliae res Poli, ticae, crimine vacua et concessa. Ac tres sunt caussae praecipuae, propter quas huiusmo: di gencalogiar! aliquanta cogniuo vnicuique  vtilis et necessaria censetur.  \pend
\vspace{0.5cm}\noindent
\vspace{0.2cm}\rule{1cm}{0.2pt}\\ 
\hspace{0.2cm}\textit{mg}
\footnotesize Gons. Gen.7. 
\normalsize| 
\hspace{0.2cm}\textit{mg}
\footnotesize Genesao. 1.Samucl.7. Psal.89. Esa.7. Luc.1. Mαιthη 
\normalsize| 
\hspace{0.2cm}\textit{mg}
\footnotesize 2Prope
\normalsize| 
\hspace{0.2cm}\textit{mg}
\footnotesize 1prope 
\normalsize| 
\hspace{0.2cm}\textit{mg}
\footnotesize -Prope 
\normalsize| 
\hspace{0.2cm}\textit{mg}
\footnotesize 5Propo 
\normalsize| 
\hspace{0.2cm}\textit{mg}
\footnotesize e Pegxe logiss regum et principum, etc. 7 Politicarum genealogiarum triplex Mc 
\normalsize| 
\section*{IN I. EPIST. D. PAVLI. }
\marginpar{[ p.15 ]}\pstart Prima est ius successionis: Scire enim operaepretium est, a quibus maioribus, et quo iure regna, prouincias, praedia, et quascunque  facultates vnusquisque  acceperit possideatque ;: tum qui rursus haeredes mortuis succedere in bonis debeant.  \pend\pstart Secunda ius respicit contrahendarum nuptiarum. Notum namque  est certis propinquitatis gradibus prohiberi coniugia. Ac de his duabus causis lurisconsulti accurate disserunt, nec magni laboris est perdiscere, quae huc requiruntur.  \pend\pstart Tertia causa sumitur ab excitatione ad virtutem, quae nascitur ex consideratione rerum, quas maiores cum laude gesserunt. Enimuero vt sciamus, a quibus doctrinam religionis acceperimus, quanta sit eius dignitas et antiquitas, vt admoneamur constantiae in confessione veritatis, fortitudinis, beneficentiae, et aliarum virtutum, que vnumquemq pro conditione et vitae suae genere decent, prodest ante oculos frequenter ponere, quae maiores nostri pro religione propugnanda, et conseruanda, pro gloria Dei illustranda, pro patria, pro vtilitate piorum et bonorum hominum fecerunt. Estquehaec caussa generalis, atque  omnibus temporibus omnibusquehominibus communis, vnde fit, vt saepe in veteri T’estamento pii describantur inuocantes Dominum patrum suorum, celebrantes varias virtutes et heroica facta eorundem. Et in nouo Testamento Apostolus in 2, ad Timoth. 1. excitare cupiens Timotheum et confirmare infide, reducit ad memoriam fidem, inprimis ipsius auiae Loidis et matris Eunicae. Et Polycrates Ephesiorum Episcopus, qui de celebratione paschatis epistolam Synodicam scripsit contra Victorem Episcopum Romanum, de se ita ait: Ego quoque minimus omnium nostium secundum doctrinam propinquorum meorum, quos secutus fum (septem siquidem fuerunt propinqui mei Episcopi, et ego octauus) semper pascha celehraui, , quando populus ludaeorum azyma faciebat. Itaquefratres, sexaginta quinque  annos ae» tatis meae natus in domino, et a multis ex toto orbe fratribus eruditus, peragrata omni Scriptura, non formidabo eos, qui nobis minantur. proferuntur haec verba ex epi£ stola illa ab Hieronymo in Catalogo scriptorum ecclesiasticorum, et Eusebio historiae Ecclesiast. liber 5. cap.4.  \pend\pstart Caeterùm pria potissimum sunt Christianis, dum genealogias suas pertexunt, cauenda. Primum est, ne studiose genus suum deducere laborent, ac nescio quas aborigines, et nimis longam ac fabulosam confictamqueantiquitatem: quomodo videmus nonnullos in priscis Troianis inaniter quaerere suae stirpis primordia. Viderint tales, ne deridiculi fiant, exemplo Arcadum, quorum stultitia etiam prouerbio locum dedit: quan do, quod haberi vellent omnium mortalium primi et antiquissimi, per irrisionem ab aliis προσέλννοι, tanquam luna ipsa antiquiores, appellati. Secundum est, ne titulentur fastus et inanis gloriae captandae libidine: semper enim fastus et ambitio hac in parte reprehenduntur, ad Roman.3.ad Philip.3. Extat et Chrysostomi homilia, qua docet, non esse gloriandum ob progenitorum virtutes. Et Satyrici poetae ementitam aciactabun dam nobilitatem perstrinxerunt. Tertium, ne quae a splendore generis bona accesserunt, quouis modo conuertant ad aliorum oppressionem, siue ad quodcunque  malum. Hoc namque  esset dignitate cum amantibus parta turpiter abuti: quomodo Catilina et nonnulli alii generis nobilitate ac potentia abusi leguntur. His atque  aliis vitiis, quae ratione circumstantiarum et causarum incidere possunt, remotis, possunt Christianige neris atque  familiarum seriem sine peccato digerere.  \pend\pstart Nunquamfiniendis.) Sic describit genealogias, tum quod perse sint institutae, sicut innumerabiles sunt familiae: tum, quod in earundem digestione ac supputatione inue niri exitus nequeat: cum vbi ad vnum aliquem quasi stipitem et primum familiae parentem perueneris, necessum sit ad alios quoque  longius progredi, a quibus ille idem originem suam detraxerit, atque  ita deinceps ad primum et communem omnium nostrum parentem Adamum ridicule deducenda erit supputatio. Ergo in hac ipsa voce argumentum latet suasorium, quo ab inutili et curiosa genealogiarum inuestigatione deterremur. Non superuacaneum modo, verum etiam immensi laboris et stultum est, rei, cuius exitus inueneri nequeat, dare operam.  \pend\pstart Quae quaestiones praebent, etc.) Duae sunt rationes coniunctae, quae docent, quare uitarioporteat omnem et fabularum et genealogiarum tractationem. Prior sumpta est a damnoso: praebent, inquit, quaestiones: ac necesse ea comprehedamus quae quodammo\pend
\vspace{0.5cm}\noindent
\vspace{0.2cm}\rule{1cm}{0.2pt}\\ 
\hspace{0.2cm}\textit{mg}
\footnotesize L Iu succes sionis. 
\normalsize| 
\hspace{0.2cm}\textit{mg}
\footnotesize Ahumpus rumt. 
\normalsize| 
\hspace{0.2cm}\textit{mg}
\footnotesize 3. Virtutum empima 
\normalsize| 
\hspace{0.2cm}\textit{mg}
\footnotesize Cautioncs ingenedogiis. 1. Cauio. 
\normalsize| 
\hspace{0.2cm}\textit{mg}
\footnotesize Arcadesxe σίληνοι. 2. Cautio. 
\normalsize| 
\hspace{0.2cm}\textit{mg}
\footnotesize 1 Cauio 
\normalsize| 
\hspace{0.2cm}\textit{mg}
\footnotesize Caussae pro pter quas de clinandae fabulɇ et gent alogiae siut. rcusse. 
\normalsize| 
\section*{AD TIMOTHEVM CAP. I. }
\marginpar{[ p.16 ]}\pstart do multitudini quaestionum inseparabiliter cohaerent, nempe dissensiones opinionum, acres contentiones, rixas, conuitia, odium implacabile, et c. Posterior ab eo, quod nihil faciunt ad pietatem, non aedificant, non illustrant gloriam Dei, non confirmant docent, vel solantur piorum conscientias. Vult Apostolus nos ab effectis iudiciun petere, vt statuamus, num doctrina aliqua sit amplectenda, vel reiicienda. Cuius do ctrinae effecta atque  fines sunt noxii, illam merito omnes vitabunt. Aedificare hebraic phrasi et per Metaphoram ductam ab exemplo Architectorum et fabrorum: trant fertur a corporali ad spiriruale: idemque  significar, quod docere, instituere, excitare, con firmare, tum verbo, tum actionibus, vt recte alii sentiant, iudicent, faciantue. Quem admodum domus et quodcunque  aedificium extruitur per fabros e lignis et lapidibus aprote concinnatis: Sic ecclesia, siue quicunque  pii ecclesiae ascripti docentur, instituuntur, confirmantur, peralios pios doctrina et exemplis, ad normam verbi seu voluntatis diuiae congruentibus  Ac valde idonea est haec phrasis, ideoquefrequenter ab Apostolis vsurpara. Vnde et ecclesia ipsa appellatur spirituale aedificium, structura, remplum, ad Ephes.2. Et pii 1.Corinth.6. vocantur templa ac habitacula Spiritus Sancti:1.Petri2. Christus appellatur lapis angularis, et primas acfundamentum iecisse, alios vero doctores superaedificasse. Et quo vnusquisque  clarius perspieiat, de spirituali aedificatione sermonem haberi, adiecit Apostolus vocem Dei, nimirum, aedificatioista non est corporalis vel humana, sed spiritualis ac diuina: Ac debet ea tantum doctrina sonare in Ecclesia, quam Deus ipse verbo suo tradidit, ac nobis omnibus commendauit, quae item ad iplius Dei tendit gloriam:nec possunt admitti fabulae uel uaniloquia.  \pend\pstart Quae est perfidem) Interpretatur ac plenius explanat, quae nam sit uera aedificatio: ut iudicemus, num doctrina aliqua aedificet, debemus in primis spectare elus aurorem nempe Deum, atque huius uerbum: Sed quia etiam falsi doctores dicere audent, Haec dicit Dominus, atque impia dogmata conantur de iactis scripturae uerbis stabilire, necesse est ulterius progredi atque considerare etiam materiam ipsius doctrinae, id est, num fides per eam commendetur. Earenus namque  doctrina aliquautiliter proponitur et recipitur in ecclesia, quatenus a Deo est, et ad Dei gloriam dirigitur, quatenus insuperfidem ac fructus ex ea pullulantes praedicat. Etenim fides iauae funda mentum est iustificationis multiplices complectitur fructus, qui iustificationem con sequuntur, dilectionem, inquam, spem, uerum timorem, pacem conseientiae, tolerantiam, mansuetudinem, etc. In iustificatis, ad quos, et de quibus Apostolus loquitur, fides otiosa nec est, nec potest sola consistere.  \pend\pstart Caeterùm decet nos memoriae commendare, quae Apostolus hocloco tradit, quo de quacunque doctrina, siue ea per pios doctores, siue per alios adferatuir, sanum et rectum feramus iudicium. Ac par sane est paulo diligentius ista excutiamus, tum propter utilicatem, tum quia eadem opera sententiae quaedam sequentes in hac epistola il lustrabuntur. Tria potissimum sunt ei consideranda, qui decernere uult, num doctrina, quae proponitur, amplecti uel reiicere expediat. Primum uerba ipsa siue uocabula, quibus dogma aliquod enunciatur. Secundum, res seu materia. Tertium, vsus, effecta atque  sinis. Ex his quomodo iudicium peti debeat, ordine ostendemus. Quod ad verba igitur attinet, notum est, verbis non rite intellectis, neminem posse quibus de rebus agatur, satis assequi Quamobrem prima semper est verborum cura, eorundeme expendi oportet significationem. Vt autem hac in parte cautus sis, ante omnia videm num sententia siue dogma nouae doctrinae pronuncietur verbis propriis et vsitaris. Nec aliunde certius possumus discere, quae propria et usitata sint uocabula in theologicis sermonibus seu dispurationibus usurpanda, quam ex sacrae Scripturae libris. leaque  qu uerbis utitur ipsius Seripturae, is merito audietur. Si quis uero noua introducar uoca bula, aut etiam nimis dura longiusue petita, suspectus habeatur, elusquedoctrinam admittere fas non erit. Ita namque  definit Apostolus 1.ad Timoth. 6.ô Timothee deposi, tum serua, deuitans profanas uocum nouitates, et appositiones falso nominatae scien te tiae quam nonnulli profitentes circa fidem aberrauerunt. Rursus 2. ad Timoth.2.  Stude teipsum probatum exhibere DeG, operarium non erubescentem et recte secan   tem sermonem ueritatis. Caeterùm profanas uocum nouitates praetermitto:ad ma   iorem enim proficient impietatem Semus autem his locis nonnullos nonvauregorre c\pend
\vspace{0.5cm}\noindent
\vspace{0.2cm}\rule{1cm}{0.2pt}\\ 
\hspace{0.2cm}\textit{mg}
\footnotesize »Cussa
\normalsize| 
\hspace{0.2cm}\textit{mg}
\footnotesize  2  nealogis et Hoctrinis. Metaphora ab Architectis. 
\normalsize| 
\hspace{0.2cm}\textit{mg}
\footnotesize Ecdt-fados mus Dei. 
\normalsize| 
\hspace{0.2cm}\textit{mg}
\footnotesize Differentia doctrinae ue rae et fase. ride fuAtus. 
\normalsize| 
\hspace{0.2cm}\textit{mg}
\footnotesize Iucum de ctrinae ex uerbis ipsis. 1.Cautiosiue consideratio. Dogmata noua ex uer bis nouis un propriis, et mnusitaus, co gnoscuntur. 
\normalsize| 
\section*{IN I. EPIST. D. PAVLI. }
\marginpar{[ p.17 ]}\pstart id est uoeum nouitates, sed πευοφωνίας, id est, uocum inanitates legere. Verùm quando et ratio sententiae, et praestantissimi quique  interpretes, Basilius, inquam, Chrysosto mus, Fulgentius, Ambrosius, Augustinus, aliique  plures priorem lectionem magis probant, nos eam sequimur ac retinemus, et cum Apostolo omnibus piis constituimus, nouas et inusitatas uoces esse in Theologicis negociis fugiendas. Notandum tamen est, Apostolum non simpliciter excludere omnes uoces nouas, sed profanas, hoc est, quarum significatio in literis sacris nusquam extat, non quas ecclesiae consensus in loquendo non reiicit, quemadmodum erudite Augustinus in lohannem tractat.97. Sunt, inquit, quaedam doctrinae religionis congruentes uocum nouitates: sicut ipsum nomen Christianorum, quando dici coeperit, scriptum est: ln Antiochia enim primùm post ascensionem Domini appellati sunt Christiani, sicuti legitur in Actis Apost. cae 11. Sicuero et ομοοσιον nouum uocabulum est, sed huius significatio deducitur ex Christi uerbis: Ego et pater unum sumus, uidelicet, eiusdem substantiae, lohan.1o.  \pend\pstart Secundon uidere oportet, an, qui proponit doctrinam nouam, uerba plura limul iuncta apte colligat: Nam in uerborum allegatione atque sermonis forma non satis apta, non obscuritas modo, uerumetiam error rectae fidei repugnans cito nascitur: unde Hieronymus in Oseam, Ex uerbis inordinate prolatis incurritur haeresis. Quem igitur animaduerteris sententiam minus apposite exprimere, quia uidelicet uerba aliter, quam decet, construit, huius dogma tantisper fugito, dum dilucidius magisqueproprie loquatur. Atque  istud faciens, praeceptum sequeris Apostoli, dicentis, 2. ad Timoth.1.   Formam habeto sanorum sermonum, quos a me audisti cum fide et charitate, quae est   in Christo lesu. Ad quae uerba Chrysostomus: mandat, inquit, Apostolus, ut si cui opus sit siue de fide, siue de Charitate, siue de Castitate aliquid loqui aut definire, exem plaria recte loquendi ex ipsius petat sermonibus. Certe Scholastica Theologia, sicut scatet impropriis et a Scriptura alienis locutionibus, ita plurimos errores inuexit in ecclesiam.  \pend\pstart Tertio: contingit interdum aliquos nouis et insolitis uerbis studiose uti, quo existimentur tanquam caeteris peritiores aut sapientiores res nouas atque  admirandas dicere: Interea autem deprehenduntur nihil diuersum adferre a caeteris, et quamuis uerbis dissentiant, re tamen ipsa conuenire. Aliquando porro nimis anxie et pertinaciter disceptant nonnulli de uerborum quorundam significationibus. Haec qui faciunt, cient tantum verborum pugnas, atque  abs re odiosa aggrediuntur certamina. Quamobrem   horum quoque  doctrinam Apostolus vult vitari, 1.ad Timoth.6. Si quis diuersam sequi   tur doctrinam, et non accedit sanis sermonibus Domini nostri Iesu Christi, et ei quae   secundum pietatem est doctrinae, is inflatus est, nihil sciens, sed insaniens circa quaestiones ac disputationum pugnas, λογομαχιας, ex quibus nascitur inuidia, contentio, maledicentia. Et 2 ad Timoth.2. Haec admone, contestans coram domino, ne λογοκαριας sequantur, ad nullam vtilitater, sed ad subuersionem auditorum. Atque  haec de expendendis verbis sufficiant. Post verborum accuratam disquisitionem operaepretium est, in materia ac rebus ipsis aestimandis iudicium exerceri. Vt autem breuiter statuamus, ad hic fieri expediat: tria possunt definiri quaestionum genera, quae ratione materiae seu rerum ipsarum vitari atque  reiici sine longa decertatione possunt. Primum est de eiusmodi rebus, quae ad Theologiam nequaquam proprie pertinent, sed potius sumptae sunt e philosophia aut poesi, aut similibus disciplinis externis. Nec enim conducunt ad expli candam doctrinam fidei, charitatis, spei, vel ad vitam in pietate recte instituendam: vel etiam ad illustrandos vtiles locos in Scripturis. Non est autem difficile ei, qui in sacris literis vtcunque  est versatus, iudicare quae nam quaestiones sint huiusmodi. Apud Epiphanium in disputationibus aduersus haereticos multa occurrunt exempla quaestionum motarum in ecclesia, quae, cum prorsus expenderentur, ex philosophia et poesi, ex quarum cognitione nonnulli haeretici aucupabantur gloriam, manasse deprehendebantur: Vnde idem Epiphanius aliquoties autores profanos Pythagoram videlicet, Platonem, Hesiodum, Homerum, et similes nominat, ex quorum libris illi sua fue rant mutuati. Et non modo ab hostibus ecclesiae quaestiones eiusmodi sunt propositae, verumetiam ab his, qui clari doctores in ecclesia extiterunt. Origenes disseruit de uετα dεκόrς, id est, animaru migratione e corporibus in corpora:secutus Prthagoricos:  \pend
\vspace{0.5cm}\noindent
\vspace{0.2cm}\rule{1cm}{0.2pt}\\ 
\hspace{0.2cm}\textit{mg}
\footnotesize vaveganta Egiende. 
\normalsize| 
\hspace{0.2cm}\textit{mg}
\footnotesize 2 mina uersio, An uerba probe sint allegatd. 
\normalsize| 
\hspace{0.2cm}\textit{mg}
\footnotesize Tbeologia scbolastica 
\normalsize| 
\hspace{0.2cm}\textit{mg}
\footnotesize 2-cautio. 
\normalsize| 
\hspace{0.2cm}\textit{mg}
\footnotesize Secundum ue dicium ex 
\normalsize| 
\hspace{0.2cm}\textit{mg}
\footnotesize Materiacir. ca quam uersantur doctrinae. 
\normalsize| 
\hspace{0.2cm}\textit{mg}
\footnotesize Tria gener quaestionum ratione ma teriae. 1Gen que stionum pe rgruaran 
\normalsize| 
\section*{AD TIMOTHEVM CAP. I. }
\marginpar{[ p.18 ]}\pstart itemquede quibusdam aliis a Christiana doctrina remotis, quae tamen, vt Hieronymus in epistola ad Anilum ait, ipse non dogmata, sed tantùm quaesita proiectaque  appellauit. Et scimus, per Augustinum, Hieronymum, Claudianum, Mamertium episcopum Viennensem in Gallia, permulta discussa esse de anima, subtiliter magis, quam pie, id est, philo sophice potius, quam theologice, etiamsi videntur huc importunitate quorundam haereticorum fuisse adacti. August. Epist. 56.refert, Dioscorum petiuisse a se, vt locos aliquos de dialogis Ciceronianis explanaret. Negari insuper non potest, varias quaestiones inductas in ecclesiam et scholas ecclesiasticasa theologis appellatis scholasticis: itemquea Canonistis: de quibus incertum est, an lureconsulti, an Theologi agnosci possint, quae tamen cum religione Christiana commune nihil habent: In summPigitur omnes huiusmodi quaestiones, quoniam, si cum verbo Dei conferantur, sunt tanquam fabulae et nugae, licebit Theologiae ac veritatis Christianae studioso reiicere, tanquam ad se nihil pertinentes. De hac re praeceptum habemus diui Pauli ad Coloss.2 Videte, inquit, ne quis sit, qui vos depraedetur per philosophiam. et inanem disceptationem, iuxta constitutionem hominum, iuxta elementa mundi, et non iuxta Christum.   Etiad Timoth.1. Denuncies quibusdam, ne diuersam sequantur doctrinam, nec attendant fabulis et genealogiis nunquam finiendis. Rursus2. ad Timoth. 4. Erit tempus, cum   sanam doctrinam non sustinebunt, sed iuxta concupiscentias suas coaceruabunt sibi doctores,   quibus aures pruriunt, et a veritate quidem aures auertent, ad fabulas vero conuertentur.    \pend\pstart Secundum genus quaestionum est, cum proponuntur specie quadam pletatis, vt plane theologicas esse, et verbo Dei comprobari, aut certe ad Scripturarum intelligentiam aliquousque  profuturas aliquis suspicetur: Attamen vbi accurate peruestigaueris, latens absurditas, vanitas, aut etiam impietas sese tandem prodit. Exemplum adfertur ab Apostolo. Doctrina legis praeclara estarque  in ecclesia perquam necessaria. Atque  sic tradi ca potest: Verùm si addatur, quod lex data sit, vt per eam iustificemur coram Deo, vel, vt etiam gentiles ipsam cum omnibus suis partibus recipiant obseruentque ́, contagium et magnum malum hinc sequetur. Atque  ita ductabunt eiusmodi doctores a principiis plausibilibus ad noxios progressus hoc est, ad mataeologiam, id est, vaniloquium. Quanobrem definit Apostolus, tales disputationes esse omni studio fugiendas.  \pend
\phantomsection
\addcontentsline{toc}{subsection}{\textit{Porro finis, inquit, praecepti est charitas ex puro corde, et conscientia bona, et fide non simulata, a quibus quod aberrarunt quidam deflexerunt ad vaniloquium, ματαολογίας, volentes esse lepis doctores, non intelligentes quae loquuntur, neq de quibus asseuerant ad perturbandam sanam doctrinam et lacerandas ecclesias: }}
\subsection*{\textit{Porro finis, inquit, praecepti est charitas ex puro corde, et conscientia bona, et fide non simulata, a quibus quod aberrarunt quidam deflexerunt ad vaniloquium, ματαολογίας, volentes esse lepis doctores, non intelligentes quae loquuntur, neq de quibus asseuerant ad perturbandam sanam doctrinam et lacerandas ecclesias: }}\pstart Quidam aliquando contenderunt, credentes vniuersos per Christum restitui primae innocentiae: hinc porro nudos debere omnes incedere et sacra peragere: hac de caussa appellati sunt Adamians referente Epiphanio: Quamuis et patrum nostrorum memoria idem sensisse ac docuisse quosdam Bohemios, Aeneas Syluius historiae Bo hemicae cap.41. narrat. Quisquis autem natura congenitum pudorem adhuc retinet, statim agnoscit, illos vana et absurda conatos tueri. Ophitae haeretici serpentem adorarunt, sed interpretantes, se in serpente colere ipsum Christum: nam in Scripturis cum Serpente a Moyse erecto confertur: ac typum sui elucere in hoc voluit Basilius, fama est, solitos tum sese, tum hospites suos religionis caussa castrare. De Rhetorianis D. Augustinus in libr. de haeresibus ad Quoduuldeum ait, Quod censeant haereticos omnes recte ambulare et vera dicere. Itaque  multa quae proponuntur et firmis existimantur nisa rationibus, absurda, vana, et nugatoria deprehenduntur, ac proinde indigna, quae percipere pii laborent. ldem iudicium esto de Enarrationibus, quibus impii ludaei libros sacros contaminant: item de plurimis quaestionibus Sophistarum scholasticorum: veluti cùm disputant, Vtrum haec propositio, Deus pater odit filium, sit possibilis. Vtrum sit aliquid in generatione diuina: Et quis enumeret omnia huiusmodi deliria ?  \pend\pstart Terin generis quaestiones sunt, quae de rebus mouentur longe supra ingenii capti positis. Frustra illis quis cognoscendis sese macerat, quae non sperat se assecuturum, Tales audiuntur nonnulle de mysterio Trinitatis, de filio Dei, de Spiritu S.de Sancto  \pend
\vspace{0.5cm}\noindent
\vspace{0.2cm}\rule{1cm}{0.2pt}\\ 
\hspace{0.2cm}\textit{mg}
\footnotesize Onues que stiones fabu losae et a uer bo Deiatiene figiende. 
\normalsize| 
\hspace{0.2cm}\textit{mg}
\footnotesize 2 Genuiqua stionum absurdarum. 
\normalsize| 
\hspace{0.2cm}\textit{mg}
\footnotesize 3.Genusquae stionum, que suutrenotab hominis 
\normalsize| 
\section*{IN I. EPIST. D. PAVLI. }
\marginpar{[ p.19 ]}\pstart rum praedestinatione, de causa mali. Vbi igitur ad haec ventum est, religiose animum submit. temus, et contenti gratia, quam Deus nobis ad mensuram donauit: ad Ephes.4. diuino Numini debitum honorem deferemus, dicentes cum Apostolo ad Roman.n1. ô pro,, funditatem diuitiarum et sapientiae et cognitionis Dei, quam inscrutabilia sunt iudi,cia eius et inperuestigabiles viae eius? Ac saepe prodest memoria repetere dictum illud sapientis Hebraei lesu Syrach cap.3. Difficiliora te ne quaeras, et quae tuas superant   vires, ne stulte exquiras. Quae tibi mandata sunt a Deo ea sancte meditare: neque  enim   opus est tibi, quae sunt arcana, oculis cernere: Quae non necessaria, sed superuacanea » sunt orationi tuae, ea curiose ne tracta. Nam plura quam humanus sensus capiat, tibi   patefacta sunt.  \pend\pstart Restat tertio loco dicamus de considerando vsu, effectis atque  finibus vniuscuiism doctrinae, quae recens pponitur. Clare autem definit: Si quis tradit ea, ex quibus non sequitur aedificatio Dei, quae est perfidem, id est, neque  doctrina fidei siue verorum dogmatum asse ueratio, siue redargutio dogmatum falsorum, neque  institutio vitae ad iustitiam, neque  correptio a recta via aberrantium, neque  consolatio afflictorum, (nam ad haec capita sicut omnis pertinet Scriptura, sta omnem quoque  explicationem doctrinae sanae par est referamus) sed potius nascuntur quaestiones superuacaneae, inutiles, imo et noxiae, tum rixae, aemulationes, obtrectationes: id esse explanandum. Et cap.6.in eandem sententiam legimus, Ex doctrina diuersa a sanis sermonibus Domini nostri lesu Christi oriri quaestiones et pugnas verborum: dehinc sequi inuidiam, contentionem, maledicentiam, superstitiones malas, superuacaneas conflictationes, hominum nempe corruptorum, et quibus adempta est veritas, quique  existimant quaestum esse pietatem. Subindicat interim diuersos subinde fines acutis disputationibus esse propositos, veluti quaestum, ambitionem, ius primatus: Ac de his diligenter esse inquirendum. 2.ad Timoth.2. ait, pugnas verborum nullam adferre vtilitatem, sed iis subuerti auditores. Ibidem addit: profanis vocum nouitatibus proficere homines ad maiorem impietatem, et longe latequemalum ser» pere, quemadmodum serpit cancer morbus. Ruisus, stultas autem et ad eruditionem ineptas quaestiones reiice, sciens eas gignere pugnas. Porro non tantum effecta atque  fines, qui ex ipsis per se sequuntur disceptationibus, vult Apostolus a nobis expendi, verum etiam affectus, quibus qui doctrinam suspectam inuehunt, mouentur, vitia quibus iidem laborant, denique  et fines, quos praestituunt sibi. Vnde Apostolus 2. ad Timoth.3. Erunt, inquit, homines sui amantes, auari, fastuosi, superbi, maledici, parentibus   immorigeri, ingrati, impii, carentes affectu, nescii foederis, calumniatores, intempe-¨   rantes, immites, negligentes bonorum, proditores, praecipites, inflati, voluptatum amantes potius quam amantes Dei, habentes formam pietates, sed qui vim abnegarint   eius: istos igitur auersare. Ex his enim sunt isti, qui irrepunt in familias, et captiuas ducunt mulierculas, oneratas peccatis, actas cupiditatibus variis: semper discentes, neque    vnquam ad cognitionem veritatis yenire valentes: Quemadmodum autem larnes et   lambres restiterunt Mosi, ita et hi resistunt veritati, homines corruptimente, reprobi   circa fidem. Et ad Roman.i5 Obsecro vos fratres, ut consideretis eos, qui dissidia et   offendicula contra doctrinam, quam uos didicistis, gignunt, et declinetis ab illis: Nam   qui eiusmodi sunt, lesu Christo non seruiunt, sed suo uentri, et per blandiloquentiam   et assentationem decipiunt corda simplicium. Ad Galat.6. Quicunque  uolunt iuxta saciem placere in carne, hi cogunt uos circumcidi, tantùm ne ob crucem Christi persecu   tionem patiantur. Sunt alibi alia effecta, alii fines, atque  diuersae falsorum doctorum notae explicatae. Operaepretium proinde est etiam in his, per quos noua doctrina euulga tur, diuersa ista effecta atque  fines, corruptos uidelicet affectus, uitiosos mores, fines inhonestos, flagitia, diligenter obseruari: quandoquidem haec omnia sunt reuera tanquam fructus impiae doctrinae, ex fraudulento pectore pullulantes. Nec potest quispiam praetexere, se non posse satis discernere, quinam tales impii doctores existant. Nam ita Deus sua iustitia atque  prouidentia hypocritas atque  omnes, qui reluctante propria conscientia falsum docent, aut ueritatem quouis modo, id est, siue per se, siue per alios, et siue aperte, siue oblique, impugnant: punitat muilctat, ut turpiter in aliqua  \pend
\vspace{0.5cm}\noindent
\vspace{0.2cm}\rule{1cm}{0.2pt}\\ 
\hspace{0.2cm}\textit{mg}
\footnotesize Teuimì dicium a fins et effectis. 
\normalsize| 
\hspace{0.2cm}\textit{mg}
\footnotesize Efeaapa gnarum. 
\normalsize| 
\hspace{0.2cm}\textit{mg}
\footnotesize Afdusde ecntiun. 
\normalsize| 
\hspace{0.2cm}\textit{mg}
\footnotesize Note diuc dae bypocri. tgun et faforum docto: rum. 
\normalsize| 
\section*{AD TIMOTHEVM CAP. I. }
\marginpar{[ p.20 ]}\pstart detestanda peccata labantur, vel certe aliis modis pudefiant, quo iudicare postea omnes de ipsis et falsa ipsorum doctrina possint. Propterea Christus Matth. 7. iubens, vt caueant sibi a pseudoprophetis, qui veniunt in vestimentis ouium, intrinsecus autem sunt lupi rapaces, ait, a fructibus ipsorum agnoscetis eos. Vbi et pluribus verbis declarat Christus, non posse fieri, quin tandem fructus illorum noti fiant omnibus. Latent et teguntur forte doli aliquamdiu, attamen progressu temporis in lucem eos protrahi necesse est. Quocirca Apostolus 2.ad Timoth.3. post enumerata hypocritarum et falsorum doctorum varia scelera, subiicit, tales homines non posse amplius proficere, sed ipsorum amentiam omnibus patefieri. Ad Philip.3. Complures ambulant, de quibus saepe dixi vobis, nunc autem et flens dico, inimici crucis Christi, quorum finis perditio est,   quorum Deus venter est, et gloria in dedecore ipsorum. Atque  omnibus aetatibus eccle -  siae fuerunt huius rei exempla in conspectu posita: et diuites et maiores, sapientes et eruditi diu restiterunt Mosi coram pharaone, tandem vero fassi sunt se victos, ac digitum siue virtutem Dei per Mosen efficaciter operatam. Exod.8. 2. Tim.3. Qui longo tempore imposuerant regibus, principibus et populo sacerdotes Baal dum se opponunt Heliae prophetae Dei, mirabilster confunduntur. 1. reg.18. Quomodo vero Arrius, cum crederetur victor in causa, et magno honore atque pompa ad Alexandriam esset deducendus ecclesiam, diuina vltione turpiter in latrina inter cacandum cum intestinis et animam simul profuderit, testantur Ecclesiasticae historiae Eusebianae liber 1o. cap.8. Tripart. histor. liber 3. cap. 1o. Nec desunt praesenti aetati illustria exempla, si ea commemorare liberet. Sed hic iam non est nostri instituti pluribus de hac re disserere. Adhortamur igitur omnes rerum sacrarum studiosos, vt quandocunque  do ctores noui irrepunt, vel potius prosiliunt in ecclesiam, prudenter expendant, an placita eorundem sequenda sint aut fugienda, eamquead rem adhibeant in consilium, quae hactenus ex sacra scriptura sunt in medium prolata, eademque  assuescant quibuscunque  argumentis accommodare. Sed et studeant, quaecunque  his cognara in scripturis occurrerint, accurate notare, atque  his adiicere. Disputant sane hoc tempore de variis rebus in Ecclesia, et hi quidem, quorum magna est propter eruditionem non vulgarem existimatio: Attamen, si libeat primùm verba ipsa, dehinc res siue materiam, postremo vsum, effecta, et fines diligenter scrutari expendereque , comperiemus praestabilius esse, vt illarum disputationum plerasque  praetermittamus, quam ut insumamus tempus in illis percipiendis. Enumerare eas ac de singulis consutam facere, nimis longum foret. Quare vniuscuiusque  diligentiae ac priuato iudicio hunc laborem committimus, Debemus pono Deo gratias agere, qui in uerbo suo huiusmodi normas siue regulas iudicandi de doctrinis extare et quasi ardentes faces praelucere uoluit, quas in rebus dubiiis, maxime cum nouis disceptationibus ac turbis prope collidi uidetur ecclesiam, tuto sequeremur. Etictum est a psalmographo Psal.119. Lucerna pedi meo uerbum tuum, et lux semitis meis. Praesidium certum ac scutum uere magnum est, audire et adhaerere doctori Spiritui sancto. Quinetiam par est Deum assidue precemur, ut facultatem diiudicandi doctrinas in ecclesia sua integram semper conseruet, acnobis etiam per suam misericordiam benignem impertiatur. Atque haec quidem hactenus generatim, quamuis breuissime, de euersione doctrinae falsae, in qua semper primo loco esse laborandum supra diximus. Sequitur nunc de plantatione doctrinae uerae ac omnium primum de lege.  \pend\pstart Porro sinis pracepii.) Apostolus partem de plantanda doctrina vera ingrediens de lege prius agit, quam de Euangelio:idquehis, ut probabile est, de causis.  \pend\pstart 1 IIophetis familiare est, in concionibus statim ex lege damnare peccata, obimr gare transgressiones, ut sic adducant ad poenitentiam, mox subiiciunt adhortationes ad poscendum ueniam et promissiones euangelicas de misericordia per fidem accipien da. Passim occurrunt exempla apud Esaiam et reliquos. Christus autem et lohannes Baptista priorem partem concionum faciunt de poenitentia: posteriorem de peccatorum remissione et regno coelorum. Vult igitur Apostolus hos antecessores imitari. 1 Siconsideramus tempora promuigationis, les primo prompigata estper ssio.  \pend
\vspace{0.5cm}\noindent
\vspace{0.2cm}\rule{1cm}{0.2pt}\\ 
\hspace{0.2cm}\textit{mg}
\footnotesize Moei auay ptiaci consan duntur. 
\normalsize| 
\hspace{0.2cm}\textit{mg}
\footnotesize Arrii intert. ue soedus.. 
\normalsize| 
\hspace{0.2cm}\textit{mg}
\footnotesize Voebz Det Lydus Lapis ac norma est omn um dogmatum, etc. 
\normalsize| 
\hspace{0.2cm}\textit{mg}
\footnotesize De lege. 1. Prophete a lege et pae nitentia suas conciones ordiuntur. LuC3.24. Matth.3. Lex prius promuιgα- la, qumeuangelum 
\normalsize| 
\section*{IN I. EPIST. D. PAVLI. }
\marginpar{[ p.21 ]}\pstart sen. Euangelium, id est, gratia et veritas postea per lelum Christum.  \pend\pstart Tertio: postulat hocordo docendi in procuranda hominis salute necessarius. Qui namque  in viam reducendus est homo, doceri ante omnia debet, vt semetipsum agnosçat atque  miserrimu esse peccatorem, aeternae damnationis reum confiteatur: id quod nul la ratione commodius, que  praédicatione legis perficit. Per legem agnitio peccati, ad Rom. 3.Vbi vero extreme delectus fuerit, eum rursus erigl oportet praedicatione Euangelii, vt ad Deum confugiat et misericordiam fide imploret.  \pend\pstart Addi potest et haec ratio. Aetate Apostoli hostes Euangelii ludaei, quorum similes permulti hoc tempore sunt hypocritae, supra modum legem vrgebant, ex ea promittentes iustiuam coram Deo: proinde quoniam necesse erat, ante confutari aduersarios, atque omnes ex omnium animis reuelli dubitationes, quam ad reliqua fieret progressio, prudente de ipsa Lege primo loco sermonem instituit. Discimus igitur et nos, quandocunque  summam doctrinae religionis tradere habemus in animo, propter casdem caussas initium face, re de lege: Secundo loco adiicere, quae sunt Euangelii. Totus vero est in explicando legis vsu atque  fine:cuius rei duae causae reddi possunt. De vsu atque  fine erant controuersie, non de aliis rebus ad legem pertinentibus. Illa autem potissimum explicari operaepretium est, de quibus dubitatur, seu de quibus oriuntur disceptationes. Estnamque  doctoris officium, pe tidere, vt difficiliores nodos dissoluat, et scrupulos omnes ex animis remoueat auditorum. 2. Maxime necessarium est cuiuscunque  doctrinae, quae proponitur, vsum atque  finem et primo et vltimo loco consi derari. Nui enim vsum probe praeuideamus, frustra in illa percipienda laborabimus. Admonemur proinde et nos hunc sequi in docendo ordinem, nempeut semper studeamus diligenter eruere et explanare quae in dubio uidentur posira et de quibus controuertitur: proxime, vt omnis doctrinae verum vsum perspicuë ostendomus.  \pend\pstart Finis praecepui est charitas.) Praeceptum siue παραγνελ ιεv intelligit ipsam legem, quod perspicuum fit ex verbis sequentibus, vbi vocabulum legis repetitur. Constituit pioinde Apostolus praeceptisiue legis finem esse Charitatem, quasi dicat: Si inquirat sinis praecipuus ob quem tota lex lata est, et ad quem cuncta praecepta respiciunt, inueniemus es se Charitatem siue dilectionem Dei ac proximi, vsque  adeô, vt qui charitatem docet, reue ra summam omnium tradat quae lege continentur, et credentibus sunt ex lege cognitu necessaria. Sic autem et Christus ipse definiuir Marci12. et Matth.22. vbi ast, in dilectione Dei et proximi legem vniuersam et prophetas contineri: Eandemqueob causam   ad Rom.13. Nemini quicquam debeatis, inquit, nisivt inuicem diligatis: Er subiicit: » Consummatio, πλ́ήρωαα, legis est Dilectio. Ac certum est, si quis per singula legis praecepta cat, cuius rei formam ibidem ostendit Apostolus, expenetur omnia proprie aut ad gloriam Dei illustrandam, aut ad proximi aedificationem tendere. Ergo paucis hisce verbis Apostolus praecidit omnes superuacaneas quaestiones siue disputationes, quae per aduersarios de distributionibus aut diuersis finibus praeceptorum legis moueri poterant. Ac docet simpliciter, homini Christiano non esse plura de lege cognoscenda vel imponenda, quam quae ad officia Charitatis in Deum et proximum conducunt. Et quicumque  hac definitione Apostoli contenti fuerint, ii se multis molestis ac perplexis quaestionibus sentient esse liberatos. Quinetiam ex hoc similibusquelocis manifestum fit, Apostolum, quoties de lege disserit, non de ceremoniis, vt quidam callide interpretantur: Sed de moralibus praeceptis, hoc est, de ipsa medulla atque ossibus legis sermonem facere.  \pend\pstart Ex puro corde, et conscientia bona:) Grauissima occupatione haec adiiciuntur. Aduer sarii calumniabantur perquam exiguum esse, quod Apostolus diceret, finem praecepti esse Charitatem. Neque  enim hanc metiri pro sua dignitate omnes poterant. Atqui Apostolus docet, se, cum de Charitate loquitur, rem intelligere multo per se amplissimam, atque  legi id tribuere, quod dignitatem ipsius mirifice illustrat. Requirit igitur Apostolus Charitatem, non quidem hominum falsa opinione rem humilem vel abiectam, sed ualde excellentem ac Deo inprimis gratam Primo namque  uult Charitatem q̃ pficiscitur ex puro corde. Cor primo loco Des pcepta excipit atque  apprehedit, et put istud se habet ac purn est uelimpurt;  \pend
\vspace{0.5cm}\noindent
\vspace{0.2cm}\rule{1cm}{0.2pt}\\ 
\hspace{0.2cm}\textit{mg}
\footnotesize icben.i 
\normalsize| 
\hspace{0.2cm}\textit{mg}
\footnotesize 3. Ordo do? ειηηα ads tutem expe diens. Ex lege ho mo ad pecca ti agnitiont ducitur. 
\normalsize| 
\hspace{0.2cm}\textit{mg}
\footnotesize 4 ludaeiex lege Mosai ca ustunan aestimant. 
\normalsize| 
\hspace{0.2cm}\textit{mg}
\footnotesize Causaecur Apostolus tam multus si in legis usu explicando. 
\normalsize| 
\hspace{0.2cm}\textit{mg}
\footnotesize 1. Causa, re controuerse prmoloco explicanda. 
\normalsize| 
\hspace{0.2cm}\textit{mg}
\footnotesize 2. Caula, fi nis et usu uniuscuiusque  doctrinae 
\normalsize| 
\hspace{0.2cm}\textit{mg}
\footnotesize semper con siderandus. Finis legis charitas. 
\normalsize| 
\hspace{0.2cm}\textit{mg}
\footnotesize Occupatio. 
\normalsize| 
\hspace{0.2cm}\textit{mg}
\footnotesize Vera Charitas qualis el se debeat. 2Corpurĩs. 
\normalsize| 
\section*{AD TIMOTHEVM CAP. I. }
\marginpar{[ p.22 ]}\pstart ita ab eo postea exeunt cogitationes et actiones bonae aut malae, quemadmodum Christus docet Matth.15. Sed enim purum ac mundum cor non efficitur, nisi ipso Deo autore: Psal.51. Act.15. Ex tali igitur corde debent actiones, quae secundum legem fiunt et charitati adscribuntur, prodire: quemadmodum ipse quoque  Deus legislator exigit, dicens: Diliges Dominum Deum tuum ex toto corde tuo, et ex tota anima tua. Ad - quae verba Apostolum respexisse et allusisse non est dubium Ergo hac particula exclusit Apostolus, primùm omnem impuritatem et malitiam, qua corda omnium non renatorum semper sunt inquinata: Dehinc reiicit omnem hypocriticam diligentiam, qua nonnulli exterius quidem videri volunt legis studiosi, interius vero plurimum a lege et vera charitate absunt: dum nimirum alii tantum seruili metu poenarum opera legis factunt, ut ludaeorum pars maxima assiduis comminationibus et infinitis suorum malis territa: alii inanis gloriae causa, vt Scribae et Pharisaei, Matth.6. alii denique  alios in fines intenti sunt. Huiusmodi autem externam charitatem siue iustitlam legalem Euan gelium non agnoscit, veluti ad legem proprie pertinentem: sed exigit charitatem manantem ex puro corde. Atque  haec demum legi respondebit.  \pend\pstart Secundo vult Apostolus charitatem profluentem ex conscientia bona. De his, que cor paulo ante apprehenderat, conscientiae est exercere iudiclum. Nihil enim aliud est conscientia quam cognitio πρεκτιηο, qua cogitationes atque actiones nostrae exami nantur, et mirabiliter modo accusantur, modon defenduntur, testante Apostolo ad Rom 2.Gentes, inquit, legem non habentes, sibi ipsis sunt lex, qui ostendunt opus legis scriptum in cordibus suis, simul attestante illorum conscientia, et cogitationibus inter se   accusantibus aut etiam excusantibus. Itaque  conscientia cuncta siue facta, siue facienda, accurate expendit, et ratiocinando disquirit, quid recte aut secus factum sit, aut fieri queat. Attamen non omnes eodem modo rationes probabiles ac certas colligunt, quarum ductu sudicium rectum ferunt. Haerent quidem in hominum animis impressae quaedam notitiae, quas tanquam regulas et leges inuiolabiles conscientia sequi deberet, estquehaec animi facultas conseruans et semper suggerens eiusmodi notitias συντήρνσις, a verbo συντηρέω, id est, conseruo, appellata: sed prout affectus praui saepenumero dominantur: Ita conscientia debilitatur ac perturbatur, quo minus ad regulas illas respiciat: Imon obtrudunt vitiosi affectus alias regulas obliquas ac distortas, per quas conscien tia decipitur, vt male iudicet. Eorum proinde, qui ita faciunt, conscientia voluntate ma la et corrupta est. Ad Titum 1.Pollutis et infidelibus nihil est purum, sed polluta est illorum et mens et conscientia.1 ad Timoth 4. Spiritus certo loquitur, quod in posterioribus temporibus deficient quidam a fide, attendentes spiritibus impostoribus ac   doctrinis daemoniorum: per hypocrisin falsiloquorum, quibus cauterio resecta est conscientia, prohibentium contrahere matrimonium, et c. iubentium abstinere a cibis. Vbi exempla proponit Apostolus conscientiae negligentis regulam in verbo Dei propositam, et sequentis peruersas suggestiones corruptorum affectuum. Ex his autem per contrarium intelligemus, conscientiam bonam, quae non discedit ab illis notitiis naturalibus, sed veritati constanter insistit, maxime autem gaudet amplecti regulam verbi Dei, quam imprimis vtilissimum est sequi, aduersus omne periculum dubitationis, erroris siue lapsus. Estquemaxime proprium bonae consclentiae, ad primam veritatem, hoc est, ad Deum ipsum, semper refetre omnia. Ex quo factum, vt qui conseientiam suam testem adhibet, Deum ipsum testem ac iudicem inuocare intelligatur. Ad Rom. 9. Veritatem dico in Christo, non mentior, attestante simul mihi conscientia mea per Spiritum Sanctum, quod dolor mihi sit magnus. Vbi conscientiam suam sistit testem, verum a spiritu sancto institutam et gu   bernatam, ne videlicet conscientiae ipsius tanquam ertanti aut deprauatae minus tribuendum existiment. Et 2.ad Corint.1.Gloriatio nostra haec est, testimonium conscientiae nostrae, que    cum simplicitate et cum synceritate Dei, non sapientia carnali, sed pergratia Dei conuersati   fuerimus in mundo potissimum autem apud uos. In quam sententiam nomnulla dicuntur ad finem   capitis secundi eiusdem epistolae. Et Actorum 23.24.1.Petri3. Eatenus ergo conscientia bona est, quatenus ad veritatem et Deum ipsum se conuertit. Quamobrem vbi Apostolus opera charitatis uult exerceri conscientia bona, omnem excludit ignorantiam, dubitationem,  \pend
\vspace{0.5cm}\noindent
\vspace{0.2cm}\rule{1cm}{0.2pt}\\ 
\hspace{0.2cm}\textit{mg}
\footnotesize 1. Conscion tia bona. Conscientiquias 
\normalsize| 
\hspace{0.2cm}\textit{mg}
\footnotesize Conscientia testis et iudex. 
\normalsize| 
\section*{IN I. EPIST. D. PAVLI. }
\marginpar{[ p.23 ]}\pstart errorem, corruptelam, malitiam, et quicquid veritati ac syncerirati aduersatur. Sunt sins gula charitatis officia tanquam in ipsius Dei conspectu facienda: alias non est quod cogitemus, Deo inspicienti corda, ac iudici conscientiarum accepta fore. Conscientia aduertens se a linea veritatis deflectere, seipsam coram tribunali Dei accusat iugulatque ́.  \pend\pstart Et fide non simulata.) Tertio, requirit etiam Apostolus fidem non simulatam: necessarion hanc particulam addidit: Nam etiansi conscientia nixa non solum naturalibus no titiis, verùm etiam verbo Dei optime iudicet: tamen sine fide ac Dei gratia non datur id vere bonas actiones progredi. Fides namque , sine qua impostibile est placere Deo, et qua accipitur Spiritus Sanctus, sola charitatis vberes fructus producit, vsque  adeo, vt frustra cor apprehendat ac proponat bona, frustra item diiudicet probetqueconscientia, nisi accedat fides, quae bene coepta sunt optime absoluens. Atque  hinc fit, qu fides et on " scientia in Scripturis aliquoties coniungantur, vt paulo post hoc ipso capite, Hoc praeceptum commendo tibi, fili Timothee nempe vt secundum praegressas de te prophetias, milites per eas bonam militiam: retinens fidem et bonam conscientiam, qua repulsa, nomnulli fidei   naufragium fecerunt. Ad Rom.14. Beatus qui non iudicat seipsum in eo, quod probat: " Qui vero ambigit, si ederit, condemnatus est, quoniam non edit ex fide: quicquid ve   ro exfide non est, peccatum est. Exigit autem Apostolus fidem non simulatam. Est sane fides quaedam plerisque  natans in ore, et ab hypocritis multum iactata:sed qualem etiam in daemonibus, vt ait lacobus 2.capite, liceat reperire, quae certe nihil vitale in se habet. Huiusmodi autem fides pro ducet fortasse aliqua opera in speciem bona. Constat namque , quantas res Pharisaei et multi etiam, gentiles facere consueuerint, et adhucfaciant. Attamen non profert opera vere bona, siue opera eius Charitatis quam Aposto lus dicit esselegis finem. A fide, quae Charitati seruit, et per dilectionem operatur, abesse simulatio et hypocrisis omnis debet, nec potest bonis operibus vacua consistere. Itaque  Apostolus tria haec, cor, inquam, conscientiam, et fidem, mirabili prudentia con iunxit, atque  optimo inter se ordine collocauit, et congruentia epitheta singulis attexuit. Cor primo excipit Dei praeceptum de agendis vel non agendis. Conscientia adiuta verbo Dei probat ac definit. Fides quoniam accipit Spiritum Sanctum, producit efficaces actiones. Cor se habet velut persona causam deferre ad superiorem volens. Con scientia est tanquam iudex examinans, cognoscens, et ferens sententiam, non modo exno titiis naturalibus, uerum etiam ex verbo Dei, veluti ex scriptis legibus. Fides postremo exequitur, per quam nisi fieret executio, abs re et cor, et conscientia laborarent. Atque  ita quasi per gradus quosdam ad ea fit progressus, per quae absoluuntur verae charitatis offi cia. Vnde et incrementi rhetorici ratio quaedam notari in hac enumeratione et collocatione potest. Ex his autem omnibus liquet, quae nam sit causa formalis charitatis, et quod non tam exiguares sit vera et perfecta charitas, atque  nomnulli suspicantur. De fide et Cha ritate multi uerbosem disserunt, et audacter tantas sibi uirtutes vendicant, cum interim rebus ipsis absint longissime. Maximae sane res sunt, fides, qualis describitur fuisse in Aprahamo ad Rom 4: et Charitas, de qua hic loquitur Apostolus, proficiscens uideli. cet ex corde puro: et conscientia bona, et fide non simulata: nec potest huiusmodi nisi in his, qui iam sunt iustificati, et Spiritu Dei praediti, inueniri: qua de re paulo ponst. Cum porro tanta res sit uera Charitas, manifestum hinc fit, qu Apostolus legem uehementer conmendet, ac multo ampliorem illi tribuat dignitatem, que  qui alios fines legis comminiscebantur.  \pend
\phantomsection
\addcontentsline{toc}{subsection}{\textit{A quibus quod quidam aberrarunt, deflexerunt ad vaniloquium, volentes esse le gis doctores, non intelligentes quae loquuntur, neq de quibus asseuerant. }}
\subsection*{\textit{A quibus quod quidam aberrarunt, deflexerunt ad vaniloquium, volentes esse le gis doctores, non intelligentes quae loquuntur, neq de quibus asseuerant. }}\pstart Post quam breuissime et tamen plenissime exposuit uerum usum legis, prout lex seruit iustificatis, eos speciatim perstringit, qui hacin parte non recte docuerunt. Estquehaec reprehensio ideo adiecta et satis acris instituta, partim ut sententiam de legis usu expositam omnes amplectantur ac retineant: partim, ut a falsis doctoribus sibi caueant, quos uident foede aberrasse, atque  ita in pericula non leuia coniectos. Ruina aliorum ante oculos posita cautiores omnes debet reddere. Itaque  haec reprehensio tacitam continet admonitiont piorum auditorum, ductam a periculoso. Siquidem iubet falsam doctrinam uitare, ne ndiserimina uaria incidant. Cumprimis elegans est, quod dicit quesdam a fineleg  \pend
\vspace{0.5cm}\noindent
\vspace{0.2cm}\rule{1cm}{0.2pt}\\ 
\hspace{0.2cm}\textit{mg}
\footnotesize a. Fides nos simmalat
\normalsize| 
\hspace{0.2cm}\textit{mg}
\footnotesize Hebr. n 
\normalsize| 
\hspace{0.2cm}\textit{mg}
\footnotesize Fides hypo critarum que  Daemonum neat. 
\normalsize| 
\hspace{0.2cm}\textit{mg}
\footnotesize Charitais uerae uirtur tes. 
\normalsize| 
\hspace{0.2cm}\textit{mg}
\footnotesize Admonitio ducta a periculoso. 
\normalsize| 
\section*{AD TIMOTHEVM CAP. I. }
\marginpar{[ p.24 ]}\pstart fam exposito aberrasse, licut sagittantes, qui aberrant a scopo siue albo. Nam ἀςοχέν est a scopo aberrare, quod qui apud sagittandi peritos, praesertim in ludis publicis facit, ab omnibus irridetur et ignominia quodammodo afficitur. Conuenit autem reprehensionem hanc in duo distinguere membra: quorum altero denotatur, quomodo in materia et reipsa doctrinae aduersarii sint lapsi: altero vero perstringuntur personae, earun demque  vitia, atque  eadem opera tum causa quaedam erroris, tum etiam poena, qua iusto Des iudicio tales afuiciuntur, breuissime attingitur. Lapsi sunt in materia et re ipsa doctre nae de lege: quonsam, vbi debuissent inculcare hanc esse a Deo, vt disceremus ex ipsa offi cia charitatis in Deum et proximum, et sciremus haec ex puro corde, et conscientia bona et fide non simulata exercenda: ausi sunt longe diuersa tradere, nempe legem esse daiam, vt per eam consequeremur iustitiam coram Deo: qua in re Christo magnam inferunt iniuriam, dum iustificationis, sanctificationis, ac redemptionis nostrae ius illi adimunt. I.ad Corinth.1. ac legi falso idem tribuunt. Deinde legis impletionem statuerunt consistere in qualicunque  externorum operum obseruatione, etiansi neque  ad charitatem opera illa dirigerentur, neque  ex puro corde, conscientia bona, et fide non simulata prouenirent: quae a et Nazaraeis et Pharisaeis inculcata fuisse constat. Qui igitur adeo turpiter a veritate aberrauerunt, nónne recte dicentur ad vaniloquium deflexisse, eosquevitandos suadebimus ?̧Admonemur proinde, ut supra quoque  diximus, diligenter expendere, in quacunque  doctrina proposita, ipsam materiam siue rem, ac videre, quid et quousque  synceriter explanatur. Porro vbi perstringit, statim causam quandam lapsus monstrat, dicens: velle eos esse legis doctores. Ambitio siue fastus, et coniuncta his malitia, malorum omnium causa sunt in ecclesiis, atque  hi ante alios ruunt in errores, qui eiusmodi vitiis laborant. Significanter autem ait, velle esse doctores legis, subindicans nimirum, non per ignorantiam esse lapsos: quo nomine excusari vrcunque  possent: sed scientes volentesque ́, aut temere erupisse ad munus docendi, ad quod non erant idonei, aut, siido nei fuerunt, studio tamen corrupisse doctrinam, et voluisse alujs rationibus, quam decebat, legem commendare: Vtrumque  vero iuxta culpandum. Caeterùm poena ambitionis et malitiae subiicitur, cum asserit Apostolus, eos ita esse excaecatos, vt, quibus de rebus, et quid asseuerent, ipsi non intelligant. Etenim tradit Deus impios doctores, praesertim qui non ignorantia, sed per malitiam peccant, in sensum reprobum, vt amentia eorum tandem nota fiat omnibus hominibus: ad Rom.i. 1.ad Corinth.2.3. et 2.ad Timoth.3. Ex his autem omnes discimus, religiose et cum timore considerandum esse admirabile iudicium Dei, quod in iis puniendis exercet: qui agnitam veritatem aut callide corrumpunt, aut aperte abnegant. Quicumque  scientes volentesque  a sana aberrant doctrina, hos necesse est vanilo quiis esse intentos, excaecari ac tandem pudefieri. Exempla non tantum aetas Mosis in Ianne et lambre, vel aetas Pauli Apostoli, in iis de quibus mentionem facit in his epistolis: uerum etiam nostra haec aetas permulta uisenda exhibet. Prae caeteris uero admonen tur eruditi, et qui uel in Scholis, uel in ecclesiis uersantur in docendi munere, ut uerita ti semel agnitae constanter adhaereant, quoque  id possint, Patrem coelestem sine intermissione orent. Interest uero et ecclesiarum, dum pastores siue doctores eligunt, et Episcoporum, dum ad ecclesias id postulantes mittunt aliquos: sapienter prouidere, quales promoueant: maximopere cauendum, ne inflati ac tumidi praeficiantur, utpote qui tam doctrina, que  actionibus offendicula praebere soleant. Doctrina quidem, quia plus eruditionis sibi uendicant, quam reuera habent: Vnde fit, ut asseuerare audeant, quae ipsi non intelle gunt. Actionibus uero, quia, dum caeteris praestantiores haberi ambiunt, mox alios con temnunt, in eosdem affectant imperium, et quo istud assequantur, tragoedias mouere nihi uerentur. Nulla pestis exitium adfert ecclesiis maius, que  si hae committuntur hominibus ar rogantibus  qui se peritiores putant, q sunt: aut, si docti sint, ad ostentationem tamen, ad quaestum, ad corruptelam ueritatis, intellectu, ingenio, et adiunctis donis abutuntur. Hucusque  autem tra dit usum legis Apostolus, quatenus ipsa seruit iustificatis: sicut mox subiicit usum eius ut seruit non iustificatis. Etenim quemadmodum uniuersum genus humanum recte diuiditur in duo genera, nempe pios siue iustificatos, et impios siue non iustificatos: lta duplicem qq usum siue sine legis  postolus explanat Priorest,ut qnfidem iusti sunt, discant ex ea, que  \pend
\vspace{0.5cm}\noindent
\vspace{0.2cm}\rule{1cm}{0.2pt}\\ 
\hspace{0.2cm}\textit{mg}
\footnotesize ετπωqdε 
\normalsize| 
\hspace{0.2cm}\textit{mg}
\footnotesize vplae pa atum bypo ritarum sicistniert ran. 
\normalsize| 
\hspace{0.2cm}\textit{mg}
\footnotesize maiiaee 
\normalsize| 
\hspace{0.2cm}\textit{mg}
\footnotesize Duplexof. fendiculum impiorum doctorum. L In doctrina. R In adionibus. 
\normalsize| 
\hspace{0.2cm}\textit{mg}
\footnotesize Duo genera hominum. Duplexusus lei. 
\normalsize| 
\section*{IN I. EPIST. D. PAVLI. }
\marginpar{[ p.25 ]}\pstart sint officia verae charitatis: Alter, vt impii et qui non sunt iusti, per eam adducantur ad agnitionem peccatorum suorum, atque  inde ad poenitentiam. Ad Rom.3.Est autem prior vsus fideliter explicandus in ecclesia, duabus maximem de causis: Primô, vt certo omni tempore sciant pii, quae nam Dei sit voluntas, et quae opera illi placeant vel displiceant. Quamuis enim hispiritum Dei habeant, tamen, quoniam, quamdiu hic uiuunt, carne sunt circumdati, atque haec cupit eos ab agnitione veritatis abducere: ideo con uenit proposita lege Dei scripta ipsos in agnitione Dei confirmari. Huc pertinent plu ra de vsulegis hancin partem exposita psalmo 119. Lex Domini, inquit, fidele testimonium, praestans sapienciam: Et, Lucerna pedi meo verbum suum, et lumen semitis meis. Secundo, vt excitentur pii ad officia charitatis praestanda, scientes, seruum, qui scit voluntatem Domini, non autem facit, dignum esse plagis multis. Atque harse causam notat Apostolus, cum de lege verba facere volens praeceptum eam vocat. Si Deus praecipit charitatem, necesse est obtemperemus. Egent ergo, qui iustificati sunt, doctrina legis, atque  ex hac discunt officia Charitatis. Quocirca Apostolus quoque in caeteris epistolis, postquam demonstrauit, homines iustificari fide absque  legis operibus, semper subiicit ac docet legis obseruantiam, obedientiam, ac bona opera dilectionis nihilominus a iustificatis requiri. Ad Rom. 13. Nemini quicquam debeatis, nisi hoc, vt inuicem diligatis: nam qui diligit alterum, legem expleuit: Siquidem illud, Non moechaberis, non occides, non furaberis, non falsum testimonium dices, non concupisces, et si quod aliud est mandatum. In hoc sermone summatim comprehendit, Diliges proximum tuum sicut teipsum. Dilectio proximi malum non operatur. Consummatio igitur legis est dilectio. Ad Galat.3.In Christo lesu neque circumeisio quicquam, neque praeputium, sed fides per dilectionem operatur. Et paulo pôst: Per dilectionem seruite inuicem. Nam tota lex in uno verbo completur, nempe hoc, Diliges proximum tuum vt teipsum. Ad eundem igitur modum et praesenti loco Apostolus vsum legis explanat, vt videlicet iustificatis ea seruit, qui ex ipsa discere debent, quae nam sint ompera Charitatis, et quae Deo placeant, et quae displiceant.  \pend
\phantomsection
\addcontentsline{toc}{subsection}{\textit{Scimus autem, quod bona sit lex si quis ea legitime vtatur. }}
\subsection*{\textit{Scimus autem, quod bona sit lex si quis ea legitime vtatur. }}\pstart Generalis sententia, in medio sermonis, de duplici vsu legis interposita: qua tum confirmat, quae de vsu legis exposuit et adhuc expositurus est: tum confutat acreiicit aduersarios, qui asserebant diuersum. Notae vero sunt ac receptae in Dialecticis, atque  disci plinis aliis regula: Cuius sinis bonus est, idipsum quog bonum est. Rursus. Cuius finis siue vsus malus est, id ipsum quog malum est. Item: Si quis abusum ponat, pro vsu non recte ratiocinatur. Itaque hinc colligit Apostolus, se recte docere de lege ac dignitate, quam debet legi tribuere. Quandoquidem ostendit verum vsum legis, quando dixit, finem praecepti esse charitatem: et mox addit: legem non esse iustis, sed iniustis positam. Ratiocinamur igitur, legem esse bonam per se, sed quatenus quis respicit ad legitimos eius fines. Atque hoc modo velut praeoccupatione quadam simul obturat os aduersariis, qui clamitabant, Paulum contra legem quaedam spargere, et imminuere eius dignitatem, maxime, qui non doceret cum ipsis, legem esse datam, vt per eam homines iustifi carentur. Sed respondet Apostolus, se scire aclibenter confiteri, quod lex bona sit, (qua de re etiam ad Rom.7.) verum hoc cumprimis requiri, ut quis ea legitime utatur. Proindee contrario significat, nullo modo audiendos esse aduersarios, qui legitimum legis usum atque  finem non docebant, sed inculcabant potius diuersa. Argumento est, quod coniugatis utitur, dicens, legem esse bonam, si quis ea legitime utatur. Vox legitime, ut a uoce legis deriuatur, ita admirabili proprietate uerum usum legis, et quaecunque  de Lege sunt scitu obseruatuquenecessaria, suggerit. Legitime enim perinde est ac si dicas, iuxta mentem legislatoris, secundum spiritum ipsius legislatoris et per omnia, sicut oportet. Prudentissime ergo Apostolus hanc generalem sententiam interponit, qua omnes, qui de lege uolunt disputare, ad considerandum uerum usum legis, quibuscunque  aliis rebus praetermissis, reuocat. lubemur ueron et nos non tantùm, cùm legem, uerum etiam quoscunque  alios locos doctrinae religionis enarrandos suscipimus, accurate pro\pend
\vspace{0.5cm}\noindent
\vspace{0.2cm}\rule{1cm}{0.2pt}\\ 
\hspace{0.2cm}\textit{mg}
\footnotesize Laeque nus seruit piis ac fide libus. 1. Ex lege Det disciti uoluntas. 
\normalsize| 
\hspace{0.2cm}\textit{mg}
\footnotesize 2. Per leg pii excitan tur ad open charitatis. 
\normalsize| 
\hspace{0.2cm}\textit{mg}
\footnotesize Regule De lechicorum. 
\normalsize| 
\hspace{0.2cm}\textit{mg}
\footnotesize Aconuz uii. 
\normalsize| 
\section*{AD TIMOTHEVM CAP. I. }
\marginpar{[ p.26 ]}\pstart uidere, ut uerum usum explanemus, et tunc modis omnibus urgeamus: necessum est hanc obseruari regulam in recta docendi methodo. Neque  probabile est, ex illorum sermonibus magnum posse prouenire fructum, qui subtilibus, uel superuacaneis occupantur quaestionibus, uerum autem usum proferre negligunt, aut obiter tantum frigideque  attingunt.  \pend
\phantomsection
\addcontentsline{toc}{subsection}{\textit{Sciens illud, quod iusto non sit lex posita, etc. }}
\subsection*{\textit{Sciens illud, quod iusto non sit lex posita, etc. }}\pstart Secundus usus legis, ut nimirum seruit iniustis et impiis, quos oportet lege adduci ac agnitionem peccatorum et poenitentiam. Quo autem dilucidius hunc usum explanet adhibet antithesin: dicit esse legem positam non iustis, sed iniustis. Tameisi uero praecedentibus uerbis diximus, legis usum esse demonstratum, prout seruit iustificatis: ue ra tamen adhuc est Apostoli sententia, quod iusto lex non sit posita. Primo, si propte eos, qui superbe et odiosem perpetuo urgebant legem scriptam eaq cupiebant grauar. uniuersos homines: intelligamus legi oppositum esse praeceptum Charitatis: quande supra admonuimus, Apostolum dara opera uocabulum legis refugisse, atque  usurpasse uocem praecepti:tunc perspicua est sententia, Legem quidem non esse iusto positam attamen praescriptum illi praeceptum de Charitate, quod legis locum obtinet ac plant sufficit. Apostolus ad Rom.2.ait, gentil.bus non esse datam legem illam scriptam in ta bulis, habere tamen eos legem seriptam in cordibus pernaturam. Quanto igitur minus lege illa scripta in tabulis eget homo iustus, qui habet praestantiorem legem in corde eau lemque  descriptam per Spiritum Sanctum, a quo plene docetur, ad bene agendum impellirur, et sapientissime regitur. Itaque  iustus ad legem illam scriptam in tabulis minime respicit, sed intentus est in eius legis finem principalem et spiritualem effectum, nempe Charitatem, quam, qui praeditus est Spiritu sancto corda purificante, (Act.15.) sponte et alacriter praestat ex puro corde, et conscientia bona, et fide non simulata. Atque  hinc fit, quod Apostolus ad Roman.8. et ad Galat.5. et alibi, crebro iustificatos hortatur, ut spi ritu ambulent, opera et fructus Spiritus proferant, haud dubie Spiritus suggestionem prae mandatis legis commendans.  \pend\pstart Secundo, lex non est posita iusto, quia datanon est, ut iustificet, et uiuificet: ad Rom3 ad Galat.3. Sed potius, ut tanquam paedagogus ad Christum deducat, a quo per fidem iustiti?̃ consequimur, quique  factus est nobis iustitia, sanctificatio, et redemptio a Deo. Itaque  qui iam est fide in Christum iustificatus, nequaquam intelligit adminiculo legis denuo justisiczri. Neque  em̃ bona opera iustorum causa sunt efficiens iustificationis perficiendae, sed tantum effecta atque  signa iustificationis iam perfectae. Huc pertinet quod Augustinus ait liber  de spiritu et litera cap.1o. Quomodo, inquit, iusto lex non est posita, Si et iusto est necessa ria, non qua iniustus ad iustificantem gratiam adducatur, sed qua legitime iustus utatur.  \pend\pstart Tertio, iusto non est lex posita: quia a seruitute ac iugo legis est per Christum liberatus, neque  ea sus aut potestatem habet, ut ipsum occidat amplius uel condemnet: ad. Rom z.8.ad Galat.5. Et hoc pacto esse a lege liberatum ampliss imum est beneficium, quoc Christus, qui solus legem exacte impleuit, attulit. Sed iis tantum communicatur qu in Christum uere credunt. Solus Christus efficit, ut occidi et condemnars a lege nor possint, qui ipsi sunt insiti.  \pend\pstart Quarto, addi potest et haec ratio: lusto non est lex posita: quia Ceremonialia et iudicialia praecepta legis sunt abrogata. Conueniebant haec namque  uni tantùm populo, suam habenti ad tempus ecelesiam et rempublicam. Caeterùm ante tempus legis per Mosen promulgatae innumerabiles ubique  uixerunt pii, qui ceremonias illas et iudiciaria plane ignorarunt: Quamobrem neque  ab aduentu Christi pii illis obstringuntur. De abrogatione ceremoniarum Actor.15. et copiose in Epistola ad Hebraeos. Reliqua igitur sunt sola moralia praecepta, quibus pii utique et semper recte utuntur, dum ex iis officia charitatis perdiscunt. Et sic iusto lex non est posita, tota uideli, cet, sed tantùm elus pars. Verùm praecedentes rationes plus habent ponderis: quando neque ab iniustis, de quibus mox sermo instituitur, ceremonialium et ludicialium obseruatio exigitur. Itaque ex his manifestum sit, recte dictum, quod prior usus legis, de quo ante uerbafecimus, d iusos proprie pertneat, et dende uertmnihdomeae.  \pend
\vspace{0.5cm}\noindent
\vspace{0.2cm}\rule{1cm}{0.2pt}\\ 
\hspace{0.2cm}\textit{mg}
\footnotesize 2. Vsus legis quo ad impi ps et iniulos. Iusto lex non est posita, quomodo accipiendum. 1. Homorenatus Spiritus regiur. 
\normalsize| 
\hspace{0.2cm}\textit{mg}
\footnotesize 2.Lexpee goous ad Christum. 1. Corintbus 1. Bona opera in uustifcatis sunt effecta tustificauonis, non causa efficiens iustitiae. 3. Alugo Le gis liberati Iimus per. Corstons. 
\normalsize| 
\hspace{0.2cm}\textit{mg}
\footnotesize 4. Lexquatenus abrogata, secundum ceremo pias uidelitet et iudi, ria. 
\normalsize| 
\hspace{0.2cm}\textit{mg}
\footnotesize Synecdotie. 
\normalsize| 
\section*{IN I. EPIST. D. PAVLI. }
\marginpar{[ p.27 ]}\pstart elie, quod iulto non sit lexposita. Caeterum propemodum ex illis ipsis ratlonibus procliue erit colligere per contrarium, quomo do lex sit posita iniustis et peccatoribus. Primo, lex illa scripta praecipueque  moralis, quam diximus etiam iustis utilem esse et pernecessariam, posita est iniustis, ut nimirum ex ea tanquam fida magistra discant generatim, quid beneuel male fiat: speciatim uero, ut assuescant propria agnoscere peccata. Est sane lex naturalis semel inscripta omnium cordibus: Sed quoniam haec multis modis obscuratur, maxime in impiis, quamuis non tota penitus extinguitur: prodest hominem intueri legem tanquam speculum, atque  ex ea percipere, quae peccata aduersus Deum et proximum designarit.  \pend\pstart Secundo, lex posita est iniusto, ut sciens se contra legem peccare, conuincatur, quod propter peccata grauissimas poenas, tam in hac, quam in futura uita sit promeritus, ideoque  aeterna damnatione dignus. Etenim lex in transgressores poenas horrendas ssatuit. De hoc usulegis ad Roman.7. Peccatum occasione accepta per praeceptum decepitme et perillud occidit. Tertio, posita est lex iniusto, eum uidet se per legem non liberari a malis, non item iustificari, sed potius opprimi, occidi, condemnari, sohcitus fit de impetranda peccatorum remissione et iustitia aliunde acquirenda. lubet igitur lex, postquam hominem extremo deiecit, confugere ad Christum, in hunc fudem coniicere, et propter hunc solum misericordiam expectare. Ad Galath 3. Si latafuisset lex quae posset uiuificare, uerem ex lege esset iustitia: Sed conclusit scriptura omnia   sub peccatum, ut promissio ex fide lesu Christi daretur credentibus. Et paulo post, lex   paedagogus noster fuit ad Christum, ut ex fide iustificaremur,  \pend\pstart Summa autem harum rationum haec est: Lex posita est iniustis, ut per eam adducantur ad ueram perfectamq poenitentiam, atque  ita salutem consequantur. Constat aenim poenitentiae, quae modo salutaris esse queat, duas esse partes, contritionem et fidem. Contritione prosternimur et morimur peccatis: fide rursus erigimur ad accpiendam remissionem peccatorum et nouam uitam instituendam. Exprimuntur autem hae partes crebro in scripturis. Matth.3.lohannes Baptista, agite poenitentiam uitae prioris, in propinquo est enim regnum coelorum. Christus Marci. I. lmpletum est tempus, instatqueregnum Dei, resipiscite et credite Euangelio. Saepe item Apostolus duas istas partes nominibus mortificationis et uiuificationis coniungit: ad Roman. 6.ad Coloss. 2. Quando igitur lex homini impio ostendit sua peccata, atq arguit esse reum mortis ac damnationis aeternae, quemadmodum duabus prioris ratiomnibus dictum est: plane permouet eum, ut ex animo doleat, conteratur, acmortificetur. Quando uero ulterius ducit ad Christum, ut dictum est ratione tertia, fidem in eo exuscitat, ac iubet misericordiam propter Christum apprehendere. Atque  hoc demum facto uera poenitenua absoluitur: frustra peccator aliquis conteritur et mortificatur, nisi insuper ei detur fide erigi et uiuificari. Tantus autem cum sit legis effectus in impiis, merito lex magnifit, ac summo studio omnibus impiis frequenter proponstur. Quando autem ad hunc modum Apostolus duplicem legis usum tradit: alterum, quatenus lex seruit iustificatis, hosquedocet et excitat ad praestandam charitatem: Alterum, ut seruit iniustis, quos adducit ad poenirentiam: abunde exponit quae de usu legis sunt cognitu necessaria. Aduersarii plura uolebant commonstrare, sed abs re, neque  in tradendo legitimo usu persistebant. Constituunt nonnulli et hoc tempore triplicem legis usum: Primum quidem, quo impii homines et iniusti mouentur ad poenitentiam. Secundum uero, quo pii seu iusti instituuntur in officiis charitatis. Tertium denique  quo qui nulla tanguntur iustitiae aut honestatis cura, coercentur aliquo modo a flagitiis, uocantquehunc usum po liticum. Atqui si bene consideramus, non est operaepretium ita partiri. Primo, hunc ier tium usum a primo disiungi, nulla cogit necessitas, quando utrunque  satis amplecti uidetur, qui simpliciter dicit, iniustis legem esse positam. Secundo, quod non uideamus ta libus impiis, qui nulla tanguntur honestatis cura, redigendis in ordinem, magis seruire praecepta iudicialia q̃ poenas ciuiles flagitiosis irrogant, atque  ita nefarios conatus reprimunt quam generalem illam legem moralem Decalogo comprehensam. Tertio: lmo ueron palam im pii atque  in malo obdurau, nulleus prorsus legis teuerentia aut metu a scelenb, deteretur.  \pend
\vspace{0.5cm}\noindent
\vspace{0.2cm}\rule{1cm}{0.2pt}\\ 
\hspace{0.2cm}\textit{mg}
\footnotesize 1. Ralio 9 lexsit pesit iniustis. 
\normalsize| 
\hspace{0.2cm}\textit{mg}
\footnotesize Lex morali utilis est in iustis. 
\normalsize| 
\hspace{0.2cm}\textit{mg}
\footnotesize Lex specult uoluntatis di uine et pec catorum ne strorum. 
\normalsize| 
\hspace{0.2cm}\textit{mg}
\footnotesize 2. Ruiole conuincitur peccati. 
\normalsize| 
\hspace{0.2cm}\textit{mg}
\footnotesize 2. Ratio. Ie deducit pee catorem profigatum ad Christan. 
\normalsize| 
\hspace{0.2cm}\textit{mg}
\footnotesize Pates pan tentiae Contritio c Fides. 
\normalsize| 
\hspace{0.2cm}\textit{mg}
\footnotesize Lae ufa est et iustificats, et initsis. 
\normalsize| 
\section*{AD TIMOTHEVM CAP. I. }
\marginpar{[ p.28 ]}\pstart Quarto addi potest, quod cum satis sit partiri omnes homines in duo genera, piorum inquam, et impiorum, satis quoque  erit duplicem legis vsum pro hominum duobus gene ribus constituere, vt intelligatur lex vniuersis sine exceptione hominibus profutura, si modo ea legitime vtantur. Quinto, sunt omnia officia omnesque  fructus legis in his tribus posita, videlicet, poenitentia, fide, et Charitate, neque  vltra haec exigi quicquam ab vllis hominibus slue impiis siue piis potest. Haec autem ipsa duplici vsu, quem Apostolus tradidit, explicantur. Ergo manifestum fit, Apostoli partitionem de duplici vsu legis sufficere, nec vllam esse necessitatem, quae aliter partiri nos cogat.  \pend\pstart Iniustis et ot sequentibus.) Enumeratio diuersorum hominum contra legem peccan tium, propter quos lex lata intelligitur, vtpote qui si recte vterentur lege, fructum nor madiocrem referre ex ea possent. Collocat autem primo loco duas voces generales, in iustos videlicet et inobsequentes, quibus statim pestringitur propria hominum mali tia et corrupta voluntas, qua ausi sunt legem contemnere, nec patiuntur, se per ean in ordinem redigi. Siquidem lex causam peccati nequaquam praebet, sed ipsorum hom num ingenita iniusticia et praua voluntas legi Dei subdinolens: ad Rom.7. et8. Quod si audire legem tandem vellent, et commouerentur per eam de vitandis peccatis, et ad bene agendum excitarentur. Elegans est ἀντανάκλασις, νόμ ω κῶται ἀνόνοις, quasl dicat. lex posita est exlegibus.  \pend\pstart Impips et peccatorieus.) Distributio, cuius priori membro redarguuntur, qui in primam Decalogi tabulam peccant, quae mandat syncerum Dei cultum, hoc est, veram in Deum pieratem. Posterius membrum eos respicit, qui peccant in secundam tabulam, et crassioribus sant vitiis inquinati.  \pend\pstart Irreuerentibus et pronsanis) Explicatio duorum membrorum proxime positorum per subiectas species. Primumn igicur recenset, quos intelligat impios, irreuerentes nimirum et profanos, qui metu Dei prorsus sunt vacui nulloqueducuntur religionis studio.  \pend\pstart Parricidis et matricidis.) lam, quos intelligat peccatores ac reos transgressionis secundae tabulae. Eo autem ordine hic enumerat diuersa hominum peccantium genera, quo conspiciuntur praecepta secundae tabulae collocata. Nominat patricidas et matricidas, propter illos omnes in genere, qui peccant contra quintum praeceptum decaIogi, quod in secunda tabula primum est, nempe, Honora patrem et matrem. Nominat homicidas, qui peccant aduersus sextum praeceptum, Non occides. Dehinc scortatores et masculorum concubitores, qui peccant in septimum praeceptum, Non moe chaberis. Tum plagiarios, qui praecepti octaui tenentur rei, quod est, Non furtum facies. Denique  mendaces et periuros, qui peccant in praeceptum nonum, quod est, de non ferendo falso testimonio. Estquehic ordo ob duas maxime causas obseruandus. Prima, ut assuescamus contra ac multi faciunt, recte digerere Decalogi duas tabulas, et in singulis certum praeceptorum numerum atque ordinem. Secunda, quandocunque  nostram ipsorum probare volumus conscientiam, et peccata coram Deo confiteri, discamus hacratione per singula Decalogi praecepta progredi, et in singulis memoria repetere, quam multa et quam horrenda peccata contra eadem admiseris mus. Hac via et permouebimur ad scelerum poenitentiam atque misericordiam a Deo postulandam, et bona Dei lege legitime vtemur. Si quis alias vult formas examinandi et probandi semetipsum praescribere, quam quae verbo Del nobis commendantur, is uehementer errat. Tutissimum est sequi ductum praeceptorum Decalogi, quando ad ea facile est reducere, et omnis generis bona opera, quae a Deo praecipiuntur, et omnis generis peccata, quae a Deon prohibentur. Contrariorum namque  cadem tenenda est haecratio.  \pend\pstart Et si quid aliud est, quod sanae doctrinae aduersetur.) Conclusio, qua signisicat eodem modo iudicandum tum de reliquis praeceptis, si quae uisus est praetermittere, tum de infinitis aliis peccatis, quae aduersus singula praecepta designantur, atque enumerare est impossibile. Estqueillud notandum, quod iudicium de peccatis uult peti ex doctrina sana. Etenim ubi doctrina sana non est, ibi rectum de bene uel male factis iudicium nor est. Corrupta doctrina con perducuntur hemines, ut quae reuera peccata sunt admode  \pend
\vspace{0.5cm}\noindent
\vspace{0.2cm}\rule{1cm}{0.2pt}\\ 
\hspace{0.2cm}\textit{mg}
\footnotesize Peccautae praua hominis uoluntate pullulat, non ex lege Dei saicta. 
\normalsize| 
\hspace{0.2cm}\textit{mg}
\footnotesize νίμω na͂τα ariuen. 
\normalsize| 
\hspace{0.2cm}\textit{mg}
\footnotesize Exigt. Det iudicium fumitur de Jactis bomimm. 
\normalsize| 
\section*{IN I. EPIST. D. PAVLI. }
\marginpar{[ p.29 ]}\pstart detestanda, pro peccatis non habeantur. Hinc uidemus superstitiones, idololatriam, cultum imaginum, et multa eius generis ibi locum obtinere, nec reprehendi, ubi doctrina nondum est repurgata. Propterea etiam supra diximus, quandocunque  ordinanda uel reformanda est ecclesia, initium a doctrina debere fieri. Non est spes, posse eradicari uitia et abusus, antequam uerbum Dei pure promulgatur. Sed neque aduerlarii ipsius Apostoli, qui legitimum usum legis non tradebant, recte poterant de uitiis et uirtutibus pronunciare. Exemplum sumi potest de Pharisaeis, scribis, Legisperius, quos Christus reprehendit, quandoquidem hypocrisin suam atque impuritatem cordium non existimabant esse peccata, et cum de externa quadam iustitia coram hominibus essent utcunque  soliciti, de interna, quam Deus potissimum requirit, ne cogitabant quidem. Causa est, quamuis haberent doctrinam et praedicarent, destituti tamen etoctrina sana, conuertebantur ad uaniloquium, ad fabulas, quibus doctrina sana opponitur.  \pend\pstart Secundum Euangelium.) Necessaria et utili occupatione adiicitur particula de Euangelio. Dicendum uidebatur: si quid aliud est, quod sanae doctrinae aduersatur secundum legem: Sed maluit Apostolus dicere, secundum Euangelium. Id autem quatuor de caussis. Prima, quemadmodum supra admonuimus, familiare est Apostolo interdum data opera uocem legis refugere, alia peraequeaut magis accommoda substituta. Secun da,ut ostendat, quanta sit dignitas Euangelii, quando in eo etiam continetur atque  abunde traditur, quicquid ad ueram pietatem et salutem tradit ipsa lex. Sic autem Christus Matth.22. et alibi mirabili compendio summam legis et prophetarum complectitur. Quin etiam Euangelium arguit peccata, adducit ad poenitentiam et fidem, non minus quam lex, quod loannis Baptistae et Christi conciones declarant. Tertio, hac uia excludit odiosas disputationes aduersariorum, qui primo atque ultimo loco senper urgebant legem, atque omnes ad eius obseruationem, quamuis legitimum usum non ostenderent, cupiebant pertrahi. Quarto, satisfacit piis, qui ipsum inique traducebant. tanquam si doceret cum lege pugnantia, essetque  legis ex professo hostis. Ego, inquit, doceo Euangelium, quod eùm non sit legi contrarium, imo uero prope ostendit legitimum usum legis: ego quidem eadem opera doceo, quae ad legis pertinent substantlam, nenpereprehenssonem ac damnationem peccatorum, poenitentiam, fidem, officia Charitatis ex fide manantia, et si quae praeterea ad salutem necessaria in illa proponuntur.  \pend\pstart Gloriae beati Dei.) Notandae sunt hic plures rationes, quae dignitatem Euangelii uehe menter amplificant. Prima est, quod doctrina sana debet aestimari ex Euangelio. Eteuim Euangelium a Christo et Apostolis traditum regula est et amussis, ad quam examinari doctrinam omnem oportet. Quicquid pugnat cum Euangelio, pro doctrina sana admitti nullo modo debet: ad Galat.1.Secunda: Euangelium appellatur Euangelium gloriae Dei. Deus quidem omnia ad gloriam suam fecit et adhuc facit, sed nulla re magis illustrata est ipsius gloria, quam ex Euangelio, hoc est, bono illo nuncio de salute humani generis per Christum parta. Itaque et Euangelium ipsum perse plurimum gloriae habet, et Euangelio ipsius Dei gloria maximopere celebratur. De gloria ipsius Euangelii 2.ad Corinth.3. De gloria Dei patris per Euangelii manifestationem, Christus lohan.17. De gloria silii lohan.1. Tertio, uocatur Euangelium gloriae beati Dei:primo, ut sciamus Euangelii autorem esse ipsum Deum, quod ipsi placuit etiam perfilium suum unigenitum Mundo patefacere: lohan.1 Hebr.. Secundo, ut agnoscamus in solo Euangelio proprie ac plene explicari, quomodo Deus solus per se et sua natura perfecte beatus sit, solus alios quoque, quos uult, beatificet atque aeterna beatitudine donet. Tertio, ut certi simus, in eodem Euangelio nobis demonstrari modum, quo beatitudinem illam consequamur. Euamngelium est potentia Dei ad falutem omni credenti: ad Rom.1. Quarto, patet etiam ex his, materiam siue subiectum Euangelii esse Deum ac beatitudinem in eodem et per eundem. Tanta proinde cum sit Euangelii dignitas, aequum sane est, ut summo in honore habeamus atque constanter amplectamur. Inprimis autem studebimus, omnem doctrinam per illud examinare, ac uera quidem dogmata suscipiemus, falsa uero, deinde uitia seu peccata quaecunque  auersabimur. Adhaec sic incumbemus sacrosancto Euangelio, idemquepraedicabimus, ut scopus noster sit, Dei  \pend
\vspace{0.5cm}\noindent
\vspace{0.2cm}\rule{1cm}{0.2pt}\\ 
\hspace{0.2cm}\textit{mg}
\footnotesize Matth.15.23. 
\normalsize| 
\hspace{0.2cm}\textit{mg}
\footnotesize Occupatio. 1. 
\normalsize| 
\hspace{0.2cm}\textit{mg}
\footnotesize 2l Euagelii dignitas. 
\normalsize| 
\hspace{0.2cm}\textit{mg}
\footnotesize Luc.3. Matth. 3.11. 3. 4. 
\normalsize| 
\hspace{0.2cm}\textit{mg}
\footnotesize Euangelium legi non est contrarium, docet enim legis usum. 
\normalsize| 
\hspace{0.2cm}\textit{mg}
\footnotesize 1. Doctrina sana aestimatur ex Euagelio. 
\normalsize| 
\hspace{0.2cm}\textit{mg}
\footnotesize Euangelio gloria Dei illustratur. 
\normalsize| 
\hspace{0.2cm}\textit{mg}
\footnotesize Euangelium gloriae Dei beati quonto do dicatur. 
\normalsize| 
\hspace{0.2cm}\textit{mg}
\footnotesize 2. In Buaugelio discitur beatitudo Dei et hominum. 
\normalsize| 
\hspace{0.2cm}\textit{mg}
\footnotesize Modus consequendae bea titudinis in Euangelio con monstratur. 
\normalsize| 
\hspace{0.2cm}\textit{mg}
\footnotesize 4. Subiectum Euangelii est beatido et Deus, per que illa obtngit. 
\normalsize| 
\section*{AD TIMOTHEVM CAP. I. }
\marginpar{[ p.30 ]}\pstart gloriam amplificare in terris, et nostram salutem aeternamquebeatitudinem quaerere in coelis. Quod si insuper dubitatio de doctrina, vel tentatio aliqua ad peccata suboria tur, vel etiam ad desperandum de Dei misericordia animos nostros moueri vllo modo contigerit, ex Euangelio certitudinem doctrinae sanae, multiplicia remedia aduersus tentationes, spem gloriae beati Dei, denique  omnis generis consolationes haurire nobis licebit. Viderint porro omnes, quam grauiter illi peccent, qui tum quouis pacto impediunt Euangelii eursum, tum Deo conantur adimere gloriam, tum hominibus institutionem et consolationem in hac vita, denique  et beatitudinem post hanc vitam inuident. Qui Euangelium inglorium vel infrugiferum arbitrantur, tota errant uia, estque ; omnibus contemptoribus durum contra stimulum calcitrare.  \pend\pstart Quod concreditum est mihi.) Addit de suo officio, idque  tribus de causis. Prima, vt au toritatem sibi diuendicet fidemqueipsi habeant auditores. Quem constat ad munus de cendi Euangelii diuinitus esse vocatum, par est hunc audiamus. Secunda, vt intelligam omnes, doctrinam, quam de lege eius quevsu exposuit, puram esse ac synceram. E diuei so autem doctrinam aliorum esse vanam et fabulis nihilo praestantiorem. Tertia, v viam sternat ad ea quae mox subiicit. Solet namque  Apostolus orationis suae partes perinde connectere, ac si ex sese inuicem fluerent. Caeterùm ob duas illas causas priores decet nos quoque  omni ope eniti, quo non modo pro virill, quae nostrae conueniunt ve cationi praestemus, verùm etiam aliis tam de vocatione nostra, quam de doctrina atque  actionibus reddamus rationem, et satis quibuscunque  hominibus faciamus.  \pend\pstart Etgratiam habeo.) Quo suam apud omnes tueatur autoritatem, nec temere quisqu? credat aduersariis, varia et sinistra de ipso spargentibus, plenius exponit, quomodo sib concreditum sit Euangelium, ostendens a quo et quomodo uocatus fuent, quomode item factus sdoneus, uariisque  donis ornatus a Deo, ut munus iniunctum fideliter exe queretur. Sed illud cum primis notandum, quod sic istud facit Apostolus, talemquefos mat orationem, ut eadem opera demonstret, quomodo legis usum duplicem a se proxi me descriptum plane sit in semetipso expertus. Atque  hoc pacto dogmata, quae tradidit, valde prudenter ad πράβιν transfert, et quomodo illorum fructum percipere queant conscientiae, patefacit: quod semper fieri a doctoribus maximopere est necessarium. Ante omnia igitur agit Deo gratias, qu dignatus sit ipsum et vocare ad sublime ministerium Apostolatus, et donare, vt in eodem summa fide constantiaque  uersaretur. Hac autem ratione non tam cauet, ne aduersarii calumnientur, ipsum quadam arrogantia munus suum inaniter iactare, quam perstringit falsos doctores, qui neque  erant a Deo uocati, neque  intelligebant, de quibus asseuerarent, quod supra dictum est. Sic autem alibi quoque pro sua uocatione et ministerio gratias uni Deo agit: Lad Corinth.3.2.ad Corinth 2. Eo autem iustius agit Apostolus Deo gratias. quo magis admirandum erat, eum, cuius uita anteacta erat impia, superstitiosa, et ab Euangelii professsone aliena, nunc pium esse idemque  Buangelium, quod pridem hostiliter persequebatur pure, fideliter, quanuis cum maximis molestiis et periculis docere. Pertinent huc, quae ab codem referuntur Act.22. 1.ad Corint. 15. ad Galath. 1.Siuero ad usum legis libet respicere, ideo quoque  agit Apostolus Deo gratias, quod lege recte uti, hoc est, adduci per legem ad agnitionem peccatorum et poenitentiam, item officia charitatis stenue exequi, nemo nisi qui speciali Dei gratia adiuuatur, possit. Discimus autem hoc loco, primum, debere nos diligenter considerare et agnoscere, quae quantaque  beneficia Deus a pueritia inde in nos contulerit. Dehinc, quomodo res ad agnitionem Euangelii, praeterea ad certum uitae genus, in quo queamus ipsi placere nullis nostris meritis praeeuntibus aut sequentibus, sed per solam misericordiam clementer uocarit. Prohis autem omnibus Deo nos semper gratias agere oportet, quinetiamut in posterum quoque  tempus eadem nos seruet ac prouehat, indesinen ter precabimur. Si ingrati suerimus, Deus a nobis dona sua iterum auferet.  \pend\pstart Qui me potentem reddidit.) Instituit agere gratias pro eo, quod constitutus sit a Deo in ministerio apostolico: praemittit autem hanc rationem ductam a Simili seu ab Effe ctu, ex quo ipsum a Deo positam in eo munere cognosci queat. Certum argumentum meApostolatus est, inquit, quond Deus me potente reddiderit, et mirabilicer sitperme te  \pend
\vspace{0.5cm}\noindent
\vspace{0.2cm}\rule{1cm}{0.2pt}\\ 
\hspace{0.2cm}\textit{mg}
\footnotesize 1. Autoritas muneris Apostolici a Deo uocante pendet. 2 Pura doctrina delege et eius usu in Euan gelio cernitur. 3. 1.Petris. Vocatio Pau li diuina et coelesus. 
\normalsize| 
\hspace{0.2cm}\textit{mg}
\footnotesize a Simili siue ab effedu
\normalsize| 
\section*{IN I. EPIST. D. PAVLI. }
\marginpar{[ p.31 ]}\pstart Euangelii praedicatione operatus. Est autem eximia et singulars potentia acuirtute o-. pus his, qui non modo in munere apostolico, uerum etiam in quacunque  functione eccle siastica aliquid laude dignum et cum fructu auditorum efficient: de qua potentia Matth.7. Marci1. Luc.4.Act.4. 1.Corinth.1.4. Atque  hanc admirabilem potentiam in praedicatione Apostolorum cum impii considerarent, permoti sunt ad Euangelium suscipiendum. Non potest autem potentia haec humana industria, arte, uel consiliis acquiri, sed solius Dei obtingit bonitate. Ideoque  Apostolus pro ea agit gratias Deo, simulque  pro certissima nota ueri Apostolatus uult haberi. Quamobrem i.ad Corinth. 2. opponit eam humanae » sapientiae et arte contextis sermonibus: Ego, inquit, non ueni ad uos cum eminentia   sermonis aut sapientiae, sed cum ostensione spiritus ac potentiae. Ac debent non Apostoli modo, uerumetiam quicuque  ministri Ecclesiae se psidio huius spiritualis potentiae in periculis et insidiis, quae per malos homines struunt, consolari et munire, nec unquam frangant̃ animo, aut in officio agant negligentius. Mundus uerbo acministris ueibi ppetuo insultat, sed inuitis omnibus consistunt incolumes, quibus Deus spiritualem hanc addidit potentiam.  \pend\pstart Christo Iesu.) Causa finalis potentiae diuinitus concelsae: ex qua limiliter iudiciu peti debet ad discernendum ueros Apostolos et doctores a falsis. Datur namque  potentia illa, tum ut de Christo et beneficiis per eundem allatis quaerant Apostoli testimonium, tum quia ad Christi gloriam referri, tam sermones quam actiones Apostolorum uniuersae debent. Qui hoc faciunt, merito legitimi et a Deo missi ministri agnoscuntur. E contrario autem pseu doapostoli et personati doctores iudicabuntur, qui non praedicant Christum purumque  eius Euangelium, sed q̃ arbitrantur accepta fore hominibus: deinde non gloriam Christi, sed suam praecipue quaerunt, priuatisque  inhiant commodis: ad Rom.6.ad Galat.5.ad Philip3.  \pend\pstart Qui side le me iudicauit ponendoin minislerium) Nunc declarat, quomodo expertus lit in semetipso usum legis, prout lex seruit iustificatis. Deus, inquit, posuit me in ministe rium Apostolatus, reddidit etiam me potentem Christo, (sunt autem haec hominis iustificati signa) hinc porro iudicabat et comperit me fidelem in officio. Etenim fidelis a Deo est iudicatus, quandoquidem cum Deus spiritu suo ipsum ornasset, potuit omni tempore ad legitimum finem praecepti intentus esse, et Charitatis officia, quae amplissimum suum munus decebant, ex puro corde, et conscientia bona, et fide non simulata praestare. Suntquehaec duo in uerbis Apostoli prudenter consideranda: Alterum, quod primo lo co constituit Deum operantem in cordibus, largientem Spiritum Sanctum, reddentem idoneos ad munus, cui ipsos destinauit: In eandem sententiam 2. ad Corinth. cap.1. ac 3.pluribus uerbis exponit Apostolus, quomodo Deus ministros Euangelii idoneos faciat. Alterum, quond secundo loco dicit, se a Deo iudicatum fidelem. Nimirum causa praecedere effectum necessario debet. Neque  enim uocatus fuit Paulus ad Apostolatum et potens redditus, quiaiam fuerat fidelis, (qui namque  probari eius fides potuisset, antequam in munere esset constitutus?) Sed ideo potius fidelis est ludicatus, quia a Do mino in ministerio erat positus et factus idoneus: Electio atque  uocatio praecedit, fides sequitur. 1.ad Corinth.4. Sicnos aestimet homo, ut ministros Christi et dispensatores mysteriorum Dei. Quod superest autem, illud requiritur, in dispensatoribus, ut fidus, fidelis sit in suo ministerio, non ferri acceptum legi, sed ipsi Deo et Spiritui a Deg accepto: lege nihilominus admoneri pios, tanquam scripta tabella memoriae caussa ante oculos posita, ut diligentes sint. Itaque  apposite usus legis in iustificatis hic traditur. Non enimut bene faciant, habent a lege, sed a Deo eiusq spiritu: est tamen lex vice monitoris, qui suo modo ad officium vnumquemque  excitat. Atque ita pariter Deg tribuitur sua gloria, et legis dignitas salua retinetur. Caeterùm vt haec et proxime sequen tia verba sentiam explicanda de vsu legis et propria ipsius Apostoli experientia, quatuor me mouent rationes. Prima, si, quae Paulus hic de se subiicit, hoc modo exponantur, optime cohaerent et consentiunt cum iis, quae de lege hactenus dicta sunt: Si autem aliter interpretemur, non perinde cohaerent. Semper autem debemus prudenter dispicere, quomodo sententiae in quocumque  scripto sint innicem colligatae. Ac Paulus, si quis alius, in con nectendis suis sententiis atque  orationis partibus mirus est artifex. Secunda:consuetudo est  \pend
\vspace{0.5cm}\noindent
\vspace{0.2cm}\rule{1cm}{0.2pt}\\ 
\hspace{0.2cm}\textit{mg}
\footnotesize Confolaic piorum ex uocationis certitudinc. 
\normalsize| 
\hspace{0.2cm}\textit{mg}
\footnotesize Caussa finc lis uocationis diuinae. 
\normalsize| 
\hspace{0.2cm}\textit{mg}
\footnotesize usdexer Ra. 
\normalsize| 
\hspace{0.2cm}\textit{mg}
\footnotesize Solu Dei est aptos et idoneos red dere ad mus nus diuinin capescendin 
\normalsize| 
\hspace{0.2cm}\textit{mg}
\footnotesize Electio et uocatio fidi pracedit. 
\normalsize| 
\hspace{0.2cm}\textit{mg}
\footnotesize Lex tabeil, memoriae. 
\normalsize| 
\hspace{0.2cm}\textit{mg}
\footnotesize Piiut bene faciant nor habent a le ge sed a deo. Constat ta men legi nihilominus Iu dignitas 
\normalsize| 
\section*{AD TIMOTHEVM CAP. I. }
\marginpar{[ p.32 ]}\pstart Apostoli, quandocunque  dogma aliquod exposuit, protinus subiicere exhortationem, vel aliud quippiam de eius dogmatis recto vsu, et quomodo ad conscientias instituendas, vel ad vitam moresquetransferri debeat. Ad Roman.6. docet efficacia mortis Chri sti aboleri peccatum, acnobis, dum baptizamur, applicari hoc Christi beneficium. Subiicit igitur exhortationem de eius doctrinae usu: dicens, Ne regnet igitur peccatum in mortali vestro corpore, etc. Ad Galath.5. plenius disserit de libertate a iugo legis» per Christum facta. Tandem vero addit de libertatis illius vsu: Vos, inquit, in libertatem vocati fuistis fratres: tantum ne libertatem datam putetis carni pro occasione, sed per charitatem seruite alii aliis. Simili igitur modo post explicationem dogmatis de legisωνsu subiicit hoc loco Apostolus propriae experientiae testimonium.  \pend\pstart Tertia ratio: Apostolus in disputatione de vsu atque effectis legis ad Roman.7.omnia ad se transfert suaeque  accommodat personae, dicens, Peccatum non cognoui nisi per le, gem, et peccatum per praecept decepitum et occidit me: et, Lex quidem spiritualis est, egovero homo sum camnalis. Et ad eum modum eruditissimi quique , tam veteres, in primis D. Agustinus liber  contra lulianum, quam recentiores interpretantur. Quid ni igitur statuamus, Apostolum hoc quoque  loco, quae dixit de vsu legis, suae accommodare personae, et probationem de propria experientia adferre?  \pend\pstart Quarta ratio: Paulon post sermonem instituit de Euangelio siue de exuberante gratia per lesum Christum: atque  ibi statim addit exemplum de semetipso, referens nimirum, quomodo vsum siue effectum Euangelii talem in se fuerit expertus: Misericordiam, inquit, sum adeprus, vt in meostenderet lesus Christus omnem clementiam. ltaque  quae hic de seipso ex-- ponit, merito accipiemus ad demonstrandum vsum legis esse adhibita, atque  Apostolum maxime cupere, vt quae ipse de vsu legis vera ac certa esse in sese comperit, ex iisdem caete ri quoque  capiant, ac de seipsis iudicium facere assuescant. Caeterum in verbis Apostoli, quibus ait, quod Dominus fidelem ipsum iudicarit ponendo in ministerlum, primon commendatur nobis doctrina de vocatione. Ac scire debemus ad ministerium nullos legitimem peruenire, aut cum fructu aliquid in eo praestare, nisi quos ipse Dominus in eo posuerit. Quamobrem e contrario vehementer improbantur, neque  sperare debent quicquam lau de dignum se effecturos, qui non vocati accurrunt, vel verius, nescio quibus artibus irron punt ad functiones per fas nefasque  occupandas. Secundon, discimus, nullos perse, sua pru dentia vel consiliis idoneos siue fideles esse, sed tales effici quatenus Deus fideles iudicat, et suo iudicio, id est, electione, vocatione, institutione ac donis necessariis additis format fideles. 1.Samuel.3. vocatur Samuel, cum nondum fuisset ei reuelatum Dei verbum, postea autem id Eli reuelatur. lerem.1.leremias vocatur ad munus propheticum agno scitis se minus aptum, sed Dominus tangit ipsius os additquevires. Sic etiam Paulus Act.9. vocatur a Christo, atque  Ananiae dicit Christus, illum elegi, siue est vas electum, vt portet nomen meum coram gentibus ac regibus. Paterfamilias commendat famulo cuipiam munus administrandum in domo, quamobrem instituit et praescribit faciem. da omnia, simulquedat potestatem illud exercendi. Istud vero dum facit, declarat quod famulum illum iudicet fidelem. Quanquam hic incidit aliquod discrimen, siquidem paterfa milias non potest ea ratione statim facere idoneum, qua Deus, sed tantùm inststuit informatque ́, subinde etiam praecedentia quaedam obseruauit, ex quibus idoneum fore suspicatus est, Deus autem quamprimum aliquos vocat, simul facit idoneos, necrespicit merita, vel vllas omnino res antecedentes. Tertio, cum Apostolus et hic et alibi passim de ecclesiastica functione sermonem faciens, utatur uocabulo διιακονιας ministerii, par est meminerint omnes Ecclesiarum ministri, se maximarum refum admoneri. Primo, ut sciant, ingentes labores, curas et pericula adiuncta esse el officio, ac diligenter expendat omnia, neque  interim grauate quae se decent faciant. Secundo, ut humiles sint, nec sese supra caeteros efferant, imo ne patiantur signa illa arrogantiae uel ambitionis in se conspici. Tertio, oculos semper intendant in potentissimum illum, atque  omnium, qui cum potestate sunt, communem Do minum, ipsum uidelicet Deum, cuius se esse ministros uniuersi non tantum eccleslarum anti stites, uerumetiam gubernatores rerumpublicarum fateri humiliter debent: ad Rom.13.Quar to et postremo ueibis Apostoli admonentur oës, qin aliqua dignitate sunt constituti,  \pend
\vspace{0.5cm}\noindent
\vspace{0.2cm}\rule{1cm}{0.2pt}\\ 
\hspace{0.2cm}\textit{mg}
\footnotesize Qui a Deo non uocantur d ministeri um, cum frus ctu id sustihere nequeBut. 
\normalsize| 
\hspace{0.2cm}\textit{mg}
\footnotesize ab aecononia siue ne fimuïari. 
\normalsize| 
\section*{IN I. EPIST. D. PAVLI. }
\marginpar{[ p.33 ]}\pstart vtfideles else in officio pergant, donis a Deo sibi collatis recte vtantur, deniq quicquid faciunt, ad Dei conuertant gloriam et hominum vtilitatem, quod quo possint, a Deo precationibus impetrare singuli contendent.  \pend\pstart Qui prius eram blasphemus et persecutor et violentus.) Altera pars de Vsu legis, pro vt ea seruit iniustis. Ex lege agnoscit Apostolus, quam grauiter ante vocationem et iusti ficationem peccauerit, et ad peccatorum confessionem adducitur. Obnciat forte aliquis, in conuersione Pauli non spectari eundem modum, qui in illis, quos suis sermonibus ad poenitentiam permouerunt Apostoli Act.2. et3 deprehenditur: vel in aliis, quos piae conciones vel meditationes de lege diuina, ad peccatorum agnitionem et confeisionem inducunt. Siquidem Apostolus subito et insolito modo, ac potius miiraculo quam confuetudine ordinaria est ad resipiscentiam ac vitam meliorem pertractus. Respondeo, verum quidem est, aliquid esse dissimilitudinis in conuersione et poenitentia Pauli Apostoli et aliorum: Sed id magis positum est in circumstantiis, que  in negocii summa. Negari enim haud potest, Paulum Vim legis, prout ea seruit iniustis, mirabiliter ac longe efficacissime esse in se expertum, primo, quando, vt scribitur, Act.9. circumsulgurauit eum subito Lux de coelo et collapsus estin terram: praeterea audiuit vocem dicentem. Saule Saule quid me persequeris? Certum est, tum audiuit Apostolus grauem de lege concionem, quin etiam sensit toto animo et corpore suo acerbam legis iram et vindictam: audiuit ipsum legislatorem ipsumqueDeum sua voce damnantem in ipso peccata ac seuerissime ea vlciscentem. Quid autem obsecro est vim legis experiri, si hoc non est? Ac propterea etiam Apostolus statim edit confessionem, qua contestatur, se agnoscere ac dolere ob peccata ante acta, et velle deinceps vitam aliter instituere, quando nimirum prostratus in terra, tremens et obstupefactus, ait,   Domine, quid me vis facere ? Conferantur igitur haec de Conuersione Pauli cum iis quae actor.2.et3.referuntur de conuersione illorum, qui auditis Apostolorum sermonibus compuncti sunt corde, et humiliter quaesierunt, quid se facere oporteret, quo sa lutem consequerentur. Et haud dubie nullum in his, quod ad confessionis modum attinet, discrimen inuenietur, nisi quod in causa Pauli circumstantiae quaedam magis sunt illustres. Ergo Apostolus ad comprobandum vsum legis, vt seruit iniustis ac peccatoribus, suum exemplum atque  experientiam perapte adhibet. Secundo Apostolus etiam a vocatione et iustificatione eundem legis vsum in seipso demonstrat, quando ante actorum criminum recordatione affectus ob ea se dolere, et nisi per misericordiam Dei clementer esset absolutus, damnationem aeternam finisse promeritum confitef. Etenim agnitio, et confessio peccatorum semper a lege est, et quotiescumque  peccata sua quisquam explorat, toties vis legis, prout ipsa iniustis et peccatoribus posita est, sese in eo exerit. Porro in hac confelsione Apostolus non tegit, non arnficiosa ambage circumloquit, non alio quouis modo extenuat sua peccata, sed simpliciter q̃ grauissima sunt quasique  capita omnium fla gitiorum, recenset. Id autem tribus de causis. Prima. hoc facto clarius ostendit effectum legis, quae hominem humiliat, psternit et corda plane conterit. Quo maiora sunt peccata, eo maior sequi debet humiliatio, maxime cum, vt Psal.51.testaf, cor contritum et humiliatum Deo inprimis placeat. Secunda, retundit aduersarios, qui cum uerum usum legis nequaquam tradebant, sed potius e contrario quasi homines legem exacte implerent, contendebant ex lege pa rari iustitiam, atque  ita tumidos reddebant et inflatos. Exemplum habemus in Phariseo, qui vendicabat sibi iustitiam, et prae se contemnebat publicanum. Tertia ad amplificandam Dei misericordiam, de qua mox uerba facit. Quanto grauiora peccata remittuntur, tanto magis illustratur Dei misericordia, per quam illa remittuntur, quemadmodum docet Christus Lucae7. quando mulieri peccatrici pecçata remisit. Sed est insuper notandus in hac confessione ordo decentissimus. Primum fatetur Apostolus, se fuisse blasphemum: id autem est peccatum grauissimum, commissum in Deum, et pugnans cum praeceptis primae tabulae. Secundo fatetur, se fuisse persecutorem et uiolentum; haec contra homines et quidem pios Christum consitentes: quamobrem sunt multon grauissi ma, committunturquecontra praecepta secundae tabulae. Constituitigit Apostolus sereum et transgressorem legis uniuersae. Nec difsicile esset demonstrare, quomo do in hacbreui  \pend
\vspace{0.5cm}\noindent
\vspace{0.2cm}\rule{1cm}{0.2pt}\\ 
\hspace{0.2cm}\textit{mg}
\footnotesize Lex quoniodo seruiat iniustis. 
\normalsize| 
\hspace{0.2cm}\textit{mg}
\footnotesize Apostolus Pallus Inusitatarationad resipiscentiam et munus pertrahitur. Diuersitas uocationis Pauli in cirsita est. sita est. 
\normalsize| 
\hspace{0.2cm}\textit{mg}
\footnotesize aguitio et cosessio pec. cai exlege. 
\normalsize| 
\hspace{0.2cm}\textit{mg}
\footnotesize Lex humiuat, proste nit et conterit cor bo minis. 
\normalsize| 
\hspace{0.2cm}\textit{mg}
\footnotesize Lucais 
\normalsize| 
\hspace{0.2cm}\textit{mg}
\footnotesize ordoconfe fious Paui ne. 2. Blasphe mus. 2. Persecutor. 
\normalsize| 
\section*{AD TIMOTHEVM CAP. I. }
\marginpar{[ p.34 ]}\pstart Enumeratione lateant horrenda peccata contra singula Decalogi praecepta. ltaque  ex his animaduertere licet, Paulum uim legis peccata redarguentis in peccatoribus, in se efficacissimam sensisse, et proprio didicisse experimento legitimum ac uerum legis usum.  \pend\pstart Sed et misericordiam adeptus sum.) Addit haec ob duas causas. Prima ubi agnoscimus et confitemur nostra peccata, deque  nobis ipsis desperamus per legem commoti: restat ur supplices confugiamus ad Deum, atque  misericordiam imploremus propter Christum. Huc nimirum lex ipsa peccatores adigit, propterea namque  ad Galat.3.Lex uocatur pae dagogus noster ad Christum, ut fide iustificemur: et ad Roman.7. Vbia lege de pecel to conuictus fatetur et exclamat, se miserum esse et infelicem, cupitquea corpore morti obnoxio liberari, ad Dei se conuertit misericordiam, ac gratias agit Deo per lesum Christum, propter quem peccata remittuntur. ltaque  haec uerba aliquo modo officium legis adhuc ostendunt. Secunda, uia sternitur his uerbis ad disserendum mox de Euangelio. Habet Apostolus in more, cum ab uno loco ad alterius loci tractationem uult progre di, in fine prioris loci paucula apponere, quae respiciunt ad sequentia, ut partes cum par tibus aptius connectantur. Sit ergo haec doctrina nobis perpetuo in conspectu. Lex est ueluti speculum, in quo, quibus maculis anima nostra sit deformata, contemplamur.  \pend\pstart Docet proinde, quomodo agnoscamus peccata: Dehine inducit ad ueram poent tentiam, et peccatorum confessionem. Denique  remittit ad Christum medicum, qui maculas nostras abstergit, et per quem misericordiam consequimur. Secundo, Deum orabimus, ut hunc usum legis recte assequamur, et ne patiatur nos sic excaecari, ut peccata nostr. neque  agnoscere, neque  confiteri queamus. Indubitatum est argumentum irae Dei, et certa uia ad exitium, non habere pro peccatis, quae tamen peccata sunt, et peccata sua nunquam agnoscere. Quos huiusmodi esse aduertimus, eos sciamus a Deo in sensum reprobi traditos. Tertio, orabimus Deum, ut ubi per legem peccata agnouerimus, non patiatur nos induci in tentationem, per quam in desperationem abriperemur: sed potius mentes nostras spiritu suo illuminet ac regat, fidem in iis accendat, quo ad Christum toti conuerl remissionem peccatorum propter eundem consequamur. Si Paulum blasphemum, et per secutorem Deus dignatus est misericordia, speremus et nos per poenitentiam et fiden posse in gratiam recipi.  \pend\pstart Quod ignorans fecerim per incredulitatem.) Occupatio. Neque  enim peccatum suum his uerbis extenuat, uel quouis modo nititur excusare, quasi aliquid in ipso esset, ob quoc Deus ad Commiserationem sui sit permotus: Sed oceurrit atque  respondet tantis quorun dam obiectionibus. Quidam habentes aliquam ueritatis cognitionem, mirari poterant, qui fieret, quod, qui confessus est, se fuisse blasphemum, receptus tamen sit a Deo in gra tiam: siquidem blasphemia eiusmodi uidetur peccatum, quod remitti nullo modo queat, Quidam impii et uitiis immersi cogitassent: Si Paulo, qui tam grauiter peccauerat noxa est condonata, futurum, ut sibi quoque  peccata omnia remittantur, atque  ita impunitatis spem concipientes animum in malis obdurassent, ac posito omni metu perrexissent inique agere. Hoc autem nihil fuisset aliud, que  Pauli exemplo et Dei bonitate tur piter abuti. Vtrique  igitur hominum generi Apostolus cupiens satisfacere, ait, se fuisse qui dem blasphemum, sed per ignorantiam et ineredulitatem: quasi dicat, blasphemiae uia mise ricordiae minimem est praeclusa, sicutnec caeteris peccatis, si modo rei agant poenstentiam et fide misericordiam apprehendant. Est uero incredulitas caussa blasphemiae, quemad modum etiam est fons et origo peceatorum omnium. Et subiicit particulam de increduli tate, quae sedem habet in cordibus:ut intelligatur, et ignorantiam, et blasphemiam, et quaecumque  hinc peccata emanant, ex hominis peruerso corde prodire, ac proinde omniz esse ipsi imputanda. Ergo hoc modo Apostolus peccata sua exaggerat et amplificat non imminuit. Sunt porro qui peccant, atque  etiam apertem oppugnant Euangelium, sed per ignorantiam. Existimant ensm se recte facere, atque , ut Christus lohan.16. loquitur, praestare Deo cultum. Haec, inquit, facient uobis, quia non nouerunt patrem, neque  me.   Atqui tales, si ex animo resipiscunt, et ueritatem monstratam amplectuntur, Deus, ue niam petentibus, ignoscit. Quocirca et Christus pro his orauit, qui cruci ipsum affige  \pend
\vspace{0.5cm}\noindent
\vspace{0.2cm}\rule{1cm}{0.2pt}\\ 
\hspace{0.2cm}\textit{mg}
\footnotesize Agnoscentes sua peccata ad Dei misericordiam confugere devent. 
\normalsize| 
\hspace{0.2cm}\textit{mg}
\footnotesize Lex ppecu' lum, in quo maculas nostras contenplamur. 
\normalsize| 
\hspace{0.2cm}\textit{mg}
\footnotesize Lex deducit     
\normalsize| 
\hspace{0.2cm}\textit{mg}
\footnotesize            
\normalsize| 
\hspace{0.2cm}\textit{mg}
\footnotesize        
\normalsize| 
\hspace{0.2cm}\textit{mg}
\footnotesize         
\normalsize| 
\section*{IN I. EPIST. D. PAVLI. }
\marginpar{[ p.35 ]}\pstart   bant: Pater, inquit, remitte illis, nesciunt enim quid faciant. Luc.23. Eo igitur modo Apostolus blasphemus et persecutor fuerat, verum cum admoneretur peccati, atq ostenderetur ei, quid facere ipsum oporteret, submisit sese, ac venia ei data est. Rursus sunt, qui peccant, atque  oppugnant Christum ipsiusque  Euangelium, idquefaciunt permalitiam Quo tempore Christus concionabatur, multi veritatem probe norant, siquidem animaduertebant, quôd doctrina Christi consentiret per omnia cum lege et prophetis, ex quibus Christus probationes sumebat: quinetiam ex signis, quae Christus aedebat, plane conuincebantur, vt honorifice de Christo eiusquedoctrina sentire cogerentur: Sed quid fecerunt? Cum nollent displicere magnatibus, vel viderent syncera doctrina Christi euerti falsa sua dogmata, protrahi in lucem ac damnari suam hypocrisin, des nique  obscurari suam gloriam, scientes volentesqueomnibus viribus Christo resistebant. Tales perstringit Christus lohan. 7.dicens, non posse eos credere, quod gloram abinui cem acciperent. In tales inuehitur Apostolus 2. ad Corinth 2. et 4. Cum dicit esse, qui Verbum Dei cauponantur ac dolosem tractant, inuertunt: ltem huius epist. cap.4. Vbi dicit orituros falsos doctores cauterio notatam habentes conscientiam. Certum autem est, si qui huiusmodi sunt, in malo perstant, nec blasphemare desinunt, aeternae dam nationis reos fieri. Habet autem et haec aetas, in quam nos incidimus, permultos, qui veritatem impugnant per ignorantiam, quos tamen postea videmus subinde in viam redire et veritatem constanter profiteri, quibus nimirum hoc datum est a Deo, quosque  ad vitam Deus praeordinauit. ad Ephes..ad Philip.1.Act.13. Sed sunt insuper non pauci, qui veritatem impugnant per malitiam, vt gloriam venentur apud imperitum populum, gratiam magnatum promereantur, opes acquirant et honores: quos omnes, si non resipiscant, sup plicia aeterna manent. Verùm quo vnsquisque  tam de praesente loco, quam de sequentibus ad finem huius capitis et initium capitis sexti, vbi iterum blasphemiae fit mentio, queat rectius iudicare, nos de Blasphemia paulo explicatius disseremus, addituri etiam nonnulla de peccato per ignorantiam.  \pend
\phantomsection
\addcontentsline{toc}{subsection}{\textit{De Blasphemia. }}
\subsection*{\textit{De Blasphemia. }}\pstart Blasphemare et blasphemia latius omnino patent, quam plersque  opinantur. Est nam θλασφυμο͂ν in genere maledicere alicui superiori, execrari, vituperare, conuitiis incessere, et quouis modo dignitati atque  honori eius derogare. Βλασφνμία est huiusmodi actio siue opus. ltaque  duplex in Scripturis notatur blasphemia: altera, quae committitur in homines dignitate praeditos: altera vero in Deum. De blasphemia in homines ad Tit.2. Admoneto illos, vt principatibus ac potestatibus subditi sint, vt magistratibus pareant, vt ad omne opus bonum sint parati: μνδιενα θλασφκαᾶν, ne de quoquam maledicant. Et 1.ad Corint 4. Paulus de se aliisque  Apostolis: Malem audientes benedicimus, persecutionem patientes sustinemus, Βλασφυμομενοι παρακαλσμεν, id est, conuitiis affecti obsecramus. Blasphemia ergo non modo in Deum, verumetiam in homines peccatur: atque  omnes appellantur blasphemi, qui non contra pares, sed aduersus superiores, in quacunque  hi sint potestate siue Ecclesiastica, siue Politica, aliquid indignem et contumeliose dicit aut facit. Hinc quoque  fit, quod aliis insuper peccatis, quibus homines inprimis of fendimus, blasphemia censeri solet, ad Ephes.4. Omnis amarulentia et tumor et ira et   vociferatio et Βλασφημία, id est, maledicentia tollatur a vobis cum omni malitia. Ad Coloss.3. Nunc deponite et vos omnia, indignationem, malitiam, θασφμμίαy, maledicentiam, ,, turpiloquium ab ore vestro. Est porro alia blasphemia proprie in Deum: qua significatione plurimum ea vox vsurpatur, quod Augustinus quoque  indicat de moribus Manichaeorum liber 2.cap.11. Blasphemia, inquit, est, cum aliqua mala dicuntur de bonis. Itaque  bla   sphemia iam vulgo non accipitur, nisi mala verba de Deo dicere: de hominibus namque  du bitari potest: Deus vero sine controuersia bonus est. Definitio veron talis colligi potest. Blasphemia in Deum est, cum quis ex incredulitate dicit aut facit aliquid, quo aut negatur de Deo, quod erat illi tribuendum, aut e diuerso tribuitur, quod ei minime conuenit, qua exre Dei imminuitur aut obscuratur apud nos gloria. ln qua definitione caufas aliquot congessimus, ad demonstrandam magnitudinem eius peccati idoneas.  \pend
\vspace{0.5cm}\noindent
\vspace{0.2cm}\rule{1cm}{0.2pt}\\ 
\hspace{0.2cm}\textit{mg}
\footnotesize Peccauma liie. 
\normalsize| 
\hspace{0.2cm}\textit{mg}
\footnotesize Caussareprobationi? gentis Iud. o. Act.28 Esae Matth. 
\normalsize| 
\hspace{0.2cm}\textit{mg}
\footnotesize Blasphemare quid? 
\normalsize| 
\hspace{0.2cm}\textit{mg}
\footnotesize Duplex Bl. sphemia 2. In Deum 2 In homines. 
\normalsize| 
\hspace{0.2cm}\textit{mg}
\footnotesize Blasphemie aduersus su periores. 
\normalsize| 
\hspace{0.2cm}\textit{mg}
\footnotesize Blaspemiat Deum. 
\normalsize| 
\hspace{0.2cm}\textit{mg}
\footnotesize Definitio blasphemiex causis. 
\normalsize| 
\section*{AD TIMOTHEVM CAP. I. }
\marginpar{[ p.36 ]}\pstart Causa efficiens blasphemiae est incredulitas in corde hominis. Liquet id ex verbis Apostoli, qui confessus se fuisse blasphemum, persecutorem, et violentum, subnectit, ita se peccasseἐυ ἀπιςία, per incredulitatem. Ac res ipsa testatur, blasphemiam omnem ita oriri. Nouit aliquis homo, quid de Deo sentire oporteat, quaequeDeo proprie conueniant: Verùm quia praeduro est corde, nec credit rem ita se habere, audet omnireligione ac metu posito contrarium dicere aut facere. Talibus nimirum initiis et progress Iu Paulus blasphemus effectus est. Audiuerat enim verbum Dei, audiuerat plures pios veritatem Euangelicam de Christo confitentes, sed credere nondum poterat, imo vero arbitrabatur Apostolos caeterosquepios errare, se vero recte sentire: vnde animatus est ad resistendum illis, et persequendum, qua ratione peccatum ipsius factum est longe grauissimum. Et quod cognitio aliqua veritatis praeire soleat, antequam ad blasphemiam veniatur, ex eo manifestum est, quod fides quoque sit ex auditu, auditus autem per verbum Dei, ideoqueper contrarium incredulitas similiter ex auditu verbi est. Verbum namque primo proponitur, consequitur deinde, vt, qui audiunt, credant aut non credant. Huc pertinent verba Christi lohan.15. Haec omnia facient vobis propter nomen meum, quia non nouerunt eum qui misit me. Si non venissem et non locutus essem iis, peccatum non haberent, nune autem non habent, quod praetexant peccato suo. Qui me odit, et patrem meum odit: Si opera non fecissem inter eos, quae nemo alius facit, peccatum non haberent: nunc autem et viderunt, et oderunt, nonso . - lùm me, verum etiam patrem meum. Ex quibus verbis colligere licet, blasphemias ludaeorum inde natas, quod verbum quidem audiuissent, verùm nequaquam credidissent. Docet etiam Irenaeus liber 5. Semper debere cognitionem aliquam praecedere, antequam quis ad blasphemiam erumpat, additqueex lustini sententia Satanam anteaduentum Domini non ausum blasphemare Deum, quod ante id tempus damnationem suam certo non sciuisset. Vtrùm autem Satanas antea per se, vel per suos, Deum non blasphemauerit, non est praesentis loci plenius excutere. Haec autem sufficiunt ad declarandum, quod blasphemia omnis ex incredulitate tanquam fonte emanet. Chrysostomus liber 3.de Lazaro et diuite, Luc.16. Blasphemiam ait oriri ex impatientia: Hierony, mus enarrans cap.4. epist. ad Ephes.extra eam collocat. Sed si bene expendimus, Impatientia, lra, et similes perturbationes ex incredulitate tanquam ex fonte ebulliunt. Qui enim fidem habet in Deum, non irascitur, non est impatiens, sed statuit eum laturum opem cum ipsi oportunum videbitur quamobrem nullum ei verbum excidit, ad Dei spectans contumeliam. lob non fuit incredulus, sed fidelis: quamobrem in tantis aduersitatibus effecta incredulitatis, qualis sunt lra, Impatientia, Blasphemia, in illo nequaquam conspiciebantur. Apostolus ad Romanos 1o. ait, corde credi ad iustitiam, ore autem edi confessionem ad salutem. Ad eundem modum dicere possumus, in corde latere incredulitatem, ore autem emitti blasphemiam, vt, sicut opponuntur fides et incredulitas: ita opposita intelligamus confessionem fidei et blasphemiam. Porro causa formalis blasphemiae sunt dicta factaquequibus il la committitur. Dictis blasphemi fuerunt Syri, 1. Reg.20. qui regi suo Benhadac persuadebant, restaurandum bellum aduersus Israelitas, quandoquidem lsraelita, rum Deus non ita posset victoriam dare in planicie, quemadmodum dederat in mon tibus. 2.Reg.18. et 19. satellites Sanherib regis Assyriorum clamitabant, non posse Deum eripere Hierosolymitanos ex ipsorum manibus, sicutnec caeteri Dii suos populos seruauerant: Sed Esaias praedixit liberationem, ac tantam blasphemiam non fe re impunitam. Facto blasphemus fuit Achasia rex lsraelis, qui cum per Cancellos caenaculi delapsus grauiter aegrotaret, nuncios misit, qui consulerent Baalzebub Deum Ecrón, vt sciret, num spes superesset recuperandae valetudinis, quasi nullus esset Deus in lsrael, qui posset opitulari: ltaque ob eam blasphemiam obiurgatus est ab Helia, et mortuus. 2. Reg.1. Inuenias etiam quosdam dictis pariter et factis blasphemos. Paulus non tantùm Sermonibus obsistebat veritati, uerum etiam saeuam mouit persecutionem. Origenes in Euangelium lohan.39. disserens de prophetia Caiaphae, ait, blasphemare eu in SpirituS, operibus et sermonibus, qui peccat Spritu S. presente in eius  \pend
\vspace{0.5cm}\noindent
\vspace{0.2cm}\rule{1cm}{0.2pt}\\ 
\hspace{0.2cm}\textit{mg}
\footnotesize Causse (i ens bla mie (uu dulitas. περηαρ duiae et dmiσiae coniuncta act22. 
\normalsize| 
\hspace{0.2cm}\textit{mg}
\footnotesize Roman.19. Esas3. Incredeline carent uerbo Dei, cuï per blasphe miam refragantur. 
\normalsize| 
\hspace{0.2cm}\textit{mg}
\footnotesize Fide uera praediti non fundunt bla sphemias in Deum. In torde latet inereduItitas. Exoremanu blasphemia. Cafafoma us blasphemiaein dictis et factis po pita est. Esa.37. 1.Reg.1. 
\normalsize| 
\section*{IN I. EPIST. D. PAVLI. }
\marginpar{[ p.37 ]}\pstart anima. Augustinus in lohan. tractat. 27. Varii, inquit, inueniuntur nuc, qui lingua blalphemant Christum, permulti qui uita. Loquitur de his, qui Christianos se esse uerbis quidem profitentur, uitam uero uiuunt Christianis plane indignam, et magna praebent offendicula: quales hodie quoque  permultos uidemus, qui inter blasphemos censerino lunt, quamdiu ore blasphemias in Deum nullas effundunt. Caeterùm uehementer augetur blasphemia uerbis commissa, quando accedunt execrationes, detestationes, periuria, conuicia aperta in Deum, atq id genus, quae ex sonte incredulitatis, et ex perturbatis per ineredulitatem affectibus solent erumpere. Ac possunt haec inter blasphemiae affinia numerari. Petrus Matth.2o.negans Christum sibi esse notum, addidit periurium et execrationem. Horrenda fuit hlasphemia luliani, qui non solùm negauit Christum esse Deum, contra confessionem antea a se editam, uerum etiam conuiciis est insectatus, maxime cùm uulneratum se ad mortem cernens de sanguine proprio quantum manu exceperat, in coelum proiiceret, et contumeliose aduersus Christum clamaret, uicisti Galisaee: Tripart: histor:libus 6.cap.47.Et Theodoreti liber 3.cap.25. Grauissime uero peccant haeretici, qui uerbis blasphemiam perpetrant, quoniam sermocinando ac docendo disseminant uirus haereseos, qui alias factis et externis moribus sic norunt se conponere, tanquam si sanctimonia insigni olerent, et aedificare magis quam destruere haberent in animo. Intellige autem haereticos, qui uere tales sunt, hoc est, qui perperam de religione et negotiis fidei dislerunt Deinde, postquam sacris literis conuscti fuerint erroris, pertinaciter adhuc defendunt impia placita, atq alios quoque  in suam pertrahunt sententiam, sectasquefaciunt. 1.ad Timoth.1. Nomnulli circa fidem naufragium secerunt, quorum de numero est Hymonaeus et Alexander, quos tradidi Satanae, ut discant non Βλασφνμο͂ν, maledicere.2.ad Timoth.3.ubi describit haereticos resistentes ueritati, mente corruptos, reprobos circa fidem, uocat eosdem sui amantes, fastuosos, blasphemos siue maledicos. Quodquealii per eos abducantur ad blasphemandum, unde exaggeratur illorum   peccatum, liquet. 2. Petri2. Quidam, inquit, falsi doctores clam inducunt sectas perniciosas et sibi ipsis accersunt celerem interitum, et plerique  sequuntur illorum exitium,   per quos uia ueritatis Βλασφυμιθήσεται, id est, maledictis incessitur. Quod si rem eo deduxerint, ut secessionis a uera ecclesia sint autores, tunc etiam factis blasphemant. Vide quomodo nomnunquam a paruis initiis in multa grauissima peccata haeretici sese praecipitent. Porro blasphemia factis commissa grauior fit, quando peccatum eiusmodi est, ut postquam in lucem protrahitur, alii occasionem blasphemandi, id est, male loquendi de Deo, deque  syncera religione accipiamnt.2.Samuel.12. Dauidi, quia adulterio suo dederat hostibus occasionem male loquendi de lsraele denunciat Nathan, puerum ex adulterio genitum moriturum. Ad Roman.2.ad ludaeos impie uiuentes. Qui de lege gloriaris, per legis transgressionem Deum de honestas. Nam nomem Dei propter uos Βλασφνμῶται, male audit inter gentes. 1.ad Timoth.6. Quicunque  sub iugo sunt serui, suos dominos omni honore dignos ducant, ne nomen Dei et doctrina Βλασφνμο͂ταε, id est, male audiat. Est proinde ualde necessarium, ut de non dando offendiculo ac de praecauenda causa blasphemiae homines frequenter admoneantur publice in ecclesiis, cum perpauei sint, qui intelligant, et quid proprium sit blasphemia, et quam graue sit peccatum, qua taque  secum mala trahat. lam ex his, quae hucusquedicta sunt, perspicua redditur postres ma particula in definitione posita, nempe quod gloria Dei dictis factisqueblasphemis apud nos obscuratur. Notatur uero hic causa finalis siue effectus blasphemiae. Deo quidem non potest de gloria quicquam detrahi: necesse em̃ est, maiestatem ac gloriam cius, quanta fuit ab aeterno, inuiolabilem semper manere. Attamen in hominum animis, qui dictis factisqueblasphemantium permouentur, tacite nascitur quidam Dei contemptus: Ex quo fit, ut de Deo numquam tam honorifice deinceps sentiant, quam debeant. Ex contemptu autem Dei et falsis opinionibus alia mox erumptunt horrenda peccata, ueluti temeritas in peccando, oppugnatio agnitae ueritatis, mentis excaecatio, perseuerantia in malo, cordis obduratio, denique  et desperatio. Sic namque uelut per gradus quosdam ad mala extrema progressio fit, et plurima consequuntur ex sese mutuo, perinde atque effecta comitantur caussas. Quinetiam punit Deus peccata peccatis, et blasphemiam  \pend
\vspace{0.5cm}\noindent
\vspace{0.2cm}\rule{1cm}{0.2pt}\\ 
\hspace{0.2cm}\textit{mg}
\footnotesize Effectus blaspemifue causs finalis. Obscuratio gloriae De inter bomis nes.  Dei gloria per se senper malet inuiolabilis: respectu aum blaspemantium in contempti du citur apuc bemnes. 
\normalsize| 
\section*{AD TIMOTHEVM CAP. I. }
\marginpar{[ p.38 ]}\pstart blasphemia. Qua de re, quib? oculi sunt a Spiritu Dei illuminati, possunt iudicare, maxime cum illorum lapsus et actiones considerant, ꝗ audent hoc tempore agnitam impugnare veritatem. His igitur modis Dei gloria per blasphemantes vehementer apud nos obscuratur. Hucusque  de blasphemiae in Deum definitione, atque  eius diuersis causis, ac finibus et effectis. Restat explicemus, quomodo blasphemia quaedam sit remissibilis, quaedam ve ro irremissibilis. Christus Matth.12. ait, Omne peccatum et conuitium, blasphemia, remit tetur hominibus: at conuitium in Spiritum S.non remittet̃ hominibus. Et, Quicunque  dixe rit verbu aduersus filium hominis, remittet illi: qui autem dixerit aduersus Spiritum S. non remittetur ei, neq in praesenti seculo neque  in futuro. Eandem sententiam Lucas cap.12. enprimit: Marcus aut cap.3.ita. Amen dico vobis, omnia remittentur filiis hominum peccata et conuitia, quibuscunque  conuiciati fuerint: at qui conuicium dixerit in Spiri tum sanctum, non habet remissionem in aeternum, sed obnoxius est aeterno iudicio. Quodnam autem sit illud peccatum in Spiritum sanctum, nunquam remittendum, diffici le est desinire. Augustinus liber  de vero baptismate, peccare, ait, in Spiritum S.qui ve ritatem contemnit, qui in fratres, quibus Deus manifestauit veritatem, malignus est, qui non agnoseit Dei in ecclesiam beneficia, nec agit gratias. Libus de Sermone Domini, dicit eum peccare in Spiritum S. qui alteri inuidet collata Spiritus S.dona. Atqui haec nun quam remitti poenitentibus non satis potest demonstrari. ltaque  idem liber de Verbis Domi ni, ex Euangelio Matthae i Sermone n.animaduertens, nulsum esse peccatum, quod poenirentia non aboleatur, et de quo non sit vel in extremo vitae poenitentia speranda sratiocinatur, peccatum in Spiritum S.irremissibile esse impoenitentiam, et quidem finalem, hoc est, ad extremum vsque  vitae durantem, quae alias vocatur desperatio obstinata. Nec mul, tum dissidet ab hoc, quod idem epist.50 ad Bonifaciu vocat duriciem cordis ad fineu vsque . Ambiosius liber de Spiritu sancto: similiter Chrysostomus ad cap.12. Matth. dicunt, peccare in Spiritu S qui dero gant dignitati personae Spiritus sancti, atque  adimunt. quod ei erat tribuendum: quemadmodum Pharisaei negabant virtute Spiritus S.eiici a Christo daemonia, sed fecisse id Christum potius ope Beelzebul. Alii alia protulerunt. A quam plurimis autem recipitur sententia illa, de impoenitentia finali, vt licet ex his cognoscere, qui enarrant Petri Lombardi sontent. liber 2. distinctionem quadragesimam tertiam: neque  nos omnino reiicimus, minus tamen satisfacit, cum Euangelistaru verba non satis commode explanet. lubent igitur verba contextus primum respicere nos generatim in blasphemiam, id est, impugnationem alquam Veritatis, idquedocet historia ipsa, in qua recitatur, Pharisaeos, cum Christus eiecisset daemonium caecum et mutum, virtute Dei, contra Veritatem inique interpretatos; a Christo id factum non uirtute Dei siue spiritu Dei, sed ope Beelzebul principis daemoniorum. Talem blasphemiam Christus rationibus aliquot refellit. Deinde subiicit generalem comminationem de peccato in Spiritum S.nunquam remittendo. Secundo, discrimen aliquod statuere oportet, inter blasphemiam siue oppugnationc̃ Veritatis, quae committitur in filium, et eam quae commit. titur in Spiritum S. Notum autem est, duobus modis impugnari veritatem vel blasphemiam committi: altero quidem per ignorantiam: altero vero per voluntariam malitiam; Qui primo modo peccarint, cum poenitentia ducuntur, iterum in gratiam accipiunt tur. Sic Paulus de se satetur, qu blasphemus et persecutorfuerit, verum quoniam ignorans fecisset, misericordiam se consecutum. Petrus act.3. Vos sanctum et iustum negastis, et postulastis vt virum homicidam donaret vobis, autorem veron vstae interfecistis. Post pauca Et nunc fratres scio, qu per ignorantiam fecistis, sicut et principes vestri. Poeniteat igitu; vos erroris, et conuertimini, vt deleant peccata vestra. Paulus 1.ad Corinth.2.Si cogno; uissent, nequaquam Dominum gloriae crucifixissent. Et pro his, qui ita peccant, aequum est o. rationes fieri, quemadmodum Christus orauit pro his qui ipsum crucifigebant: Similites Stephanus pro lapidantibus. Generatim mandauit Christus, vt pro his oremus, ꝗnos persequuntur. Quando enim per ignorantiam aliqui oppugnant Veritatem, non est de conuersione et salute eorum spes abiicienda. Qui autem secundo modo, id est, uolur tarie impugnant veritatem agnitam et blasphemant, his peccatum non remittit. Ad Hebrao.S volentes peccauerimus post acceptam cognitionem veritatis, non vltra pro peccati?  \pend
\vspace{0.5cm}\noindent
\vspace{0.2cm}\rule{1cm}{0.2pt}\\ 
\hspace{0.2cm}\textit{mg}
\footnotesize Blaspemia quatenus remittatur, aut inexpiabilis permaneat. 
\normalsize| 
\hspace{0.2cm}\textit{mg}
\footnotesize Peccatum in Spiritum Sanctum irremissibile diffict Ie definita. DefintioD. Augustini. Impoenitentia finalis peccatt red ait irremissiAugustinum. SententiAmbresit, et Chrysostomi. Lucai. Marci.s. Matth.i.. Petrus Lom bardus. 1 Discrimen blasphemie. 1 filium et Spiritum Sanctum. 
\normalsize| 
\hspace{0.2cm}\textit{mg}
\footnotesize Duobus mo dis oppugna tur ueritas, Ignorantia et Malitia uoluntaria, Blasphemia ignorantiae remissibilis si sequatur poententia Act.7. Matth.5. Luc.6. Blasphamia ex malitia et uoluntate peruersa est urenissibi
\normalsize| 
\section*{IN I. EPIST. D. PAVLI. }
\marginpar{[ p.39 ]}\pstart   reliqua est hostia, sed formidabilis quaedam expectatio iudicii, et ignis vehementia, qui   deuoraturus est aduersarios. Et quia huic peccato negatur remissio, ideoκατ' ἐξογην vocatur peccatum ad mortem, nec est pro eo orandum.. lohan.5. Si quis viderit fratrem suum peccare peccatum non ad mortem, petet et dabit ei ueniam: peccantibus dico non ad mortem. Est peccatum ad mortem, non pro illo dico vt roges. Omnis iniquitas peccatum est, sed est peccatum, quod non est ad mortem. Ex his proinde licet colligere blasphemiam, quae remittitur, nempe quae fit ex ignorantia, recte statui blasphemiam in filium: Eam vero, quae non remittitur, quaeque  fit voluntate siue malitia voluntaria, blasphemiam esse in Spiritum Sanctum. ltaque  qui veritatem nondum tenent, neque  Spiritu Sancto sunt illuminati, et quoad possunt veritati resistunt, maledicunt, conastur eam interuertere, tales siue in personam patris, siue in personam filii, adde etiam in perso nam Spiritus Sancti conuitia coniecerint, dicuntur per ignorantiam peccare: quandoquidem iudicare rectius non possunt, facturi procul dubio aliter, si melius edocti essent. Quare spes est remitti eiusmodi blasphemiam posse. Addidi etiam de persona Spiritus Sancti, qu in hac parte, in blasphemia, inquam, per ignorantiam, non sit necesse discrimen inter personas praecipue notare: sed potius spectanda est ipsa blasphemiae origo, igno rantia videlicet. Probatur id autoritate Stephani Actor.7. qui pro illis generatim orat, quos paulô ante dixerat spiritui sancto semper restitisse: in quorum numero erat etiam Paulus, qui postea est conuersus. Resistebant illi spiritui sancto, et peccabant in spiritum sanctum, sed adhuc per ignorantiam, alias nequaquam pro illis orasset. Si autem aliqui veritatem tenuerint, atque  a spiritu sancto fuerint satis edocti, et interim pergunt, reluctante propria conscientia, traditam de Deo rebusquediuinis veritatem scientes ac volentes impugnare, hi proprie peccant in spiritum sanctum, neque  his spes veniae vlla superest. Ideo autem huiusmodi peccatum dicimus appellari proprie peccatum siue blasphemiam in spiritum sanctum, quia omnis cognitio veritatis a spiritu sancto est, et sine spiritus san cti occulta reuelatione atque  inspiratione veritatem non possumus confiteri. lohan.16. 1.ad Corinth.2. 1.lohan.2. Qui non fuerint ita edocti, non possunt eadem ratione culpari velut peccantes in spiritum sanctum. Est porro tale peccatum irremissibile ob duas caussas. 1. Quia reuera est per se longe grauissimum. Quis enim non iudicet indignissimum esse, quod, qui semel fuerunt illuminati ac participes facti spiritus sancti, ac gustauerunt bonum Dei verbum, tanta Dei beneficia aspernantur, Deum afficiunt con tumeliis, eiusqueveritatem modis omnibus impugnant. 2. Tales iusto Dei iudicio plurimum etiam in hac uita excaecantur, traduntur in sensum reprobum, indurantur in malo, vt conuerti ad Deum, poenitentiam serio agere, veniam humiliter et cum fide postulare nunquam velintaut possint. Quin potius pergunt indesinenter et ad extremum vsque  oppugnare veritatem. Pharao, lannes et lambres, Exod.7.2.ad Timoth.3.audiunt veritatem videntquesignis consirmari, vt conscientia ipsorum conuicta fateri cogerentur, digito Dei seu virtute spiritus geri omnia: Sed audebant adhuc obsistere, ideoque  excaecabantur, indurabantur, ac persequi non desinebant, donecmiserabiliter perirent. Esaiae 23. populus rebellis et conturbans Spiritum sanctum habet Deum aduersantem, qui et facit illos errare a via recta, et permittit corda illorum timore Dei esse vacua. Paulus ad Roman. I1o. pluribus demonstrat, ludaeis annunciatam fuisse veritatem Euanges lii, sed eos noluisse amplecti, imo porius contradixisse perpetuo acrebellasse: idqueprobat ex Esaiae cap.53. Quamobrem capite 1.statim subiieit ex Esaiaeó. et Psalmo 69. illos iusto Dei iudicio excaecatos et a Deo reiectos. His ita expositis, peccatum in Spiritum sanctum proprie dictum licet hoc pacto definire.  \pend\pstart Poccatum in Spiritum S. est, cum quis veritatem de Deo rebusque  diuinis, Spiritus S.beneficio cognitam sciens ac volens, reluctante propria conscientia, abnegat vel oppugnat, dumque  id excaecatus vel induratus constanter facit, remissionem tantae malitiae nunquam accipit. Est uero notandum, q, quinam ita grauiter blasphement ut peccatum eorum sit prorsus irremissibi le, non possumus, quandiu blasphemi in uiuis sunt, certo iudicare, nisi forte reuelatione quadam singulari id cuipiam manifestaret̃. Olim qdem Ecclesia uariis donis erat illustrata, et passim inueniebant nomnullis, q discemere ac probare spiritus et uidere, quos Deusa  \pend
\vspace{0.5cm}\noindent
\vspace{0.2cm}\rule{1cm}{0.2pt}\\ 
\hspace{0.2cm}\textit{mg}
\footnotesize Peccatumad mortem. 1Ioan.5. 
\normalsize| 
\hspace{0.2cm}\textit{mg}
\footnotesize Blasphemiin Filiumre missibilis. Blasphemi in Spiritum Sanctum ii renssibilie 
\normalsize| 
\hspace{0.2cm}\textit{mg}
\footnotesize Case Peccatum in Spiritums duabus de eausis irremissibile. 2. Caussa eur nou remittatuerpec catum in Spiritums. 
\normalsize| 
\hspace{0.2cm}\textit{mg}
\footnotesize Definitio peccati in Spiritum s. 
\normalsize| 
\hspace{0.2cm}\textit{mg}
\footnotesize Donx aifge tonis fpirituum in primuua eccle sia 
\normalsize| 
\section*{AD TIMOTHEVM CAP. I. }
\marginpar{[ p.40 ]}\pstart se omnino reiecisset, possent.1.ad Corinth.12. 1.loan.4. Quemadmodum ita fieri tunc expediebat, eratquenecessarium. Nunc vero iudicium certum ferre non possumus, nisi a poste riori, hoc est, dum animaduertimus manifeste effecta blasphemiae. Etenim qui palan excaecatus in sensum reprobum traditus, obduratus, atque  in blasphemia, veritatis oppugnatione, gratiae Christi contemptu perseuerat ac moritur, de hoc pronunciare audebi mus, quod peccati in Spiritum Sanctum irremissibilis reus fuerit: Ante vero que  vitae absoluent cuilum, non est plane desperandum de eius conuersione. Quis namque  scit, an Deus Illi velit rursus aperire oculos, vt peccatum agnoscat, veniam ex animo postulet, eandemque  consequatur. Quocirca verba Christi ad ludaeos recte intelligent generalem quandam comminationem continere, vt sit sententia: Si qui ex ludaeis tunc praesentibus ita senper interpretarentur, perfecta esse potentia Beelzebul, quae diuina virtute ipse fecisset, futurum, vt propter tale peccatum in Spiritum S. excaecarent, in sensum traderent reprobum, denique  damnarentur in aeternum: Nec dubium est, quin aliqui ex iis resipuerint: quidam vero eiusmodi effecta blasphemiae fuerint experti. Stephanus quoque  non vniuersaliter, sed indefinite dixit: Vos semper resticistis Spiritui Sancto: significans nimirum, non omnes, qui ibi aderant, sed tantum aliquos scientes volentesquein Spiritum sanctum peccasse, quorum peccatum foret irremissibile: quosdam vero alios illi coetui interfuisse, qui in Spiritum Sanctum peccarent quidem, attamen per ignorantiam, quibus peccatum, quandocunque  poemtentiam agerent, remitteretur: ac pro talibus, in quibus et Paulum adolescentem fuisse constat, orationem ad Deum fecit. Si, reuelante Spiritu Sancto cognouisset, oës sine exceptione ita restitisse Spiritui Sancto, vt peccatum eorum foret simpliciter irremissibile, nequaquam pro illis generatim orasset, nec ecclesia postea vidisset Paulum factum Apostolum. propterea 2.ad Timoth.2.ita praecipit Apostolus: Seruum Domini non opore tet pugnare, sed placidum esse erga omnes, aptum ad docendum, tolerantem malos: cum ce lenitate erudientem eos, qui contrarie sunt animo affecti experturum, ecquando det illis Deus poenitentiam, ad agnoscendam veritatem, et resipiscant e diaboli laqueo capti ab' eo, ad ipsius voluntatem. Ad Tit.3. Haereticum hominem antequam habeatur pro aeuerso atque  damnato, vult semel iterumqueadmoneri. Ad Rom.11. ludaeos excaecatos ait posse iterum inseri, si in incredulitate non permanserint. Quare nos quoque  de blasphe mantibus et Spiritui Sancto resistentibus, quamdiu in viuis sunt, non protinus despe rabimus, sed quemadmodum iubet Apostolus, frequenter admonebimus, docebimus sedulo, quin, etiam exemplo Stephani pro illis orabimus, fiduciam habentes, quod Deus dignabitur illorum misereri. Hinc quoque  fit, quod ecclesia pro inimicis et persecutoribus, pro ludaeis et haereticis orationem nonnunquam ad Deum facere solet Si aute certa reuelatione patefieret, aliquos ita resistere Spiritui Sancto, vt poenitentiae acremissionis nulla esset amplius spes reliqua, pro nis, quemadmodum loannes docet, non orabimus. Haecad iudicandum, quae nam blasphemia in Spiritum Sanctum sitirremissibilis, sufficiunt.  \pend\pstart Nune de peccato ignorantiae nonnulla sunt adiicienda. Quae autem maxime sunt ter peccascitu necessaria, nos propositionibus aliquot complectemur.  \pend\pstart 1. Peccatum aliud iudicio voluntario perficitur, aliud ignorantia. De priori genere satis notum est: omnes, qui legis diuinae habent cognitionem, sed tamen transgredi au dent, reos eius peccati statui: perlegem cognitionem peccati: ad Roman.3. De altere genere Leuit 4. et Numer.15.ubi Deus praescribit, qualia sacrificia offerri ad ostium t? Dernaculi et cornua altaris debeant, pro eo, qui per ignorantiam peccauerit, siue sacei dos, siue tota multitudo filiorum Israel siue priuatus quispiam ita deliquisset.  \pend\pstart 2. Quo autei uncernere polsimus, quae nam ignorantia, et quomodo, uel quando pro crimine statuatur: operae pretium est, diuersas ignorantiae formas siue species expli cari. Primum est ignorantia quaedam uniuersalis, et omnibus hominibus communis, quae est effectus peccati originalis, qua fit, ut nunquam tam plenem ea, quae Deiue salutis nostrae sunt, cognoscamus in hac uita, uti debebamus, ueluti post hac uita dabig cognoscere. Atque haec omnis erroris in hac uita fons est et origo, nec quamdiu uiui mus, prorsus a nobis tollitur. Secundo, est ignorantia particularis, atque  alia quider  \pend
\vspace{0.5cm}\noindent
\vspace{0.2cm}\rule{1cm}{0.2pt}\\ 
\hspace{0.2cm}\textit{mg}
\footnotesize udicium a bosteriori et postmor tem. 
\normalsize| 
\hspace{0.2cm}\textit{mg}
\footnotesize Ante mortem de hominis salutenon est desperandum, 
\normalsize| 
\hspace{0.2cm}\textit{mg}
\footnotesize De peccato per ignoran tiam. 
\normalsize| 
\hspace{0.2cm}\textit{mg}
\footnotesize 1. Propositio. Duplicitur, exignorantia et uo luntatae. Sagificia pro peccatis ignorantiae. 2. Propositio. 
\normalsize| 
\hspace{0.2cm}\textit{mg}
\footnotesize 1. Genusignorantiae uniuersalis. Effectum pec cati originalis. 
\normalsize| 
\hspace{0.2cm}\textit{mg}
\footnotesize 2. Ignoranuaparticils 
\normalsize| 
\section*{IN I. EPIST. D. PAVLI. }
\marginpar{[ p.41 ]}\pstart eorum, quae sciri a nobis nullo modo est necesse:alia eorum, quae sciri est necesse. Ex his duabus formis prior illarum usque  adeo omni culpa uacat, vt, si scire, quae nos scirenon oportet, studeamus, vitio verti nobis possit, atque  crimem incurramus, si de rebus sublimibus temeritatis, si de rebus humilioribus, ineptae curiositatis. Huc illa pertinent Ecelesiastis cap.8. Aspexi omne opus Dei: sed nemo est, qui potest inuenire rationem huius " operis, quod fit sub Sole, propter quod homo laborat, vt quaerat, sed non inuenit, quinetiam si dixerit sapiens, se dare operam ut sciat, non poterit inuenire. Et in prouerbiis, in libris, Hiob et lesu Syrachi, multa inuenias, contra quarumcunque  rerum non necessariarum scrutatores dicta.  \pend\pstart Caeterùm ignorantia rerum scitu necessariarum nunquam sine crimine est, quamyis alia quidem mitius, alia vero seuerius punienda existimetur. Qua de re iudicium utfer ripossit, alias partitiones oportet considerari: Est igitur quaedam ignorantia, cum glexit. Hanc vocant ignorantiam crassam ac supinam, quae spectatur in pueris et multis hoquis quae scire merito debebat, aut omnino nesciuit sibi addiscenda, aut addiscere neminibus recta educatione destitutis. Alia est ignorantia, cum quis quae sciebat sibi addiscenda, tamen etiamsi voluit addiscere, non potuit: a quibusdam appellatur ignorantia inuincibilis siue probabilis, qualem apparet fuisse in Eunucho Actor. 8. et Centurione, Actor.1o. multisque  aliis, praesertim gentilibus, qui Deum vtcunque  timuerunt. Alia est, rentur addiscere: Hanc nominant affectatam et simulatam, estqueomnium perniciosissima. cum quis quae scire debuerat et poterat, noluit tamen, quamuis facultas et occasio daNon pauci hoc pacto tres ignorantiae formas distinguunt, et sicuti libet explorare cau sas diuersas, ob quas ignorantia prouenit, competit distinctionem hanc non esse reiiciendam. lam quomodo, et quare oës hae formae cedant ad peccatum, non difficile est explicare. Dam nari ignorantiam supinam multi Scripturarum loci comprobant. Psal.32. Ego prudentem te reddam, et docebo te in via, per quam ambulabis, consillumquetibi dabo oculo meo. Non eritis sicut equus et mulus, qui nihil intelligunt. Psalm.49. Homo in honore non perdurabit, cum sit assimilatus pecudibus, quae intereunt. Haec ipsorum via stultitia   eorum est. Et Paulus Actor.7. ait, Deum tempora ignorantiae in gentilibus dissimu» lasse, nunc autem annunciari hominibus, vt omnes vbique  resipiscant. Est autem hominibus vniuersis insita quaedam de Deo atque  honestis actionibus cognitio: quod planius docet Apostolus ad Roman.1. et 2. Quare quicunque  lumen sibi concessum non agnoscunt nullamquedant significationem cognitionis, qua sunt praediti: sed in ignorantia ac tenebris persistunt adeoquegaudent persistere: tales, eodem Apostolo teste, sunt inexcusabiles, et sicut sine lege peccant, ita sinelege et peribunt. Ad Rom.1. et 2. In quam sen tentiam D. Augustinus de gratia et libero arbitrio cap.3.ait, ignorantia eorum, qui tan quam simpliciter nesciunt, nullum sic excusari, vt aeterno igne non ardeat, sed fortasse vt minus ardeat. Non enim sine caussa dictum, Effunde iram tuam in gentes, quae te non nouerunt: Psal.76. Et illud quod ait Apostolus 2.ad Thessal.1. Cum venerit Christus cum incendio flammae, qui infligit vltionem iis, qui non nouerint Deum, et qui non obediunt Euangelio. Estquetalis ignorantia in adultis quidem poena peccatorum, tum originalis tum actualium: in paruis vero est poena peccati originalis. Atque  idem Augustinus de Ciuitate Dei lib 21 cap.14. quomodo ignorantia siue insipientia sit in paruulis poena eleganter declarat. Similiter epist. 14o. ad Sixtum presbyterum. Itaque  nulla est tam barbara natio, nulli homines tam imperiti rerum, qui ignorantiam possint praetexere, quod Deum non rite colant. Lex et Euangelium sunt doctrinae vniuersales ab initio inde propositae generi humano: Deinde ab Apostolis per vniuersum orbem diuulgatae: Psal. i9.ad Rom.1o. Vtfieri non possit, quin vbique  aliquid de utroque  genere doctrinae sit auditum perceptumque .  \pend\pstart Omnium vero minimem excusari possunt ignorantiae praetextu, qui annumerantur Christianis. Frustra causantur quidam sufficere ad salutem fidem implicitam, hoc est, vt profiteantur, se ita credere, quemadmodum credit ecclesia, nec ipsi interim persese possunt fidei suae vllam reddere rationem. Frustra praetexunt simplicitatem vel rusticitatem. Nemo tam rudis est, qui non possit in ecclesla Deisaltem prae cipua capita legee  \pend
\vspace{0.5cm}\noindent
\vspace{0.2cm}\rule{1cm}{0.2pt}\\ 
\hspace{0.2cm}\textit{mg}
\footnotesize Res scitu ne estrreo non necesse rie. 
\normalsize| 
\hspace{0.2cm}\textit{mg}
\footnotesize Ignorantia trassa et su pina. 
\normalsize| 
\hspace{0.2cm}\textit{mg}
\footnotesize ignorantie probabilis 
\normalsize| 
\hspace{0.2cm}\textit{mg}
\footnotesize Ipercr affectata
\normalsize| 
\hspace{0.2cm}\textit{mg}
\footnotesize Bpere supina dam nata
\normalsize| 
\hspace{0.2cm}\textit{mg}
\footnotesize lrpαι adultispoenpeccatorum actualium: Il Q0 peue en endis. 
\normalsize| 
\hspace{0.2cm}\textit{mg}
\footnotesize Fides implicita uana et inutilis ad salutem 
\normalsize| 
\section*{AD TIMOTHEVM CAP. I. }
\marginpar{[ p.42 ]}\pstart et Euangelii, siue fidei et charitatis addiscere. Sicut artificium vel alias res agendas:ita religionem et res credendas facile assequi quilibet Christianus, si velit, potest. Ac generale illud praeceptum 1.Petri 4. sint parati ad respondendum cuilibet petenti, et loquamini de ea, quae in vobis est spe, cum mansuetudine et reuerentia, conscientiam habentes bonam. Quod ad alteram speciem ignorantiae attinet, qualem diximus in Cornelio et Bunucho fuisse, certum est, et haec ipsa, si non remouetur, damnationis causam coram Deo adfert. Sicut namque  omnes intelligimus insigne esse opus misericordiae Dei, si quado dignatur Deus talem ignorantiam tollere, missis videlicet doctoribus, qui imperitos doceant, vel aliis accommodis remediis adhibitis: Ita per contrarium iudicium atque  iustitia Dei in eo elucet, quando in ignorantia, eorum, quae ad salutem sunt necessaria, haerêre aliquos permittit. Neque  enim Deus obstrictus est, vt quomodo vel quando homines putant aut cupiunt, ex tenebris eos eripiat lucemqueaccendat in ipsorum cordibus. Quamobrem inter poenas peccatorum antecedentium, iudice Augus stino epist.1o5. haecignorantia non abs re numeratur. Huc pertinet, quod Lycaones Act.14. quamuis gustum aliquem cepissent Euangelii, et Apostolos admirarentur, ac cuperent honore persequi, non tamen dignatus est Deus pleniori cognitione eos illustrare: sed Apostoli ab illis discesserunt, relictiquesunt in ignoratione quanta et quam   diu Deo placuit. Nos annunciamus, inquiunt Apostoli, vt ab istis vanis conuertamini ad   Deum viuum, qui secit coelum et terram et mare, et omnia quae in eis sunt, qui in praeteritis aetatibus sinebat omnes gentes ingredi viis suis, quanquam non expertem   testimonii seipsum esse sinebat, dum beneficia confert, de coelo nobis dans pluuias ac   tempora fructifera. Sunt item Apostoli interdum prohibiti, ne quibusdam veritatem patefacerent.  \pend\pstart Quod insuper puniatur hae cignorantia, et sit ipsa partim poena peccatorum, partim etiam per sese culpa quaedam, manifestum est, tum quia omnis ignorantia corum, quae ad salutem sunt necessaria, in vitio est, et peccandi adfert causam: tum quia regula generalis a Christo descripta est de ignorantia punienda: etiamsi quandam mitius, quandam vero seuerius puniri verosimile est. Sic autem Christus Luc.12.ait, Seruum qui cognouit voluntatem Domini, et non se praeparauit, nec fecit secundum voluntatem eius, plagis vapulabit multis. Qui autem non cognouit, et fecit digna plagis, vapulabit paucis  \pend\pstart Tertiaspecies ignorantia affectata maxime vitiosa est, et grauissime a Deo punitur Vnde sepe contra ea in Scripturis verba fiunt. Omnibus temporibus inuenti sunt permul. ti homines improbi, qui veritatem ac voluntatem Dei addiscere cum potuissent, noluerunt Sed haec quoque  apte distinguitur. Quaedam enim prouenit ex errore seu falsa persuasione, cum videlicet quis arbitratur doctrinam, quam semel hausit, esse omnino veram ac certam: quam vero audit ab aliis pro veritate proponi et commendari, falsam esse existimat. Etsi vero haec ignorantia subinde pariat alia grauissima peccata in hominibus  nempe vt resistere audeant Spiritui sancto blasphement, pios persequantur, et violen tiam in hos exerceant: quia tamen non nisi ad tempus durat, atque  doctrinam sanam illi tandem recipiunt, aguntquepoenitentiam, Deus ipsorum miseretur, et concedit veniam. Comprobantur haec pluribus Scripturae locis. Christus lohan.16.Alienos a Synagogis facient vos: Et veniet tempus, vt quisquis interficiat vos, uideatur cultum praestare Deo. Et haec facient uobis, quia non nouerunt patrem neque  me. Petrus Actor.3. Vos sanctum et iustum negastis, et postulastis, ut uirum homicidam donaret vobis, autorem uero uitae interfecistis. Et paulo post: Et nunc fratres, scio, quod per ignorantiam fecistis, sicut et principes populi uestri: Poeniteat igitur uos erroris, et conuertimini, ut deleantur peccata uestra.1.ad Timoth.1.Qui prius eram blasphemus, et persecutor, et uiolentus: sed misericordiam adeptus sum, quod ignorans fecerim per incredulitatem.  \pend\pstart Quaedam porro ignorantia affectata habet adiunctam malitiam et obdurationem, quando nimirum, cum saepius audierint doctrinam sanam elare atque  euidenter explicari, sic insuper eam percipiunt ut negare non possint solidam ac salutarem existere: pergunt nihilominus scientes ac uolentes ueritatem agnitam modis omnibus oppugnare. Talem proinde ignorantiam, si manere contigerit, nulla unqua subsequente poe\pend
\vspace{0.5cm}\noindent
\vspace{0.2cm}\rule{1cm}{0.2pt}\\ 
\hspace{0.2cm}\textit{mg}
\footnotesize ignorantia et causa pecatorum, et Cce ce sese. 
\normalsize| 
\hspace{0.2cm}\textit{mg}
\footnotesize ignorantia affectata gra - uiter a Deo punitur. 
\normalsize| 
\section*{IN I. EPIST. D. PAVLI. }
\marginpar{[ p.43 ]}\pstart nitentia, decet pro peccato grauissimo reputare. Erumpit namque in blasphemiam contra Spiritum sanctum, quam Christus Matthaei i2. irremissibilem pronunciat. Sic Pharao, cum nollet audire Mosen, permansit in diuinae voluntatis ignoratione. excaecatus est, et obduratus, poenitentiam non egit, ideoqueexitio aeterno destinatus est. Exod. 14. Et de ludaeis agnitam veritatem contemnentibus et a Deo reiectis plura occurrunt passim testimonia. Ierem.4. Stultus est populus meus, et me   non cognouerunt filii insipientes, et non intelligentes sunt: sapientes quidem sunt ad male faciendum sed bene facere nescierunt. Vbi mox plura subiiciuntur de corundem perpetua internecione. Psalm. 81. Populus meus non audiuit vocem meam,   neque lsrael acquieuit mihi. Dimisi igitur eum in grauitate cordis sui, et ambularunt   in consiliis suis. Ad Roman.1o. Zelum Dei habent, sed non ex cognitione. Nam   ignorantes Dei iustitiam et propriam iustitiam quaerentes constituere, iustitiae Dei   non fuerunt subditi. Tandem eo rem deducit Apostolus, vt demonstret ex prophetarum oraculis, ludaeos propter obstinatam ignorantiam et peruicaciam reprobatos a Deo et excisos.  \pend\pstart Porro admonendi sunt pii, ut cautem circumspecteque  agant in omnibus, ac diligenter prouideant, ne propter ignorantiam mandatorum Dei, uel etiam constitutio. num humanarum, praesertim quae uerbo Dei nituntur, atque ob aequitatem merito commendantur, in uerbis aut factis, adde etiam in causis aut circumstantiis siue praecedentibus siue comitantibus, siue adiunctis in negociis delinquant, aut etiam offendiculum praebeant reliquis. Siquidem late ignorantiae malum patet, et cùm minime credimus nos, omni aetate, quouis tempore, per ignorantiam tam nostram, quam alienam, in peccata prolabimur. Frequenter succurrat nobis apostolica admonitio ad Ephes. 5. Ne sitis imprudentes, sed intelligentes quae sit uoluntas Dei. Quin etiam crebro debemus Deum precari, ut ignorantiam nostram, et quae per hanc admiserimus, nobis condonet. Non enim de nihilo est, quod Deus in lege praecipit, ignorantiae peceata debere sacrificiis expiari.  \pend\pstart Quia uero perdifficile esset, regulis complecti modos omnes, quibus per ignorantiam peccatur, aut ubi excusatio a peccato inueniatur: ideo consultum est, quotiescunque  scrupus oboritur, uel proponitur quaestio, ad exempla, quae in Scripturis passim occurrunt, confugere atque ex illis certum petere iudicium. Psalm.25. postulat propheta remitti sibi peccata, quae in adolescentia ( nam haec ignorantiae mor» bo plurimum obnoxia est) perpetrarat: Peccatorum adolescentiae meae atque scelerum meorum ne recorderis. Genes.19. Loth ebrius tanta obruitur ignorantia, ut cum propriis filiabus committat incestum. Verùm ignorantiam hanc cum vltrô sua culpa sibi ipse accersiuerit, (debuerat namque sibi cauere ab ebrietate, quae peccandi frequentissime subministrat materiam) excusari nequaquam potest. Genes.12. Pharao grauiter punitur ob Saram inductam in domum suam, quam tamen uxorem Abrahami fuisse ignorarat: Debuerat nimirum diligentius prius omnia explorasse. Genes.2o.rex Abimelech a Deo praemonetur in somnis, ne cum eadem Sara per igno rantiam designet adulterium. 1.ad Corinth. 8. acio. Si pius his, quae immolata fuerunt idolis, uescitur, ignorans tamen fuisse immolata, absoluitur a crimine. Si autem   postea resciscit oblata fuisse, atque  tunc uel propria uel alterius conscientia infirma laeditur, culpa non uacat. Actor.23. Paulus durius respondet pontifici, uocat parietem dealbatum: admonitus autem non debere maledicere pontifici, excusat se, dicens, se ignorasse, quod Pontifex esset. Sed quis enumeret exempla eius modi omnia? Interest ergo uniuscuiusque pii diligenter causas et circumstantias omnes in eiusmodi exemplis perscrutari. Quod si deinde libeat nostra facta cum factis, quae expressa sunt sacris literis, conferre ac singula accurate expendere, promptum erit, ubi peccetur, uel ubi excusatio locum habeat, definire. Hactenus de prima parte doctrinae Christianae, nempe delege. Plura quidem potuissent de lege, de eius partibus, de abrogatione legis, in medium proferri, sed arbitramur, nos, quae ad explicandam mentem Apostoli sufficiunt. dixisse.  \pend
\vspace{0.5cm}\noindent
\vspace{0.2cm}\rule{1cm}{0.2pt}\\ 
\hspace{0.2cm}\textit{mg}
\footnotesize Exod.5.7 
\normalsize| 
\hspace{0.2cm}\textit{mg}
\footnotesize Confulende sunt exen pla Sorpe2 
\normalsize| 
\section*{AD TIMOTHEVM CAP. I. }
\marginpar{[ p.44 ]}
\phantomsection
\addcontentsline{toc}{subsection}{\textit{Exuberauit autem supra modum gratia: }}
\subsection*{\textit{Exuberauit autem supra modum gratia: }}\pstart Sequitur altera pars doctrinae Christianae, quae est de Euangelio: ac paribreuita te, qua pars praecedens, excutitur. Cur uero doctrinam Euangelii explanari opor teat post doctrinam legis, supra exposuimus, minimequeopus est repeti. Nullo au tem exordio, nulla transitionis forma apposita, statim aggreditur Euangelii negocium. Ac separari haec uerba debere ab antecedentibus, nouamquepartem hic institui, ex eo patet, quod proxime tractauit exemplum de sua persona priuatim: quae uero nunc sequuntur, doctrinam continent generalem, atque accipi ut uniuersaliter dicta de nouo loco communi conuenit. Item sicut praecedentia omnia accommodasa fuerunt ad demonstrandum usum legis: ita hic sine inuolucris proponuntur quae propria sunt Euangelii: Ac paulon post Apostolus iterum quaedam de sua persona annectet, ad usum Euangelii aperiendum idonea. Breuiter, nihil ex praecedentibus ad ista respicit, nisi quod narrando se adeptum misericordiam, uiam utcunque  strauit ad hunc de Euangelio locum. Diserte uero hic agere Apostolum de natura Euangelii, perspicuum fit non una ratione.  \pend\pstart Primo exuberans gratia Domini nostri, de qua hic instituitur sermo, Legi in Scripturis opponitur. Iohann.1.de plenitudine eius nos omnes accepimus, et gratiam pro gratia: quia lex per Mosen data est, gratia et ueritas per lesum Christum exorta esti Ad Roman.6. Non estis sub lege, sed sub gratia. Ad Galat.5. Quicunque per legem iustificamini, a gratia excidistis. Ergo Apostolo placet summam doctrinae Christianae hac methodon tradere, ut postquam absoluit sermonem de lege, pergat nunc dicere de Euangelio legi opposito.  \pend\pstart Semcundon, gratia est ueluti subiectum Euangelii, ita ut ubicunque de gratia Der aliqua occurrunt, ibi commendari negocium Euangelii omnes interpretentur. Vnde Apostolus ad presbyteros Epheslos uerba faciens Actor.20. Euangelium appel lat Euangelium gratiae: ut consummem, inquit, cursum meum cum gaudio, et mi nisterium, quod accepi a Domino lesu, ad testificandum Euangelium gratiae Dei. Ad Ephes.2. Audistis dispensationem gratiae Dei, quae data est mihiin uos, etc. Atq hinc fit, quod Apostoli passim in suis scriptis primo et ultimo loco gratiam incul, cant. In prophetarum autem sermonibus, imo in toto ueteri testamento talis loquendi ratio non obseruatur.  \pend\pstart Iertio, in his uerbis Apostoli cumprimis elegans et foecunda continetur Euange. lii desinitio seu descriptio, quod mox plenius ostendemus. Itaque nouam hic institu partem de doctrina Euangelii, ex his liquido constat. Proinde ad uerba Apesto. li ueniamus.  \pend
\phantomsection
\addcontentsline{toc}{subsection}{\textit{Exuberauit autem. }}
\subsection*{\textit{Exuberauit autem. }}\pstart Continetur his uerbis luculenta Euangelii definitio, quam hoc pacto crassius licet exprimere. Euangelium est doctrina laetum proponens nuncium de gratia et remissione peccatorum gratis et abundanter in nomine lesu Christi a Deo oblata omnibus hominibus, quorum interest deinceps in uera fide ac dilectione pie sancteque  ad Dei et lesu Christi gloriam uiuere. Estquedefinitio illa Apostoli uehementer commen danda. Primo, quia breuissima est: fiquidem in definitionibus breuitas probatur. Secundo, quia maxime propria est. Nihil enim aliud continet, quam quod spectari in primis in Euangelii natura debet. Nonnulli definiunt Euangelium esse doctrinam de poenitentia et remissione peccatorum. Neque id quidem incommode, ue rùm non satis proprie de Euangelio loquuntur, sed magis de doctrina uniuersa, quam Christus proponi in ecclesia mandauit, quae legem pariter cum Euangelio complectitur. Etenim praedicatione legis excitantur homines ad poenitentiam: quandoquidem lex adfert peccatorum cognitionem: Comminationes autem legis incutiunt dolorem ob peccata. Sed quando haec duo, cognitio, inquam, peccatorum et dolor ob peccata locum habent eciam in impiis, et non semper ubi sia duo sunt, peccatore  \pend
\vspace{0.5cm}\noindent
\vspace{0.2cm}\rule{1cm}{0.2pt}\\ 
\hspace{0.2cm}\textit{mg}
\footnotesize t. Pars dotrinae Christianae est de Euangelio. 
\normalsize| 
\hspace{0.2cm}\textit{mg}
\footnotesize l.Gratiaex. uberans opponitur le. *. 
\normalsize| 
\hspace{0.2cm}\textit{mg}
\footnotesize 2. GratiaE  ubiectum. Actor.20. Euangelium gratiae. Dispensatic praue De 
\normalsize| 
\hspace{0.2cm}\textit{mg}
\footnotesize 1. Euangelii definitio hic traditur. 
\normalsize| 
\hspace{0.2cm}\textit{mg}
\footnotesize Euangelium quias 
\normalsize| 
\hspace{0.2cm}\textit{mg}
\footnotesize Definitio Etagelii propria. Euangelii definitio impropria. 
\normalsize| 
\hspace{0.2cm}\textit{mg}
\footnotesize Cognitio peccatori et dolor ob peccata non sunt propria Euangelii. 
\normalsize| 
\section*{IN I. EPIST. D. PAVLI. }
\marginpar{[ p.45 ]}\pstart remissio sequitur: ideo non sunt propria Euangelii, quod simpliciter agit de gratia et peccatorum condonatione. Tertio, commendatur haec definitio, quia con gesta est ex Euangelii effectis, sicut supra quoque disseruit de vsu et quibusdam effectis legis. Estquehoc genus definitionum in Theologia vsitatissimum, ac frequenter cogimur de Deo, de esus voluntate iudicium colligere a posteriori et effectis acsignis. Aptissime autem huiusmodi definitiones verum cuiuscunque  doctrinae Spiritualis usum expla nant et augent in conscientiis experientiam, vt maiori iudicio assuescant in rebus diuinis et spiritualibus omnia considerare: Sine qua experientia et sedula obseruatione nemo serio in theologia rebusque ́, ad salutem pertinentibus proficit. Sic porro distributa sunt membra in hac definitione Euangelii, vt in priori parte statuantur consideranda effecta in Deo, siue ex parte Dei, ut modus et causa efficiens iustificationis erque  salutis humanae, siquidem hanc constat solius Dei gratia perfici.  \pend\pstart Exuberauit, inquit supra modum gratia Domini nostri. Iln posteriori autem parte subiiciuntur effecta consideranda in hominibus, qui gratiae Dei participes facti fuerint: quando ait, cum fide et dilectione, quae est per Christum Iesum. Etenim in iustificatis fides et dilectio coniunctim elucere perpetuo debent, quae ipsa sunt Dei dona ac gratiae Dei praeclara effecta. Ac notandus est hicordo, quo primo generatim vt origo et fons omnis boni commendatur gratia Dei. Secundo, fides, quae etiam requiritur ante iustificationem: quanquam et ipsa tunc Dei donum est, Ephes.2. et sequi iul stificationem debet. l’ertio, Dilectio, quae est tanquam regula ac norma bonorum operum, quae a fide producuntur. Ad Galath.5. Haec vero simul omnia inuicem coniunguntur et cohaerent, vt mutuo sese ponere quodammodo videantur. Si qui uere adepti sunt gratiam Dei seu remissionem peccatorum, ac sunt iustificati, hi quoque veros Euangelii fructus ac dignos gratia Dei, fidem, inquam, et dilectionem producent: Et in quibus fides est et dilectio, in iis gratiam Dei et prius esseoperatam, et adhuc operari necesse est. Gratia hoc loco significat misericordiam, qua Deus gratis remittit peccata propter meritum lesu Christi, qui morte sua pro pec çatis nostris satisfecit. Sic autem in negocio iustificationis ac salutis nostrae considerari gratia in Deo et ex parte Dei debet: vt ex parte nostra neque opera, quantumlibet legi diuinae consentanea, neque merita locum habeant. Huc pertinent hi loci: Ad Roman3. lustificamur autem gratis per illius gratiam, et quae sequuntur: Quod si per gratiam, non iam ex operibus. Ad Galat.5. Quicunque  per legem iustificamini, a gratia excidistis. Ad Tit.3. Postquam bonitas et erga homines amor apparuit seruatoris nostri Dei: Non ex operibus iustis quae fecerimus nos, sed ex sua misericordia seruaustnos, etc. His locis clarem et per plures circumstantias in negocio salutis nostrae considerandas, explanatur, quid praesenti loco Apostolus nomine gratiae velit intelligi.  \pend\pstart Exuberauit, ὑπερεπλεόνασι.) Elegans compositio ad significandum insignem, admirabilem, et superabundantem gratiam. Atque hoc loco mirifice amplificatur Dei gratia, quae Amplificatio alibi tribus potissimum ex causis colligitur atque illustrat̃. 1. Quod per sese amplissima sit in ipso Deo. Ad Ephes.1. dicitur, quod in Deo sint diuitiae gratiae. Et Deus pater 2. ad Corinth.1. pater misericordiarum appellatur. De filio autem lohan.1. Conspeximus gloriam eius gloriam uelut unigeniti a patre, plenus gratiae et ueritatis.  \pend\pstart 2. Quôd Deus largissime impendat suam gratiam hominibus eandemquesemper au geat. ad Ephes.I. Per quem habemus remissionem peccatorum iuxta diuitias gradiae suae, de qua ubertim impartiuit nobis in omni Sapientia ac prudentia. Iohan.1. De plenitudine eius nos omnes accepimus, et gratiam pro gratia. 1. Petri 1. Perfecte sperate in eam, quae ad uos defertur, gratiam, dum patefiat lesus Christus.  \pend\pstart 3. Amplificatur gratia Dei, quando confertur hominibus, qui ob peccata aeterne damnationis omnes rei sunt. ad Roman.5. Si per unius delictum mors regnauit per unum, multo magis ii qui exuberantiam gratiae et doni iustitiae accipiunt, per uitam  \pend
\vspace{0.5cm}\noindent
\vspace{0.2cm}\rule{1cm}{0.2pt}\\ 
\hspace{0.2cm}\textit{mg}
\footnotesize Definitio a effectis. 
\normalsize| 
\hspace{0.2cm}\textit{mg}
\footnotesize Fides et d lectio donDei in san ctis. 
\normalsize| 
\hspace{0.2cm}\textit{mg}
\footnotesize Dilectio et regula bont rum operum quae ex fid manant. 
\normalsize| 
\hspace{0.2cm}\textit{mg}
\footnotesize Gratia Des fidem et di lectionem ho minis praec dit. 
\normalsize| 
\hspace{0.2cm}\textit{mg}
\footnotesize Gratia et m fericordisent idedn aSντα. 
\normalsize| 
\hspace{0.2cm}\textit{mg}
\footnotesize mn uegotio iustificatio nis gratid 
\normalsize| 
\hspace{0.2cm}\textit{mg}
\footnotesize Dei non ad mittit mert ta hominum 
\normalsize| 
\hspace{0.2cm}\textit{mg}
\footnotesize Amplificatio gratiae Dei 
\normalsize| 
\hspace{0.2cm}\textit{mg}
\footnotesize 11. Gratia ii Deompii simal 
\normalsize| 
\hspace{0.2cm}\textit{mg}
\footnotesize 2. Deus lar giter suam il credentes e fundit gratiam. 
\normalsize| 
\hspace{0.2cm}\textit{mg}
\footnotesize 2Grdliabe iis, qui dan nationis rel sunt, confe ur. 
\normalsize| 
\section*{AD TIMOTHEVM CAP. I. }
\marginpar{[ p.46 ]}\pstart Lauberaust gratia. Tluiumour ignur Ampiicanones gratiae Der notantur abApostolo, dicente uno uerbo, quod gratia Domini nostri super modum exuberarit.  \pend\pstart Proinde obseruanda nobis est haec methodus, qua, cum agitur de negocio Euangelii et Salutis nostrae, initium fit ab exuberante Domini nostri gratia, imo et finis in eadem fieri debet. Hanc methodum par est nos quoque in docendo sequamur. Atque omni studio retinendum praedicandumqueest in ecclesiis dogma, quod gratia Deisit non modo iustificationis nostrae, uerum etiam electionis, uocationis, et glorificationis fundamentum ac causa efficiens. Vult namque  Deus hac ratione suam celebrari bonitatem, gloriam et potentiam. Ad Roman.3. Ad Ephes. 1. Nec ferendi funt qui tanquam contemptores et hostes gratiae Dei plus iusto deferunt hummis viribus, operibus et meritis, siue ludaeisint, propriam quaerentes statuere iustitiam, ad Roman. 9. Siue pelagiani, siue quocunque alio appellati nomine. Pii interea agnoscant, se plurimum Deo debere pro accepris beneficiis, in primis spiritualibus, ideoque  agant indesinenter gratias, atque omnibus modis bonitatem Dei celebrent. Verum in hoc incumbent praecipuë, vt fructus gratiae Dei respondentes, hoc est, actiones ueris Christianis dignas per omnem uitam proferant: ne, quemadmodum 2. ad Corinthios on. dicitur, gratiam Dei frustra recepisse existimentur: quinpotius cunctis uidentibus conuersationem nostram subministremus occasionem glo rificandi coelestem patrem. Matth.5. 1.Petri 2. Assuescamus insuper in tentationibus, in pauoribus conscientiae, denique in quibuscunque  perieulis, consolationum materiam colligere ex gratiae Dei exuberantis consideratione.  \pend\pstart Iamad secundam partem definitionis, de fide, inquam, ac dilectione, quae ex parte ipsius hominis considerari in negocio salutis debet. Oritur primo loco fides, quandoquidem hanc Deus ante alia quaecunque in homine requirit: ad Hebr.11. Sine fide fieri nonpotest, ut quis Deo placeat. Nam qui accedit ad Deum, hunc credere oportet esse Deum, et esse remuneratorem quaerentibus se. Significat autem hocloco fides non sidem ciuilem, quam detiniunt philosophi ac iuriseonsulti dictorum factorumqueconstantiam: non item fidem, qua eduntur miracula, cuius fit mem1o1. Corinth. 13. neque fidem historicam, qua uidelicet credit quispiam simplicitez uera esse, quae de Deo, uel a Deo facta referuntur: qualem lacobus ait inueniri etiam in Daemonibus: lacobus 2. Sed firmam certamquecognitionem innixam diuinis promissionibus de gratia seu misericordia nobis impendenda propter Chustum. Quae descriptio uariis Scripturae locis explicatur: ad Hebrae.11. Fides est earum rerum, quae sperantur, substantia, argumentum eorum, quae non uidentur. Ad Romanos 4. Qui praeterspem sub spe credidit. Niri autem fidem promissionibus de misericordia a Deo factis, liquet ad Romanos 4. ldcircon ex fide datur haereditas, ut siu per gratiam, ut firma sit pmissio uniuerso semini Rursus de Abrahamo: Ac minime imbecillis fide, non considerauit suum ipsius corpus iam emortuum, cùm centum circiter na tus esset annos, nec emortuam uuluam Sarae: Verùm ad promissionem Dei non haesitabat ob in credulitatem, sed robustus factus est fide, tribuens gloriam Deo, certa persuasione concepta, quod is qui promiserat, idem potens esset etiam praestare. Quapropter etiam imputatum est illiad iustitiam. Et promissam gratiam conferre propter semen benedictum, uidelicet lesum Christum, crebro explicatur. Ad Galat.3. Vt in gentes benedictio Abrahae ueniret per Christum lesum, ut promissionem spiritus acciperemus per fidem. Item, Porro Abrahae dictae sunt promissiones et semini eius. Non dicit, Et Seminibus, tanquam de multis, sed tanquam de uno: Et in semine tuo, qui est Christus. Nascitur uero huiusmodi fides, simul, quan do auditur Dei ueibum, simul quando Spiritus Sanctus eum informat in cordibus Ad Roman.10. Fides est ex auditu, auditus autem per uerbum Dei. Actor.1e.Audientes autem gentes gauisae sunt, et glorificabant sermonem Domini, et credide. runt quotquot erant ordinati aduitam aeternam. 2. Thessal. 3. Non omnium est fi des, sed fidelis est Dominus, qui stabiliet uos. Ad Philippenses 1. Vobis donatum est pro Christo, non solum vt in eum credatis, verumetiam vt pro ilo patiamiie  \pend
\vspace{0.5cm}\noindent
\vspace{0.2cm}\rule{1cm}{0.2pt}\\ 
\hspace{0.2cm}\textit{mg}
\footnotesize Gratia Dei fundamentum et caussacf. ficiens est tust ficationis, electionis, uocationus et glorificationis. Pelagianorie et Iudaeorum eror. 
\normalsize| 
\hspace{0.2cm}\textit{mg}
\footnotesize 4 Parsdestnitionis: de side et dileclione quae in homine, qui iustifica Flee Pesse FIdes politica. Fides miraculosa. Fides histo. rica. Fidesususi cans. perces dei ex epist. ad Hebr. capne Rom.4. 
\normalsize| 
\hspace{0.2cm}\textit{mg}
\footnotesize Gratia prohissa confer erprepier Christum sehen bencdi, Eue Fides exuer Iegcere mpa. 
\normalsize| 
\section*{IN I. EPIST. D. PAVLI. }
\marginpar{[ p.47 ]}\pstart Ergo lic in homine spectanda est fides, ut tamen sollus Dei dono contingere, non parari studio uel industria humana intelligamus. Simulata et hypocritarum est fides, quae a Deo non inditur. Debet porrô haec inuicta ac certa fides in hac breuitate uerborum dupliciter considerari.  \pend\pstart Primo, ut necessaria est ante iustificationem: quandoquidem fide absque operibus homo iustificatur. Ad Roman.3. lustitia Dei per fidem lesu Christi in omnes et super omnes qui credunt. Non enim est distinctio: omnes enim peccauerunt ac destituuntur gloria Dei: iustificantur autem gratis per illius gratiam, per redemptionem quae est in Christo Iesu, quem proposuit Deus reconciliatorem per fidem, interueniente ipsius sanguine, ad ostensionem iustitiae suae propter remissionem praete.   ritorum peccatorum. Ibidem. Arbitramur fide iustificari hominem absque opcribus legis. Ad Galatas 2. Scimus non iustificari hominem ex operibus legis, nisi per fidem lesu Christi: et nos in lesum Christum credidimus, ut iustificaremur ex fide Christi, et non ex operibus legis. Ad Ephes.2. Gratia estis seruati per fidem, idquenon exuobis, Dei donum est, non ex operibus, ne quis glorietur. Etenim in negocio iustificationis fides sola, quemadmodum Sancti Patres loqui solent, locum habet.  \pend\pstart pore elucere debet adiuncta ei dilectio, proueniuntqueex fide bona opera secundum  Secundon, debet considerari fides, ut necessaria est etiam a iustificatione, quo temdilectionis regulam facta: qui enim iustificatus est, fidem non abiicit, neque in ipso minuitur, sed potius per Spiritum sanctum in iustificatione acceptum augetur, ideoque  studet per dilectionem bonis operibus abundare. Ad Galat.5. Nos spiritu ex fidespem   iustitiae expectamus. Nam in Christo lesuneque  circumcisio quicquam valet, neque  praeputium, sed fides per dilectionem operans. ad Ephes.3. Det uobis iuxta diuitias gloriae suae, » ut fortitudine corroboremini per spiritum suum in internum hominem, ut inhabitet Christus   per fidem in cordibus uestris, in charitate actis radicibus. ad Colos.1. Reconciliauit uos cor   pore carnis suae per mortem, ut exhiberet uos sanctos et irreprehensibiles et inculpatos in conpectu suo : Siquidem permanetis in fide fundati ac stabiles, nec dimouemini a spe Euangelii quod audiuistis. lacobus 2. fides absque  operibus mortua est. Ita namque  ubi aliqui fi de absque  operibus fuerint semel iustificati, in his necesse est post iustificationem fructus fidei, qui potissimum consistunt in operibus charitatis, perpetuô elucere. Etenim omnes loci scripturarum, qui de fide ac dilectione, uel de fide et operibus coniunctim loquunt, ad eos intelliguntur pertinere, qui iam sunt fide iustificati, quemadmodum non obscure docet Augustinus liber  de fide et operibus.  \pend\pstart Quae est per Christum Iesum.) Addunt̃ haec uerba ad discernendam ueram fidem ueramquedilectione a salsa, simulata et hypocritica: Notum namque  est utramque  ab hypocritis iactari uerbis, nomnunquam etiam specie quadam externa ostentari. Verum quando Deo, qui serutat corda, non probant, inanes et emortuae plane sunt. Matth.7. Non omnis, qui dicit mihi, Domi ne Dñe, intrabit in regnum coelorum. Et paulo ante dictum est de Charitate ex puro corde, et conscientia bona, et fide non simulata, prodeunte. Caeterùm tribus de causis dicunt, tam fides, quam dilectio esse in Christo siue per Christum. Primo, si spectar̃ earum causa efficiens atque  origo, oportet uni Christo accepta feramus. Fides enim Dei donum est, et propter Christum donat, ut paulo ante ostendimus. Vera dilectio ebullit ex fide, et accendit in cor dibus per Spiritum S.qui dat̃ in iustificatione. Secundo fides et dilectio in Christo sunt, ratione causae formalis, siue propter modum quo exercentur: In ueris Christianis multo aliter sese exerunt fructus utriusque uirtutis, quam in hypocritis et gentilibus: semper spirant in illis aliquid eximium et singulare, quod in aliis minimem deprehenditur. Hinc Christus distinguit iustitiam Christianorum a iustitia Pharisaeorum et seribarum: Matth. 5. Et Petrus in priori Epistola varias colligit circumstantias, ex quibus discrimen notari vult inter actiones hominis Christiani et non Christiani. Capite secundo iubet propter Dominum, propter voluntatem Dei, propter conscientiam Dei, tanquam seruos Dei, secundum gratiam Dei, pro gratia apud Deum beneagere. Capite tertio: Sanctificate, inquit, Deum in cordibua vestris: Sitis autem parati semper adrespondendum cuslsbet petenti, yt loquamini de  \pend
\vspace{0.5cm}\noindent
\vspace{0.2cm}\rule{1cm}{0.2pt}\\ 
\hspace{0.2cm}\textit{mg}
\footnotesize rides ex d qo peco. tingil 
\normalsize| 
\hspace{0.2cm}\textit{mg}
\footnotesize 1. Fides ne cessara es ait iaitr cationem 
\normalsize| 
\hspace{0.2cm}\textit{mg}
\footnotesize 2. Fides necessaria est post iustificationen. 
\normalsize| 
\hspace{0.2cm}\textit{mg}
\footnotesize Fidei fruct' in iustificatis debent clucere. 
\normalsize| 
\hspace{0.2cm}\textit{mg}
\footnotesize Discrimen uerae et fa se fidei. 
\normalsize| 
\hspace{0.2cm}\textit{mg}
\footnotesize Quomodo fides et di lectio sit pel Christum, 
\normalsize| 
\hspace{0.2cm}\textit{mg}
\footnotesize Cusac udls fas L aQ ne.** Chr. sto. tgens 
\normalsize| 
\hspace{0.2cm}\textit{mg}
\footnotesize 4. Christus est, proptes quem fides donatur a Deo. 
\normalsize| 
\section*{AD TIMOTHEVM CAP. I. }
\marginpar{[ p.48 ]}\pstart ea, quae in vobis elt, spe, cum mansuetudine et reuerentia, conscientiam habentes bonam, vt in hoc quod vobis detrahunt tanquam scelerosis, pudefiant hi, qui incessunt   vestram bonam in Christo conuersationem. Item: bona conscientia bene respondeat apud Deum. Capite quarto, charitatem credentium vult esle vehementem et eximiam, vult vnumquemque agere secundum dona accepta a Deo, vult ea um, qui loquitur, loqui, vt eloquia Dei: vult eum, qui ministrat, ministrare tanquam ex virtute, quam suppeditat Deus, idque ́, vt in omnibus glorificetur Deus per Iesum Christum.  \pend\pstart T'ertio, sunt fides et dilectio in Christo, ratione causae finalis: Sempernamque  inteni tae sunt ad celebrandam apud omnes, etiam hostes, gloriam Dei.1. ad Corinth.10. Omnia in gloriam Dei facite. Arque operibus bonis vacare nos oportet, vt glo, rificetur Deus: Marth. 5. ad Philip.L.1. Petri 4. Ac sancti Patres diligentes fuerunt in demonstrando discrimine inter virtutes bonasqueactiones Christianorum et virtutes gentilium et quorumcunque incredulorum. Vide Tertulian. Apolog. capite 25. Origenem contra Celsum, Augustinum contra lulianum libro quarto: libro de ciuitate Dei 19. capite 25. Dignum autem est obseruatu, quod Apostolus tam paucis summam doctrinae Euangelii expressit. Suntquehi loci, in quibus breuiter comprehenduntur capita vniuerlae doctrinae Christianae, tam de lege quam de Euangelio, notandi ac seligendi. Nam saepe occurrunt in nouo testamento, frequentissime vero in Paulinis Epistolis. Studiosi adolescentes excerptos in libellis, in chartis, slue tabulis ac parietibus merito saepius describerent: Imo publice in ecclesiis in loco digniori deberent exarari. Aequum autem est, assiduë inculcari in ecclesiis doctrinam de iustificatione per fidem absque  operibus, cum gloriam Dei vehementer commendet, et ad illustrandam illius bonitatem ac potentiam, rursumquead detegendam nostram infirmitatem, indignitatem, et bonorum operum, momentum magnum magnum adferat: ldcirco etiam paulo post gloria Dei immortalis et inuisibilis ab Apostolo praedicatur. Consimile est iudicium de fide ac dilectione, bonisqueoperibus, quae requiruntur coniuncta in iustificatis. Atque  interest piorum seipsos imô et alios perpetuo excitare, vt bonis operibus intenti sint, atque ita suam fidem ac iustificationem, dehinc fidem cum charitate a iustificatione declarent. Viua, efficax, et operans esse debet fides per fructus charitatis.  \pend\pstart Certus sermo et dignus quem modis omnibus amplectamur :) Priusquam absoluat summam doctrinae Euangelii, interponit autoritate apostolica et magno affectu paucaverba, quae certitudinem ac dignitatem eiusdem doctrinae mirifice commen dant. Adduntur autem quadam praeoccupatione. Nam cum Apostolus in tanta breuitate nullas admisceat probationes ex scripturis, facile aliquis reiiceret omnia, et causaretur scripturarum testimonia deesse. Verùm Apostolo, et quidem scribens ti ad Timotheum istarum rerum iam egregie peritum, satis est hoc pacto duobus verbis, Euangelii cùm certitudinem tum excellentiam asserere. Quo pacto iterum praesentis Epistolae cap33. et 2.ad Timoth2.ad Titum3. eadem asseuerationis formula vtitur, nec puto alibi eum id facere. Etiamsi protinus subiiciet probationem admodum efficacem, sumptam videlicet ab experientia siue de exemplo suae personae. Latenter tamen suggerit Apostolus his paucis verbis non Timotheo modo, verum etiam cuicunque  Christiano recordationem omnium locorum, quibus haec doctrina de Euangelio in scripturis comprobatur. Nam ad scripturas potissimùm remittit auditores: Et qui agnoscunt conceduntque Christianis doctrinae suae principia, hoc est, sacram scripturam, sufficit huiusmodi asseueratio per Apostolum facta. Sciunt namque de scripturis aperta haberi omnium dogmatum firmamenta. Si qui autem principia scripturarum Christianis non concedunt, his nullae rationes sufficient neque conueniet cum eis disputare. Non est autem operae pretium, hic per nos repeti locos seri, pturae, quibus  doctrina de gratia, fide, et dilectione per lesum Christum, item Christi aduen tu in camne, vt peccatores saluos feceret, confirmat Si quide ab initio mundipromisaio:  \pend
\vspace{0.5cm}\noindent
\vspace{0.2cm}\rule{1cm}{0.2pt}\\ 
\hspace{0.2cm}\textit{mg}
\footnotesize 3. Causa finalis fidei et dilectionis in Christu gloria Dei. 
\normalsize| 
\hspace{0.2cm}\textit{mg}
\footnotesize Praeoccupatio. 
\normalsize| 
\hspace{0.2cm}\textit{mg}
\footnotesize Doctrinae Christiane principia ex scriptura sacra deriuanda. 
\normalsize| 
\section*{N L. EPIST. D. PAVLI }
\marginpar{[ p.49 ]}\pstart nes de gratia et de Christo gratiae autore mittendo ad Adamum, deinde ad Patriarchas, postea ad prophetas et Dauidem factae sunt. De hac re oracula, uaticinia, conciones diuersas prophetae ediderunt. Hanc rem praefigurarunt multae Sanctorum patrum in ueteri testamento actiones, huius rei umbrae et uypi in sacrificiis et plurimis ceremoniis legalibus uisebantur. Recte ergo dici potest, quod hic sermo, haec doctrina Euangelica certa sit. Acuide optimum ordinem seruatum in uerbis Apostoli. Prius dicit, sermonem esse certum, postea addit, dignum omni commendatione. Expriore sequitur posterius, iusta forma Enthymematis. Nam illa demum doctrina commendarl mertur, quae solida accerta est  \pend\pstart Discimus autem hinc in docendi methodo conducere nonnunquam, ut tractationi dogmatum interseramus breues aliquas asseuerationes, non tam ad confirmandos, quam excitandos auditorum animos. Conuenit hoc interdum eorum intelligentiae uelnegligentiae, quos docendos suscipimus. Quanquam prudenter dispicere oportet, quando huiusmodi utisorma expediat. Neque enim eam decetadhibere passim, sed tùm cùm exposita est alias satis diffuse eadem doctrina, ut non necesse sit, probaticnes taediose repeti: uel auditoribus, quandoquidemin sacrarum literarum lectione exercitati sunt, iam ante nota esse potest:uel tania est eius, qui docet, autoritas, et de doctrina ipsius opinio, ut etiam sua non probanti ex sacris literis fides habeatur. Interea uerô auditores pii asseuerationes eiusmodi tempestiuë adhibitas benigne accipient, doctrinam propositam altius infigent memoriae, et pro certa indubitataqueeam retinebunt. Nec difficile est iudicare, qui nam doceant cum uirtute et tanquam autoritatem habentes. Aut si quit scrupulus reliquus est, conferemus ipsi locos scripturarum, exempla patrum, tantis per dum nobis de docrinae certitudine constiterit, nisi malimus doctiores consulere. Est uero haec magna laus doctrinae Christianae, quod certa sibique  constans ubique sit. Longe aliter schola Christiana, quam Academia illa Socratica, quae in hoc laborans ut aliena refellat dicta, et probabilia tantùm aliquo usque defendens, nihil definit, sed iudicium omne relinquit auditori. Schola Christiana proponit quae certa sunt, atque ab omni dubitatione procul remota. Certa autem sunt, quia diuinitus reuelata sunt, et eruditos nos possunt reddere ad salutem 2. ad Timotheum 2.2. Petri1.Aequissimum igitur est confirmemus animos, ut eam doctrinam constanter retineamus, neque ullis disputationibus, contentionbus, offendiculis, temtationibus, persecutionibus, insidiis patiamur nos auelli. Hominum quidem commenta dies abolet:uerbum autem Domini in aeternum manet. Ac potest fieri, ut interdum obscurius tradatur doctrina ac negligentius, praeterea hominum caecitas malitiaque  inducit subinde errores: Attamen ueritas doctrinae Christianae semper apud aliquons integra seruatur sine corruptelis et fermento Pharisaicarum traditionum, ac tandem patefit: Christus nunquam in totum deserit suam Ecclesiam, sed semper aliquam seruat, cui aderit ad consummationem usque seculi. Potest itaque ad tempus, et quibusdam in locis oppugnari ac labefactari ueritas, expugnari uero et prorsus delerinon potest. Efficax profecto mentium Christianarum confirmatio, praeterea ad resistendum tentationibus, quas suggerunt Diabolus et caro nostra, incitantes nos ad dubitandum de ueritate doctrinae Euangelicae, uel etiam de Dei in nos gratia, consolatio praesentissima et praesidium inexpugnabile.  \pend\pstart Qu'od Christus Iesus venit in mundum, ut peccatores saluos faceret.) Nonsine Occupatione haec asseruntur. Breuissime dictum est in negocio salutis humanae, ex parte Del exhiberi gratiam, ex parte uero hominum fidem ac dilectionem in Christo necessariam esse. Non abs re autem postularent aliqui etiam modum omnem, quo gratia illa Dei sic exuberauit, plenius cognoscere. Ergo quo satisfiat omnibus, clarius nunc exponit Apostolus, quamuis iterum mira breuitate, quomodo salus nostra ex gratia a Deo perfecta sit.  \pend\pstart Chrislus lesus, inquit, venit in mundum.) Tria autem sunt obseruanda membra, quibus vt generatim tantùm et aliquo modo obscurius dictum videri poterat de gratia Domini exuberante, distinctius explicatur.  \pend
\vspace{0.5cm}\noindent
\vspace{0.2cm}\rule{1cm}{0.2pt}\\ 
\hspace{0.2cm}\textit{mg}
\footnotesize Eoe Crei, ta ae solida 
\normalsize| 
\hspace{0.2cm}\textit{mg}
\footnotesize Pict ctrinae Chril stiane et So craticae siut philosophil cae Hominum con. menta et coera, se impfecta e caduca: Ver tmmnase lidum este eternum. 
\normalsize| 
\hspace{0.2cm}\textit{mg}
\footnotesize Quomodo Paice mibere 
\normalsize| 
\section*{AD TIMOTHEVM CAP. I. }
\marginpar{[ p.50 ]}\pstart Primum membrum est de persona atque  autore illius gratiae, nempe Christo lesu seu filio Dei. Vt enim nobis applicaretur Dei gratia, oportuit filium Dei humillari et sua Obedientia nobis diuinam gratiam promereri: oportuit illum fieri lesum siue Seruatorem, et Christum, id est, unctum siue regem, pontificem et prophetam, qui pro nobis apud patrem intercederet. Non est aliud nomen, in quo remissionem peccatorum consequamur: et gratia acueritas per lesum Christum exorta est. Actor.4.1o. Iohan.1. Atque hocnomine distat doctrina Euangelica a doctrinis aliarum sectarum: quod illa lesum Christum agnoscit, et perpetuo praedicat: caeterae autem abnegant et reiiciunt. Valde igitur necessarla est haec particula ad recte intelligendum, quid sit Euangelium.  \pend\pstart »Secundum membrum est de modo, quo gratia Dei nobis per filium parta est Venit, inquit, in mundum. In coelo erat gloriam habens cum patre communem: ul. tro inde descendit nullis illectus nostris meritis aut uirtutibus. Et uenit in mundum, ubi formam serui suscipiens est extreme humiliatus. Itaque hac particula compre, henduntur omnsa, quae Christus in sua humiliatione nostri caussa perfecit, ut incarnatio, mors, sepultura, resurrectio. Neque hoc membrum doctrinae Euangelle? ignoscunt ludaei, Turcae, et similes. Est ergo doctrina de humiliatione Christi, eiusque  pro salute nostra actionibus, Euangelii propria.  \pend\pstart I’ertium membrum est de Effectu gratiae, quod nimirum per eam salui fiunt peccatores. Late alioqui patet Dei gratia.' Etenim quod sumus, uiuimus, et mouemur, item, quotidiana beneficia, fertilitas terrae, pluuia tempestiua, ut ctus, uestitus, caeteraqueaccepta serri gratiae Dei debent: quod intelligunt et praedicant omnes homines, ettam ludaei ac Turcae et quicunque hypOcIIae. E uangelium autem sigillatim concionatur de gratia, qua ad salutem aeternam adducuntur homines peccatores. Deus namque salutem aeternam gratis et per gratiam, et in gratiam filii peccatorsbus confert. Quod qui non credunt, ab’E, uangelio sunt alieni, siue ludaei sint, qui propriam iustitiam quaerunt constituere, atque ex suis uolunt iustificari operibus: Ad Romanos. 9. siue titulotenus Christiani hypocritae, siue alii. Peccatorum nomine intellige primo quoscunque sine exceptione homines. Omnes enim peccauerunt: Ad Romanos.3. et1. Secun do, peccatores hoc loco sunt, qui tales se esse humiliter confitentur, ac perpetuon orant: petunt autem in nomine Christi sibi remitti peccata. Nam his tantum peccatoribus beneficium salutis impenditur. Ad peccatores, qui simpliciter tales sunt, et manent, siue confiteantur siue non confiteantur, concionatur lex, atque ante oculos ponit damnationem. Ad peccatores autem poenitentes dirigitur uoxE uangelii. Euangelium namque poenitentiam expetit, quemadmodum Lucae 24 mandat Christus, praedicari in nomine suo poenitentiam et remissionem peccato rum, in omnes gentes. Et loannes Baptista, Christus et Apostoli, a poenitentiase. cerunt docendi initium. Marth.3. 4. Marci 6. Actor. 2.3. 5.11. Mirifice uero amplificatur Dei potentia ac bonitas ex eo, quod peccatores potissimum uelit sal uos facere. Superba ratio hoc non credit, uerum statuit primum, quod homines aliqui per se iusti sint, sibique iustitiam quandam comparent: Deinde, quod ho rum tantùm qui uidelicet iusti hoc modo effectisint, habere velit rationem. Sec turpiter hallucinatur, quandoquidem falsae sunt hae duae hypotheses, quibus ipst nititur. Vera autem et certa est doctrina Euangellii, cui et acquiescere nos de cet, quod videlicet Christus idcirco in mundum uenerit, vt peccatores saluos faceret, quod ipsemet asserit Marci 2. qui bene habent, non egent medico: non veni ad vocandum iustos, sed peccatores ad poenitentiam. Matth.11. Venite ad me omnes qui laboratis et onerati estis, et ego refocillabo uos. Itaque gloria Dei haec maxima est, quod non iustos ( qui tamen nulli reuera sunt) sed peccatores saluos facit.  \pend\pstart Porrô salutis nomine comprehendit Apostolus electionem ab aeterno factam, vocatio nem in hac vita ad Euangelium, et in ecelesia iustificationem cum bonis et fructibus, qui  \pend
\vspace{0.5cm}\noindent
\vspace{0.2cm}\rule{1cm}{0.2pt}\\ 
\hspace{0.2cm}\textit{mg}
\footnotesize 1. Christus author gra ne 
\normalsize| 
\hspace{0.2cm}\textit{mg}
\footnotesize 2. Christi hus hiliatio. Philip.2. debr.2 
\normalsize| 
\hspace{0.2cm}\textit{mg}
\footnotesize 3 Effectun xσAι lus peccato; luin. Ag Euangelĩe de gratia Dei erga peccutores qu cutre Ialuantur, praedicat. 
\normalsize| 
\hspace{0.2cm}\textit{mg}
\footnotesize Modus acce fionis ad Ch stunt et alic hationis in Pide et intredulitate sita. 
\normalsize| 
\hspace{0.2cm}\textit{mg}
\footnotesize Peccatores, quircasiv. faluentur. Nebiete catoribus lex et quibus Euangeliñ praedicetur. Rauo buma nanon perci. pit mysterium lafcat. stificatiouis, etc. lohan3. Matth.48. Dugiere patefitex gratuita pee Fatoruures P mpere ne sificatone. 
\normalsize| 
\section*{IN I. EPIST. D. PAVLI. }
\marginpar{[ p.51 ]}\pstart comitari ipsam etiam in hac uita solent, denique et glorificationem, quae post hanc uitam aeterna apud Deum in coelis succedit. Vide, quam magna et multa Apostolus uno uerbo suggerat expendenda. Haec uero omnia per Christum contingere creden tibus luculenter docet Apostolus cum alibi, tum ad Ephes.1. et 2. Porro hoc membium de saluandis peccatoribus, postremo loco positum esse videtur, vt eo commodius conuertatur sequens probatio de exemplo suae personae.  \pend\pstart 1. Obseruandum proinde est dogma Euangelicum, quoscunque  sine vlla personarum exceptione peccatores salutem per lesum Christum consequi. Haec summa est siue generalis propositio, et Status doctrinae vniuersae, quam Euangelistae et Apostoli suis scriptis inculcant, sicut et loannes hoc fere pacto concludit suum Euangelium cap.20. Haec scripta sunt, inquit, ut credatis, quod lesus est Christus ille fillus Dei, et vt credentes vitam habeatis per nomem eius. Similiter lex et prophetae perpetuo sunt de vno Seruatore Christo concionati, Matth.5.Luc.24-  \pend\pstart 2 Redarguuntur omnes haeretici, qui de Christi descensu de coelis in terras, de eius incarnatione, morte, caeterisque , quae idem nostrae salutis caussa in assumpta natura humana perfecit, quouis modo dubitant, aut quaestiones praeter ea et contra ea, quae de his aperte sunt in Scripturis podita, curiose mouent et periculose contendunt. Sed illi potissimùm relmquantur, qui iubent homines salutem aliunde, quam a solo Christo expectare. Non opera propria quamlibet speciosa: non merita aliorum honinum, qui bonarum actionum suarum communicationem aliis promittunt, ne dicam pretio accepto communicant: non denique sanctorum hominum, qui e viuis excesserunt, inuocationes et intercessiones, salutem nobis adferunt, sed solus lesus peccatores, quemadmodum hic Apostolus diserte asserit, saluos faciat. Christus intercessor et mediator noster apud patrem, qui ideo in mundum venit, ut  \pend\pstart Eorum vero, qui salutis per Christum allatae sacti sunt participes, interest pie sanctaeque  viuere. Id autem quo animaduerti possit, tum quod ipsi vere in Christum credant, memores et grati sint pro beneficiis per Christum acceptis, tum quod Christi mors in illis sit efnicax atque mirabiles fructus spirituales producat. Ad Tit. 2.Illuxit gratia Dei salutifera omnibus hominibus, erudiens nos, vt abnegata impierate et mundanis concupiscentiis sobrie et iuste et pie viuamus in praesenti seculo, expectantes beatam illam spem et illustrationem gloriae magni Dei et Seruatoris nostri lesu » Christi, qui dedit semetipsum pro nobis, vt redimeret nos ab omni iniquitate, et pu,rificaret sibi ipsi populum peculiarem, sectatorem bonorum operum.  \pend\pstart Abs re autem iactant se esse Christianos, se vere in Christum credere, se agnoscere, quanta bencficia a Christo in omni vita acceperint, qui sic interim viuunt, vt barbaris hominibus, quibus nihil vnquam de Christo et doctrina Christiana fuit propositum, nihilo sint mellores: Imo propter eos nomen Dei et Euangelium blasphematur apud aduersarios.  \pend\pstart Caeterùm pii haec verba Apostoli altissime infigent memoriae, et ex singulis caplent vberem materiam consolationum, quandocunque se vel alios tentari animaduertent. Nec enim fieri potest, quin aliquando aut nos ipsi aut alii metu iudicii atque irae Dei, et propriorum scelerum conscientia animo consternemur, atque etiam in hac vita circumuallemur mortis et inferni doloribus. Operae pretium igitur est, nos hac sententia Apostoli et aliis similibus tanquam valido praesidio atque armis muniri atque sustentari. Ac colligent pii plures tales locos diligenter, qui ad manum sint quandocunque postulabit necessitas. Matth. 9. Non egent medici opera qui sani sunt, sed aegroti. Matth. 11. Venite ad me omnes qui laboratis et onerati estis, et ego refocillabo uos. Matth.1. Non sum missus, nisi ad oues perditas domus lsraelis. Marci 2. Non ueni ad uocandum iustos, sed peccatores ad poenitentiam. Matth.18. et Lucae 19. Venit filius hominis quaerere et saluum facere quod perierat. Lu» caeI5. Gaudium erit in coelo super uno peccatore resipiscente magis, quam super nonagin   ta nouem iustis, qui non indigent poenitentia. loan.3. Sic Deus dilexit mundum, ut filium suum   unigenita daret, ut omnis qui credit in eu, non pereat, sed habeat uita aeternam. Non enim  \pend
\vspace{0.5cm}\noindent
\vspace{0.2cm}\rule{1cm}{0.2pt}\\ 
\hspace{0.2cm}\textit{mg}
\footnotesize Propositio Euangelii e doctrinae el Cesicseer ralis est de s. lute homini et iustifica tione pecca torum 
\normalsize| 
\hspace{0.2cm}\textit{mg}
\footnotesize Horeticori parsaios Pontificiort traditione. En8d bo Dei et Euangelii de ctrinae aduersantur 
\normalsize| 
\hspace{0.2cm}\textit{mg}
\footnotesize 2. Vita Chi Repe aee uu esse debs at. 
\normalsize| 
\hspace{0.2cm}\textit{mg}
\footnotesize 4. Impiorl mores fide pressier Indigne 
\normalsize| 
\hspace{0.2cm}\textit{mg}
\footnotesize 1cosctee E LLhe 1eAie. los conscie tiae. 
\normalsize| 
\section*{AD TIMOTHEVM CAP. I. }
\marginpar{[ p.52 ]}\pstart misit Deus silium suum in mundum, vt condemnet mundum, sed vt seruetur mundus per cum. Maxima est foelicitas, quanta bona nobis in Christo proposita sint, agnoscere.  \pend\pstart Ex omnibus autem hucusque dictis constat, quid sit Euangelium, et quae illius dis e. gnitas atque prestantia. Decet autem pios etiam in aliis locis eiusmodi breues descriptiones diligenter notare, atque  inter se conferre. Extant autem passim, Matth.9. 11.24. Marci 1. 2.16.Luc.2.4.6.loann.1.3.5.6.17. Ad Roman.11o. 1.ad Corinth.1.2.5. 2. ad Corinth.5. ad Ephes.1.2. ad Coloss.1.2. ad Timoth.1. ad Tit.1.2. Actor. 1o.1. Petri 1.1. Iohann.2.5. etc.  \pend\pstart Quorum ego primus sum.) Comparatio descriptionis Euangelii ante positae, sumpta devero Euangelii vsu, quem declarat a suae personae exemplo et propria experientia perinde ac si dicat: verum esse, quod exuberarit gratia Domini nostri, quodq lesu Christus venerit in mundum, vt peccatores saluos faceret, id possum iure optimo asseuerare: nam in me ipso id sum reuera expertus. Maxime autem fide digna est pro batio, quae sumitur tum ab experientia, tum ab alicuius proprio iudicio atque sensu. Sic lohannes quoque ab experientia et sensuum iudicio confirmat Christi diui-- nitatem: conspeximus, inquit, gloriam velut vnigeniti a patre. Et 1. Iohann. 1.. Quod erat ab initio, quod audiuimus, quod vidimus oculis nostris, quod perspexie mus, et manus nostrae contrectauerunt, de sermonevitae, et vita manifestata est, et   vidimus, et testamur, et annunciamus vobis vitam aeternam, quae erat apud pa-- trem, et manifestata est nobis. Loquitur de persona Christi, in qua duae confideran; tur naturae contunctae: atque asserit certissimum et omnibus sensibus sibi comper. tum, quod duae illae naturae in Christo verae et perfectae existant. Huiusmodi autem probationes sumptas ab experientia et sensuum manifesto iudicio, et libenter proferent hi qui docent, et magnifacient auditores: Id autem partim, quia sunt perse certissimae, partim quia mirabiliter erudiunt conscientiam, ac demonstrant, quis sit germanus dogmatum, quae traduntur, vsus, siue quomodo doctrinae fructum applicare nobis ipsis debeamus. Ait igitur Apostolus, se compertum habere, quod Christus Iesus venerit in mundum, vt peccatores saluos faceret: quandoquidem ipse vnus exhis sit, quos Christus seruarst. Primum se vocat, id est, insignem et quasi principem, quomodo apud Latinos quoque vox primus non semper numerum, sed saepe excellentiam significat, veluti cùm dicitur primus Romana Liuius in historia: non quod ante Liuium historiam nemo scripserit, sed quod nemo id fecerit excellentius. Consulto autem Apostolus se primum peccatorem hoc pacto appellat. Prima causa est, vt testatam faciat suam humilitatem, quae in omnibus, qui peccatorum ducuntur poenitentia, considerari primo loco debet. Non dignatur Deus sua gratia, nisi qui seipsos extreme humiliarint: Vnde olim poenitentes ciner bus aspergebant capita, in cinere sedebant, sacco aut cilicio induebantur, ieiunabant, flebant. lob42. Plalm. 102. Esaiae 58. lerem.6. Et verus ecclesia Christianorun poenitentibus certum locum in sacro coetu, certas obseruationes, et signa quaedan externa poenitentiae internae praescribebant.  \pend\pstart Secunda, quia peccata ipsius reuera fuerunt maxima, nempe, blasphemia, perse cutio, violentia, vti semel confessus est. Etsi veron semel dixerit, se per ignorantian peccasse, qua ratione videbatur nonnihil extenuari peccatum: tamen quia illi non sa tis fuit veritatem ab Apostolis et a Stephano fideliter propositam non admittere (qua potissimum in parte ignorantia locum habet, quemadmodum ipse declaraus dicens, se ignorantem, fecisse per incredulitatem) verum etiam adiunxit manife stam blasphemiam, atque, vt verbis ipsius vtar Actor. 26. putauit se Hierosolymis aduersus nomen lesu Nazareni repugnando multa facturum: Et multos sanctorum inclusit carceribus, a principibus Sacerdotum potestate accepta: Et cum occiderentur, detulit sententiam, et peromnes Synagogas frequenter puniens eos, compellebat blasphemare, et amplius insaniens in eos, persequebatur etiam in exteras usque ciditates. Talia, inquit, et tam multa flagitia cum ipse de se commemoret, not\pend
\vspace{0.5cm}\noindent
\vspace{0.2cm}\rule{1cm}{0.2pt}\\ 
\hspace{0.2cm}\textit{mg}
\footnotesize Ab exemplo proprio prt bat gratis serueri peccatorem a Christe, etc. 
\normalsize| 
\hspace{0.2cm}\textit{mg}
\footnotesize n persona Christi dunaturae contuncte. 
\normalsize| 
\hspace{0.2cm}\textit{mg}
\footnotesize Ecιsfeα primum se uo cet peccato. rem. Hamilitss Apostolipec catoris; Ceremonipoenitentia agentium, lonae3. Mattbus 12. Luc.11. 2Caussa. Pecrata Paus li maexima. Acta.2e. 
\normalsize| 
\section*{IN I. EPIST. D. PAVLI. }
\marginpar{[ p.53 ]}\pstart abs re hoc loco confessionem istam edit, atque se uel primum peccatorem fatetur. Sola certe incredulitas maximum per se peccatum est:siquidem loan.16. Christus eam appellat peccatum per Antonomasiam: Spiritus sanctus, inquit, arguet mundum de   peccato, quia non credunt inme. Quid igitur statuetur, ubi aperta blasphemia, persecutio, uiolentia accesserint? Tertia, vocat Apostolus se primum seu maximum peccatorem, ut Dei, qui dignatus est tantum peccatorem in gratiam recipere, amplificet atque  extollat misericordiam. Quanto grauiora sunt peccata quae remittuntur, tanto illutrior fit Dei misericordia. Ad Rom.3. Iniustitia nostra Dei iustitiam commendat. Ad Roman.5. ubi exuberauit peccatum, ibi magis exuberauit gratia.  \pend\pstart Sunt autem quaedam nobis hic obseruanda. Primum ingenue confiteamur nostra peccata, praesertim quae nota sunt, et offendicula fratribus praebuerunt, nec vnquam dissimulemus illa, vel quouis modo extenuemus. Atque  hoc eos cum primis decet, qui contestari volunt, quod serio et ex animo ducantur poenitentia, quodquede Dei misericordia suapte natura desperent. Secundo, notabunt haec verba Apostoli, qui praesenti tempore veritati diuinitus manifestatae credere et se submittere nolunt. Sola incredulitas maximum peccatum et sufficiens aeternae damnationis causa est. Quid vero facient, qui non solum non credunt, verùm etiam palam blasphemant, oppugnant, persequuntur? Tertio, comprobatur dogma hic, quod qui serio egerint poenitentiam, et vera fide apprehenderint misericordiam propter Christum, certam fiduciam de Dei gratia debeant concipere, et credere, Deum etiam in hac vita sibi esse propitium Paulus fatetur se quidem fuisse maximum peccatorem: nunc autem videtur asseuerare, se esse ex eorum numero, qui misericordiam sint per lesum Christum consecuti. Sophistae quidem dicunt, nos non posse statuere, quod Deum habeamus placatum, et simus in gratia, nisi ex coniectura morali, prout rursus se dignum illa homo inuenerit. Nos veron Paulo adhaerebimus, et ad iudicandum, quod simus in gratia, oculum fidei non ad opera nostra, vel aliquam nostram dignitatem, sed ad exuberantem Dei gratiam, ad promissiones Dei, ad dignitatem lesu Christi, ad Spiritus sancti operationem in cordibus nostris dirigemus. Ad Roman.5.per Christum contigit nobis, vt fide perduceremur in gratiam hanc, in qua stamus et gloriamur sub spe gloriae Dei: et quae ibi sequuntur. Ad Roman.8. Accepimus Spiritum adoptionis, per quem clamamus Abba pater: Idem Spiritus testatur vna cum Spiritu nostro, quod simus filii Dei.1.Co rinth.2. Nos non Spiritum mundi accepimus, sed spiritum, qui est ex Deo, vt sciamus, quae a Christo donata sint nobis. 2. Corinth.13. Vosipsos tentate, num sitis in fide. Vosipsos probate. Annon cognoscitis vosmetipsos, quod lesus Christus in nobis est ? I.loann.2. Per hoc scimus, quod manet in nobis e spiritu, quem nobis dedit. Itaque  qui iubent credentes dubitare, remouent id, quod fidei proprium est: siquidem lubitatio et fides inter sese pugnant. Paulus autem sicut veram habuit et fidem et cha ritatem, ita audet dicere, se, quamuis maximus fuerit peccator, salutem tamen in Christo esse consecutum. Vehementer autem commendat haec doctrina Dei bonitatem et misericordiam: deinde pios excitat ad diligentiam in producendis bonis actionipus, quae sunt Spiritus sancti et fidei fructus: quinaetiam sunt testimonia gratitudinis erga Deum pro beneficiis acceptis: denique  mirifice consolantur in omni dubitatione et quibuscunque tentationibus. Quare merito retinenda est in ecclesia ac piis iuxta methodum Scripturae proponenda. Plurimum igitur aberrant, qui doctrinam hanc reiiciunt. Nam si verum vsum eius intelligerent, longe aliter sentirent. Causantur forte, hac doctrina induci in animos auditorum quandam stultam confidentiam, segniciem ad bene agendum, excuti Dei timorem: sed falluntur. Tantum namque abest, vt haec certitudo de gratia Dei tardos ad bonas actiones reddat, aut timorem Dei expellat, vt haec ipsa uel maxime inserat credentium pectoribus. Si quidem vbispes Dei est, ibi quoque  necessario est studium bene agendi, vera fiducia in Deum, verus Dei timor, non quidem seruilis et degener, sed qualis in filiis requiritur, quique  cum charitate est coniunctus, de quo Esaiae 11. et 1.loann.4. et alibi passim in scripturis. Quisquis autore Christo in gratiam perductus in ea stat, Deo agat gratias: interea autem ne  \pend
\vspace{0.5cm}\noindent
\vspace{0.2cm}\rule{1cm}{0.2pt}\\ 
\hspace{0.2cm}\textit{mg}
\footnotesize 3. Caussa. Dei misericordia et gratia erga peccatore maxima. 
\normalsize| 
\hspace{0.2cm}\textit{mg}
\footnotesize 1.Confessid ingenua de cet prccato res. 
\normalsize| 
\hspace{0.2cm}\textit{mg}
\footnotesize 2. Incredu tas peccatum maximum: Huic siBl. sphemia, pl fecutio, uio lentia, acce fruetons fit maeximum 
\normalsize| 
\hspace{0.2cm}\textit{mg}
\footnotesize 1. Poeniten tibus propt sita est pec catorumt re. missio. Sophistarum error. 
\normalsize| 
\hspace{0.2cm}\textit{mg}
\footnotesize De gratia Dei non est dubitandun piis et cre dentibus. 
\normalsize| 
\hspace{0.2cm}\textit{mg}
\footnotesize Certitude gratiae nonnobis, sed a Deo pendet. Rom.8 1.Corinth. 2.Corin. B. 1.Ioann.2. Fides et dil bitatio pugnant. Fidei fructe Doctrinafi dei necessaria in ecdesia. 
\normalsize| 
\section*{AD TIMOTHEVM CAP. I. }
\marginpar{[ p.54 ]}
\marginpar{[ p.55 ]}\pstart efferatur, sed timeat: Ad Roman.4. et u. fides in Deum, charitas in Deum, Timor Dei, coniuncta sunt, atque  his efficitur, vt pii clament, Abba pater.  \pend\pstart Verùm ideo misericordiam sum adeptus:) Duas adiicit et palam praedicat insignes omni busquenotandas causas, ob quas Deus dignatus est ipsum misericordia: quaȩ ita sunt digestae et e pluribus electae, vt vsum siue effecta Euangelii declarent, simulqueverum es se doceant, quod proxime dixerat lesum Christum venisse in mundum, vt peccato res saluos faceret. Licet autem, imo et prodest causas operum Dei nonnunquam inus stigare, et in lucem protrahere: quandoquidem earum consideratio in nobis fidem au cuit, vel ad timorem Dei, vel ad praedicandum eius nomem, ad gratiarum actionem, a stud um bene agendi aut omnino ad alia bona inuitat. Quod licet vidêre, tùm in mul tis disputationibus  in libro Hiob, tum in plerisque  psalmis, vt 19.104, 105,106.107. in qui bus excitantur homines ad consideranda admirabilia opera Dei, et canendas propter haec Deo laudes. Variae autem sunt eausae actionum Dei: Quaedam enim sunt, quas serutari humanae menti plane est impossibile, de qualibus ad Rom.11. dicitur: Quam inserutabilia iudicia et imperuestigabiles viae eius. Quaedam sunt, quas etiamsi huma nus intellectus posset assequi, tamen nosse minime expediret: in quibus explorandis qui vehementer laborat, stultae curiositatis incurrit notam. De negotiis hominum particularibus multae huiusmodi quaestiones moueri possunt:ut, Cur Deus huichomini longum vitae tempus, illi vero breue praestituerit: Cur Deus voluerit, vt ex duobus fra tribus vnus esset maior, alter minor: et maior interim minori seruiret: ad Rom.9. Harum et similium actionum Dei causas omnes nosse ne expediret quidem: sed satis audire, hoc ita erat decretum, Deus cuius vult miseretur. Cur autem velit, scire non est necesse. Relinquitur ergo, has demum causas in operibus Dei recte a nobis explorari, quas intellectum siue potius fide possumus apprehendere, et quas scire nobis nunc expedit, et has ipsas eatenus diligenter expendamus: proponemus insuper aliis expendendas, quatenus instituto de quo agimus, sunt accommodatae Nam omnes perquirere causas, et alias quoque  quae instituto non seruiunt, neque  vtile est, neque  necessarium, vt non dicam saepenumero etiam impostibile. Ad hunc proinde modum Paulus praesenti loco grauissimas causas ante oculos ponit, ob quas Deus ipsum dignatus est misericordia.  \pend\pstart Ptin me primo ostenderet Iesus Chrislus omnem clementiam.) Prima causa spectans a illustrationem gloriae Dei:quae et ideo primo loco posita est, quod in cosideratione causarum decet nos in ea esse ante omnia intentos, quae ad Deum ipsum pertinent: Siquidem De us sui ipsius et suae gloriae causa perficit vniuersa: ad Ephes.1.Hominum vero officium est, de creatoris ac Domini sui gloria inprimis cogitare, quemadmodum Christus quoque  iubet prius quaeri regnum Dei, quam quaecunque  alia: Matth.6. Vsus est autem Aposto lus verbis valde appositis et ad institutum accommodatis. Primon simpliciter respondet illis, quae de vsu Euangelii ante sunt posita, Venit Christus Iesus in mundum, vt peccatores saluos faceret. Secundo, singulae voces quadam Emphasi enunciari debent, vtpote quae mirabilem vim obtinent amplificandi. Ostendit, inquit, Christus, id est, euidentur et multo manifestissime demonstrat, vultqueDeus ab omnibus animaduerti. Ey Ataçis est indicatio rei cum primis animaduertendae: qua voce etiam medici vtuntur, ad significandum, rem praeuertendam et accuratem considerandam in morbs alicuius cutatione. Et item dicit, in me primo id est, maximo et insigni peccatore. Nec dicit simpliciter κακροθυμίαν, id est, clementiam seu potius longanimitatem, lenitatem animi sed omnem clementiant siue lenitatem, id est, multo maximam amplissimamque ́. Estque ; vocabulum hoc exquisitum et negotio Pauli aptissimum. Vere enim Deus μακροθυuos est, qui patienter expectat peccatorum conuersionem, nec quaerit eorum mortem, sed magis vt conuertantur et viuant. Atque  in conseruatione Pauli, qui asperrimus fuerat Euangelii hostis et persecutor, μακροθυμια Dei insignis eluxit. Celebratur autem hoc pacto ista Dei longanimitas ad Roman.2.1. Petri3.2.Petri3. Bonitatem vero et clementiam suam Deum velle, cum peccatores seruat, ostentare, ad Roman. 9. 11. ad Ephes.1.ad Tit.3. et passim alibi praedicatur.  \pend\pstart Adexpremendum exemplar je qui erediur essent in ipsa in vitamatemam.) Secunda çur  \pend
\vspace{0.5cm}\noindent
\vspace{0.2cm}\rule{1cm}{0.2pt}\\ 
\hspace{0.2cm}\textit{mg}
\footnotesize Quae causae ob quas mi. sericordian adeptus est Paulus. 
\normalsize| 
\hspace{0.2cm}\textit{mg}
\footnotesize vule est operum Dei causas inuestigare, ad fidem, adtimorem Dei, ad gratiarum. actionem, ad studium uita piae. 
\normalsize| 
\hspace{0.2cm}\textit{mg}
\footnotesize Adiones Dei quatenus intestigande, 9 contra. 
\normalsize| 
\hspace{0.2cm}\textit{mg}
\footnotesize LCaussa Il. uspos repes. 
\normalsize| 
\hspace{0.2cm}\textit{mg}
\footnotesize Emphaseis. 
\normalsize| 
\hspace{0.2cm}\textit{mg}
\footnotesize toeses 
\normalsize| 
\hspace{0.2cm}\textit{mg}
\footnotesize Deus auts duuos, 
\normalsize| 
\hspace{0.2cm}\textit{mg}
\footnotesize Excch.1n 
\normalsize| 
\section*{N L. EPIST. D. PAVLI }
\marginpar{[ p.56 ]}\pstart la ad aeumcatione et vmtate proxmiipectat. Deus quae facit, primo lacitpropterIele, et ad bonitatis, potentiae, iustitiae, et gloriae suae ostentationem. Secundo, propter homines, quos vult illa ipsa experiri ac notare, et propterea ab eisdem coli, metui, inuocari, praedicari. Quin etiam quaecunque  videt aliquis erga vnum hominem a Deo fieri, eadem debet et ipse a Deo expectare, siue bene agat Deus per suam misericordiam, siue puniat per suam iustitiam. Badem namque  est Deo potestas in omnes. Hoc igitur respiciens Apostolus ait, se expertum esse in seipso Dei clementiam, quo caeteri quoque  peccatores exemplo ipsius cle mentiam sperare auderent. Vocabulum ὐποτυπώσιως accommodatur tam doctrinae, quam actionibus. De doctrina 2. ad Timoth μυποτύπωαιν, id est, formam siue exemplar habeto sa norum sermonum, quos a me audisti. De actionibus vero loco praesenti. Quanquam aliis quoque  rebus, vt reipublicae, aedificiis, etc-accommodatur. Estigitur ὕποτύπασις exemplar bene expressum et euidens, quod quis anilno facile apprehendit ac sibi imitandum studiose proponit. Significat igitur Apostolus, se esse certissimum exemplar, et quassignum in publico expositum omnibus peccatoribus, blasphemis, persecutoribus, vt de Dei misericordia non desperent. Ne porro aliqui in impietate et flagitiis persistentes hoc Pauli exemplo proposito abuterentur, ac temere spem misericordiae ad se transferrent, prudenter adiicit Apostolus, quibus demum vsui esse exemplum ipsius queat, et quid in his, qui misericordiae Dei cupiunt fieri participes, requiratur. His, inquit, ego ero exemplar sperandae misericordiae, qui exemplo meo credituri sunt, qu per lesum Christum vita aeterna ipsis sit parata. Respicit autem ad id, qu paulo ante dixit in descriptione Euangelii, gratiam Dñi nostri exuberare cum fide et dilectione. Etenim Euangelium proprie offert gratiam et remissionem peccatorum per fidem in lesum Christum, ad Rom.1. Itaque  Apostolus plensssime, quauis mira breuitate, comprehendit hic, quicquid ad doctrinae Euangelicae summam pertinet: atque  ita absoluit, quae ad comprobationem descriptionis Euangelii antem positae adiicere visum fuit. Caeterùm ex his discimus, primô officii nostri esse, in omnibus rebus humanis considerare bonam ac iustam Dei voluntatem, et quo ad fieri potest, eiusdem rationes religios e inquirere. Physici dicere solent, Naturam nihil frustra agere.At Theologi multo rectius dicent, Deum nihil frustra agere. Nam in res vniuersas Dei prouiden tia perpetuo intenta est, easque ; sapientissime moderatur. Secundo, illas causas explora bimus diligentissime, quae conducunt ad celebrandam Dei potentiam, bonitatem, iustitiam, et gloriam: tùm ad hominum aedificationem et rectam institutionem conscientiarum. Hoc enim pacto de omnibus, quae quotidie in hac vita eueniunt, siue in ecclesiis, siue in rebus publicis, siue in priuatis nostris actionibus, pium et salubre iudicium ferre poterimus, acfructus spirituales vberrimos decerpere. Quod alicubi recte ordinantur ecclesiae vel respublicae, qu tyramni imperant et perturbant omnia, qu cientur bella, rursumquepax componitur, quod aliquis se conuertit ad studium sacrarum literarum, alius uero aliud amplectitur vitae genus, qu vni officium facienti bene cedit, alteri nequaquam: certum est, haec omnia sic diuinitus ordinari, ac partim quidem ad ipsius Dei gloriam, partim uero ad hominum institutionem. Ac considerant qdem illa omnia homines que  plurimi, sed longe aliter et multum diuer sas causas expendunt philosophi, quiquemundo sunt sapientes: alia vero Christiani uereque ́ pii. Illi considerant externa, carnalia, quaeque  secundum mundum sunt vtilia velnoxia, et ad aliorum hominum commendationem, aliorum vero oppressionem tendunt. Pii vero considerant interna et spiritualia bona, quaeque  prosunt ad illustrandam gloriam Dei, et ad aedificationem conscientiarum, hoc est, ad doctrinam, redargutionem, institutionem rectae vitae, correptionem vitiorum et consolationem. His certe qui seripturas legere consueuerunt, et aliquo suns Spiritus S.lumine adiuti, non est difficile plures et salutares causas diuinae voluntatis in omnibus, quae hic eueniunt, obseruare. Et concionatores in commonstrandis talibus cau us rudi populo diligentes esse debent. Porro, prima causa, quam Apstolus de Dei voluntate erga se recenset, suggerit dogma, nempe, qu Deus nos seruat, non propter nos vel nostra merita:quae em̃ esse huiusmodi queant in peceatoribus et filiis irae, quales omnes sumus etiam natura? Sed propter seipsum, sola sua bonitate et gratuita misericordia. Placet namque  Deo in his, quos ad salutem adducit, μακροθυμιαν omnem ostendere. Secundô, damnat iasuper hec ratio Apostoli omnes hypocritas et fallos doctores, quistatuuntali\pend
\vspace{0.5cm}\noindent
\vspace{0.2cm}\rule{1cm}{0.2pt}\\ 
\hspace{0.2cm}\textit{mg}
\footnotesize nifgoitordiae Dei erga Pailum, estaedificatio proximi Applicatio fadtori De 
\normalsize| 
\hspace{0.2cm}\textit{mg}
\footnotesize Ruangelj proprietas in oblatione gratie et re missionis peccatorum posi ta est. 2. Hominis est perpendere Dei a δηύαp. ProuidentiDei omnia sunt subiesta. a. Omnia Deigloriae debent inse nire. 
\normalsize| 
\hspace{0.2cm}\textit{mg}
\footnotesize vauaga mae cause. 2. Axioma. Deus non ser uat nos propter merita: nostra, sed Propter seipsum. 2 ustifari et mpocriudnar. 
\normalsize| 
\section*{AD TIMOTHEVM CAP. I. }
\marginpar{[ p.57 ]}\pstart ter, quam per solam Dei clementiam, homines saluos fieri. Tertio, admonemur, vt nostram semper agnoscamus indignitatem, coram Deo prostrati humiliemur, peccata nostra confiteamur, veniam propter Christum petamus, vtpote quos necessum sit perire, nisi Deus subonitate nos toleret ac succurrat. Quarto, quoram conscientia premitur ob peccatoru multitudinem vel magnitudinem, hi consolationem accipient intuentes in Del lenitatem et clementiam, qua peccatores ad poenitentiam et salutem multo benignissime atque  vltro reuocat. Secunda causa ab Apostolo proposita, primon docet nos omnia quae in Sanctis vel peceatoribus diuina ordinatione fiunt, esse proposita nobis pro exemplis, quae nos tum ad omne bonum excitent, tum a quouis mali genere deterreant. Secundo, discimus ex cadem causa, qu vnusquisque  nostrum quaecunque  exempla proposita prudenter et cum iudicio ad st debeat transferre, suisque  rebus accommodare. Nisi em istud faciamus, quomodo il lanobis proderunt? Legere vel audire nullius est momenti, sed ex applicatione ad nostra negocia commoditas spiritualis nascitur. Tertio, agemus Deo gratias, simulquenos sub, inde consolabimur, qu dignatus sit nobis in Paulo et plerisquealiis ut Dauide, Matthaeo, Zachaeo, muliere peccatrice, Petro, latrone, exempla monstrare, ex quibus consolationem hauriamus in tentationibus, et per quae excitemur ad veram poenitentiam, ad inuocationem ex fide, ad veniam et misericordiam petendam et expectandam propte Christum, ad vitae emendandae studium.  \pend
\phantomsection
\addcontentsline{toc}{subsection}{\textit{Regi autem seculorum immortali, inuisibili, soli sapienti Deo, honor, gloria, in secula seculorum, Amen. }}
\subsection*{\textit{Regi autem seculorum immortali, inuisibili, soli sapienti Deo, honor, gloria, in secula seculorum, Amen. }}\pstart Vt finem faciat huius primae partis de doctrina legis et Euangelii fideliter in ecclesia tradenda, orationem abrumpit, et tanquam per Epiphonema subiicit breuem, sed tamen lu culentam gratiarum actionem. Ac prudenti sane consilio istud facit. Quod em̃ postremo loco de gratia et gratuita per Christum misericordia disseruit, quid rectius agat, quam proti nus omnes vna secum inuiter ad gratias agendas Deo? In Euangelio gratia et veritas per lesum Christum proprie annunciata est vniuersis, et communicata credentibus. Ioan.1. Aequum igitur est, omnes pro manisestato Euangelio, proque  tantis beneficiis simul collatis Deum praedicare, et gratum exhibere animum.ltaque  optime cohaeret haec gratiarum actio cum antecedentibus, utpote quae ex illis ipsis deruatur consequiturque ́. Debemus autem istam gratiarum actionem accipere ueluti profectam ex magno affectu, et animo propter uitae anteactae et beneficiorum a Christo acceptorum recordationem, vehementer commoto, quomodo alias solent sermones quodam vehementi affectu abrumpi. Et constat ipsum Paulum in omnibus acrem et vehementem fuisse: quae res in caussa est, ut nonnunquam ali quanto perturbatior eius appareat oratio. Itaque  hac gratiarum actione subiecta plane innuit Apostolus, se non posse oratione exprimere, quanta sit Euangelii maiestas, et quam praeclara bona omnibus in Christum credentibus per Euangelium proponantur: ideoque  se uelle dicendi finem facere, nec habere quod adiiciat, nisi qu cupiat Deo omnes agere gra tias. Et huiusmodi abruptiones tacite monent auditores ad complectendum animo res lon ge maiores, quam dicuntur. Sic uero ad Rom.11. et alibi orationem abrumpit: ô profunditatem diuitiartt, et sapientiae, et cognitionis Dei:que  inscrutabilia sunt iudicia eius ? Ibidem: .  Quonia exillo, et perillum, et in illum omnia, ipsi gloria in secula, Amen. Caeterum haec uerba, ut Chrysostomo, Theodoreto, Ambrosio, Theophylacto placet, ad Patrem reseruntur a quo filius missus est, et qui sic dilexit mundum, ut filium suum unigenitum daret. Quanquam haud multum interest, siue ad Patrem et filium communiter haec uerba referantur, quemadmodum et Chrysostomus demonstrat, maxime cum Christus ipse dicat loans. Qui " non honorat filium, non honorat Patrem, qui misit illum:siue potius, ut Theophylactus monuit, ad Deum simpliciter et ad totam Trinitatem, hoc est, ad patrem et filium et SpirirumS. communiter. Nam potentia personarum in diuinitate non distinguitur, suntqueopera Trinitatis ad extra, ueluti in creatione mundi, et in restitutione humani generis fiunt, indiussa, semper interim seruata uniuscuiusque  personae proprietate. Vnde frequenter, cum de gratia Euangelica sermo instituitur, trium personarum fit mentio. ad Rom.1. Segregatus in Euangelium Dei, quod ante promiserat per prophetas de filio suo, qui declaratus fuit flius Dei cum potentia, secudum Spiritu Sanctificationis, ad Rom-s. Quod sispiritus eius,  \pend
\vspace{0.5cm}\noindent
\vspace{0.2cm}\rule{1cm}{0.2pt}\\ 
\hspace{0.2cm}\textit{mg}
\footnotesize 3. Homofe Iipsum debet noscere. 
\normalsize| 
\hspace{0.2cm}\textit{mg}
\footnotesize 4. Consolauio peccatorum, quibus conscienua lauxia est et ffficta. Vlitas secundae cαιses 1. Vsus. Exempla proposita. 2. Vsus. Applicano usums habet Spirituelem. 8. Vsus Deo agenda graia. 
\normalsize| 
\hspace{0.2cm}\textit{mg}
\footnotesize Oranioatrupta conti. nens gratia rum actionem. 
\normalsize| 
\hspace{0.2cm}\textit{mg}
\footnotesize Interpretae tuio parure etc. 
\normalsize| 
\hspace{0.2cm}\textit{mg}
\footnotesize Regulade operibus Trinitatis. E Personaerum proprietates fontincorfese. 
\normalsize| 
\section*{IN I. EPIST. D. PAVLI. }
\marginpar{[ p.58 ]}\pstart qui excitauit lesum a mortuis, habitat in vobis, etc. 2. Corinth.13. Gratia Dñi nostri lesu Christi, et charitas Dei, et communicatio Spiritus S.sit cum omnibus uobis. Singulae personae in bonis spiritualibus, quae per Euangelium accipimus, suo modo cooperantur. Sed et quae in hac gratiarum actione numerantur attribura, proprie vni Deo, tamen singulis in diuinitate personis conueniunt. Siquidem pater regnat: Psal.47.1.ad Corinth.15. filius regnat: Psal.2. Matth.28. lohan.18. Et Spiritus S.regnat, utpote qui eos, qui alios regunt, siue in ecclesia, siue in politia, constituit et gubernat, Num.27.ludic.5.6. Actor. to.ltem ducitacregit filius Dei:ad Rom.8.ad Gal.5.Destruit regnum diaboli et Antichristi Matth.12.2.Thes.2. Sic insuper et pater est conditor seculorum: et filius condidit secula. ad Hebr.t.ad Ephes...loan. 1.Et Spiritus sanctus est creator. Psal.46. Summa, in unum uerum et solum Deum competunt, quae hic recensentur singula, neque  ulla personaunius substantiae diuinae eximitur.  \pend\pstart Seculorum:) Quia Deus aeternus est, qui et secula ipsa condidit eademqueperpetuo gubernat, nullum facturus regnandi et gubernandi finem alludit ad Plal.145. Regnum tud̃ regnum in omnia secula, et potestas tua in omnem generationem et generationem. Homines praepotentes uocant reges terrae, et principes huius seculi, qui et abolent. Psal.2.1 ad Corin. 2.ltem diabolus uocat princeps mundi huius, Iohan. 12. et Deus huius seculi, excaecans sensus incredulorum:2. ad Corin.4. Atqui uerus Deus habens potestatem transferendi omnes hos reges et regna, quae tantum sunt simulachra inania, regnum habet aeternum, et tam in coelis, quam in terris positum. Ergo hoc attributo notatur immensa Dei potentia, per quam mirabiliter toti mundo patefactum est Euangelium.  \pend\pstart Immortali:) Quae sequuntur attributa plenius et euidentius explicant Dei naturam. ἐὔφεαρτω, id est, incorruptibilis, quod de Deo etiam dicit ad Rom.1. Significatur autem ahi quid amplius, quam si simpliciter diceretur, immortalis, nempe quod non solùm mors nullum in Deo habeat locum, sed ne corruptio quidem ulla, quae tanquam initin et aditus quidam est ad mortem. Sunt et Angeli incorruptibiles atque  immortales: similiter animus hominis: adhaec corpora nostra incorruptibilitatem induent, iad Cora5. Sed tamen in his ipsis corruptibilitas quaedam spectaf: In Angelis quidem ex lapsu, ex inquinatione eorundem, ex aeterna item damnatione, quam illi et similiter animi hominum impiorum subibunt cum maximis cruciatibus: corpora insuper nostra multam corruptione et mortem ipsam experiuntur. Atqui Deus singulari quadam ratione incorruptibilis existit, utpote in quem nec lapsus nec inclinatio uel imperfectio, nec mors uel damnatio cadit. Vere igitur natura et substantia sua Deus incorruptibilis omnibusque  rebus superior hic praedicatur.  \pend\pstart Inuisibili:) pertinet et hoc ad naturam siue subsfantiam Dei explicandam. Deus non est substantia corporea, sed Spiritus: quo nomine rursus agnosci debet omnibus rebus creatis superior et praestantior.1.loan.4. Deum nemo uidit unquam, ad Colos.1. Qui est imago Dei inconspicui. loan.4. Deus est Spiritus. Hinc uero sequitur, Deumuere, et, uti est neque  uisuneque  ullo sensu humano comprehendi a nobis posse. Ac solus Deus, inquit Ambrosius, per se incorruptibilis siue immortalis et inuisibilis est, si quae sunt creaturae immortales uel inuisibiles, eae non ex se, sed a Deo id habent. Habentautem per gratiam, quod uni Deo inest per naturam.  \pend\pstart Soli sapienti Deo:) Et sapientiam, et ea quae ante dicta sunt, in solum Deum quadrare satis perspicuum fit ex haclocutionis forma. Singulari autem iudicio sapientiae additur particula soli: quia uerba fecit de insigni illa sapientia, quae est doctrina Euangelii, in qua sapientia Dei insignis elucet, ad quam collata sapientia huius mundi, planem stultitia est apud Deum:. Cor.1.et 2. Roman.1. Nec contingit vllis haec vera sapientia, nisi quibus  Deus eam reuelat, a quo solo impetrani oratione debet.1.Reg.3. lacobus 1. Matt.11. Honor et gloria:) Quae amplissima sunt, quibusque ; maiora concipere mens humana haud potest, Deo tribui precatur. Apud homines honor et gloria sunt res dignissimae praeferunturque  opibus et pecuniis: vnde excelso animo homines gloriam prae omnibus rebus expetiuerunt. ltaque  et Deo, quae summa sunt, impendi debent. Vult autem ea a nobis praestari ex animo, et ex tota mente, ut ipse est καρδιιογνωσης. Deinde, non ad praescriptum modon tempus, sed in omnia secula, sicut ipse quoque  seculerum omnia est  \pend
\vspace{0.5cm}\noindent
\vspace{0.2cm}\rule{1cm}{0.2pt}\\ 
\hspace{0.2cm}\textit{mg}
\footnotesize Atribaus quae comme niter conue niunt fingu lis personts 
\normalsize| 
\hspace{0.2cm}\textit{mg}
\footnotesize Deus summus rex. 
\normalsize| 
\hspace{0.2cm}\textit{mg}
\footnotesize Acompara iis. 
\normalsize| 
\hspace{0.2cm}\textit{mg}
\footnotesize Deus immortalis: toclaρrω. 
\normalsize| 
\hspace{0.2cm}\textit{mg}
\footnotesize Adisparais Differunt inter se creator et cred turae, natur. et substantia. 2.Petri 2. 1ob 4.er13 Genes.3. Rom.5. Deus inuisi. bilis dega10. 
\normalsize| 
\hspace{0.2cm}\textit{mg}
\footnotesize Deus solus sapiens. 
\normalsize| 
\hspace{0.2cm}\textit{mg}
\footnotesize Deo honor e gloria de betur. 
\normalsize| 
\section*{AD TIMOTHEVM CAP. I. }
\marginpar{[ p.59 ]}\pstart dominus. Par namque  est honor et cultus, qui Deo defertur, quique  naturae Dei, quoad fieri potest, respondeat. Oportet veros adoratores adorare in spiritu et veritate, quemadmodum Deus Spiritus est, lohan. 4. Eadem comprecationis forma legitur ad Rom. 11.16. ad Galat... 2.ad Timoth.4.  \pend\pstart Amen:) Vox hebraea confirmans, quae ante dicta sunt, precationibus, vt saepe in Psamis, vel etiam imprecationibus ac deuotionibus addi solita. Numer.5. Deuter.23. Con sueuit proinde et Apostolus eam precationibus subiicere. Admonent porro nos ista omnia: Primùm, ut Deum talem agnoscamus et praedicemus, qualis hic describitur, Nam haec attributa sicreuera in solum Deum competunt, vt alteri nulli accommodari pari modo queant. Secundo, confirmant haec attributa Deum, quoniam talis tantusque ́ est, vere et abunde praestare posse ea bona omnia, quae per Euangelium exhiberi nobis propter Christum, proxime dictum est. Qui rex est potentissimus, incorruptibilis, inurisibilis et solus sapiens, potest ac vult amplissima, incorruptibilia, et spiritualia bona in nos conferre, Beneficia omnis generis quae absunt, a solo Deo debemus petere: si quae vero adsunt, soli Deo debemus accepta ferre. Tertio, Deum quem tantum agnoscimus, et a quo beneficia longe maxima accipimus, oportet nos modis omnibus celebrare, cultu interno ac spirttuali colere, gratias agere perpetuo, atque  omnia nostra, quae tam animo, quam corpore a nobis fiunt, ad gloriam Dei in terris amplificandam comparare.  \pend\pstart Hocpraeceptum.) Conclusio totius primae partis, qua summam doctrinae Christianaeput in ecelesiis recte constitutis tradi debet, breuiter repetiuit. Sic autem conclusio haec est contexta, ut permouere aculeosquerelinquere in animo Timothei possit: quandoquidem suasoria aliqua argumenta habet non parui ponderis, quemadmodum et fieri decet in conclusionibus. Optime vero respondent haec iis, quae ab initio dicta fuêre, sicut conuenit conclusiones semper respicere ad antecedentia. Vt enim in initio dixit, Timotheum debere denunciare siue praecipere quibusdam, ne diuersam sequantun doctrinam:lta hoc quoque  loco, Hoc praeceptum, inquit, commendo tibi. Singula autem uerba εμφατηῶς debent accipi, ut argumenta suasoria ex eisdem eliciantur. Primum igitur argumentum ducitur ab affectu: quando ut efficacius moueat Timothed appellat filium. Quo quidem pacto Apostolus primùm tuetur suam autoritatem, quam in eo esse decet, qui aliis maxime de ordinanda ecclesia praescribit. Dehinc demonstrat se studio optimo atque  affectu plane paterno dicere ac facere omnia. Adhaec testaf, se Timotheum plurimi facere, ac uelle insuper, vt a caeteris magno in honore habeat̃.  \pend\pstart Commendo praeceptum:) Particula.n! cum accusatiuo aliquando significat tempus, aliquando ad, vel in, vel de:quae res in causa est, vt diuersis modis horum verborum reddatur sententia. Qui igitur ad tempus referre malunt, ita vertunt: secundum praeceden tes vsque  ad te prophetias: Hoc autem pacto colligitur hic secundum argumentum suasorium ductum a perpetua consensione doctrinae in Ecclesia. Doctrina proxime explanata de lege et Euangelio consentit cum doctrina omnium prophetarum, qui Timothei aetatem praecesserunt. Quare scire debebat, sibi ab ea veritatis linea per prophetas monstrata minime discedendum. Plurimum autem ponderis habere haecratio debet: quandoquidem basis est, cuiomnes ministrorum labores inniti oportet, ideóque  passim a Christo et Apostolis inculcatur, cum demonstrant, omnem doctrinam suam legi et prophetis respondere Matth.22.26.Luc.1o.24. loan.5. Petrus Act.3.10. 1. Epist. cap.1.Paulus ad Rom.13.ad Ephes2. et3. Alii interpretantur το ἐπι notare, quod prophetiae, id est, praedictiones, reuelationes variae, et honorifica multorum testimonia de Timotheo praecesserint, quinetiam diu ante in ipso eluxerint certa signa, quae magni aliquid de ipso promittebant, quorum consideratione decebat ipsum ad fidem et diligentiam in omni officio praestandam vehementer excitari: Atque  ita interpretantur plerique  veteres scriptores, I heodoretus, Theophylactus, Chrysostomus, Ambrosius. Secundum hos igitur tertium argumentum hic habebitur, sumptum a signis antecedentibus atque  ab omnium de ipso expectatione.  \pend\pstart Vtmihites in cus bonam militiam.) Quartum argumentum a genere functionis eccle\pend
\vspace{0.5cm}\noindent
\vspace{0.2cm}\rule{1cm}{0.2pt}\\ 
\hspace{0.2cm}\textit{mg}
\footnotesize Ma 
\normalsize| 
\hspace{0.2cm}\textit{mg}
\footnotesize Concusio partis primae. 
\normalsize| 
\hspace{0.2cm}\textit{mg}
\footnotesize LArgumaï ub effectu. 
\normalsize| 
\hspace{0.2cm}\textit{mg}
\footnotesize 2. Argumen tum a perpe tua consenfione doctrine in ecdesia 
\normalsize| 
\hspace{0.2cm}\textit{mg}
\footnotesize 3. Argumen tum a signis antecedentibus  4 Argume. 
\normalsize| 
\section*{IN I. EPIST. D. PAVLI. }
\marginpar{[ p.60 ]}\pstart siasticaeatque  difficultatibus eidem adiunctis. Vtitur Metaphora siue Allegoria sumpta are militati, quae reuera multo aptissima est ad demonstrandum quantae molestiae, quot pericula eos maneant, qui ad munus docendi in ecclesia sesc accingunt, vel iam in eo versantur. Intelligi autem debent molestiae et pericula non tantùm externa, qualia multa perferet, qui docebit ignaros et rudes, qui in uiam reducet errantes, qui con solabitur afflictos: verumetiam perpetuo certare debet cum falsis fratribus, cum haereticis, cum philosophis, denique  cum ipso satana, qui mille modis ecclestas interturbat.  \pend\pstart Habens fidem et bonam conscientiam:) Quintum argumentum, quo indicatur modus et causa formalis recte tradendae doctrinae sanae. Fassi doctores, hypocritae, haeretici, aggredi quidem possunt docendi munus et tractare scripturas: verum quando non ha bent fidem, errant in dogmatibus: qui praeterea non habent bonam conscientiam, Hi doctrinam sanam corrumpunt, et saepenumeron scientes volentesqueabducunt in errores. Excitare igitur debent hae rationes eos qui ministerio ecclesiae sunt addicti ut in omnibus strenue laborent, sese non tantum milites, sed etiam imperatores praebeant: sed inprimis prouideant, vt doctrinam omnem pure explanent, acfide et bona conscientia proponant auditoribus: qui si similiter fide et bona conscientia verbum auditum amplectantur, fieri non potest, quin admirabiles sequantur fructus, tam ad gloriam Dei, quam ad hominum aedificationem referendi. Respiciunt haec verba ad illa semel posita in initio, vbi dixit Apostolus, Finis praecepti est charitas ex puro corde, et conscientia bona, et fide non simulata.  \pend\pstart 4 Qua repulsa nonnulli circa fidem naufragium fecerunt.) Ratio praecedentis argumenti a periculoso. Propterea, inquit, necesse est eos qui docent in ecclesia, ex fide et bona conscientia negocium suum facere, quoniam qui alia ratione ausi sunt progredi, neque  bona conscientia perfecerunt quod sui erat muneris, in pericula inciderunt longe maema. Respondent haec verba illis in initio, a quibus quod aberrarunt quidam deflexerunt ad vaniloquium: pulchre igitur vltima cum primis congruunt. Apostolus vt graphicem depingeret periculorum magnitudinem, vsus est multo aptissima metaphora. Naufragium, inquit, fecerunt, et quidem circafidem, hoc est, omnium bonorum spiritualium et falutis aeternae fecerunt iacturam. Perdunt salsi doctores et se, et auditores suos omnes, saepenumero etiam vniuersam perturbant ecclesiam. Nullum igitur grauius discrimen excogitari potest, quam error siue naufragium circa fidem.  \pend\pstart Quorum de numero sunt Alexander et Hymenaeus:) Comprobatio rationis ab exemplis eorum, qui cum non fuissent bona conscientia amplexari doctrinam euangelicam, neque  etiam pure synceriterquedocuerint, infide turpiter lapsi sunt. Plura verba addit argumento quinto, ac diutius explicationi eius immoratur, quod in ecclesia praecipua cura esse debet de fideli et pura docendi ratione. De Hymenaei falso dogmatez. ad Timoth.2. Hymenaeus, inquit, et Philetus circa veritatem aberrarunt, dicentes, resurrectionem iam esse factam, et subuertunt quorundam fidem. Alexandrum probabile est in eadem fuisse sententia, de quo et2. ad Timoth.4. Alexander faber aerarius mul tis malis affecit me, reddat illi Dominus iuxta facta ipsius. Quem et tu caue: vehementer enim restitit sermonibus nostris. Notat igitur Ephesi tunc ipsum fuisse et veritati solitum aduersari, quod hic moxuocat Βλασφηυε͂ν. De codem Actor.19. refertur, quod ludaeus fuerit, atque tumultum studuerit sedare exortum Ephesi.  \pend\pstart Quos tradidi satanae.) Secunda ratio, cur enitendum, ut fide et bona conscientia doctrina religionis pure semper tradat, sumpta a genere vindictae quae autoritate ecclesiae falsis do ctoribus infligitur: quae omnibus, qui non agunt bona conscientia, praesertim in his, quae ad doctrinam pertinent, merito proponitur metuenda. Tradere aliquem satanae, est quodammodo excludere a regno Christi, quod est ecclesia, et redigere in regnum ac porestatem satanae, ut ab hoc affligatur, molestetur, et ita extreme humilietur. Sed de huius ecclesiasticae uindictae ratione mox plenius disseremus.  \pend\pstart Vi discant non maledicere:) Caussa finalis siue effectus uindictae, quae per ecclesiam irrogatur: Coercentur hac falsi doctores, vt non amplius blasphement, hoc est, de rebus diuinis male loquantur. Debent namque, cum traditur quispiam satanae, ne\pend
\vspace{0.5cm}\noindent
\vspace{0.2cm}\rule{1cm}{0.2pt}\\ 
\hspace{0.2cm}\textit{mg}
\footnotesize tumn agene re functionis ecclest sticae eiusdemquediffi cultatibus. Metaphora are militar 
\normalsize| 
\hspace{0.2cm}\textit{mg}
\footnotesize Doctorum u rorum labd res. 
\normalsize| 
\hspace{0.2cm}\textit{mg}
\footnotesize 5. Argumen tum a modo et causafos mali doctri ne sanae tri dendae. 
\normalsize| 
\hspace{0.2cm}\textit{mg}
\footnotesize Heretici e faisi docto res, aliigd bus mala e conscientie et errant, in errores abducunt. A pericule fo. 
\normalsize| 
\hspace{0.2cm}\textit{mg}
\footnotesize Meuphor a nauigant pus. 
\normalsize| 
\hspace{0.2cm}\textit{mg}
\footnotesize Fassi doctores et see auditores su os in perni ciem ducunt. Ab exemplis. 
\normalsize| 
\hspace{0.2cm}\textit{mg}
\footnotesize Hymenaei e Philetierror de resu rectione. 
\normalsize| 
\hspace{0.2cm}\textit{mg}
\footnotesize Alexandri fabri aerari mditia. 
\normalsize| 
\hspace{0.2cm}\textit{mg}
\footnotesize Tradere di quem satanae quid. Ab effectu. Cause, proi pter quas si tanae tradi impios con uent. 
\normalsize| 
\section*{AD TIMOTHEVM CAP. I. }
\marginpar{[ p.61 ]}\pstart cessarii et vtdes proponi fines. Primo spectari oportet honorem Deidsiquidem ecclesiae interest prouidere, vt Deo suus tribuatur honos, ac nemo eius nomen blasphemet. Secundo, quaerenda est vtilitas eius qui peccauit, nempe vt culpam agnoscat,atque  poe nitens redeat in viam, sicqueecclesiae Dei rursus inseratur. Peruinet huc quodi ad s-orinth.5.mandat Apostolus, eum qui incestum commiserat, Satanae tradi, ad interitum carnis, quo spiritus saluus fieret in die Domini lesu: et2. ad Corint.2 vult. eundem recipi in gratiam. Tertio aedificatio et vtilitas aliorum. Fit enim hoc pacto, vt caeteri et tuti reddantur a contagio falsae doctrinae, quae per falsos doctores inuehif, et ipsi peccata conmittere abhorreant, obqueea parem vindictam sciant in se exercendam. Caeterum ex his discimus, primo maximopere necessarium esse, vt ministri eccleslarum, et quicunque alii, quos Deus dignatus est dono prophetiae siue vberiori cognitione verbi sui exornare: semper illi sint in tenti, quo exfide et bona conscientia doctrinam legis atque Euangelii pure synceriterqueexplanent. Erunt autem hac in parte diligentes, si frequenter cogitarint hoc esse praeceptum Dei, ac omnino Deum velle, vt ita faciant: Nemo qui sapit, contra Dei praeceptum facere quicquam praesumet. Secundo id perpetuo ante oculos nostros versetur, quod maxima eos maneant pericula, qui doctrinam sanam quocunque  modo audent corrumpere: Siquidem blasphemi fiunt, et in Deum ipsum iniurii: sicut verba praesentia declarant. Dehinc naufragium in fide faciunt: qua ratione non se modo, verumetiam alios innumeros in exitium pertrahunt spirituale et aeternum. Postremo, in iudicium et condemnationem incidunt Sanctae ecclesiae Dei, quando nimirum traduntur satanae. Tertio, cum audimus etiannum aetate Apostolorum quosdam foede lapsos, admirari ac metuere debemus Dei iudicium. Datum fuit illis audire ipsos Apostolos: conspexerunt in eisdem inaudita Spiritus sancti dona: viderunt miracula longe maxima: experiebantur Euangelii et ecclesiarum insperatam propagationem : viderunt Sanctorum fortitudinem et constantiam in confessione veritatis ad mortem vsque: quin etiam sen serunt in semetipsis occultam quandam vim supernaturalium et spiritualium bonorum: et tamen adhuc praecipitati sunt in errores, abnegarunt, imô et impugnarunt agnitam ueritatem, Dei nomensblasphemarunt. Nimirùm haec sunt inscrutabilia Dei iudicia, quibus alios eligit, uocat, et seruat: alios uero reprobat, excludit, ac perdit. Qui igitur stant, timeant Deum, acuideant ne cadant. qui lapsi sunt, Deum orent, ut lerigere ipsos rursus uelit. Secundo prouidebit interea unusquisque, ne causas lapsus ipse in se foueat aut stabiliat, quae sunt incredulitas et conscientia non bona: quamdiu in nobis firmae sunt fides et bona conscientia, tamdiu sperare debemus successus et constantiam. Verùm quando hae nutant, debilitantur uel concidunt, errori et abnegationi sumus proximi. Tertio, quemac modum tempore Apostolorum multi a ueritate Euangelii defecerunt, multi falsa dogmata muexerunt atque  sectas, neque  tamen propterea fas erat iudicare, quod Apostolorum doctrina siue Euangelium non esset ueritati innixum: ita et hoc tempore quanuris audiamus dissidia et offendicula quotidie oriri, nonnullos ex his etsam, qui ueritatem Euangelii profitentur, mouere certamina: non tamen suspectam habebimus aut damnabimus protinus doctrinam uniuersam: sed eam quae pura est, atque cum lege, prophetis et Apostolis consentit, prudenter distinguemus a fermento Pharisaeorum, atque illam quidem constanter retinebimus, hoc uero iubebimus ualere. Offendicula atque  haereses euenerunt et euenient omnibus  temporibus, quo nimirum probetur fides electorum, et qui probati fuerint, manifesti fiant: I.ad Corinth.11. 1.Petri.1. Nec caussa ulla erit sufficiens, ut nutemus ac doctrinis uariis et peregrinis circumferamur: ad Hebr. B. Quarto obseruare nos decet: primo episcoporum et eorum qui praesunt ecclesiis, officium esse, falsis doctoribus et quibuscunque  auditoribus sectarum in omni tempore obsistere, antequam plures per eos inficiantur. Secundo, diuersa esse quaerenda remedia, quibus tantam rem prudenter aggrediantur: benigne audire, docere arguere, increpare, et nihil non facere oportet, quo malum resecetur: Postremo  \pend
\vspace{0.5cm}\noindent
\vspace{0.2cm}\rule{1cm}{0.2pt}\\ 
\hspace{0.2cm}\textit{mg}
\footnotesize Honor deis 
\normalsize| 
\hspace{0.2cm}\textit{mg}
\footnotesize 2. Vilit.s peccaius 
\normalsize| 
\hspace{0.2cm}\textit{mg}
\footnotesize AcdiEcatia cierunt. . Tioth.3. 1. Doctrina legis et Euam ge.is pure tra dendi in ecdoia. 
\normalsize| 
\hspace{0.2cm}\textit{mg}
\footnotesize I. Periculolo quuecer4 rumpere. LElaspbemia, 2. Naufragil̃ circa fidem. 3. Iudicium ecdesie. Ill Exemplo et malo alie n0 debemus lapere. 
\normalsize| 
\hspace{0.2cm}\textit{mg}
\footnotesize Causa erroris et abnegaionis.. 
\normalsize| 
\hspace{0.2cm}\textit{mg}
\footnotesize 311. Osseii e piscoporuς 
\normalsize| 
\section*{IN I. EPIST. D. PAVLI. }
\marginpar{[ p.62 ]}\pstart autem vbi propem omnia videntur frustra tentata, licebit contumaces satanae tradere. Tertio, conuenit interea publice de omnibus his rebus agere, sicut qui falsa dogmata sparserunt, publice peccarunt. Propterea. n. Apostolus etiam nomina falsorum doctorum exprimit, idq tùm quo ipsi vehementius promoueantur ac pudefacti aliquando resipiscant: tum quo caeteri sciant, quinam ipsis fugiendi sint. Quarto: semper veron huc tendant consilia omnia et conatus, vt in omnibus demonstremus, nos ex animo quaerere Dei gloriam, et non de cuiuspiam pernicie, sed potius de communi salute omnium esse so licitos. Decet nos exemplo seruatoris nostri Christi quaerere oues perditas easdenque ad caulas reducere, et siquas morbidas deprehenderimus, sanare, arundinem com minutam non confringere. Math1o.loan.2.  \pend\pstart Porro, quid sit tradere Satanae: tum quas ob causas: item quomodo id ab ecclesiamiat: quae deinde sint effecta atque  fines, paulo accuratius disquiramus. Satanan est hebraeis odio habere, inimicum esse, aduersari: Genes.24. Esau ipsum lacob odio habuit: Vsurpatur eadem significatione verbum Satan, quod verbum occurrit i.Reg.5 Zachar. 3. Psalm. 39. 109. Hinc descendit et nomen satan, id est, aduersarius, quod generatim de omnibus accipi potest, qui nobis aduersantur et laedere student. Dauid 2. Sam. 19. ad Abisai filium Zeruiah quid mihi et vobis filii Zeruiamh? Estis enim mihi hodie in satan, id est, mihi aduersamini. Suadebant namque  vt inter ficeretur Semei, qui regem conuiciis affecerat. Atqui Dauid interpretatur, si id in reditu ac restitutione sua face, ret, odium et periculum inde sibi nasciturum: quamobrem sentit consilium hoc sibi aduersari. Et Matth.16.Christus ad Petrum cupientem vt Christus sibi parceret, nec adiret Hierosolymam, vbi praedixerat se interficiendum: Abi post me satana, obstaculo es mihi, quia non sapis ea quae sunt Dei, sed ea quae sunt hominum. Quo pacto et diaboli nomen legimus generali significatione etiam homini attributum. Christus loan.6. de luda, vnus ex vobis diabolus est. Cùm porro in diabolo odium, inuidia et inimicitia aduersus genus humanum maxima sit et perpetua ( siquidem inuidia diaboli peccatum introiit in orbem terrarum, et ab initio, vt Christus ait loan.8. homicida est) per Antonomasiam diabolus satan passim in scripturis appellatur. Et si quid memoratur factum ad hominum perniciem, id plurimùm dicitur a satana procuratum 1. Paralip. 21. surrexit satan contra Israel, et incitauit Dauid, vt numeraret lsraelem. EtI. ad Thessal. 2. queritur Apostolus, obstitisse sibi satanam, quo minus, cum tamen semel et iterum decreuisset, ad ipsos peruenerit. Itaque  tradere aliquem satanae simpliciter est permittere, eum aliquo modo ab illo malo spiritu molestandum, affligendum, diuexandum, prout Domino visum fuerit. Si autem a Domino traditus est satanae pius Hiob, vt primum affligeretur in facultatibus, deinde etiam in corpore: Hiob 1. et 2. Item Paulus 2. ad Corinth. 12. fatetur, datum sibi fuisse stimulum per carnem, nuncium satanae, qui ipsum colaphis caederet. His affine est, quod impium Saulem exagitauit spiritus Domini malus: 1. Samuel. 16.18. Item satanas in ludam introiisse dicitur, lohan.33. Ex quibus locis colligere licet. Primùm, infestari a satana pariter pios et impios. Dehinc alios plus, alios vero minus, nempe in fa cultatibus, in corpore, vel etiam in anima et corpore simul. Sed hoc obseruare opor tet discrimem inter pios et impios, quod illi quidem infestantur a satana, ad exercitationem, probationem, humiliationem, ad illustrandam gloriam Dei in ipsorum liberatione et conseruatione, vel ad alios eiusmodi fines, qui ipsis sunt salutares, et gloriam parant: Hi vero ad vindictam, ad poenam peccatorum, frequenter etiam ad aeternam damnationem. Cùm proinde Deus singulari iudicio propia sua autoritate consue uerit permittere interdum satanae, vt modo pios, modo impios, et alios mitius, alios acerbius diuexaret, placuit Deo eandem potestatem Ecclesiae suae communicare, et quidem vt pars disciplinae adeoquevindictae ecclesiasticae esset, qua redigere in ordinem posset improbos Etenim in nouo testamento peculariter data est piis potestas vt exemplo Chisti et virtute seu nomine Christi, de q pmissio facta fuerat ad primos parentes et contriturus esset antiqui serpentis caput, Genes.3. satanae insultus reprimant, eundem  \pend
\vspace{0.5cm}\noindent
\vspace{0.2cm}\rule{1cm}{0.2pt}\\ 
\hspace{0.2cm}\textit{mg}
\footnotesize Fradere sa nae quid? Eomonse n4. 
\normalsize| 
\hspace{0.2cm}\textit{mg}
\footnotesize Diaboli no men. 
\normalsize| 
\hspace{0.2cm}\textit{mg}
\footnotesize 2 Pii et im phparti festantur a s. tana, alii pl'. alii minus. Piorum et im piorum tentationes et ex ercitationes fine differunc 
\normalsize| 
\hspace{0.2cm}\textit{mg}
\footnotesize c lrcache ta ecclesi
\normalsize| 
\section*{AD TIMOTHEVM CAP. I. }
\marginpar{[ p.63 ]}\pstart doment, prosternant, vel aliquid eidem permittant. Huc pertinent hiloci: loan.12. Nunc iudicium est mundi huius: nunc princeps mundi huius eiicietur foras. Ioan.18. Spiritus sanctus arguet mundum de iudicio: quia princeps huius mundi iam iudicatus est. Et Matth.o. Marci 16. Luc.9.clare dicitur, Christum fecisse diseipulis hanc potestatem, vt in nomine ipsius daemonia eiicerent. Ad Rom.16. Deus pacis conteret Sa tanam sub pedes vestros. Ex his et similibus locis perspicuum fit, in nouo testamento piis datam esse longe aliam potestatem in Satanam, quam vnquam antea data fuerat. Verum quidem est, habuit ecclesia in veteri testamento nomnullas imprecationes, quibus peccantes deuouentur ac vindictae atroci subiiciuntur, vt Numer.5. Mulieri, quan maritus suspectam habet de adulterio, porriguntur cum execrationibus aquae amarae et haledictae. ltem Deurer.27. et 28.variae execrationes recensentur. Sed in illis non exprimitur aliquid de potestate in Satanam, quemadmodum in nouo Testamento: siquidem nondum venerat Christus, per quem Satanas est superatus, et subiectus cre dentibus. Legimus 2. reg. I. Elia dicente ad quinquagenarium: si vir Des sum, descendat ignis de coelo, et deuoret te, et quinquaginta tuos: bis quinquaginta subito absumptos. Et 2.reg.2. Elisa maledixit pueris qui illudebant ei, statimque  a duabus  vrsis e sylua prosilientibus  discerpti ex eis quadraginta duo sunt. Verùm haec vindicta non erat ordinaria, neque  ecclesiae in lege proposita, sed potius numeranda inter miracula quibus Deo placuit demonstrare, quantus honor veris ministris debeatur. Iden iudicandum de eo est, quod Act.5.refertur de Anania et vxore elus saphira, quod, cùm redarguerentur a Petro ob mendacium et fraudem de pretio agri diuenditi, subito expirarint. Adhaec incertum, an tales punitiones cesserint tantummodo ad interitum carnis, ad spiritus verô salutem. Siquidem illis subito morientibus non est concessus resipiscentiae locus, vri his, quos ecclesia in nouo Testamento Satanae tradit, quorum certe salutem ecclesia quaerit, non exitium: et medicinam illis potius exhibet, quam occidit: sicut 1. ad Corinth.5.licet videre. Ergo, vt concludamus, tradere aliquem Satanae est pars ecclesiasticae disciplinae ac seueritatis, qua nimirum declarat ecclesia, talem hominem propter flagitia sua non agnosci inter viua membra ecelesiae seu regni Christi: sed pos tius redactum in potestatem Satanae, a quo et possit ad tempus diuexari, grauius vide licet aut mitius, prout Deo visum fuerit ecclesiae iudicium comprobare:id autem, quo castigatus et humiliatus in viam redeat: quandoquidem de salute eius nondum despe ratur. Qua de causa hoc ipsum vindictae genus mitiori nomine excommunicationen nomnulli appellarunt: siquidem denunciantur homines propter peccata a communiont ac societate credentium tantisper exelusi ac separati, dum resipiscant. In quam sententiam Christus Matth. 18.ait, eum, qui non audierit ecclesiam, habendum pro ethnico et publicano. Et Paulus 1ad Corinth. 5. prohibet piis, ne cum flagitiosis commisceantur: ac ne cibum quidem cum illis capiant. 1.ad Corinth.16. et ad Galat.1.fiunt anathe mata. Itaque  nolle commiscêri cum aliquo, siue excommunicare, habere aliquem pro ethnico et publicano, tradere aliquem Satanae, pronunciare esse anathema, fere pro eodem accipiuntur, et simpliciter significant grauissimum vindictae seu poenae genus, quod ecclefia postomnia tentata ad emendandos et redigendos in ordinem homines infligere potest.  \pend\pstart Nunc dicamus, quas ob causas tradi aliquem Satanae iuxta verbum contingat. Cum Ape stolus 1.ad Timoth.. Satanae tradat Alexandrum et Hymenaeum, propter corruptam ab i psis doctrinam: Rursus 1.ad Corinth.5. Satanae eum tradat, qui rem habuerat cum nouerca: mi nifestum hinc fit, duo recte distingui hominum genera, in qs ecclesia suo illo iure vtetur, Alterum eorum est, qui peruertunt doctrinam. Alterum eorum, qui manifestis obnoxii sunt fla gitiis Ac de corrumpentibus  doctrinam plures scripturae loci aperte agunt. Ad Romae ꝗ dissidia et offendicula contra doctrina semel recte traditam gignunt. Ad Galat. 1.si quis praedicauerit Euangelium preter id, quod Apostolus praedicauerat. 1.ad Timoth.4. spiritus impostores tradentes doctrinas daemonsorum. Act.2o. lupi falsa loquentes. Ad Tit.3. sectarum autores: L.loan.4. et2. Epist. loan. qui doctrinam de Christo non adferunt. Hi, inquam omnes illis locis iubentur vitariab omnibus, vt nec Aue illis dicatur, et habeantur a\pend
\vspace{0.5cm}\noindent
\vspace{0.2cm}\rule{1cm}{0.2pt}\\ 
\hspace{0.2cm}\textit{mg}
\footnotesize Imprecationes ueteris te lanenti. 
\normalsize| 
\hspace{0.2cm}\textit{mg}
\footnotesize Tradere Satanae est ecclesiasticae di ciplinprrs. 
\normalsize| 
\hspace{0.2cm}\textit{mg}
\footnotesize Vindictaextraordinaria, Tmiraculo Ia 
\normalsize| 
\hspace{0.2cm}\textit{mg}
\footnotesize Excoemnuencatio. 
\normalsize| 
\hspace{0.2cm}\textit{mg}
\footnotesize Ancbenae 
\normalsize| 
\hspace{0.2cm}\textit{mg}
\footnotesize Cause, propter quas ati quis tradi s. lanae debeat, Duo genera eorum, in quos eccleI4 exercere aisciplinamt pus potse 
\normalsize| 
\section*{IN I. EPIST. D. PAVLI. EPIST: D. PAVLI }
\marginpar{[ p.64 ]}\pstart res in scripturis recensentur. 1.ad Corinth.5. Qui cum vxore patris, hoc est, cum nouernathemata. Caeterùm manifesta flagitia latius patent, et qui ob ea traduntur Satanae plu ca consueuerat. Vbi mox adiiciuntur, scortatores, auari, simulachrorum cultores, conuiciatores, ebriosi, rapaces. Repetuntur eadem.1.ad Corinth.6.ad Galat.5. vbi de operibus carnis agitur. Ad Ephes.5. Item 2. ad Thessalon.3. nominantur qui nolunt operari et inordinate ambulant. De his omnibus diserte dicitur, quod pii non debeant ipsis commisceri, imô ne cibum quidem cum illis capere, quin potius eosdem ex suo grege profligabunt. Porro generaliter dictum est a Christo Matth.18. Quod si ecclesiam non audierit sit ille velut ethnicus et publicanus. Similiter generale est illud Pau li.1.ad Corinth.16. Si quis non diligit Dominum lesum Christum, sit anathema, marana tha. ltem 2.ad Thessal.3. particula illa, inordinate ambulantes, generaliter accipi potest Li buit horum locorum separatim meminisse, vt animaduertat vnusquisque , sub his plura con prehendi peccata, quam in locis antecedentibus exprimebantur. Ac fieri potest, vt quis initio ob peccatum non adeo grande moneat ab ecclesia: si tamen contingit eum admo nitionem aliquoties factam contemnere aut etiam ludibrio habere, satis significat is, se ecclesiam audire nolle, et Christum ecclesiae caput non diligere: quamobrem meriton seueriore disciplina, hoc est, excommunicatione coërcebitur. Hucusq ex scripturis de cau sis, ob quas tradere aliquos Satanae ecclesia iure potest.  \pend\pstart Nunc videamus, qua via ad tradendum aliquem Satanae siue ad excommunicandum progressus fiat. Duo sunt in scripturis loci, qui prae caeteris de hacre clarius loquuntur: Alter est ipsius Christi Matth.18. vbi, postquam verba fecit de remouendis offendiculis, et reducendis ad ouile perditis ouibus, ita ad Petrum ait: Si peccauerit in te frater tuus, vade et argue ipsum inter te et ipsum solum: Si te audierit, lucratus es fratrem tuum: sin vero te non audierit, adhibe tecum adhuc vnum aut duos, vt in ore duorum aut trium testium con sistat omne verbum: quod si non audierst eos, dic ecclesiae: quod si ecclesiam non   audierst, sit tibi velut ethnicus et publicanus. Amen dico vobis quaecunque alligaue * ritis super terram, erunt ligata in coelo: et quaecunque solueritis super terram erunt soluta in coelo. Alter locus est Apostoli 1. ad Corinth.5. lam decreui tanquam   praesens, vt is qui hoc sic patrauit, in nomine Domini nostii lesu Christi congregatis   vobis et meo spiritu, vnam cum potestate Domini nosti lesu Christi, tradatur Satanae, .  ad interitum carnis, quo Spiritus saluus sit in die Domini lesu. Expendamus quaenam ordine hic dicantur: siquidem Christus gradus quosdam, qui ad exercendum grauissimum hoc iudicium pertinent, constituit. Primo gradu loquitur de peccato, quod, qui fra ter appellatur, in fratrem committit, dequeeiusdem priuata correptione. Sunt vero haec obseruanda. Primo distinguunt peccata vt alla sint occulta alia manifesta: vtrique  generi quomodo remedium adhiberi debeat, Christus praescribit. Ac de occultis primo loco agif. Occulta vero appellant, quae vni tantùm vel duobus  sunt cognita, vel etiam aliquot alujs, nondum tamen publicata et quam plurimis siue omnibus  manifestata. Secundon: quod si proinde frater ita peccauit, vt tu solus sis conscius: satis erit, si solus et priuatim eum commonefacias, vt resipiscere uelit. Tertio: cum Christus dicat de hoc officio, te debere fratrem lucrari, ita nimirum, ut perficias, quo peccatum suum agnoscat, desistat ab eodem, ac uitam libenter emendae: perspicuum hinc fit, summam charitatem, fidelem ex uerbo Dei institutionem, modestiam, oportunitatem, prudentiam, in toto hoc negocio requiri. Quarto: Necesse est, peccatum tibi soli notum diligenter celes. Nam, si euulgaueris, primo charitatem, modestiam, prudentiam, caeteraque ́, de quibus dictum est, exuis, ac testaris te potius esse inimicum proditorem, quam amicum monitorem. Secundo, laedis, perturbas, et exacerbas eum, qui lapsus est, sentitque pro remedio sibi infligi uulnus. Tertio: aliis fratribus infirmis et ignorantibus praebes offendiculi materiam et occasionem. Quarto: futurum est, ut temporis progressu peccatum alterius etiam sine teinnotescat: Deus non patitur crimina manêre occulta: Dies os mnia reuelat. Quinto quod de unoquoque homine priuato hic dicimus, accipietur etiam de ministro ecclesiae, qui criminis occulti, ab aliquo perpetrati, siue quia ipse uidit, siue quia unus alterue patefecit, siue ipse reus uluo est fassus, conscius suerite  \pend
\vspace{0.5cm}\noindent
\vspace{0.2cm}\rule{1cm}{0.2pt}\\ 
\hspace{0.2cm}\textit{mg}
\footnotesize agradus cor rectionis. 
\normalsize| 
\hspace{0.2cm}\textit{mg}
\footnotesize Quomodo quis trada tur Satanae. Locus Mat th.18. 
\normalsize| 
\hspace{0.2cm}\textit{mg}
\footnotesize RCortart 
\normalsize| 
\hspace{0.2cm}\textit{mg}
\footnotesize 1. peccata o custa et ma nifesta. 2'Adnoniuie priuata. 3. Charitas, institutio mo destia, opon tunitas, prut dentia hoc il negotiorequiritur. 4. Peccat occulta sun celanda fratri. 
\normalsize| 
\hspace{0.2cm}\textit{mg}
\footnotesize 2Miuisbric clesiaeɇ officis in corrigen us lapsis. 
\normalsize| 
\section*{AD TIMOTHEVM CAP. I. }
\marginpar{[ p.65 ]}
\marginpar{[ p.66 ]}\pstart Hiccerte neque  aperiet aliis, neque  quasi ad ecclesiam res esset delata, excommunicare la psum attentabit, sed tantùm vt priuatus priuatum corripiat. Rationes sunt hae. Primo .. cauere debet, ne detegendo alterius occultum crimen, ipse iis modis, quos iam enu-2- merauimus, peccet grauius Secundo, interest ministrorum ecclesiae, non tantùm publice vniuersos, verùm etiam priuatim singulos, et per singulas domos docere, arguere, hortari, solari, etc. quemadmodum se facere solitum refert Apostolus Actor.2o. I'ertion in conspectu est exemplum Christi qui cum sciret ludam esse auarum et proditorem, nequaquam tamen eum a caena exclusit: Quamuis et hoc verum est quod eum aliquoties non obscure monuit, vt ludas se pungi a Christo satis intelligeret: non tamen eum Christus caeteris Apostolis nominatim designauit. Quamobrem minister quoque  eum, quem occulte nouit peccasse, admonebit sane diligenter, ac consulet, a sacra caena vt abstineat: Si tamen adire ausus ille fuerit, nequaquam eum publice arcebit vel pudefaciet.  \pend\pstart lam de secundo gradu correptionis ecclesiasticae: Si te non audierit, inquit Christus, adhibe tecum adhuc vnum aut duos, etc. Haec porro conuenit in praesenti obseruare. Primo: adhuc de peccatis occultis sermo est, nempe his, quae aliquot tantùm hominibus, non autem omnibus innotuêre, vt publica dici nondum possint. Secundon: qui, adhibebuntur, ex iis erunt, qui iam ante sunt conscii. Nam ignaris non temere peccatum, alterius detegetur. Tertio decer et hic cuncta fieri, vt fratrem lucremur, hoc est, eo mo 4. do, quem proxime exposuimus. Quarto: Sed neque  statim, si non audiat qui monetur, deferenda causa est ad ecclesiam: quandoquidem peccatum eius nondum pro publico habetur. Fieri autem non potest, quin illius crimen, si, quemadmodum coepit, perre xerit, tandem manifestetur omnibus. Tum verô Ecclesia proculdubio cognoscet. Quinto: quod si inter conscios perpetrati criminis aliquis fuerit minister Ecclesiae ad* priuatam istam admonitionem adhibitus, ne is quidem ad ecclesiam rem deferet, multo minus ius ecclesiae in excommunicando a se exercendum putabit. Vnus minister, vel etiam duo: tresue ecclesiam non constituunt, nec ecclesiae ius habent: Recte tamen maiori, quam antea, seueritate lapsum corripiet, atque  ne sacris participet cum caeteris, stu debit persuadere. Sexto: quicumque  autem ad hanc admonitionem adiunguntur, ii non tantùm vt otiosi testes, verùm etiam vt amici consultores aderunt, et pari simul diligentia hortabuntur, obsecrabunt, vrgebunt, atque  omnem mouebunt lapidem: quo peccatorem reducant in viam: Id subindicant haec Christi verba, Si eos non audierit.  \pend\pstart Tertius gradus, quo fratres peccantes ad salutem addueuntur, is est: Quod si eos non audierit, dic ecclesiae. Vbihaec notantur. Primo, si praecedentibus duobus gradibus quicquam profectum est, atque  lapsus peccatum suum agnoscens vitam emendalrit, nequaquam opus est, ad hunc tertium gradum conscendamus. Secundo:agitur hoc lo-2.  co proprie de peccatis publicis, et a quibus offendiculum longe latequead omnium ordi num homines serpit. Tertio: vbi quis pertinax est in malo, atque  impietas eius paulatim? innotescit omnibus, fieri non potest, quin ecclesia aliq modo resciscat. Nec multum refert quomodo primùm reseiscat ecclesia, dum tamen certo et legitima ratione cognoscat. Siue autem id renunciet ecclesiae per eos, qui primi sunt alterius crimine offensi, vel qui iam ante aliquoties ipsum monuerunt, sed frustra se laborasse aduertunt: vel per propinquos siue amicos: vel per aliquos e ministris ecclesiae satis apte id factum iudicabimus, si modo constet, eosdem id facere ex charitate atque  salutis illius desiderio. Non magis improbari debet istorum fides ac diligentia, quam medici vel chirurgi, qui cùm videt, se aegrum non posse curare, aliorum opematque  consilium implorat, quos nouit esse peritiores. Quarto: ecclesia, vbi renunciat fratris constans malitia, Chrysost. interprete, homil.61.in Matth. sunt presules et presidentes ecclesiae. Certem Apostolus 1.ad Timoth.4. congregationem siue consessum eorum, qui manus imponebant designatis ad ministerium, reliquasqueecclesiae pu blicas actiones exequebant, vocat πρεσθυτεριοy. Et Ambros.ca.. eiusde Epist ad Timoth. testat̃, in vnaquaque  ecclesia olim fuisse aliquot seniores sine quorum consilso nihil in ecclesia agebat. Et capa.in ecclesia. inqt, septem diaconos esse oporret et aliquot pres, byteros, vt digni sint per ecclestas, et vnus in ciuitate Episcopus Sed et ipse Aposto.  \pend
\vspace{0.5cm}\noindent
\vspace{0.2cm}\rule{1cm}{0.2pt}\\ 
\hspace{0.2cm}\textit{mg}
\footnotesize Il. gemus cor reptionis etclesiasticae. 
\normalsize| 
\hspace{0.2cm}\textit{mg}
\footnotesize III. gradus correctionis. 
\normalsize| 
\section*{IN I. EPIST. D. PAVLI. }
\marginpar{[ p.67 ]}
\marginpar{[ p.68 ]}
\marginpar{[ p.69 ]}
\marginpar{[ p.70 ]}\pstart lus non obscure dutinguit duo presbyterorum genera in ecclesia: alios, qui docent et Sacramenta dispensant: de quibus sunt illa verba ad Tit.I.Presbyterum, siue vt mox loquitur, Episcopum oportet esse tenacem eius, qui secundum doctrinam est, fidelis " Sermonis, vt potens sit exhortari per doctrinam sanam, et contradicentes conuincere: » alios, qui adiuuant Episcopum et hosce presbyteros in ecclesiae gubernatione: quibus Lad Corinth.12. dicitur dininitus collatum donum gubernationis, quod profecto de magistratibus politicis, qui nulli tunc erant inter Christianos, accipi non poterat. De vtroq autem genere 1.ad Timoth.5. Qui bene praesunt presbyteri, duplici honore digni habeantur, maxime ii, qui laborant in Sermone et doctrina. Nimirum et Episcopi et presbyteri omnes simul administrant Ecclesiam. Ex presbyteris autem cum aliisglùm in gubernatione versentur atque  praesint, alii pariter in gubernatione et doctisna ecclesiae praesunt, atque  eo nomine duplici honore digni sunt. lraque ; ad hos omnes simul quae deferuntur, ea non male dixerimus ad ecclesiam deferri, siquidem hi communi omnium ordinum nomine, quae ecclesiae sunt curant. Per hos quoque  fit, vt postea plebs sessum ecclesiasticum caussa producta est, illius interest diligenter de omnibus inquirere, quo rem totam certo cognoscat:nec licet ad cuiusquam condemnationem festinare, nisi diligenti inquisitione et maturo consilio praeeunte. Huc pertinent illa 1.ad   Timoth.5. Aduersus presbyterum accusationem ne admiseris, nisi sub duobus aut tri* bus testibus. Et paulo post: obtestor, vt haec serues sine praecipitatione iudicii, nihil faciens iuxta propensionem animi. Sexto: Quia vero ad consessum ecclesiasticum contingit varia de peccatis publicis deferri, operae pretium est, vt Seniores omnibns benem exploratis genera criminum diligenter distinguant. Quaedam sunt leuiora possuntque ́ nuncupari delicta:quaedam sunt grauiora et plane scelera. Adhaec, quidam labuntur qui antea modeste et inculpatem vitam semper transegerunt: quidam vero nunquam anea visi sunt frugi esse. Cum illis igitur, et qui ob prius illud genus peccatorum monitione acreprehensione egent, merito agetur mitius, iuxta ulud ad Galat.6. Fratres eti,amsi occupatus fuerit homo in aliquo delicto, vos qui spirituales estis, instaurate huiusmodi Spiritu mansuetudinis. Cum posteriori hominum genere, et qui flagitiorum sunt rel, seuerius licebit agere, sic tamen, vt peccatores semper intelligant, ex charitacunque  autem fuerint et cuiuscunque  criminis rei, qui a consessu ecclesiastico reprehenduntur, si obiurgationibus, exhortationibus, et authoritate Seniorum permouentur, excommunicationis medicinam procedatur: Sed illis praescribetur publica poenitentia, qua videlicet stato tempore coram tota ecclesia confitebuntur peccatum, humiliter postulabunt sibi condonari, vitae promittent emendationem in omne tempus posterum. Atque  hoc pacto ecclesia fratrem lucratur, ouem perditam reducit in ouile. De modo vel tempore poenitentiae non est quod in praesenti agamus. Octauo: Quod si, qui reprehenduntur, nihili faciunt salubria et mitia Ecclesiae remedia, sed contemnunt quicquid ad reducendos ipsos in viam tentatur, tùm solum restat extremum remedium Excommunicationis, ad quod inuita ac gemens Ecclesia progreditur: sicut medicus postquam omnia frustra conatus est, cum dolore parat sese, vt membrum de corpore abscindat. Hic iam constituamus quartum et vltimum gradum disciplinae ecclesiasticae, qui in ipsa actione Excommunicationis consistit. in qua verba Apostoli inprimis expendere et sequi nos oportet. Primo: Quando igitur admonitiones priuatae, admonitio in consessu ecclesiastico nihil profuerunt: sed frater peccatis tanquam Sus luto adhuc inuoluitur: ante omniarecte fiet, si, quibus inlocis vel Episcopus ipse non est praesens, vel ꝗ sedent ad ecclesiae gubernacula presbyteri, arbitrantur aliquas ob caussas expedire, ad Episcopum, superiores, aut doctiores homines alibi agentes scribat, atque  vel mandatum ab Episcopo, vel certe consilium a praestantioribus petatur. Hacratione declarabunt, se nihil ex vitioso affectu, vel parum considerate agere, pariterque  maiore fiducia in publico ecclesiae negocio deinceps procedent. Hucpertinet, quodexrad te non ex odio cuncta fieri, et salutem suam quaeri, non perniciem. Septimo: Qualesagnoscunt peccata, petunt veniam, submittunt se ecclesiae Dei, tùm non est, quod ad si quando ad Excommunicationem proceditur, resciscat. Quinto: Vbi porro ad con\pend
\vspace{0.5cm}\noindent
\vspace{0.2cm}\rule{1cm}{0.2pt}\\ 
\hspace{0.2cm}\textit{mg}
\footnotesize Dioges presbytero nen. 
\normalsize| 
\hspace{0.2cm}\textit{mg}
\footnotesize Quomodo tractandi l poenitentes Excomunt catio extremum ecclesiremedium. A simili. Agradus di sciplinae ecclesiastice est Excommu ricatio. 
\normalsize| 
\hspace{0.2cm}\textit{mg}
\footnotesize 4 Quailo et aduersu quos exerch da excommunicatio. Ordie etconsiliores gerenda. 
\normalsize| 
\section*{AD TIMOTHEVM CAP. I. }
\marginpar{[ p.71 ]}\pstart Corinth.s, colligitur, Paulo Apostolo a nonnullis Corinthiis fuisse significatum, quou apud ipsos aliquis patris sui haberet vxorem, Et 2.ad Thessal.3. quod si quis non obedit sermoni nostro, per Epistolam hune indicate, et ne commercium habeatis cum illo, vt pudore suffundatur. Sin vero caussae subsunt, ob quas non est necessarium ad illum aliosqueperscribi, ecclesia facere istud omittat suoque vtatur iure. Secundo Expediet porro, vt, quoniam peccatum iam est manisestum, is qui peccauit argu atur etiam publice, atque ad resipiscentiam omni studio inuitetur. Sic enim A postolus praecepit 1. ad Timoth. 5. Eos qui peccant coram omnibus argue, vt et caeteri timorem habeant. Hacratione plebi innotescit actio ecclesiae, atque ad plebem defertur, vt testimonio atque consensu suo approbet. Nec enim esse conscia et adfliberi debet, vt actionem moderetur: verùm vt obseruet, testis sit, et significationem det, sibi iudicium ecclesiae placere. Ac fieripotest, vt qui intelligit, se publice perstringi ob suum peccatum, incipiat moueri ad remittendum sceleratae vitae nuncium. Quod si vel hoc tempore fiat, bene erit, atque lapsus dum se ecclesiae subiiciens iudicio publice peragat poenitentiam, ac vitae promittat emendatio nen, abioipetente.  \pend\pstart Tertio, Fratre lapso nullam adhuc prae se ferente resipiscentiam iudicat ecclelia extremum restare remedium, nempe, vt membrum putridum a reliquo corpore resecetur: Itaque  presbyteri siue seniores conuenientes decernant, ecclesiae publicum esse indicendum luctum et moerorem. Esse enim fratrem, qui in peccatis pertinax manc at, qui nolit audire ecclesiam, quem necesse erit, quando aliter saluti ipsius con suli haud potest, a communione excludere fidelium: lstud ex verbis Apostoli discimus I. ad Corinth.5. Et vos inflati estis, ac non potius luxistis, vt tolleretur de me, dio vestrum, qui facinus hocperpetrasset. Et 2. ad Corinth.12. Metuo ne iterum, v bivenero, humilem faciat me Deus meus apud vos, et lugeam multos eorum, qui ante peccauerunt, nec eos poenituit immundiciae, libidinisque , et impudicitiae quam patrarunt. Decet autem hic publicus luctus totam ecclesiam diuersas ob cam sas. Primo, ostendere vult, se dolere vicem fratris, qui adeo est excaecatus et induratus, vt monitiones adeoque vniuersam ecclesiam aspernetur ac propte rea sit abscindendus a corpore ecclesiae: omnia membra corporis dolent, quando v num abscindi debet. Secundo, ecclesia declarat studium suum pro asserenda Del gloria, et vindicando nominis diuini contemptu. Tertio, vt agnoscatur peccati magnitudo. Significat autem ecelesia eiusmodi esse, propter quod Deus vniuersam multitudinem puniret, nisi in illud animaduertatur, simulqueorat ecclesia, vt Deus parcat, et mala impendentia auertat. Quarto, commendatur hac ratione ecclesiastici iudici seueritas. Oportet iustis et grauibus  causis submittatur iudicium, propter quod toti ecclesiae luctus indicitur. Et ecelesia dum paret, testatur se approbare. Quinto, qui lugen peccata aliena, eadem opera deflent sua: suam ipsi consitentur indignitatem ac deplo. rant, petunt veniam criminum ante patratorum, petunt seruari ne in posterum ruant in peccata similia. Possunt et aliae commoditates ex hac totius ecclesiae humiliatio. ne sequi.  \pend\pstart Quarto, publico luctu frustra consumpto indicitur publice, vt stato die omnes ii nomine Domini conueniant, sicut Paulus loquitur, ad perficiendam formidabilis ex communicationis actionem. Atque  his tot modis, nempe publica redargutione fact, indicto omnibus luctu, nouo isto mandato de conuentu agendo in ecclesia, denun ciatur multitudini, quid agere consessus ecclesiastieus siue ecclesia velit. Ex pio rum autem omnium prompta obedientia intelligitur eorundem consensus et appre batio. Interea fratri lapso omnibus hisce diebus interpositis conceditur spacium de liberandi, vtrum velit resipiscere. Nec omittunt pii fratres, siue ex senioribus, siue ex grege aliorum, illum exhortari, ad cogitandum de sua salute. Ad eam usque  horam, qua excommunicationis sententia coram tota ecclesia fertur, resipiscenti patere aditus debet, ut libere se ecclesiae sistat, atque agens poenitentiam reconcilietur. Quinto, ubi tandem dies ille constitutus illuxit, conuenit statim in loco tota ecclesia ad excon  \pend
\vspace{0.5cm}\noindent
\vspace{0.2cm}\rule{1cm}{0.2pt}\\ 
\hspace{0.2cm}\textit{mg}
\footnotesize 2. Publice an gui a bet is, qui publie. peceaut. 
\normalsize| 
\hspace{0.2cm}\textit{mg}
\footnotesize 8. Extremum remedium in membri putri di resectione ponitur. 
\normalsize| 
\hspace{0.2cm}\textit{mg}
\footnotesize Cum luctu Ecclesiae fiert debet exconmuricatio. 
\normalsize| 
\hspace{0.2cm}\textit{mg}
\footnotesize 1. Cor. 5. 2. Cor. 12. 
\normalsize| 
\hspace{0.2cm}\textit{mg}
\footnotesize Caisseob quas ecclelia lugere de bet amputa tionem mem gri. 
\normalsize| 
\hspace{0.2cm}\textit{mg}
\footnotesize A raione membrorum; condolentite, 2. Zelus pro gloria Dei. 3. Peccati magnitudo. 4. Causa, Jeueritas iudicii ecclestastici. 5. Causa, Peccata pro pria deplorandaindie nis. venie deprecatio. 4. Dics indicendus excon municationi. Spaciumt resipiscentiae. 5.Conuentus ecclesiae. 
\normalsize| 
\section*{IN I. EPIST. D. PAVLI. }
\marginpar{[ p.72 ]}\pstart municationis (nili statim resipiscentiae dederit significationem frater lapsus) iudicium peragendum. Itaque  absolutis omnibus, quae alias fieri in sacro coetu solent, tria ordine perfieiuntur. Primon causas actionis institutae aliquis e ministris coram multitudine bre uiter exponit, mox autem precatio ad Deum fit, qua Deiopem atque praesidium postulant, vt ipso totam actionem gubernante, quod iustum est exequantur, et quod sta tuent, ratum habere ipse velit. Hoc pacto declarant, se ex praescripto Apostoli esse in nomine Domini nostri lesu Christi congregatos, et velle in causa instituta cum eiusdem potestare progredi. Secundo, commemorat aliquis e ministris peccatum hominis lapsi, diligentiam ab Ecelesia adhibitam pro illius emendatione, pertinaciam illius in malo, demonstrat omnia, quae scripttira de his rebus praecepit, esse tentata, ideoque  mandato et authoritate ipsius Dei tradi peccatorem satanae, atque  a coetu fidelium separari publi ce denunciat. Hoc nimirum est cum potestate Domini nostri lesu Christi hominem sa tanae tradere. Tertio, monet atque  obtestatur omnes, vt fratrem illum vitent, quo pudefa ctus occasionem habeat ce gitandi de poenitentia. Atque his demum partibus viuersum excommunicationis negocium absoluitur. Vitare autem vnulquisque  pius excom municatum debet hoc nimirum modo atque  fine. Primum, ne illum in malo audaciorem reddat confirmetque ́: dehinc, ne consentire eius pertinaciae vlla ratione videatur, atque  ita alieni criminis particeps fieri: tum, ne quis talis hominis consuetudine vitium contrahat. Denique  ne illum a moeroris atque  tristitiae proposito, in quo persistere ipsum oportet, donec emendetur, abducat. Sed neque  permittitur excommunicato, commu nibus interesse precationibus, aut Dominicam adire mensam, ne videlicet ecclesia existimetur cum illo sentire, aut flagitium eius probare:ne item sacri coetus vilescat au toritas. Huc pertinent illa praecepta Math. 18. Sit tibi velut ethnicus et publicanus. I.ad corinth.5. Ne commisceamini cum illts, acne cibum quidem cum eiusmodi capia tis. Ad Ephess. Ne sitis consortes, συμμέτογοι, illorum. Ad Tit.3. Sectarum autorem hominem fuge. ln secunda Epistola Ioannis: ne recipiatis eum in domum:nec Aue ei dixeritis. Qui enim dicit illi Aue, communicat operibus eius malis. Attamen adire, salutare, alloqui eum, vt admoneas, doceas, corrigas, inducas ad poenitentiam, cuiuis pio integrum est. Quinetiam excommunicato licebit, vt emat et quaerat sibi necessaria, atque id genus ciuilia negocia priuata exequatur. Adire insuper sacrum coetum, vbi verbi Dei audiat enarrationem, illi permittitur: sicut etiam Ethnicis et publicanis: quam doquidem spes est, posse ea ibi percipi, per quae flectantur et emendentur.2. ad Thes. 3. Ne commercium habeatis cum illo, vt pudore suffundatur: neque velut inimicum habeatis, sed admonete vt fratrem. Sed haec ipsa, si excommunicationis considerentur fines, intelligi clarius possunt.  \pend\pstart Agamus igitur nunc de finsbus, ob quos aliquis satanae traditur siue excommunica tur. Primus sinis est, vt remoueatur ab ecclesia publicum offendiculum, ne videlicet male audiat flagitiosos impunes ferendo, ac iudicent aduersarii, omnes qui in ea congres gatione sunt, flagitiosis illis nihilo esse meliores. Ad Ephes.5. Ne commercium habueritis cum operibus infrugiferis tenebrarum: quin ea potius arguite.1.ad Corinth.5. Etiam profligate istum qui malus est, ex vobis ipsis. Secundus finis. Ne caeteri eius commercio reddantur deteriores. 1.ad Corinth.5. An nescitis, quod paululum fermenti totam conspersionem fermentat? Expurgate itaque  vetus fermentum. Tertius finis, poenitentia, emendatio, atque  salus spiritus eius qui peccarat.1.ad Corinth.5.Tradatur satanae ad interitum carnis, quo spiritus saluus sitin die Domini lesu2 ad Thessal.3. Admoncte il illum vt fratrem. Ergo excommunicatio ad vtilitatem, cum totius ecclesiae, tum corum qui lapsi sunt: denique  ad remouenda offendicula, item ad auertendas calumnias hostium ecclesiae, diuinstus instituta est. Cumque tales sint eius fines, fieri non potest quin vbi illa recte exercetur, frugifera sequantur effecta finibus respondentia. Quod ad effecta attinet in ipso peccatore, sciendum quod uerba illa Apostoli, quibus ait, tra di hominem satanae, ad interitum carnis: Chrysostomus quidem homil.15.in 1. Epist. ad Corinth. et qui ipsum sequuntur, interpretatur, quasi olim ecclesia cum potestate Christi concesserit satanae, et quasi paedagogo castigatori iniunxerit, ut eum, qui audi\pend
\vspace{0.5cm}\noindent
\vspace{0.2cm}\rule{1cm}{0.2pt}\\ 
\hspace{0.2cm}\textit{mg}
\footnotesize Triamfaert coeuperus? da. 
\normalsize| 
\hspace{0.2cm}\textit{mg}
\footnotesize 1. Causaex ponenda a ctionis ustitute. 
\normalsize| 
\hspace{0.2cm}\textit{mg}
\footnotesize 2. Peccatiex positio, in cu Ius emendatt, one eccclesia nihit fecitr liquum. 
\normalsize| 
\hspace{0.2cm}\textit{mg}
\footnotesize 8. Exbortg tio ecclesid de uitando ex communiea to. 
\normalsize| 
\hspace{0.2cm}\textit{mg}
\footnotesize Quomodo et quamob rem excom municatus 
\normalsize| 
\hspace{0.2cm}\textit{mg}
\footnotesize debeat uitar Sacrorum in terdictio. 
\normalsize| 
\hspace{0.2cm}\textit{mg}
\footnotesize Vsus rerun cuilumm excommunica
\normalsize| 
\hspace{0.2cm}\textit{mg}
\footnotesize tis relinqui tur integer. 
\normalsize| 
\hspace{0.2cm}\textit{mg}
\footnotesize 1. Finss ex communica tionis est o fendiculi declinatio. 
\normalsize| 
\hspace{0.2cm}\textit{mg}
\footnotesize a. Fiuis ex communica tionis. Ne alii reddantus deteriores. 
\normalsize| 
\hspace{0.2cm}\textit{mg}
\footnotesize 3. Finis ex. communica tionis salus uidelicet pet cantis. Effecta ext communica tionis. Bipartitain terpretatio traditionis sataue. 2De corpo rali diuexi nont satane. 
\normalsize| 
\section*{AD TIMOTHEVM CAP. I. }
\marginpar{[ p.63 ]}\pstart re ecclesiam recusarat, diuexaret in corpore, morbum videlicet, vel aliud malum (exe plo Hiob, quem Deus voluit eo pacto probars) quo caro, quae petulantior fuerat, domaretur, immittendo. Alii vero per traditionem satanae intelligunt simpliciter declarationem separationis a Christo et a coetu fidelium, ad tempus, ob quam peccator vsque  adeo angeretur vt videns se contemni et vitari ab omnibus, absorberetur propemodum prae moerore: quo pacto caro illius humillata afflictaquequodammodo interiret. Vtraque  autem interpretatio, si primae ecclesiae statum consideremus, locum merito habebit. Etenim consentaneum vero est, quod quo tempore expediebat veritatem doctrinae Euangelicae et ministerii ecclesiastici dignitatem comprobari miraculis, nomnunquam corpore visibiliter sint diuexati, quos ecclesia Dei propter peccata satanae tradidit. Haec conueniebat sane tunc multas ob caussas ita fieri, Caeterùm ex quo tempore tam doctrina, quam ministerlum ecclesiae satis est confirmatum, nequaquam pereque vt ante, necessarium est, eos qui excommunicantur, in corpore a satana cruciari. Magnum est, si pudore, vt ad Thessalonicenses dicitur, suffundantur, atque ita emendentur:ideoque erit excommunicatio legitime facta, iuxta metuenda et efficax, atque si externa afflictio in corpore accideret. Rata vsque  manet Christi sententia et promissio dicta ad ecelesiam, sontes iustis de caussis et seruato ordine excommunicantem, Matt.18. Amen dico vobis, quaecunque  alligaueritis super terram, erunt ligata in coelo: et quaecunque  solueritis super terram, erunt soluta in coelo. Neque  vero eos, qui excommunicati ac traditi satanae sunt, existimabimus sic penitus exclusos a regno Christi, quasi iam simpliciter damnati, nulla spe salutis relicta, essent: sed statuemus esse admodum aegrotos, quibus tamen medicamem adhuc porriga, et quos seruari posse in vita confidamus. Ne habeatis, inquit Apostolus, velut inimicum, sed admonete vt fratrem: 2. Thessal. Quapropter his, qui aliquamdiu excommunieati fuerunt, si modo egerunt poenitentiam, reditus patet ad ecclesiam: atque  sano corpori membra sanata rursus copulantur.ltaqueApostolus 2. ad Corinth.2. Scripsi vobis, inquit, non vt moerore afficeremini, mox addit : Sufficit istiusmodi homini increpatio haec, quae facta est a pluribus  adeo, vt e diuerso magis condonare debeatis, et consolars, ne quomodo fiat, vt immodico dolore absorbea tur huiusmodi. Quapropter obsecro vos efficite, vt valeat in illum charitas. Quando igitur in homine lapso dolor, seria poenitentia, vitae emendatio, apparent, suadet charitas noxam ei condonari, atque  ecclessaeiterum adiungi. Porro apud diuersas ecclesias diuerlum fuilse olim morem lapsos et poenitentes in gratiam recipiendi, ex patrum scriptis licet cognoscere. Vide Tertullian. in liber  de poenitentia, Cyprianum liber  3. Epistol. Epistola 14. et 16. Augustinum in Enchiridio ad Laurentium capite 65. Chrysostomum homil. 4. in 2. Epist. ad Corinthios. Et haec de altera parte doctrinae episcoporum quae est defide.  \pend
\phantomsection
\addcontentsline{toc}{subsection}{\textit{Sequitur de doctrina εὐταξίας, seu quarundim rerum pertinentium ad bonum ordinem in ecclesia. }}
\subsection*{\textit{Sequitur de doctrina εὐταξίας, seu quarundim rerum pertinentium ad bonum ordinem in ecclesia. }}
\endnumbering\beginnumbering\section{CAPVT II}
\phantomsection
\addcontentsline{toc}{subsection}{\textit{\huge\textbf{A}\normalsize Dhortor igitur, vt ante omnia fant deprecationes, obsecrationes, interpellationes, gratiarum actiones pro omnibus hominibus, pro regibus, et omnibus in eminentia constitutis, etc vsque  ad cap.3. }}
\subsection*{\textit{\huge\textbf{A}\normalsize Dhortor igitur, vt ante omnia fant deprecationes, obsecrationes, interpellationes, gratiarum actiones pro omnibus hominibus, pro regibus, et omnibus in eminentia constitutis, etc vsque  ad cap.3. }}\pstart \huge\textbf{S}\normalsize Equitur de doctrina εὐταξίας, seu quarundam rerum pertinentium ad bonum ordinem in ecclesia: quae doctrina semper necessaria est post doctrinam fidei. Etsi enim non sit ad salutem necessaria, quippe apud perfectos ea nequaquam opus est, vtpote qui omnia huiusmodi agunt sponte, neque oportet eos vllis ceremoniis regi aut gubemnari: tamen quia hinc accedit non\pend
\vspace{0.5cm}\noindent
\vspace{0.2cm}\rule{1cm}{0.2pt}\\ 
\hspace{0.2cm}\textit{mg}
\footnotesize De traditione satanae spirituali. 
\normalsize| 
\hspace{0.2cm}\textit{mg}
\footnotesize An corporaliter uexati sint excommunicati. 
\normalsize| 
\hspace{0.2cm}\textit{mg}
\footnotesize Excommunicatio legitina aeque  efficax, ac corporalis affuictio. 
\normalsize| 
\hspace{0.2cm}\textit{mg}
\footnotesize Ande exconmunicatorum salute prorsus sit desperandum. Lapsis et excommunicatis reditus patet ad ecclesuam per patentiam. 
\normalsize| 
\hspace{0.2cm}\textit{mg}
\footnotesize Εὐταξία necessaria est in ecclesia. 
\normalsize| 
\section*{IN I. EPIST. D. PAVLI }
\marginpar{[ p.64 ]}\pstart nihil adiumenti rudibus, et per hanc doctrinam opera quaedam Charitatis promouentur, et seruatur in conuentibus publicis bonus ordo, ideo non est temere contennenda. Ceremoniae et huiusmodi ordinationes habendae sunt in vita, sicut fabro habentur praeparamenta ad aedificandum, quae non in hoc parantur, vt aliquid sint, aut aliquid maneant, sed quod sine iis aedificari non possit. Nam perfecto aedificio tolluntur, et stultus esset, qui semper illa magnifaceret et ampliaret, neglecto interim principali aedificio, id est, omissa doctrina fidei. Vnde eatenus adhibendae sunt ceremoniae, quatenus sunt vtiles, id est, quatenus significant aliquid et figurant, et quatenus sunt necessariae, id est, non possunt commode negligi. Alias enim periculum foret, ne doctrina fidei per doctrinam vrαξιας obscuraretur et laederetur: Sicut pueris opus est et vtile, vt foueantur mulierum et puellarum gremiis. quod tamen periculosum foret paulon prouectioribus: Ita iudicandum de ceremoniis. Et sicut paupertas spiritualis periclitatur in diuitiis, fidelitas in negociis, humilitas in honoribus, abstinentia in conuiuiis, castitas in delitiis: lta doctrina fidei periclitatur interdum in doctrina εὐταξίας. Haec de doctrina εὐταξίας in genere.  \pend\pstart Caeterùm non hic praescribit Apostolus omnia, quae ad bonam ordinationem pertinent, quae in coetu ecclesiastico agi conueniat: sed vnam tantumodo rem vrget, vel ante non satis commendatam, vel parum diligenter obseruatam. Ea est de orando pro omnibus hominibus, potissimumque pro magistratibus. Quae quidem ἐυταξία siue bona constitutio adeo est necessaria in ecclesia, vt nulla magis: Et doctrinae fidei sic adiuncta, vt diuelli aegre queat: et ad charitatem declarandam nulla magis idonea, sic itaque  ait Apostolus: Adhortor, vt ante omnia fiant deprecationes:) Adhortor, inquit, Dignissima res de qua inprimis episcopi admone antur, et cuius prima cura omnibus Christianis esse debet, nimirsi vt oretur. In genere autem oratio dicitur constans et ex intimo affectu alicuius rei necessariae a Deo petitio: Etiamsi alsj malunt, esse mentis ad Deum eleuationem. Dico petitionem, quia oportet in ea agnoscamus aliquid nobis deesse, quod cupimus a Deo impetrare: vt videlicet nostram fragilitatem vel miseriam intelligentes humiliemur contra Pharisaeorum superbiam, et cum semper egeamus, semper ad orandum excitemur. A Deo autem, quia ab illo solo omne bonum descendit, et a nullo alio debemus petere vel expectare. Necessaria res petenda est, quia stulta vota vel inutilia Deus non audit, et petita iuste negat. Constans praeterea esse debet oratio: quia Deus vult saepe rogari, et fidem probare, et saepe rogatus tandem concedit, oportet semper orare et nunquam defatigari. Ex intimo affectu: quia exanimo, non fucate, et ardenti affectu praecipuem fide, loann.4. Veri adoratores in spiritu et veritate orant. Iacobus sine haesitatione et cum fide vult nos orare. Haec vero requiruntur in oratione, quam vnusquisque  pri uatim facit quotidie et vbiubisit locorum. Sed multo magis requiruntur in oratione publica, quae fit ab omnibus in publico coetu et congregatione: quam credi par est semper etiam efficaciorem fore, quando a quamplurimis et concordibus animis fit, Vult autem fieri orationes ante omnia. Id dupliciter accipi potest: vel vt significet, nihil sic necessarium, atque est oratio: ac propterea eam prae quibuscunque  alujs charitatis officiis exercendam: vel vt significet circumstantiam temporis, id est, prima statim diei parte conueniendum ad orationem, priusquam aliquid operis faciamus. Id quod etiam conuenientissimum est, vt primum opus et prima diei portio rebus diuinis dicetur: et quia nihil prius Deo habere debemus, non est quod putemus bene successurum in relsquis rebus, nisi Deum prius orauerimus: quam opinionem commodissimum est teneris et rudibus animis inculcare. Atque hunc morem olim seruatum scriptores testantur: inprimis Tertullianus, apud quem saepe antelucanorum coetuum mentio est. Et Plinius ipse Epistol. libro decimo, ad Traianum Imperatorem idem refert. Testantur eundem morem certi hymni, vt ille Prudentii, Alius diei nuncius et alter, iam lucis orto sydere, etc. Atque ita faciendo nihil detrahebantur diurnis operis. Etanimus sobrius, ieiunus, a quiete corporis mane maxime compositus est ad res diuinas peragendas, magis quam alia diei parte. Tametsi vtnec eleemosyna, nec ieiu\pend
\vspace{0.5cm}\noindent
\vspace{0.2cm}\rule{1cm}{0.2pt}\\ 
\hspace{0.2cm}\textit{mg}
\footnotesize A fimili. 
\normalsize| 
\hspace{0.2cm}\textit{mg}
\footnotesize Ceremoniae quatenus in ecclefia retirendae et coIendae. A ratione aetatis. 
\normalsize| 
\hspace{0.2cm}\textit{mg}
\footnotesize Ordtionis de finitio. 
\normalsize| 
\hspace{0.2cm}\textit{mg}
\footnotesize 1. In oratione homo aguoscit, et confitetur suam ipsius fragilitatem. 
\normalsize| 
\hspace{0.2cm}\textit{mg}
\footnotesize 2.A Deo solo res bonae patendae. 
\normalsize| 
\hspace{0.2cm}\textit{mg}
\footnotesize 3. Res necessariae petendae, uota inutilia praeter mittenda. 
\normalsize| 
\hspace{0.2cm}\textit{mg}
\footnotesize 4. Oratio debet esse constans et ardua. 
\normalsize| 
\hspace{0.2cm}\textit{mg}
\footnotesize 5. Oratio proficisci debet ex animo et corde. lacobus  1. Preces labores et studia hominum debent praecedere. Coetus antelucani. Preces non impediunt labores et operas homium. 
\normalsize| 
\section*{AD TIMOTHEVM CA P. II. }
\marginpar{[ p.55 ]}\pstart nium: ita nec oratio temporibus aut diebus  est alliganda: sicut nec tunc tantum modo, sed praeterea quacunque alia hora orare, oportunum est. Porro, quo magis excitet omnes ad orandum, vtitur congerie, plures orationum species coaceruans: quo nimirum declaretrem esse maxime necessariam et plurimi faciendam, vt pro eadem signifcatione liceat omnia illa accipere, et saepe contunduntur, sicut in exordio Philippensium et alibi Apostolus precationibus includit gratiarum actionem. Nonnullis tamem magis placet distinguere, vt Charitatis officia in orando plenius intelligantur, sic autem legimus: Fiant deprecationes, obsecrationes interpellationes, gratiarum actiones Deprecatio quae Graece δὲησις, significat orationem, qua petimus liberari a malis vrgentibus vel imminentibus: vnde dicimus quotidie, libera nos a malo. Obsecratio, Graece προσευχή, est qua precamur nobis dari bona: quo pertinet illud, Panem nostrum quotidianum da nobis. Interpellatio, ἔντευξις, qua querimur de iniuria nobis illata: quo in genere sunt multi Plalmi, quibus petimus tales in viam redire et resipisscere: quo videtur pertinere nllud. Sanctificetur nomen tuum. Gratiarum actio ευχα- ριστία, cum pro bonis accepris, imo et pro malis bona voluntate Det nobis inflictis, Deo agimus gratias: quo illud, Fiat voluntas tua. Et haec omnia non solum nostra causa facienda sunt, verum etiam propter alios. Deprecandum enim, vt alii quoque a malis liberentur:obsecrandum vt aliis quoque bona contingant:interpellandum pro aliorum quoque aduersariis, vt Deus eos emolliat: gratiae agendae, tùm etiam, cum aliis bene fuerit. Ideo additur, pro omnibus hominibus: Charitas siquidem Christiana nullam personam debet ex cludere. Nec solum pro nobis orare debemus, verùm pro aliis quoque. Nec pro illis tantûm, qui amici sunt, qui bene volunt, quiqueconiuncti sunt religione: verum etiam pro inimicis et pro aduersariis nostrae fidei. Pro omnibus, inquit, ἐμφατικῶς. Deus facit solem suum oriri super malos perinde ac bonos. Et si non nisi amicos diligimus, nihilo meliores sumus gentibus incredulis. Christus orauit pro erucifigentibus, et Stephanus pro lapidantibus: et quondie oramus, Remitte nobis debita sicut nos remittimus debitoribus, et iis qui nos offendunt. Et pro bonis ac coniunctis nobis in fide orandum est, vt tales perseuerent, et quantum fieri potest, in melius proficiant: pro malis orandum, vt conuertantur, resipiscant. et tandem Deus per eos et in eis glorificetur. Pro regibus et omnibus in eminentia constitutis, Cum pro omnibus orandum esset constet, at multo iustiisime orandum pro magistra tibus. Atque id eo tempore voluit Apostolus fieri quam dillgentilsime in ecclesia cûm magistratus passim erant ethnici, religionis Christianae hostes: sed ideo multon maxime erat pro omnibus orandum, quia illorum nutu atque arbitrio multa fieri poterant per quae Euangelii propagatio velimpediretur, vel propagaretur. Vt autem animi eorum immutarentur ad fauendum potius Euangelio, quam ad persequendum indurarentur, aequum erat orationibus id a Deo flagitare. Quae doctrina etiam a Sanctissimis viris olim diligenter fuit tradita. Sic enim Hieremias capite 29. iubet orari pro salute regis et regni Babylonis, vnde apparet Pauli verba quaedam huc translata. Quaerite, ait ille, pacem ciuitatis, in quam migrare vos feci, et orate pro ea Dominum, quia in pace eius erit pax vestra. Similiter Esdras pro Cyro rege Persarum, Baruch. cap.1. pro salute Nabuchodonosor orari voluit. Tertullianus Apologetici sus capite 30. hanc Christianorum pietatem celebrat. Nos oramus, inquit, pro omnibus Imperatoribus, vitam illis prolixam, imperium securum, domum tutam, exercitus fortes, senatum fidelem, populum probum, orbem quietum, et quaecunque hominis et Caesaris vota sunt. Sunt autem multae caussae, quae nos moueant ad orandum pro magistratibus, quas exposuimus in libello nostro de honorandis magistratibus. Apostolus etiam hic argumenta quaedam suasoria attexuit. Et placidam ac quietam vitam agamus cum omnt pietate et honestate.) A caussa finali siue ab vtili. Ideo orandum est pro magistratibus, quia per eos potest nobis contingere,ut placide ac quiete agamus, aclibere Deum laudemus sine ullis molestiis. Certe si nunquam datur piis quiete agere, necessario omittunt multa bona opera:imo  \pend
\vspace{0.5cm}\noindent
\vspace{0.2cm}\rule{1cm}{0.2pt}\\ 
\hspace{0.2cm}\textit{mg}
\footnotesize semper et omni loco orandum est. Congeries συναθροισμός. 
\normalsize| 
\hspace{0.2cm}\textit{mg}
\footnotesize Deprecatio dέησις. Objecratio προσευχή. Interpellatio ἔντευξις. 
\normalsize| 
\hspace{0.2cm}\textit{mg}
\footnotesize Gratiaram action ευχαριστία. 
\normalsize| 
\hspace{0.2cm}\textit{mg}
\footnotesize Emphasis. 
\normalsize| 
\hspace{0.2cm}\textit{mg}
\footnotesize Pro inimicis etiam orandum. Math.5. Luc. 6. Matth.27. Actor.7. 
\normalsize| 
\hspace{0.2cm}\textit{mg}
\footnotesize Pro magistratu orandum. 
\normalsize| 
\hspace{0.2cm}\textit{mg}
\footnotesize Vota Christianorum pro Imperatoribus. 
\normalsize| 
\hspace{0.2cm}\textit{mg}
\footnotesize ab utili. 
\normalsize| 
\section*{IN I. EPIST. D. PAVLI }
\marginpar{[ p.66 ]}\pstart quidam etiam malorum magnitudine perculsi, Euangelio malunt renunciare, quam diutius propter veritatem affllgi. Qua ex re cùm obscuratur gloria Dei, merito est orandum pro magistratibus. Et credibile est etiam ethnicorum principum iras remissas, imminuisseque  nonnihil suam truculentiam, vbi resciscerent Christianos pro ipsorum salute orasse.lta nimirum tum oraruat Christiani pro impüs principibus, et quietem sibi a Deo impetrarunt, et beneuolentiam saltem principum conciliarunt aliquam. Proinde non eo line optanda est placida et quieta vita, vt luxu et licentia perdamur, et quasi eiecto freno excurramus in quaeuis vitia, sed vt agamus cum omni pietate et honestate, id est, ex nostra libertate incrementum accipiat vera pietas et promoueatur itidem uitae honestas. Emphasis in nota uniuersali. Nomine pietatis intellige quicquid ad cultum Dei, ad ueram religionem pertinet, siue per dogmata, siue per externa opera:ut liberi conuentus ad orandum, libera praedicatio uerbi, libera exercitatio operum Charitatis, quae proprie pertinent ad propagationem Euangelii et ad incrementum gloriae Dei. Nomine honestatis comprehenduntur quae commeudant nos etiam ciuiliter apud homines, et per quae homines aedificantur, instruunturquein bonis exemplis, et occasionem habent bene et sentiendi et loquendi de nobis. Ad hunc ergo finem, ut uidelicet liberem perficiamus ea, quae pertinent ad gloriam Dei et aedificationem proximorum, optanda est placida ac quieta uita. Quod si hic finis desit, certe metuendum, ne quies et orium potius caussa sint malorum, quam bonorum. Certe uix unquam ecclesia fuit in meliori statu, nunquam credentes fuerunt tam puri, simplices, integri, incorrupti, quam tùm, atque  cum ryranni saeuirent. Contra ubi ecclesiae coeperunt otiari, quiescere, ditari, dominari: diuitiae et quies extinxerunt penitus fidem et charitatem:filia deuorauit matrem, inquit Bernardus.  \pend
\phantomsection
\addcontentsline{toc}{subsection}{\textit{Nam hoc bonum est et acceptum coram seruatore nostro Deo, qui cunctos homines vult saluos fieri et ad agnitionem veritatis venire. }}
\subsection*{\textit{Nam hoc bonum est et acceptum coram seruatore nostro Deo, qui cunctos homines vult saluos fieri et ad agnitionem veritatis venire. }}\pstart Aliud argumentum a Pio. Ait orare pro omnibus, inprimisque ́, pro magistratibus, bonum esse et acceptum coram Deo. Ac videtur subesse occupatio. Nam vbi fortassia quorundam sragilitas aut coecitas iudicasset absurdum orare pro iis, qui persequuntur pios, et impediunt gloriam Dei promoueri : respondet, imo id pium esse et Deo in primis gratum. Quando igitur audimus id per se esse bonum, et acceptum Deo, nequaquam debemus hic esse cessatores. Et vere id bonum ac pium dici potest, quia inde sequuntur plerunque  multa bona, et ea inprimis per quae syncera pietas promouetur. Cum addit, Qui cunctos homines vult saluos sieri, et ad cognitionem veritatis venire, Latet argumentum ab exemplo: quasi dicat: Deus vult omnes homines saluos fieri, quicunque illi sint: igitur et tu debes idem velle:ac proinde pro omnibus, etiam impiis, orare. Est ergo hic exemplum insignis bonitatis diuinae, quam conuenit et nos tùm affectu tum etiam diligentia omni imitari. Possumus etiam hic discere, quid maxime debeamus impiis magistratibus comprecari, videlicet vt ad cognitionem veritatis perueniant. Veritas autem hic significat Christum, qui est via, vita, veritas. Sed hic quaestio oritur. Si Deus velit omnes homines saluos fieri, qui fit, quod non omnes saluantur ? Respondet paucis verbis Ambrosius. In omni locutione sensus est, conditio latet: Vnde secundae Petri primo dicitur, omnis scriptura indiget interpretatione. Vult Deus omnes saluos fieri, sed si accedant ad eum. Non enim sic vult, vt nolentes saluentur, sed vult illos saluari, si et ipsi velint. Nam vtique qui legem omnibus dedit, nullum excipit a salute. Numquid item medicus idcirco proponit in pubhco, vt omnes ostendat se velle seruare: si tamen ab aegris requiratur ? Non est enim vera salus, si nolenti tribuatur, nec gaudere potest in percepta salute, qui inuitus, si  \pend
\vspace{0.5cm}\noindent
\vspace{0.2cm}\rule{1cm}{0.2pt}\\ 
\hspace{0.2cm}\textit{mg}
\footnotesize A pio. Occupatio. 
\normalsize| 
\hspace{0.2cm}\textit{mg}
\footnotesize Ab exemploA maiori ad minus. 
\normalsize| 
\hspace{0.2cm}\textit{mg}
\footnotesize loan. 14 Questio. Cur non omnes saluentur. Responsio ex Ambrosio. 
\normalsize| 
\hspace{0.2cm}\textit{mg}
\footnotesize a simili. 
\normalsize| 
\section*{AD TIMOTHEVM CAP. II. }
\marginpar{[ p.67 ]}\pstart E tamen fieri potest, accepit medicinam, ut non dicam, quod medicinae effectum habere non potest, nisi ad illam aeger animum accommodauerit, quia haec medicina non est corporalis sed spiritualis, quae neque dubiis proficit, neque inuidis. Fides est enim quae dat salutem, quam nisi mens tora susceperit uoluntate, non solum nihil proderit, sed et oberit. Fidei etenim gratia hanc habet potestatem, ut deuotis sibi diuinam infundat medelam, indeuotis uero conferat morbum, per quem totus homo intaereat. lta Ambrosius. Nec obest, quod Roman.9. dicitur: Deus cuius uult miseretur, et quem uult indurat. Sicut saluatio hominis non est humani meriti. Ita perditio siue induratio non est simpliciter diuinae uoluntatis. Deus peccatorum non est autor. Sed indurare sic quodammodo dicitur: quia eos ea facta perpetraturos praeuidit, per quae et ob quae merito indurari eos permittit. Quare sic statuendum est in tota hac causa de hoc loco et similibus: quod, ut nostris uiribus nequaquam tribuendum est, quod debetur soli gratiae et potentiae ipsius Dei: lra e diuerso diuinae bonitati nequaquam est imputandum, quod nostra contingit propria malitia. A nobis igitur excaecatio est et indusratio, a Deo gratia et salus. Excaecatio contingit nobis per Dei iustitiam, quae merito et iure punit nostra peccata: saluatio contingit per Dei bonitatem et misericordiam, quam nullis nostris meritis promeremur: Sicque  Deus excaecare et illuminare, misereri et indurare, sed respectu diuerso dicitur. Nonnulli clarius uolentes explicare et simplicius, dicunt, in hac propositione, Deus uult cunctos homines saluos fieri, in nota illa uniuersali esse Hebraismum, quo, Omnes, accipitur, pro Omnis generis:ut sit sensus, Deum uelle omnis generis homines saluos fieri, siue gentiles, fiue ludaeos, siue magistratus, siue subditos. Sic quoque accipiendum, ubi mox dicitur, quod Christus dedit semetipsum redemptionem pro omnibus. Et facit haec ratio non parum ad commendationem diuinae bonitatis, utpote cùm ostenditur, Deum non esle acceptorem personarum, neque agnoser uelle tantummodo a ludaeis, aut in sola ludaea, uerùm ab omnibus et ubique terrarum.  \pend
\phantomsection
\addcontentsline{toc}{subsection}{\textit{Vnus enim Deus, unus etiam reconciliator Dei et hominum, homo Christus lesus, qui dedit semetipsum pretium redemptionis pro omnibus, vt esset testimonium temporibus suis. }}
\subsection*{\textit{Vnus enim Deus, unus etiam reconciliator Dei et hominum, homo Christus lesus, qui dedit semetipsum pretium redemptionis pro omnibus, vt esset testimonium temporibus suis. }}\pstart Plenior probatio adiecta, quod Deus uelit omnes homines saluos sieri : et adhuc persequitur argumentum ab Exemplo, Vnus est, inquit, Deus omnium, qui non respicit personam. Ergo uos quoque respicere personam non debetis, sed pro omnibns, etiam aduersariis orabitis. Praeterea unus est reconciliator Dei et hominum, uidelicet Christus, qui dignatus est semetipsum dare in prerium redemptionis pro omnis generis hominibus. Ergo uos quoque saltem orando intercedere pro quibuscunque aequum est. Proinde dicit, Vnus Deus ommum, etiamsi ab omnibus nequaquam agnosceretur. Verùm cum unus Deus omnes condiderit, idemque  omnes uoluerit in gratiam recipere, et ad agnitionem sui nominis uenire, nullos spernetis uel abiicietis: merito omnium et unus Deus dicitur. Vnde Roman.3.An ludaeorum Deus tantùm? Annón et gentium? Certe et gentium. Quandoquidem unus Deus, qui iustificabst circumcisionem ex fide, et praeputium per fidem. Ergo quando unus Deus omnium praedicatur, mirifice hic nobis commendatur charitas: ac iubemur inuicem, imo et inimicos et tyrannos diligere, ac pro eorum salute orare, utpote quorum idem est nobiscum Deus, idem conditor et communis quasi pater, adeoque  ipsi quasi communes sunt nostri fratres. Quos Deus ipse non despicit, iniquum est a te contemni.  \pend\pstart Vnus etiam reconciliator Dei et hominum, homo Christus Iesus.) Cum diceret, unum esse Deum, non exclusit spiritum sanctum uel filium: quippe substantia unum sunt, qui personis distinguntur. Sicut itaque unus est Deus atque omnium, eo mouert debemus, ut neminem aspernemur. lta unus quoque est omnium reconcliator, unus  \pend
\vspace{0.5cm}\noindent
\vspace{0.2cm}\rule{1cm}{0.2pt}\\ 
\hspace{0.2cm}\textit{mg}
\footnotesize Deus quomodo indurare dicatur. 
\normalsize| 
\hspace{0.2cm}\textit{mg}
\footnotesize A contrario. A nobis est excecatio et a Deo gratia et salus.. 
\normalsize| 
\hspace{0.2cm}\textit{mg}
\footnotesize Senentia fimplicior et dilucidior. 
\normalsize| 
\hspace{0.2cm}\textit{mg}
\footnotesize Actor.16. Coloso.3. Galat.3. 
\normalsize| 
\hspace{0.2cm}\textit{mg}
\footnotesize Vuus Deus 
\normalsize| 
\hspace{0.2cm}\textit{mg}
\footnotesize Vnus reconciliator. 
\normalsize| 
\hspace{0.2cm}\textit{mg}
\footnotesize Dei dilectio nos ad mutuam hortatur dilectionem. 
\normalsize| 
\hspace{0.2cm}\textit{mg}
\footnotesize Tres personae, una substantia, unus Deus. 
\normalsize| 
\section*{IN I. EPIST. D. PAVLI }
\marginpar{[ p.68 ]}\pstart est qui omnes cum Deo patre reducit in gratiam, acpropterea non est quod quilquam caeteros vilipendat, quasi non habeant eundem reconciliatorem omni tempore ad subue niendum ipsis paratum. Proprie autem reconciliator iste seu mediator (ita enim graece μεσίτης) Christus est, qui solus ad munus hoc fuit idoneus, et se paratum obtulit. Dissidium erat inter duas partes, videlicet inter Deum patrem, et homines, qui patrem peccando legesqueeius transgrediendo offenderant: ita ut iam Deus pater per suam iustitiam merieo vniuersum genus humanum erat perditurus. Nec interim homines vlla ratione iratum patrem potuissent placare. Christus itaque  solus reconciliandae multitudinis patrocinium ad se recepit, solus uoluit et potuit. Proinde ut commodius id faceret, sicut diuersa erat partium conditio: ita placuit ei, ut se vtrique  parti accommodaret, vtriusque  partis personam induere: sicque  mediator inter Deum et hominem dum vult fieri, sicut ipse semper fuit Deus, ita homo quoq in tempore factus est, aptissimeque  hominem cum Deo reconciliauit, quando ipse ut Deus erat, apud Deum patrem quiduris potuit efficere, et ut homo potuit hominum necessitatem et miseriam agnoscere: et sic utrique  parti visus est aequus arbiter, quando vtriusque  partis naturam in se habuit. Ille ergo homo Christus lesus, idemqueDeus dependit patri Deo quicquid pro hominum salute erat dependendum. Vnde sequitur,  \pend\pstart Dedit semetipsum pretium remdemptionis pro omnibus.) Caussa dissidii erat peceatum. Peccatum autem quoniam non nsi sanguine potest expiari, et quidem non nisi preriosissimo sanguine, non taurorum aut hircorum, dignatus est Christus suum ipsius sanguinem fundere, imo totum sese dare in pretium redemptionis, et quidem pro omnibus, ita ut id pretium esset sufficiens ad satisfaciendum irato patri, pro peccatis non nostris tantum, vt dicitur 1. Iohan.2. verum etiam totius mundi, et nullum uoluerit a satisfactionis illius participatione excludi: nisi qui semetipsum suis peccatis et perfidia excludit. Singula autem uerba habent suam Emphasin, illa praesertim semetipsum et pro omnibus. EstqueMetaphora in uoce, pretium redemptionis. Solent enim captiui redimi: nam et graece legitur ἀντιλυτρον, ubi nos legimus pretium redemptionis. Significat proprie id uocabulum pretium, quo redimuntur captiui ab hostibus, eamque  commutationem, qua capite caput, et uita redimitur uita. Proinde est egregia laus Christi, quaque  insignis illius charitas erga humanum genus quod perierat, celebratur, cum ipse uoluerit agere reconciliatorem, et quidem impendendo semetipsum pretium redemptionis, et pro cuiuscunque  conditionis hominibus. lohan.1o. Maiorem charitatem nemo potest exhibere, quam ut animam suam ponat quis pro amicis suis. Roman.5. Commendat suam charitatem erga nos Deus, quod cum adhuc essemus peccatores, Christus pro nobis mortuus est. Ac proinde nos hoc exemplo Christi ad charitatem moueri debemus. Praeterea ingens consolatio conscientiis nostris esse debet, quando nobis sic in scripturis proponitur Christus homo aut reconciliator: ut nunquam profecto sit desperandum, cum hominem habeamus Christum, qui uelit intelligere nostras miserias, et quem audeas compellare: (nam propterea Apostolus eum hic hominem diserte appellat:) et eundem Deum, qui possit peccata nostra    tollere. Huc pertinent illa Roman.5. Iustificati perfidem pacem habemus ad Deum,    per Dominum nostrum lesum Christum. Roman. 8. Christus est qui mortuus est, et suscitatus est, et sedet ad dexteram Dei, et intercedit pro nobis. 1.Iohan.2. Filioli mei    haec scribo uobis ne peccetis: et si quis peccauerit, aduocatum habemus apud patrem, lesum Christum. Et ipse est propitiatio pro peccatis nostris. Hebr.9. Ob id    noui testamenti conciliator est, ut morte intercedente ad redemptionem earum praeuaricationum, quae fuerant sub priori testamento, si qui uocati sunt, promissionem accipiant aeternae haereditatis. Hebr.4. Christus dicitur thronus gratiae. Romanor.3. Hunc proposuit Deus reconciliatorem per fidem, interueniente ipsius sanguine. In hos ergo locos et similes passim obuios decet perturbatas conscientias intueri. Postremo discimus unum duntaxat reconciliatorem agnoscere Christum, ne plures nobis reconeiliatores statuamus. Verum quidem est, orant sancti in terris positi pro se mutuo et exaudiuntur, sed propter illum vnum reconcillatorem Christum, qui  \pend
\vspace{0.5cm}\noindent
\vspace{0.2cm}\rule{1cm}{0.2pt}\\ 
\hspace{0.2cm}\textit{mg}
\footnotesize Vnus media tor Dei et hominum Christus Iesus. Dissidium inter Deum et homines. 
\normalsize| 
\hspace{0.2cm}\textit{mg}
\footnotesize Christus utriusque , partis personam suscipit. Christus Deus. Christus Homo. Peccatum caussa dissidii mter Deum et homi nem. 
\normalsize| 
\hspace{0.2cm}\textit{mg}
\footnotesize Hebr.9. 
\normalsize| 
\hspace{0.2cm}\textit{mg}
\footnotesize Peccata hominum non nisi sanguine Christi possunt expiari. Metaphora a captiuis. ἀντίλυτρον quid. 
\normalsize| 
\hspace{0.2cm}\textit{mg}
\footnotesize Christi charitas erga homines ineffebilis. 
\normalsize| 
\hspace{0.2cm}\textit{mg}
\footnotesize Ioan.10. Rom.5. 
\normalsize| 
\hspace{0.2cm}\textit{mg}
\footnotesize Vsus charitatis Christi in nobis. 
\normalsize| 
\hspace{0.2cm}\textit{mg}
\footnotesize Ron.8. 
\normalsize| 
\section*{AD TIMOTHEVM CAP. II. }
\marginpar{[ p.69 ]}\pstart ipssemet intercedit pro nobis, exaudiuntur. Iniuriam facit Christo, qui alium sibi mediatorem fingit. Sequitur:  \pend\pstart Vt esset testimonium temporibus suis.) q' d. Quod Christus dederit se in pretium redemptionis pro omnibus, non pro ludaeis tantùm, sed pro cuiuscunque  conditionis hominibus: id Deus testatum fieri uoluit toti mundo, omnibusqueretro temporibus, vt homines intuentes in tantam charitatem Dei, similiter ad Charitatis officia inter sese praestanda excitarentur. Addit, sibi iniunctum id munus, testandi uidelicet mundo, quod ita esset.  \pend
\phantomsection
\addcontentsline{toc}{subsection}{\textit{In quo positus sum ego praeco et Apostolus, veritatem dico in Christo, non memtior, doctor gentium cum fide et veritate. }}
\subsection*{\textit{In quo positus sum ego praeco et Apostolus, veritatem dico in Christo, non memtior, doctor gentium cum fide et veritate. }}\pstart Confirmatio praecedentis assertionis de vno reconciliatore Christo, a ratione officii sui, et sura ipsius fide integritateque  quam iuramento probat. Ostendit proprium suum officium esse testari et palam facere eximiam Dei erga nos charitatem, praecipueque  quod Christus sit vnus reconciliator Dei et hominum, qui semetipsum obtulerit in pretium redemptionis pro omnibus. Pronide hic notatur, hanc propriam esse doctrinam Apostolicam de uno reconciliatore Christo. Quare grauiseime peccant, et a doctrina apostolica longissime discedunt, qui alios reconciliatores, uel alias reconciliandi formas meditantur.  \pend\pstart Veritatem dico in Christo, non mentior.) Confirmatio est ueluti iuramento adhibito et interpretatio per contrarium. Doctorem gentium se uocat, simulut hanc doctrinam de reconciliatore Christo gentibus magis commendet, simul ut insinuet plenius Christum pro omnibus satisfecisse, non tantùm pro ludaeis, sed etiam pro gentibus: neque fastu ullo uel arrogantia, sed ministerium suum significat.  \pend\pstart Cum siue et vertate ) A sua side et integritate. Cum fide, inquit, non cum Syllogismis aut uana eloquentia, non ut fallam aut decipiam. Discant proinde hinc ministri uerbi, quibus rationibus doctrinam suam confirmare debeant, dentqueoperam, ut uere his uerbis uti queant, uidelicet, se ueritatem dicere in Christo, se docere in fide et ueritate. Hactenus, quod sit orandum pro omnibus.  \pend
\phantomsection
\addcontentsline{toc}{subsection}{\textit{Volo igitur orare viros, in omni loco, sustollentes puras manus absq ira et disceptatione. }}
\subsection*{\textit{Volo igitur orare viros, in omni loco, sustollentes puras manus absq ira et disceptatione. }}\pstart Addit nonnulla de orandi modo. Orare autem uult uiros omni loco. Nam orandum quidem tunc maxime, cum in coetum conueniunt, ubi orari, docert uerbum solet: sed non ibi tantùm, uerumetiam ubique  et quocunque  tempore. Itaque  excitat uiros ad orandi diligentiam et frequentiam: et praeterea superstitiones omnes de loco aut tempore orandi resecat. Orandum est omni loco, non solùm ubi omnes conueniunt publice, sed etiam domi. Est domus Dei domus orationis: sed nihilominus bene etiam erat, qui, ut habetur Matth.6. Petit secretum cubiculi sui, et ibi orat, humanam deuitans gloriam. Est etiam orandum foris, quia, ut Matth. 5. lucere debet lux nostra coram hominibus, ut uideant opera nostra bona, et glorificent patrem qui est in coelis: Sed sic orandum in propatulo, ut absit Pharisaica ostentatio: sie in cubieulo, ut animus foris non uagetur, sed ad solum Deum feratur. Superstitiosum autem erat, quod ut habet loan, 4. ludaei dicebant, orandum praecipue in Templo Hierosolymitano, Samaritani uero in montibus. Christus ubique  uult orari, modo in spiritu et ueritate oretur. Adsit spiritus et ueritas, nec multum refert, quo sis loco. Non locus orationem commendat, sed affectus fidei, id quod etiam uerba sequentia declarant.  \pend\pstart Sustollentes puras manus, abque  ira et disceptatione.) Metaphoris sumptis et a sacris deorum, in quibus omnia exterius munda esse debebant, et ab habitu corporis ueteribus inter orandum usitato, ostendit quomodo animos nostros ad orandum componi deceat. Vult sustolli manus: sic enim ueteres manibus eleuatis in altum orabant, id quod et picturae quaedam testantur. Et Exod.17. manus leuante et orante Mose uincit lsrael Amalechitas. Significatur ergo per manuum eleuationem eleuatio mentis: sicut et manibus  \pend
\vspace{0.5cm}\noindent
\vspace{0.2cm}\rule{1cm}{0.2pt}\\ 
\hspace{0.2cm}\textit{mg}
\footnotesize A ratione officii. 
\normalsize| 
\hspace{0.2cm}\textit{mg}
\footnotesize A contrario. Paulus doctor gentium. 
\normalsize| 
\hspace{0.2cm}\textit{mg}
\footnotesize Vbi orandum sit. 
\normalsize| 
\hspace{0.2cm}\textit{mg}
\footnotesize Onni loco orandum. 
\normalsize| 
\hspace{0.2cm}\textit{mg}
\footnotesize Quomode et domi et foris sit orandum. 
\normalsize| 
\hspace{0.2cm}\textit{mg}
\footnotesize In spiritu et uertate orandum. Metaphorae. 
\normalsize| 
\section*{IN I. EPIS T. D. PAVLI }
\marginpar{[ p.70 ]}\pstart sursum sublatis coelum demonstramus, Deum significamus. ltaqueut totus in Deum feratur animus per fidem, est necessarium, si exaudiri volumus. Sic Thren3. Leuemus corda nostra cum manibus in coelum. Debent praeterea manus esse purae, id est, nullis sordibus vitiorum inquinatae. Nam cum manibus multa peccata perpetrantur, vt furtum, caedes, rapina, vis illata, pugnae, et c. Cum exiguntur manus purae, exigitur ab omnibus sceleribus immunitas, ita vt Synecdochice per manuum puritatem intelligatur puritas totius hominis, praesertim animi siue cordis. Fide purificantur corda. Sic vero et impuritas manuum damnatur Esa.1.Cum extenderitis manus vestras, auertam oculos meos a vobis, et cum multiplicaueritis orationem, non exaudiam. Manus enim    vestrae sanguine plenae sunt. Sed vt plenius exprimat Apostolus, quid per manuum impuritatem intelligi queat, quasi per interpretationem adiicit, Absque ira et discgtatione. Ira praecipue manus polluit, et per eam omnes affectus vitiantur. Qui irato est animo, plerunque erumpit ad perpetranda scelera. Sed et ideo ira abesse debetab oraturo, quia oratio debet fieri tranquilla et quieta mente. Quies autem non contingit, vbi ira vel odium, vel appetitus vindictae habet locum. Huc pertinet illud Eccles.28. Homo homini seruat iram et a Deo quaerit medelam? Matth. 6. Si non remiseritis errata sua, nec pater vester remittet errata vestra. Item quod quotidie oramus, Remitte nobis debita nostra, sicut et nos remittimus debitoribus nostris. Disceptationem autem interpretantur nonnulli haesitantiam, qua videlicet quis secum disceptat, dubitans, num sit exaudiendus. Oportet ergo hominem oraturum non disceptare secum, non fluctuare: sicut lacobus docuit. Nam qui haesitat, is similis est fluctui maris, qui ventis agitur, et impetu rapitur. Neque existimet homo ille se quicquam accepturum a Domino. Habemus ergo hic non obscura documenta, quomodo nos ad orandum componere debeamus.  \pend
\phantomsection
\addcontentsline{toc}{subsection}{\textit{Consimiliter et mulieres in amictu modesto cum verecundia et castitate ornare semetipsas, non tortis crinibus, aut auro, aut margaritis, aut vestitu sumptuoso, sed quod decet mulieres profitentes pietatem per bona opera. }}
\subsection*{\textit{Consimiliter et mulieres in amictu modesto cum verecundia et castitate ornare semetipsas, non tortis crinibus, aut auro, aut margaritis, aut vestitu sumptuoso, sed quod decet mulieres profitentes pietatem per bona opera. }}\pstart Eadem sine dubio puritas affectuum requiritur a mulieribus quae a viris: sed est praeterea, cuius mulieres sunt admonendae: quomodo videlicet etiam in vestitu sese gerent, vbi ad ecclesiasticum coetum conuenerint. Consimiliter et mulieres, id est, orent, sustollentes puras manus, absque  ira et disceptatione.  \pend\pstart Sed praeterea volo mulieres ornare semetipsas in amictu modesto, cum verecundia et castitate.) Solent mulieres huic vitio esse obnoxiae, quod magnificentius comptae ambiunt prodire in publicum: qua ex re dum arroganter captant gloriam, et quemadmodum ille ait, videre cupiunt ac videri: sine dubio verecundiam deponunt, quae praecipuum est mulierum castarum ornamentum, et periclitatur earundem castitas. Quae enim uerecundiam exuerit, quaeque  uidere et uideri ambit, certe seipsam quodanmodo exponit aliis concupiscentibus expetendam. Vnde Ecclesiastici26. Fornicatio mulieris excellentia oculorum, et in palpebris illius agnoscetur. Et prouerbus 6. mulier ornatu meretricio describitur, garrula, uaga, quietis impatiens, quae apprehendens iuuenem deosculatur, ac procaci uultu blanditur, decipiens animam illius. Ecclesiastici 19. de uestitu. Amictus corporis et risus dentium, et ingressus hominis annunciant de illo. Itaque  non indecorum modo, sed saepem periculosum est, quod mulieres se uel supra morem patriae, uel supra facultates suas, uel secus omnino, quam decet matronas, praesertim profitentes pietatem, comunt exornantue. Debent ornatae esse amictu modesto, qui satis est, ut placeant suis uiris. Graece est, ἐν καταστολῇ κοσμίῳ. Significat autem καταστολὴ stolam quandam et amictum simplicem, qui undiquaque tegit ac uelat, ne quid sit conspicuum ceu fallax, ut Theophylactus annotauit. Cum humili animo sit orandum, nequaquam decet ueste ea uti, quae excitet fastum et superbiam. Et quae uenit in coetum, ut agat cum solo Deo, nequaquam debet in ueste prae se ferre, quo uidetur uelle agere cum quibusuis  \pend
\vspace{0.5cm}\noindent
\vspace{0.2cm}\rule{1cm}{0.2pt}\\ 
\hspace{0.2cm}\textit{mg}
\footnotesize Esa.1. 
\normalsize| 
\hspace{0.2cm}\textit{mg}
\footnotesize Irae effecta mala. 
\normalsize| 
\hspace{0.2cm}\textit{mg}
\footnotesize Iacobus 1. 
\normalsize| 
\hspace{0.2cm}\textit{mg}
\footnotesize Omamenti in muliebris sexus uerecundia est et castitas. 
\normalsize| 
\hspace{0.2cm}\textit{mg}
\footnotesize Qualis esse debeat mundicies muliebris. 
\normalsize| 
\section*{AD TIMOTHEVM CAP. II. }
\marginpar{[ p.71 ]}\pstart Hiris. Certem quae ueste operosa excedit reliquas, et suam ipsius conscientiam grauat superbia, et aliorum animos irritat ad ministerium libidinis. In summa, deformat matronam quicquid laedit verecundiam aut suspectam reddit castitatem. Nam etiam suspicionem omnem euitare consultum est. Sed quid sit amictus modestus, Apostolus per contrarium voluit explicare, dicens:  \pend\pstart Non tortis crinibus, aut auro, aut margaritis, aut vestitu sumptuoso.) Cum orandi gratia adeatur locus orationis, nequaquam se sic component, quasi adirent nuptias aur theatrum, aut lupanar, caussa saltandi aut libidinandi. Vestitum magno sumptu paratum merito damnat Apostolus, quem etiam damnarunt leges Vestiariae quorundam ethnicorum principum. Nam et Suetonius in Caesare scribit, qu conchiliatam vestem et margaritas, nisi certis personis et aetatibus, perquecertos dies, permisit. Sed multo iustissime damnantur etiam illa superuacanea et operosa, cuiusmodi est torquere crines, et id genus quae non nisi ad lasciuiam sunt excogitata. ldeoque  Petrus i. cap.3. Vxorum ornatus sit non externus, qui situs est in plicatura capillorum et additione auri, aut in palliorum amictu, verum oecultus, qui est in corde homo, si is caret omni corruptela, ita vt spiritus pla cidus sit ac quietus, qui spiritus in oculis Dei magnifica et sumptuosa res est. Nam ad eum modum olim et sanctae illae mulieres sperantes in Deo, comebant sese et subditae erant suis uiris. Sed de omni hac re copiosissime Tertullianus, qui unum librum conseripsit de habitu mulierum, in quo potissimum ornamentorum et monilium usum reprehendit: alterum, de cultu foeminarum in quo fucum, lenocinium et nitorem operosum reprobat: EundemqueCyprianus imiratus est libro de habitu uirginum.  \pend\pstart Sequuntur ouo racita argumenta, quibus ad modestiam uestitus hortatur. Sed quod decet mulieres: Argumentum a Decoro. Ideo contentas esse decet uestitu simpliciore, quia id decer, maxime eas quae conueniunt ad orandum. Non conuenit matronis euntibus ad locum orationis idem habitus, qui uix conueniat adeuntibus nuptias aut theatrum: sed omnium minime decebit Profitentes pietatem per bona opera: Argumentum ex circumstantia a conditione religionis. Vt talis sumptuosus habitus alias mulieres fortasse deceat, at nequaquam decet eas, quae profitentur religionem Christianam, quae pietatem suam testari debent per bona opera, id est, in omnibus moribus externis, non. solum dictis et factis, sed etiam ipso uestitus genere: Ita namque  religionem suam atque  innocentiam testari debent Christiani, in omni uita, in dictis, factis, uestitu, incessu, denique  in omnibus rebus, quae ad aedificationem seu bonum exemplum aliorum queant fieri.  \pend
\phantomsection
\addcontentsline{toc}{subsection}{\textit{Mulier in silentio discat cum omni subiectione. Ceterùm mulieri docere non permitto, neq autoritatem vsurpare in viros, sed esse in silentio. }}
\subsection*{\textit{Mulier in silentio discat cum omni subiectione. Ceterùm mulieri docere non permitto, neq autoritatem vsurpare in viros, sed esse in silentio. }}\pstart Quoniam hic praescripsit, quomodo in coetu ecclessastico mulier sese geret in uestitu, inter orandum, adiicit per occasionem quomodo se habebunt in coetu ecclesiastico in discendo:atque haec ordinatio ualde quoque necessaria est et utilis. Proinde uult mulieres in coetu ecclesiastico, cum docctur, nihil aliud quam humiliter silere et benignis auribus amplecti quod proponitur: non aliqua a doctoribus inquirere, etiam si fortassis sit de quo dubitent, multo minus ferendum ut ipsae doceant. Clarius explicat̃ haec ordinatio siuae εὐταξία 1. ad Cor.14. ubi uiris quidem permittit singulatim omnibus prophetare, ut omnes discant, et omnes consolationem accipiant: Sed seruato interim ordine, et ne plures simul loquantur: At mulieres uult prorsus silere: Mulieres vestrae in ecclesits sileant: inquit: nec enim permissum est illis ut loquantur, sed ut subditae sint quemadmodum et lex dicit. Quond si quid discere uolunt, domi suos uiros interrogent. Nam turpe est mulieribus in coetu loqui. En quod, etiam si quid discere uellent, ipsarum tamen non est in ecclesia interrogare doctores. Caussam addit, indecorum esse mulierem quoquo pacto loqui in coetu, Est enim superbum animal mulier, et sicut ex uestitu gloriam inanem captat, ita multo magis captaret gloriam ex quaestionibus. Et praeterquam quod fortassis parum conuenientia interdum proponerent, tum doctorum contemptus inde nasceretur, praesertim si absurdis quaestionibus muliercularum pro palato omniu non satisfaceret. Sed potuisset eadem res ridicula uideri ethnicis, si permissum  \pend
\vspace{0.5cm}\noindent
\vspace{0.2cm}\rule{1cm}{0.2pt}\\ 
\hspace{0.2cm}\textit{mg}
\footnotesize A decoro. 
\normalsize| 
\hspace{0.2cm}\textit{mg}
\footnotesize A circumstantia et condinone religionis. 
\normalsize| 
\hspace{0.2cm}\textit{mg}
\footnotesize Quomodo mulieres sese gerent in coetu ecclesiastico in docendo. 
\normalsize| 
\hspace{0.2cm}\textit{mg}
\footnotesize Mulierin in ecclesia silaet. 
\normalsize| 
\hspace{0.2cm}\textit{mg}
\footnotesize A turpi siue in decoro. Mulier animal seperbum. 
\normalsize| 
\section*{IN I. EPIST. D. PAVLI }
\marginpar{[ p.72 ]}\pstart fuisset mulieribus Christianorum in coetu loqui aut disputare. Minimem itaque  decet mulierem aliquid curiose aut fastuose inquirere, cum deceat ipsam in omnibus humilitatem prae se ferre. Ideo addidit Apostolus, discat, cum omni subiectione: Nequaquam se geret vt subiectam, quae in publico conuentu cupit videri vel audiri loquens vel disputans. lam si ne interrogare quidem licet mulieri aliquid in coetu, iuxta Apostoli praescriptum, sed potius domi debet discere ex marito:at multo minus permittendum est, vt in ecclesia doceat: Vnde sequitur,  \pend\pstart Caeterùm mulieri docere non permitto, neg autoritatem vsurpare in viros, sed esse in silentio.) Vnum inconueniens ex altero sequeretur. Si mulier permittatur docere. et praescribere etiam viris credendi et viuendi normam: lam plane sequitur, quod mulier arroganter vsurpabit autoritatem in viros: quod valde est absurdum et prorsus contra nasturae ordinem. Nam mulieri dictum est a Deo Genes.3. Sub viri porestate eris, et ipse dominabitur tibi. Quod praeceptum cùm de obedientia mulieris etiam in priuatis rebus sit intelligendum: nequaquam ferendum erit, vt illa aliquid pro suo arbitrio constituat in rebus publicis siue in coetu ecclesiastico. Si non potest suo nutu aliquid agere domi in causis priuatis, multon minus id licebit ei in ecclesia in negociis religionis. Est aunt turpissimum si qui viri tam degenerant, qui sibi permittant vxores licenter imperare, et videntur in naturam ac Dei ordinationem peccare. Eccles.25. Mulier si primatum habeat, contrario est viro suo. Tametsi non abutentur viri sua hac potestate sibi diuinitus concessa. Aliquid etiam concedent mulieri intra parietes. Nam et docere domi, mulieri non est negatum, siquidem Priscilla Actor.18. vna cum marito Aquila, docuit quaedam Apollon Alexandrinum. Et alibi fidelis mulier iubetur virum infidelem commonefacere, et si possit, Christo lucrifacere. Et exigitur ab omnibus matribus  familias, vt prima elementa seu Catechismum doceant suos liberos. Imo aliquando permittit̃ in aliqua causa mulieri etiam publicem adferre sanum, si quod habeat, consilium, tametsi non temere. ludic.5. Debora erudiuit populum lsrael. Non autem quond publice doceret, sed dedit quibusdam in rebus consilium, quoniam afflata erat Spiritu sancto, qui interdum et per sexum fragilem dignatur res magnas efficere, et occulta vel admiranda indicare. Vnde et Ioel.2. legitur: Et prophetabunt filii vestri et siliae vestrae. Et Maria virgo cecinit spiritu Sancto afflata canticum, Magnificat:Item Actor.21. Philippi quatuor filiae uirgines prophetabant: Rursus 1.Corinth.11. Mulier orans aut prophetans iubeturvelare caput. Sed tamen ex his non sequitur, quod publice docere permittentur, cur autem non sit istud ferendum, sequitur:  \pend
\phantomsection
\addcontentsline{toc}{subsection}{\textit{Adam enim prior formatus est, deinde Eua: et Adam non fuit deceptus sed mulier seducta fuit per praeuaricationem, }}
\subsection*{\textit{Adam enim prior formatus est, deinde Eua: et Adam non fuit deceptus sed mulier seducta fuit per praeuaricationem, }}\pstart Duo argumenta, cur mulieri non sit permittendum, vt doceat uel usurpet autoritatem in uirum. Prius est ex ratione creationis. Aequum est, ei praeeminentiam concedere in docendo et imperando, qui primùm est a Deo conditus. Sed Adam formatus est prior Eua. lgitur debetur ei praeeminentia et imperium. Intelligendum autem est, idem obseruari debere ius dignitatis in posteris, quod contigit primis parentibus. Et per Adamum conuenienter intelliguntur omnes uiri, per Euam omnes mulieres. Quod autem prior conditus fuerit Adam, Geneseos liber clare docet. Alterum argumentum est a ratione lapsus, negatiue. Illi non est permittendum munus docendi, qui ruinae caussa extiterit. Sed mulier extitit caussa ruinae. Igitur ei non est permittendum docendi munus. Significanter dicit mulserem seductam, non Euam, quia id malum in uniuersum genus mulierum debet redundare. Nequaquam autem tutum, ei credere, quae semel in praeeundo sic decepit. Nec est, quod ambiat docendi aut imperandi munus, quae in ipso exordio, ad leuissimam persuasionem impegit et in errorem induxit. Quomodo autem mulier prius assenserit fallaci serpenti, et prius fuerit praeuaricata gustando fructum ueti» tum, uide Genes.3. Sic ibi legitur: Vidit mulier, quod bonum esset lignum ad uescen» dum et pulchrum oculis aspectuquedelectabile, et tulit de fructu illius et comedit: de„ ditqueuiro suo, qui comedit. Et post pauca, Adam dicit, Mulier qua dedisti mihi sociam.  \pend
\vspace{0.5cm}\noindent
\vspace{0.2cm}\rule{1cm}{0.2pt}\\ 
\hspace{0.2cm}\textit{mg}
\footnotesize Ab inconue nienti. 
\normalsize| 
\hspace{0.2cm}\textit{mg}
\footnotesize A minori. 
\normalsize| 
\hspace{0.2cm}\textit{mg}
\footnotesize Ab exemplo. 1.Corin.7. 
\normalsize| 
\hspace{0.2cm}\textit{mg}
\footnotesize 1.Arg. A raione creationis siue ordine naturae. Synecdoche. 2.arg.a ratone lapsus. 
\normalsize| 
\section*{AD TIMOTHEVM CAP. II. }
\marginpar{[ p.73 ]}\pstart dedit mihi de ligno, et comedi. Et mulier. Serpens decepit me, et comedi. En ibi se priorem et seductam mulier confitetur. Ratio ergo est euidens, cur mulieri non sit permit tenda facultas docendi. Nunc viderint qui committunt mulieribus gubernandas respublicas et amplas ditiones: Sine dubio ira Dei est manifesta, cum mulieres alicubi admitrant. Amazones quidem maudito exeipio tenuerut regnu, ieu tantum inter mulieres, in regno viris vacuo. Neque  exemplum id est probatum, neque  diuturnum fuit earum regnum.  \pend
\phantomsection
\addcontentsline{toc}{subsection}{\textit{Salua tamen fiet per generationem liberorum, si manserint in fide ac dilectione, et sanctificatione, cum castitate. }}
\subsection*{\textit{Salua tamen fiet per generationem liberorum, si manserint in fide ac dilectione, et sanctificatione, cum castitate. }}\pstart Emollit, si quid fortasse durius dixisse uideatur: et ostendit aliud mulierum esse officium, nempe quod educare et docere debeant suos liberos et sic si dicantur caussa ruinae extitisse, saluae esse queant etiam: quasi dicat: Non sit vobis molestum, ô mulieres, si vobis cum sitis fortassis ingenio pollentes, negatur tamen munus docendi. Potestis alia ratione vos exercere in pietatis operibus ad salutem vestram, et si ita vultis, etiam docere. Etucate et docete diligenter, tum rudimenta religionis, tum caetera quae pertinent ad morum honestatem, vestros liberos, in his docendis operam impendite, et haec opera erit vobis salutaris, maxime sivos ipsae semper in fide et charitate perstantes sancte et caste vitam transigatis. Haec conditio vere honorifica est, et satis habens in se dignitatis, etiamsi publice docendi in ecclesia ius vobis adimatur. Et recte adiecit ista Apostolus, post inustam muliebri sexui notam, de caussa ruinae humani genenis. Quamuis mulier fuerit occasio damnationis, tamen potest illa quoque  esse occasio conseruationis. Ipsa inprimis potest salua fieri, si manserit in fide et dilectione, et aliquo modo dici caussa salutis filiorum, quos ipsa bene instruxerit, et instrus ad gloriam Dei rite curauerit. Sicut sine dubio multum iuuit Timotheum, quod recte ab infantia, procurante matre, in sacris literis fuerat institutus. Itaque hic cum ait, per generationem liberorum, intelligi conuenit, quicquid ad puerorum generationem et educationem pertinet, nempe partus dolores, labores in educando, diligentiam et curam in formanda infantia et pueritia, quae sine dubio laboriosa sunt, et Deo inprimis grata officia, maxime si ex animo et longanimiter a matribus perficiantur. Hinc discimus, non satis esse, si parentes curent, vrliberi viuant, sed oportet eos curare, vt bene viuant, vt in fide et timore Domini etiam a teneris vnguiculis instituantur. Eccles.7. filii tibi sunt, erudi illos, et incurua illos a pueritia illorum.  \pend\pstart St manserint in fide ac dilectione et sanctificatione cum castitate.) Graece legitur, ἔαν μένωσιν, plurali numero, si manserint. ἑτέρωσις numeri per hebraismum, qui est Paulo frequens, vt hic numerus respiciat ad id, quod principio habebatur. Consimiliter et mulieres. Quidam vitantes hebraismum malunt legere in singulari numero. Perspicuum autem est, verbum istud referendum ad mulieres, non ad liberos. Quis enim sensus fuerit: Mulieres saluas fore, si liben manserint in fide et dilectione: quasi vero in potestate mulieris esset, efficere, vt tales maneant, aut ipsa alioquin seruari non possit, nisi illi integri in fide perstent? Significatur itaque mulieres ipsas saluas fore, non quidem ratione educationis prolium, quasi ea re salutem promererentur, sed (quod per Correctionem uidetur addi) si quoque  diligentiam omnem expenderint in educandis prolibus, manserint in fide et dilectione, id est, si a fide et Charitate non excidant Et si manserint in sanctificatione, id est, ea praestent, per quae eas sanctas esse agnosci queat: denique  cum castitate, id est, si uerbis, factis, denique  uniuersa uita sobrietatem et castitatem, et modestiam prae se ferant. Proinde hic locus est admodum consolatorius, praesertim iis mulieribus, quae arbitrantur, se prolium turba grauari. Scire debent, suam conditionem Deo probari, maxime si diligentes sint in liberorum suorum honesta educatione. Et matronarum singulare decus esse liberos, constat ex Corneliae Gracchorum iudicio. Apud hanc enim cum Campana quedam matrona diuertisset, quae ostendit ei monilia sua, qualia illo seculo maximi fiebant, illa remorata est Campanam sermone, donec liberi e schola redirent. Tum, haec, inquit, ornamenta mea sunt: sentiens matronae nihil esse pretiosius, quam liberos recte educatos. Item com  \pend
\vspace{0.5cm}\noindent
\vspace{0.2cm}\rule{1cm}{0.2pt}\\ 
\hspace{0.2cm}\textit{mg}
\footnotesize vbi  et quos docere debeuut mulieres. 
\normalsize| 
\hspace{0.2cm}\textit{mg}
\footnotesize Mulier quomodo saluari queat. 
\normalsize| 
\hspace{0.2cm}\textit{mg}
\footnotesize ἑτέρωσις numeri per Hebraismum. 
\normalsize| 
\hspace{0.2cm}\textit{mg}
\footnotesize Correctio. 
\normalsize| 
\hspace{0.2cm}\textit{mg}
\footnotesize 1. Thess. 4. Haec est uoluntas Dei, Sanctificatio uestra, etc. 
\normalsize| 
\hspace{0.2cm}\textit{mg}
\footnotesize Liberis bene educatis mulieri nihil pretiosius, nihilquelaudabilius. 
\normalsize| 
\section*{IN I. EPIST. D. PAVLI }
\marginpar{[ p.74 ]}\pstart mendantur hic praecipua ornamenta matronarum, per quae salutem queant consequi. fides, inquam, et charitas, quae in omnibus requiruntur: Deinde sanctificatio et castitas, quae tanquam monilia sunt et decus praecipuum omnium matronarum. Ethaec quidem de doctrina τῆς εὐταξιας  breuiter delibata.  \pend
\phantomsection
\addcontentsline{toc}{subsection}{\textit{Posterior pars Epistolae, in qua ostendit qualis episcoporum vita esse debeat, ac primo quomodo se gerent priuatim in domo sua et erga familiares. }}
\subsection*{\textit{Posterior pars Epistolae, in qua ostendit qualis episcoporum vita esse debeat, ac primo quomodo se gerent priuatim in domo sua et erga familiares. }}
\endnumbering\beginnumbering\section{CAPVT III}
\phantomsection
\addcontentsline{toc}{subsection}{\textit{\huge\textbf{I}\normalsize NDVBITATVS sermo. Si quis episcopi munus appetit, honestum opus desiderat. Oportet igitur episcopum irreprehensibilem esse, vnius vxoris maritum, vigilantem, sobrium, etc. usque ad cap.4. }}
\subsection*{\textit{\huge\textbf{I}\normalsize NDVBITATVS sermo. Si quis episcopi munus appetit, honestum opus desiderat. Oportet igitur episcopum irreprehensibilem esse, vnius vxoris maritum, vigilantem, sobrium, etc. usque ad cap.4. }}\pstart \huge\textbf{Q}\normalsize Vemadmodum in argumento diximus, distribuitur haec epistola in duas partes principales: quarum prior de doctrina episcoporum agit, posterior de vita. De hac incipit apostolus in praesenti disserere. Cum autem in vita episcopi spectandum sit, primo quomodo se gerat priuatim in domo sua et erga familiares: Deinde vero quomodo publice et in ecclesia erga cuiuscunque  ordinis homines: hoc loco depingit, qualem episcopum domi suae esse deceat, et quibus virtutibus ornatus conspiei debeat.  \pend\pstart Indubitatus sermo.) Statim autem, et velut lege Attica ἄνευ προοιμίου καὶ πάθος rem aggreditur, vti saepe alias. Indubitatus sermo: quasi dicat: Quae nunc dicam, certa acrata habeantur, fide digna sunt. Sic enim graece πιστός ὅ λόγος. Vno itaque  uerbo praefatur et attentos reddit. Nonnulli existimant ueluti per occupationem praefixum esse. Nam cum diuersi munus episcopi appeterent, atque  alii quidem spe quaestus aut gloriae, alii sine dubio synceriter, et ad promouendam pietatem, atque  illi nihilominus praetexerent honestatem perinde atque  isti : videtur Apostolus se velut arbitrum voluisse huic quaestioni interponere, dicens: Indubitatus sermo: quasi diceret: Dicam quodres est: Si quis episcopi munus appetit, honestum opus desiderat: Si tamen is talis esse velit, qualem in ea functione esse oportet. Theophylactus autem secutus suum Chrysostomum, haec verba, epilogi vice adiungit praecedentibus quae dicta sunt de officio mulieris, vt sit sensus: Ea quae dicta sunt de mulieribus per generationem et educationem liberorum saluandis, uera esse, indubitata, fide digna, ac minime de illis dubitandum. Magis consentaneum videtur, connectenda iis, quae sequuntur de vita et moribus episcoporum, et ad renouandam attentionem exordii loco praefixa. Proinde sic definit Apostolus.  \pend\pstart Si quis episcopi munus appetit, honeslum opus desiderat.) Quoniam episcopi munus recte administratum, per se bonum est et gratum Deo, definit Apostolus eum bene agere, qui id munus appetit, ea videlicet intentione vt recte administret. Obseruanda autem vocis proprietas ἐπισκοπὴ, ita enim graecem vna vox est, vbi nos habemus, episcopi munus. Significat igitur ἐπισκοπὰ diligentem inspectionem, speculationem, obseruationem, visitationem: a verbo ἐπισκοπὲω, quod est considero, superintendo, recenseo, recognosco, etiam uiso aegrotum, uel, ut uulgo dicunt, uisito. Vnde ἐπίσκοπος obseruator, speculator, explorator, custos, uisitator. Et erant episcopi dicti quidam magistratus apud Athenienses et Lacedaemonios, quos et αομοςάς uocabant. Praeerant ii reficiendis seruandisquearcibus et publicis aedificiis. Etiam iudices in publicis certaminibus dicti sunt episcopi. Hinc nomen episcoporum accommodatum est iis, qui praesunt ecelesiis, quorum interest iuxra nominis significationem, considerare, superintendere, uigilare, uisitare, denique  omnibus modis prospicere et cauere, ne quid gregi credentium incommodi in doctrina fidei uel morum accidat. Est episcopus magis officii.  \pend
\vspace{0.5cm}\noindent
\vspace{0.2cm}\rule{1cm}{0.2pt}\\ 
\hspace{0.2cm}\textit{mg}
\footnotesize Deuita episcopi. 
\normalsize| 
\hspace{0.2cm}\textit{mg}
\footnotesize Auentio. 
\normalsize| 
\hspace{0.2cm}\textit{mg}
\footnotesize ἐπισκοπή quid. 
\normalsize| 
\hspace{0.2cm}\textit{mg}
\footnotesize ἐπίσκοπος. 
\normalsize| 
\hspace{0.2cm}\textit{mg}
\footnotesize Episcopis Atheniensium et Lacedaemoniorum. Episcopi munus. 
\normalsize| 
\section*{.AD TIMOTHEVM CAP. III. }
\marginpar{[ p.75 ]}\pstart diligentiae, et laboris nomen, quam dignitatis. Multa cura; multi labores, muntum taedum ei erit deuorandum, qui hoc munus recte obibit. Quod cu isa fit, recre ait pauius, Si quis episcopi munus appetit, bonenstum opus desiderat.) Episcopi munus opus vocat  \pend\pstart Nam reuera non est quod in otio et per lusum poisit obiri: sed valde operosa res et ple na molestiarum functio, quam nemo ambire debeti nisi qui prius metitur, quid hume ri eius possint, quid ferre recusent. Qui sua metitur pondera serre potest. Recte qui dam interpretantur opus Hebraico more accipi pro labore: Sic enim opus manuarium et laborem manuarium pro eodem dicimus. Et honestum opus vocat episcop munus, quia si in virtute honestas sita est, merito ea functio honesta appellabitur, in qua oportet quamplurimas virtutes relucere. Neque vlla plus honoris meretur, quam quae tota ad Dei honorem augendum expenditur. Hoc igitur pacto definit Apostolus omnes in genere quicunque sint, qui episcopi munus appetunt, honestum opus desiderare. Si quis, inquit. Neque enim ita ad Timotheum ista scribit, vt non velit eadem cognosci ab omnibus, et vsui esse omnibus piis, et omnibus temporibus. Ista itaque debemus ad nostra tempora, sicut et quamcunque doctrinam Apostolicam accommodare. Videre quidem est interdum multos, qui episcopi munus appetunt, sicut Apostolus appetendum intelligit. Nam appetere, uti Apostolus hicsignificat, nihil aliud erat, quam vitam suam et doctrinam per omnia ita instituere et cognoscendam omnibus piis proponere, vt decernere omnes possent, eu qui sic viuebat et docere poterat, idoneum esse, cul tantum munus committeretur. Haec erat appetendi tantum munus ratio, hic erat ambitus, et desiderium honesti operis. Nam tunc quidem eligebantur in vnaquaque ciuitate ab vniuersa credentium multitudine, vel saltem ab episcopo, seu Apostolo aliquo, qui illic docuerat et alio proficisci debebat, populo interim consentiente et approbante. Hoc vero pacto non erat eligendus nisi qui prorsus talis erat doctrina et vita, qualis hic ab Apostolo describitur. Itaque eniti quemquam, vt talis fieret, et idoneus ad tantum munus iudicaretur id orat episcopi munus appetere At nunc, ꝗđ dolendum, longe alia ratio est. Nulli magis appetunt, et nulli eriam citius obtinent, quam qui omnium minime sunt idonei:cum tamen nihil sic optandum foret nihilque conducibilius reformandis ecclesiis, quae passim iam sunt corruptae et desolatae, et doctrina sana orbatae, quam vt ad veterem illam eligendi episcopos formulam rediretur. Certum est, si vel husus rei haberetur ratio, id est, si illi duntaxat praeficerentur ecelesiis, qui vere sunt idonei, qui doctrina et vita sua se commendant, qui non dignstatis sed officii, non quaestus sed labo ris, non arrogantiae aut propriae gloriae, sed vtilitatis publicae et gloriae Dei desiderio appetunt: sperari posset sanam docttinam ecclesiae restituendam, et abusus omnes et sectas prope; dic̃ cuertendas. Non est itaque temere et a quibusuis id munus ambiendum, ne que quibusuis ambientibus deferendum. Hebus 5. Nemo sibi vsurpet honorem, nisi a Deo uocetur, quem admodum et Aaron. Percussus est lepra Ozias rex luda de quo 2. Paralip..6. et e domo Domini eiectus, quod officium sacerdotum sibi arripere uoluit, et thus superaltare thimiamatis adoleust, quod tamen factum per se non uidetur malum. Sic Nadab et Abihu filiiAaron, quod alienum ignem intulissent in Templum Domini, ut habetur Leuit. 1o. ab igne deuorati sunt.Item Numer. 16. Chore, Dathan et Abiron, quia impu denter aspirarunt ad dignitatem, quae ipsis non debebatur, sed Most et Aaroni, a tetra dehiscente absumpti sunt, et descenderunt in infernum uiuentes. Sicut autem pericu losum est appetere munus episcopi, cum a Deo non uocaris, ita rursus periculosum est, si uocatus id recuses, uel negligas. Vnde Exod.3.4. Dominus iratus est Mosi, qui recusabat, caussans suam balbutiem. et Hierem. 1. Hieremias reprehenditur a Deo dicente, Noli dicere, quia puer ego sum, quoniam ad omnia, quo mittam te, ibis, et uni. uersa quaecunque mandauero tibi, loqueris. Ideo laudatus Esaias, qui, ut habetur Esaiae 8. intelligens se uocari, lubenter se obtulit, dicens. Ecce ego, mitteme. Et Paulus Act. 9. promptum sese exhibuit, quando ait: Domine, quid me uis facere. Foelix ergo ille, qui cum sit idoneus, se etiam idoneum ostendit, et sic modeste taciteque episcopi munus ambit. Foelix item ille, qui ad munus id a Deo per muleitudinem eligentem uotatur, et dum uocatus est, non negligit, sed strenue et cum laude sludobit.  \pend
\vspace{0.5cm}\noindent
\vspace{0.2cm}\rule{1cm}{0.2pt}\\ 
\hspace{0.2cm}\textit{mg}
\footnotesize Soiscopim nus in laboribus et soucitudine uersatur, non AAuse etfeStuitate. 
\normalsize| 
\hspace{0.2cm}\textit{mg}
\footnotesize Episcopine aus non temere estanbiendum. 
\normalsize| 
\hspace{0.2cm}\textit{mg}
\footnotesize Exέρμεστι qui ob temere affectatum ntunuis, diuinius sunt pu niti. Munusa beomposilum non redisnum 
\normalsize| 
\section*{IN I. EPIST. D. PAVLI. }
\marginpar{[ p.76 ]}
\phantomsection
\addcontentsline{toc}{subsection}{\textit{Oportet igitur episcopum irreprehensibilem esse, et vnius vxoris maritum, vigilantem, sobrium, modestum, hospitalem, aptum ad docendum, non vinolentum, etc. }}
\subsection*{\textit{Oportet igitur episcopum irreprehensibilem esse, et vnius vxoris maritum, vigilantem, sobrium, modestum, hospitalem, aptum ad docendum, non vinolentum, etc. }}\pstart lam qualem eum esse deceat, qui episcopi munus ambiet, quique ad id omnium suffragiis admittetur, describit, requirens non solùm in eius persona uirtutes plurimas, uetumetiam a crimine immunes vult vxorem, liberos, ministros. Ac primùm quidem de episcopo in genere ait,  \pend\pstart Oportet episcopum irreprehensibilem esse:) id est, inculpatae vitae, tamque ; immunem ah omni vitio, vt ne Momus quidem, si fieri possit, aliquid in eo reprehendat, saltem iuste. Nam a malis iniuste uituperari, id etiam in laude ponendum est: sicut ab illaudato viro laudari illaudabile est. Quis enim omnibus placeat ? aut quis malis placere gaudeat? Et quomodo refrenentur linguae malorum hominum, qui in quemuis vel integerrimum linguae suae spicula venenata iaciunt? Christus ipse, Apostoli, caeterique sancti non potuerunt peruersorum obloquia et detractiones effugere. Tunc autem erit apud bonos irreprehensibilis, ubi uita ipsius doctrinae responderit. Alias enim si innocentiam multum commendare velit, qui ipse nocens reperitur, si de puritate multum disi serit ipse impurus, si vitia in aliis damnat, quibus ipse est obnoxius: hunc ne boni quidem satis ferant. Quantum enim is docendo aedificat pridie, tantum non recte agendo destruit postridie. Istud quidem in genere, quond irreprehensibilis esse debeat, et vacare omni crimine episcopus, abunde satis dictum videbatur: sed placuit Apostolo etiam diuersas virturum species, quibus ornari uult episcopum, necnon per contrarium uitiorum species, a quibus episcopus abesse debet, subiicere.  \pend\pstart Vnius vxoris maritum.) Quoniam, si episcopus scortatione vel adulterio sese polluat (quod peniculum fere semper iis imminet, qui propriae fragilitatis immemores se obstringere volunt coelibatui:) non potest in aliis eadem vitia reprehendere: imo et malo suo exemplo fenestram aperit et aliis ad omnia impudicitiae genera: ideo episcopus, vt etiam suspicionem illorum vitiorum effugiat, vnius vxoris maritus esse debet. Sane si cui donatum est a Deo, vt possit vxore carere et pudice viuere, non est coniugii vinculis irretiendus. Bonum est enim mulierem non tangere, ob incommoda quae solent coniugatos ministros verbi, maxime vbi perturbantur imminente persecutione, vel ad fugam compelluntur, consequi: Verùm qui non potest ita se continere, satius est, eum vxorem ducere quam vri. Honorabile enim connubium apud omnes, et cubile impol » lutum: Hebr.13. Caeterùm vnius vxoris erit maritus. Nam plures simul uxores habere aut gregem concubinarum, licet ludaeis sub lege erat permissum, tamen incontinentiae notam incurreret apud Christianos, qui idem uellet, maxime autem si id atrentaret episcopus, qui coniugalis pudicitiae exemplum omnibus esse debet. Neq hinc sequitur, episcopum mortua priore uxore alteram haud posse ducere. Haud enim dicendus est plures habere uxores, qui defuncta priore ducit alteram. Vnde 1.Corinth. 7 de uiduis et inconiugatis loquens Apostolus, Quod si se non continent, inquit, contrahant ma trimonium. Nam satius est matrimonium contrahere, quam uri. Et rursus. Solutus es ab uxore? ne quaeras uxorem: quod si duxeris uxorem, non peccasti. Dixerat enim pri ,us. Propter stupra uitanda, suam quisque uxorem habeat, et suum quaeque uirum habeat. A qua lege uiduus episcopus non est excludendus, non magis quam natura liberat eum ab illicitis affectibus. Sed neque conueniens esset, eum, florenti adhuc aetate, liberos, familiam, domum deserere, sine ulla domestica oeconomia, quae per mulierem optimem solet administrari. Certe nisi episcopus ducta iam altera uxore uideatur uel a charitate abduci, uel segnior reddi in suo officio, non uidetur reprehendendus, multo minus reiiciendus. Vnius itaque uxoris maritum esse uult episcopum, id est, excludi uxorum seu concubinarum turbam, sanctum uero et iustum matrimonium inter duos (iuxta praescriptum Christi Matth.19.) stabiliri, et pudice coli non prohibet.  \pend\pstart  Vigilantem, νuφαλιου.) Videtur alludere ad etymologiam nominis episcopi: nam et iπlmor o lgnsieatuig le custodem. Ary oporiet epscopu quoue qrgs uic lantere.  \pend
\vspace{0.5cm}\noindent
\vspace{0.2cm}\rule{1cm}{0.2pt}\\ 
\hspace{0.2cm}\textit{mg}
\footnotesize vitutes EPiscopt. 
\normalsize| 
\hspace{0.2cm}\textit{mg}
\footnotesize Qgatenus bonum mulie remnontan gere. 
\normalsize| 
\hspace{0.2cm}\textit{mg}
\footnotesize vigilantia ciscoρι 
\normalsize| 
\section*{AD TIMOTHEVM CAP, III. }
\marginpar{[ p.77 ]}\pstart esse, dareque operam, ne ipso dormiente et negligente, inimicus hominis in agro Domini Zizania serat, ut est in parabola Matth. 13. Eam ob causam dicti quoque sunt episcopi speculatores, quoniam pro grege excubias agere debent: vndique circumspicere ( more eorum qui excubias agunt in aeditis speculis ciuitatum vel arcium) ne alicunde irrumpant aduersarii! quos si vel procul prouideat, annunciabit statim infla1o cornu uerbi Des, ut omnes ad depellendum hostem sint parati. Sic Dominus Ezechiel. capite tertio aiebat: Fili hominis speculatorem dedi te domui lsracl. Eódem pertinent quaedam Ezechiel.33 et Actor.2o. Attendite uobis et cuncto gregi, in quo, uos spiritus Sanctus posuit episcopos ad regendam ecclesiam Dei. Ego enim nout hoc,, quos ingressuri sint post discessum meum lupi graues in uos, non parcentes gregi.Se quitur. Propter quod Vigilate, etc.  \pend\pstart Sobrium: σώφρονα.) id est, temperatum, quique nullo calore animi accensus res Ieni affectu transigit, et valde composite se gerit. Actor.26. Non insanio optime Feste, sed ueritatis et sobrietatis uerba eloquor, id est, composite, mature, et deliberato animo loquor. Qui nulla in re excedit modum, sed temperate et modeste agit omnia, talis sobrie agere dicitur. Accepta significatio ab his, qui parce cibum et po tum sumunt. Hi enim non facile excedunt in dicendo uel agendo, sed modeste, quiete, leni animo agunt omnia:sicut contra qui ingenium, mentem obtundunt sese ingur gitando immodico cibo aut potu, plerunque uel prorsus stupidi sunt, uel quid quomodoue agant, non sat uident, ubique modum excedentes.  \pend\pstart Atodestum, κόσμιον,) id est, facientem ea, quae ipsum ornant ac decent, non deformant, non dehonestant: qui nihil facit praeter decorum. Resertur ad habitum: cultumque corporis, etiam ad gestuum rationem. Multum moments habet, qua dex. teritate gestuum quis quid faciat. Et in uestitu pariter sordes, pariter immoderatus sumptus et cultus dedecent. Mundicies honesta et mediocritas uestitus, neque adhypocrisin neque ad fastum, neque ad lasciuiam sit compositus. Denique in omnibus actionibus candorem, puritatem, frugalitatem, simplicitatem suam testatam faciat. Absit ab eo et nimia seueritas et nimia leuitas. Neque quosuis Timonico misanthropo supercilio auersetur:neque quibusuis histrionica leuitate sese accommodet. Medium ubique tenendum, ut constet modestia.  \pend\pstart Hospitalem quλesyou) Est insigne officium Charitatis, bene mereri de egenti. bus, praesertim iis qui uel propter iustitiam, uel alia de caussa fortunis suis spoliati coguntur peregre agere, Exigitur autem eximia haec charitas ab omnibus Christianis. Nam 1. Petri 4. Hospitales estote erga uos inuicem, sine murmuratione. Et Hebr. 13. Hospitalitatis ne sitis immemores: per hanc enim quidam inscientes exceperunt Angelos hospitio. Sed inprimis exigitur ab episcopis. Atque hanc praecipue ob caussam amplae facultates ecclesiis sunt superioribus seculis donatae, ut uidelicet hospitalitatem episcopi exercerent itemut inde alerentur pauperes uiduae, curarentur aegroti, instituerentur puen, alerentur sacrarum literarum studiosi:ob eandem hospitalitatem erecta multis in locis Xenedochia. Sed quem in usum nunc conuertunt episcopi eas facultates? Extant autem praeclara hospitalitatis exercendae exempla in scripturis. Notae sunt historie Abrahae, Loth, mulieris Sunamitidis, Marthae, Coriarii, apud quem Petrus Act. 1o, diuersatus fuit: Publii, Actor.28. lasonis. Gentiles quoque hospitalitatem Gilliae Tarentini aliorumque celebrarunt. Vide et Ciceronem in officiis.  \pend\pstart Aptum addocendum duf'axlixop:) Inuenias, qui uita sint laudata et probis moribus praediti: inuenias etiam qui recondita quadam doctrina sint condecorati: sed raro iidem sunt apti ad docendum. Qualis doctrina esse debeat episcoporum, prioribus capitibus et in aliis Epistolis diligenter est ab Apostolo traditum. Sed alia quaedam facultas requiritur, ut quis aptus sit ad docendum. Etiam doctissimos uiros com pertum est aliquando infoeliciter attentasse docere populum. Nonnulli foeliciores sunt in exercendo stylo, alii rectius docent in Schola:pauci id dextre agunt in templo apud imperitam multitudinem, in qua quot capita, tot sententiae. Ad Tit. 1. Vult  \pend
\vspace{0.5cm}\noindent
\vspace{0.2cm}\rule{1cm}{0.2pt}\\ 
\hspace{0.2cm}\textit{mg}
\footnotesize Ezech.3.33. 
\normalsize| 
\hspace{0.2cm}\textit{mg}
\footnotesize Sobrietas et temperantia. 
\normalsize| 
\hspace{0.2cm}\textit{mg}
\footnotesize Modestia Ad Philip.4 Modestia ue stra nota sit omnibus homninibus Decet et status, et color et res. 
\normalsize| 
\hspace{0.2cm}\textit{mg}
\footnotesize Hospitalitas episcoporum 
\normalsize| 
\hspace{0.2cm}\textit{mg}
\footnotesize Caussa, propter quas ecclesia facultatibus auctae et locupleta tae sunt. 
\normalsize| 
\hspace{0.2cm}\textit{mg}
\footnotesize Διδασκαλεις 
\normalsize| 
\section*{IN I. EPIST. D. PAVLI. }
\marginpar{[ p.78 ]}\pstart Apostolus episcopum, qui potens sit exhortari per doctrinam sanam, et contradicentes conuincere. In hoc genere excelluit apud veteres scriptores Chrysostomus, id quod Homiliae eius testantur: ab hac facultate videtur traxisse id nominis. Eam ob caussam Valerius Hipponensis episcopus, homo Graecus et haud dubie doctus, asciuit nouo exemplo tum Augustinum, qui magis aptus erat ad docendum, egitque ille Deo ob talem virum diuinitus sibi concessum gratias, eundemque se viuo episcopum ordinari cu rauit. lraque qui formantur vt praesint ecclesiis, diligenter exercendi sunt in artibus dicendi, dandaque opera, vt ad oratoriam etiam dignitatem, quantum ministro verbi con uenit, assurgant.  \pend\pstart Non vinolentum μὴ πέροινοg.) Nox et amor vinumque nihil moderabile suadent, inquit poeta, proinde ebrietas cum omnem aetatem et sexum dedeceat, et aliena esse debeat ab omnibus, qui in aliqua publica functione versantur, maxime tamen defugienda est episcopis. Non modo minus ideonei redduntur ad docendum, qui uino sunt dediti et obnoxii, verum multis aliis malis sunt expositi, vt taceam foedissimum inde exemplum ebrietatis ad alios deriuari.ldeo lex Mosaica interdicebat Leuitis, imo et summo pontifici, cùm ad altare essent accessuri, omnem vini vsum et cuiuscumque potus qui inebriare posset. Recte dicitur prouerba3. Ne intuearis uinum quando flauescit, cum pellucet in uitro color eius. Ingreditur blande, et in nouissimo mordebit vt coluber, et sicut regulus venena diffundet. Moderatus vini vsus laudatur, et Paulus eum Timotheo praescribit: Et Augustinus vino semper usus legitur, quod acueret ingenium.  \pend\pstart Non percussorem.un πλήκτιο.) Manibus pulsare quempiam aut uim inferre, in nullo vel plebeio tolerabile est. Quare ab episcopo requiritur, ut non solùm manibus temperet, uerùm ne uel lingua quenquam procacius impetat. Atque ita hic accipiendam percussionem de ea quae per linquam fit, satis euidens est ex eo, quod postea iterum subditur, alienum esse debere a pugnis. Ergo ne lingua quidem aliquem laedet, etiam in legitima obiurgatione peccatorum, eam seruans modestiam, ut paternum ubique affectum ostendat, et hortari aut condolere potius uideatur, quam obiurgare, et menderiuelle, non exacerbare. Et quemadmodum dedecet, antequam eligatur ad tam splendidum munus, percussorem esse:ita ne in domum quidem propriam saeuire aut grassari iniquius conuenit. Cum ratione domus est gubernanda et mulieri uasi infirmiori sapi enter cohabitandum.  \pend\pstart Non turpiter lucri cupidum, uη αἀσχοκερδῆ.) Stipendio suo contentum esse decet:stipendium autem tale esse conuenit, unde et familiae suae honeste prospicere, et hospitalitatem, eleemosynas, similiaque praestare queat. Caeterum ab omni auaritia debet se ostendere alienum, ne existimet turpiter inhiare lucro Lucrum autem est, ubi plus accipit quam dat. Ab hocaffectuoportet episcopum esse alienum, dareque per operam, ne lucrum uel sperare uideatur ex aliquare, multo minus ne ex lucro pendêre eius fortunae putentur. ldeoque  mercatus, opisicium, similiaque ex quibus lucrum speratur, iudicantur episcopo nequaquam conuenire, ne uidelicet ulla lucri cupiditate ardere uideatur. Nec munera etiam accipiet, id quod et in iudicibus grauiter reprehenditur. Ezech.34. contra lucri cupidos: Vae pastoribus Israel, qui pascebant semetipsos, etc.1. Petris. Pascite quantum in uobis est, gregem Christi, curam illius agentes non coacte, sed uolentes: non turpiter affectantes lu crum, sed propenso animo, neq ceu dominium exercentes aduersus cleros, sed sic, ut *sitis exemplaria gregis. En quomodo ne affectare quidem lucrum decet.  \pend\pstart Sed aquum:) qui in omnibus aequabilitatem iusta trutina seruet, personas non respiciat, nusquam fouens partes, aut uni magis addictus, quam alteri: in puniendis peccatis iustitiae adiungens misericordiam, et misericordiae iustitiam. Nam haec duo semper in aequabilitate comiungi debent. De hac aequabilitate uidetur illud intelligendum Deut.  \pend
\vspace{0.5cm}\noindent
\vspace{0.2cm}\rule{1cm}{0.2pt}\\ 
\hspace{0.2cm}\textit{mg}
\footnotesize Vocatio At gustini epi scopi Hipp uensis. 
\normalsize| 
\hspace{0.2cm}\textit{mg}
\footnotesize Vmolentie incommoda 
\normalsize| 
\hspace{0.2cm}\textit{mg}
\footnotesize Leuit.1e. Legem Na caraeorum lege Num.6. 
\normalsize| 
\hspace{0.2cm}\textit{mg}
\footnotesize Linguauer berare quie Reprebensio moderas ta et paer na. 1.Pet.3. 
\normalsize| 
\hspace{0.2cm}\textit{mg}
\footnotesize Auaritia fugienda epl scopo. 
\normalsize| 
\hspace{0.2cm}\textit{mg}
\footnotesize Lagutones corrumpun et excaeca sapientes. Exod.23. 
\normalsize| 
\hspace{0.2cm}\textit{mg}
\footnotesize Aequitas. eneuan. 
\normalsize| 
\hspace{0.2cm}\textit{mg}
\footnotesize 
\normalsize| 
\section*{AD TIMOTHEVM CAP. III. }
\marginpar{[ p.79 ]}\pstart nihil per iram aut praecipitantiam designet. Ac verum quidem elt ab his quoque alienilsimum esse debere episcopum Altercari, rixari, nulli probo conuenit: et nimium altercando veritas amittitur.  \pend\pstart HρλασνοLAlienum ab auaritia ἀφιλάργυρον,) nequaquam pecuniae studiolum: Vix vnquan possunt rebus diuinis vacare, quorum animus rerum terrenarum cupidate impeditur. Nemo potest seruire Deo et Mammonae. Sed neque illis committenda pe Luc.6. cuniarum ecclesiasticarum, e quibus pauperes alendi sunt, cura, qui vitio φιλαργυριας sunt notati. Dum autem alienum vult ab auaritia, eadem ratione alienum vult ab omnibus, quae auaritiae et turpis lucri caussa committuntur, cuiusmodi sunt mendacium pergrium, dolus, etc.  \pend\pstart Qui suae domui bene presit.) Sicut auaritia detestanda est: ita et prodigalitas seu rei fa Prudentuain Edministran miliacis profusio. Dilapidator seu decoctor, aut alioqui incurius rerum domesticard̃ la oeconoparum idoneus iudicatur qui administrare possit ecclesiam. Fere enim solent illam in nia. scitiam aut ignauiam alia vicia comitari. Proinde ex administratione rei priuatae disce, Sane qui domi suae negligens fuerit, non est quod expectetur ab eo magna diligentia in rebus ecclesiae: et qui domi suae iracundus, litigiosus, immodestus, auarus, talis quoque  erit in ecclesia. Circumfert inter Chilonis apophthegmata, quod admonuerit vaumquemque bene praeesse suae domui. Alias enim nunquam vtilem fore reipublicae, qui pri uatam recte non gubernarit. Cùm domus quandam speciem paruae reipublicae praese ferat. pophibeChilonis a\pend\pstart Qui liberos habeat in subiectione cum omni reuerentia.) Ex liberorum educatione lolet de parentum prudentia et probitate iudicari: id quod Plato quoque in Lachete siue de fortiudine admonuit. Nec verosimile est, eum effecturum in ecclesia vt tum ob doctrinam, tum ob mores ab omnibus honoret, qui domi suae non tantum effecerit prius, ve, pprii filii ipsum audiant et reuereant. Pueris seu filiis quidem incumbit sedulo iuxta praeceptum Dei, honorare patrem et matrem, ut sint longaeui super terram: Atqui a parentibus etiam exigit Deus vt sic se gerant apud liberos, eosque sic forment, vt sciant quomodo praeceptum illud sit implendum. Pariter autem vitio vertitur parentibus et nimia austeritas, et nimia indulgentia. Tametsi hac sicut plures perdunt̃ liberi, ita grauius a Deo punit, vtpote quae fenestram aperiat ad scelera. Nota est historia Heli summi sacerdotis, a quo exemplum sumere debent Episcopi 1.Sam.3.4. Contra austeritatem parentum est illud Colos.3. Patres ne prouocetis vel irritetis liberos vestros, ne despondeant animum. Tobias recte instituebat fi lium, dicens: fili omnibus diebus vitae tuae in mente habero Deum, et caue ne aliquando pec cato consentias, et pretermittas praecepta dei nostri. Ethnicos in hac parte diligentes fuil se testant libelli diuersi, continentes precepta moralia iuuenturi pponi solita:vt est lsocratis paraenesis ad Demonicum:item Epicteti stoici libelle, Xisti philosophi, Catonis disticha. dare familiri. gma de curam uferitag indulgentia nimia perinde in uitio est.  \pend\pstart Quod si quis propriae domui praeesse non nouit, quomodo ecclesiam Dei curabit? Argumentum a minori. Non tam parua res est bene prçesse domui atque vulgon existimat. Prospicere omnia ad victum necessaria vxori ac liberis, alere ancillas aut ministros, pecuniam honestis parare artibus commutare, locare, conducere, emere, vendere, denique quacunque alia ratione contrahere, haec saepe magnam molestiam habent adiunctam, neque sine cura et prudentia recte peragunt. Res ipsa docet, quam paucis bene cadat haec alea. Vnde non abs re eruditi homines putarunt sibi quoque libros de oeconomia pscribendos, qualis extat Xenophontis. Sed multo alia res, multis interuallis superior, praeesse ecclesiae... Sicut se habet flumen ad mare, et exiguum munus in republica ad totam magistras tuum collectionem et summum principem: sic se habet priuata domus ad ecclesiam. Vt ergo qui flumen nauigare non audet, ei non est sulcandum mare: et qui in 2. minutioribus offieiis in ciuitate administrandis deprehensus fuerit minus satisfecisse, non promouebitur sapientum iudicio ad gubernacula totius reipublicae, siue ad supremum magistratum: Ira qui paruae familiae suae non sufficit, quomodo par erit gubernandae toti simul congregationi ? In qua cum administratio praecipua sit rerum spiritualium, sanem in iis errare vel negligentem esse, non vacat periculo. Illi ergo duntaxat qui domum suam bene administrandon specimen ediderit suae prudentiae, cura Ammonia baus. Simuuales.  \pend
\vspace{0.5cm}\noindent
\vspace{0.2cm}\rule{1cm}{0.2pt}\\ 
\hspace{0.2cm}\textit{mg}
\footnotesize ε Matth.6. 
\normalsize| 
\hspace{0.2cm}\textit{mg}
\footnotesize Tobus 4. 
\normalsize| 
\hspace{0.2cm}\textit{mg}
\footnotesize 
\normalsize| 
\hspace{0.2cm}\textit{mg}
\footnotesize 
\normalsize| 
\hspace{0.2cm}\textit{mg}
\footnotesize 
\normalsize| 
\hspace{0.2cm}\textit{mg}
\footnotesize 
\normalsize| 
\hspace{0.2cm}\textit{mg}
\footnotesize 
\normalsize| 
\hspace{0.2cm}\textit{mg}
\footnotesize 
\normalsize| 
\section*{IN I. EPIST. D. PAVLI. IN I. EPIST. D. PAVLI. }
\marginpar{[ p.80 ]}\pstart eccleliae committenda est. Huc pertinet illud Matth.25. Euge serue bone et fidelis, su per pauca fuisti fidelis, super multa te constituam intra in gaudium Domini tui. loseph prius probatus fuit in administratione domus Potipharis summi sacerdotis, quam el committeretur administratio regni Aegyptii: Genes. 4i.  \pend\pstart Nouitium reddidit interpres, vbi Graecis est κεόφυτα, quod, si etymon inspicimus, si gnificat nuper natum siue nuper insitum. Non ergo de eo qui iunior est annis, sermo est, quandoquidem interdum iuueniles anni habent adiunctam senilem prudenttam. et Timotheus adhuc iuuenis aetate fuit, vnde postea dicit:luuentutem tuam nemo despiciat: sed de iis id accipiendum est, qui nuper ex gentibus vel ludaeis per baptismum ecclesiae Dei insiti sunt, ac nondum in doctrina religionis satis sunt confirmati. Talem itaque nondum satis expertum, vel vt doctus sit, at non satis solide doctum Aposto lus in episcopum non vult eligi. Oportet prius essebonum et exercitatum discipulum quùm siat bonus praeceptor. Quidam Gregorius presbyter scripsit vitam Nazanzeni, ac narrat Cynicum philosophum Maximum nomine, qui curauit se eligi in epiAreNene scopum Constantinopolitanum. Sequitur ἀτιολογια: ne inflatus in condemnationem incidat calumniatoris. Solent tales subito ad dignitatem euecti inflari, et confidentes esse, sibi plus iusto tribuere, ad res nouas inducendas propensi: quasi ipsi soli in tanta multitudine idonei inuenirentur, et ipsi prae caeteris longo tempore exercitatis omnia forent mysteria consecuti: Vnde alios incipiunt contemnere, suo arbitrio omnia velle statuere, de quibusuis rebus statim definire: quo fit, vt in errorem subinde incidant. Qui se putat esse aliquid, seipsum seducit. Hinc verô, dum tales sunt, facile incidunt in condemnationem calumniatoris, id est, facile maleuoli sinistre de iis iudicant, imo et increduli habent quod rideant totam congregationem, vtpote in qua non sit alius idoneus ad tale munus nisi qui nuper primùm in societatem receptus e st. Et Synecdochice calumniator em dicit pro multis calumniatoribus. Velut alii, Theohylalactus, et Ambrosius, per calumniatorem intelligitur diabolus, qui vt est astutus artifex, facile potest neophytum irretire et adducere in errores. Dicitur diabolus αυτονοuxsuös Calumniator, et proprium estetymon.  \pend
\phantomsection
\addcontentsline{toc}{subsection}{\textit{Non nouitium, ne inflatus incidat in comdemnationem calumniatoris. }}
\subsection*{\textit{Non nouitium, ne inflatus incidat in comdemnationem calumniatoris. }}
\phantomsection
\addcontentsline{toc}{subsection}{\textit{Oportet autem illum et bonum habere testimonium ab extraneis, ne in probrum incidat et laqueum calumniatoris. }}
\subsection*{\textit{Oportet autem illum et bonum habere testimonium ab extraneis, ne in probrum incidat et laqueum calumniatoris. }}\pstart Vult prorsus inculpatum cuique de anteacta vita nihil probrosi possit obiici:ne extranei, id est, increduli, qui nondum Christo dederunt nomen, habeant quod in dedecus religionis exprobrent. Esset autem totius religionis ingens dedecus, si talis cerneretur ecclesiae praefectus, qui antea vel lenonem egerit, vel lapidator fuerit, aut aliud quiddam commiserit, ob quod ignominia foret notandus. Eam ob caussam Christiani olim, teste Aelio Lampridio in vita Alexandri Seueri, cum sacerdotes ordina re vellent, prius ipsorum nomina palam proponebant, hortantes et petentes, vt si quis haberet, quod posset illis publici criminis exprobare, in medium proferret et probaret. Vtinam hic mos laudabilis in hunc voes diem perdurasset: et haud dubie magna et multa scandala ab ecclesiis suissent depulsa. Oportet legitimum episcopum aeque apud hostes atque amicos bene audire. Vnde nunquam legimus Apostolos alicubi accusatos, tanquam fornicatores, adulteros, fures, vel homicidas, sed tan tùm seductores, quia videlicet docebant, quae hominibus, praesertim malis, parum placebant, abolentes omnem superstitionem ludaicam, et Ethnicam idololatriam. Caeterùm si danda opera, vt etiam ab exteris et aduersariis bonum habeatur testimonium: quanto magis aduigilandum vt bonum testimonium habeamus a domesticis et ipsis fidelbus? Profecto autem res magnas Apostolus ab episcopo ecigt  \pend
\vspace{0.5cm}\noindent
\vspace{0.2cm}\rule{1cm}{0.2pt}\\ 
\hspace{0.2cm}\textit{mg}
\footnotesize Quinam νεο ρυτοιnon sin admittendi ad episcopatum. 
\normalsize| 
\hspace{0.2cm}\textit{mg}
\footnotesize Galat6. 
\normalsize| 
\hspace{0.2cm}\textit{mg}
\footnotesize Suecdoche 
\normalsize| 
\hspace{0.2cm}\textit{mg}
\footnotesize Malus epr scopus eccle siɇ est dede cori. 
\normalsize| 
\hspace{0.2cm}\textit{mg}
\footnotesize Act.24. Lucaa-. 
\normalsize| 
\hspace{0.2cm}\textit{mg}
\footnotesize 
\normalsize| 
\section*{AD TIMOTHEVM CAP. (III. }
\marginpar{[ p.81 ]}\pstart sed si bene introspicimus, veras virtutes et vera ornamenta, quae sola magnum faciunt. Parum est alicubi episcopos quosdam ornari gemmatis mitris, auratis pedis, purpureis palliis, si haec omnameta eis desunt.  \pend
\phantomsection
\addcontentsline{toc}{subsection}{\textit{Ministros itidem compositos, non bilingues, non multo vino deditos, non turpiter lucri auidos, tenentes mysterium fidei, cum pura conscientia. Atque hi probentur prius, deinde ministrent, sic vt nemo possit illos criminari. }}
\subsection*{\textit{Ministros itidem compositos, non bilingues, non multo vino deditos, non turpiter lucri auidos, tenentes mysterium fidei, cum pura conscientia. Atque hi probentur prius, deinde ministrent, sic vt nemo possit illos criminari. }}\pstart Hactenus formauit episcopum: Nunc çadem opera praescribit, quales esse debeant qui ab episcopo ad ministrandum publice in ecclesia constituuntur. Rectissime au tem vbi de episcopo egisset, haec de ministris adiecit: tùm quia mnistrorum officia as finia sunt officio episcopi: tum quia a ministrorum conditione simul admonetur epi scopus, qualiter se gerere debeat. Si enim huiusmodi virtutibus praeditos esse decet ministros, at multo iustius est talibus ornatum conspici virturibus episcopum: quasi tacitum hic subesset argumentum a minori. Ac reuera sicut episcopis incumbit diligenter videre, quales ministros ecclesiis praeficiunt: ita prodest nonnunquam cos intueri in probam et sanctam vitam bonorum ministrorum, vt horum exemplo admoniti suam quoque vitam ipsi componanr. Caeterùm hoc loco animaduertendum, qualis in primitiua ecclesia, quae erat integerrima, fuerit eorum, qui praeerant eccles siis, ordinatio. Erantitaque episcopi, et erant eorum ministri, qui Graecis vocantur διάκονοι a verbo διακονίω, quod significat inseruio, ministro, obsequor, operam alicui nauo. Et Mercurius vocatus est διιακονα, quia minister Deorum est. Presby terorum nomen, vt Ambrosius ait, tam episcopis quam diaconis commune erat. Omnis episcopus, inquit, presbyter est, non tamen omnis presbyter episcopus est. Hic enim opiscopus est, qui inter presbyteros primus est. Apparet autem eos, qui ecclesiis praeerant, ab aetate dictos presbyteros, propterea quod olim non tere alii creabantur ab episcopis eccclessarum ministri, quam, qui sieut doctuina et moribus, ita aetate quoque erant venerandi, et ab omni suspicione mali immunes. Sic namque Graecis ωρευθήτης significat senem, vnde πρδνθυτσet senior, et propter aetatem honoratior. Porro quoniam Apostolus hic bis loquitur de diaconis siue ministris, colligere licet, quod ministri, slue διάκονοι, siue presbyteri, ab coiscopis instituti, in duos ordines distinguebantur. Atque id etiam ex aliis locis deprehendi potest. Primo enim, vt legitur Actor. 6. constituti sunt diaconi spectatae probitatis viri, et pleni spiritu Sancto ac sapientia, qui minsstrarent mensis, siue qui facultates ecclesiae distribuerent. Huiusmodi septem nominantur, qui et per impositionem manuum spiritum Sanctum accipiebant: inter quos celeberrimus Stephanus: tametsi hic quoque et disputasse et praedicasse Euangelium, per oc casionem et oportunitatem legitur. Atque hos quidem Apostoli ordinarunt, vt ministrarent mensis, ipsi vero minus impediti deprecationi et ministerio verbi liberius vacarent. Sed quia praeterea Apostoli non sufficiebant, vt omnibus locis semper concionarentur, vt baptizarent, et coenam dominicam peragerent: ideo opus erat alios quoque substituere, partim episcopos, partim δδακουςς qui ministerio verbi et sacramentorim quoque dispensationi intenderent, et hos quidem episcopis pares, nisi quôd ordinandi alios potestatem non habebant, vt Theophylactus att. Hi autem διιάκονοι, priusquam ad munus illud admitterentur, debebant in facris literis probe esse eruditi et exercitati, ut paulatim idonei fierent ad tantum docendi et sacramenta administrandi munus obeundum. Erudiebantur autem et exercitabantur et praeparabantur etiam a iuuentute in veteri testamento apud summos sacerdotes, vt indicat historia Samuelis, item apud Prophetas, sieut Heliseus in lericho suae fidei concreditos habebat, qui tùm vocabantur Nazaraei et filii prophetarum, item in Scholis, quales illae synagogae erant, de quibus Actor. sexto Libertinorum, Cyrenensium, Alexandrinorum, etc. In nouo autem Testamento exercitati sunt tales ab ipsis Apostolis, quod postea obseruatum est et ab  \pend
\vspace{0.5cm}\noindent
\vspace{0.2cm}\rule{1cm}{0.2pt}\\ 
\hspace{0.2cm}\textit{mg}
\footnotesize Oudes esse debeant ecelesiae miniMtri. 
\normalsize| 
\hspace{0.2cm}\textit{mg}
\footnotesize Aminori. 
\normalsize| 
\hspace{0.2cm}\textit{mg}
\footnotesize Ecclesiepn nitiuae integritas. 
\normalsize| 
\hspace{0.2cm}\textit{mg}
\footnotesize Mercurius λάινω dι. ctus. 
\normalsize| 
\hspace{0.2cm}\textit{mg}
\footnotesize Madanul tri ecclesia blim creati fint apudue teres. 
\normalsize| 
\hspace{0.2cm}\textit{mg}
\footnotesize Diaconis in ecllesia Hie rosolymitana 
\normalsize| 
\hspace{0.2cm}\textit{mg}
\footnotesize Diaconort nstitutio. 
\normalsize| 
\hspace{0.2cm}\textit{mg}
\footnotesize Scholae prophetarum 
\normalsize| 
\section*{IN I. EPIST. D. PAVLI. }
\marginpar{[ p.82 ]}\pstart episcopis. Vnde Tychicus Coloss.4 πιςος δάκονσκὰ σμύδιολω in domino ab Apofte lo nominatur: eundem ait se mississe Ephesum, 2. Timoth.4. Sic alii quoque quorum opera ipse subinde vtebatur, quique labores eius subleuabant alii describentes Epistolas, alii perferentes casdem, et consolantes ac confirmantes eos qui crediderant, sicut hic Timotheus saepe hac de caussa ab Apostolo emissus fuerat: et ipse alios commemorat tales discipulos suos Titum, etc. qui ministrabant ei incarcerato, qui indi; uulsse ei adhaerebant. In Bpistolis quae Ignatio ascribuntur dicitur, Stephanum fuis se diaconum lacobi, Timotheum et Linum Pauli, Anacletum et Clementem Petri. Ac praeclarum exemplum illud est, vt nunc quoque talis ratio in ecclesiam reuocetur, iuuenes idoneos in sacris literis instituendi, qui ecclesnjs postea praeficiantur. Ac primum de his, qui informant ad gubernandas ecclesias, qui sacris literis instituuntus vt verbi et sacramentorum dedantur ministerio, Oportet, inquit, esse ministros itidem compositos. σεμνος Graece legitur. Multa autem hoc nomine complectitur. σεμνος enim significat castum, venustum, seuerum, venerandum, pudicum, grauem, et apud Lucianum, religiosum, unde et σεμνα loca sacrata et ob religionem inaccessa. ltaque  uult Apostolus, eos qui ita instituuntur et ordinantur in presbyteros διιακόνος, ut pon pulum doceant in ecelesia, et sacra illa munia obeant, esse optimis integerrimisque moribus praeditos, uerecundos, pudicos, instar probarum uirginum, graues, candidos, ue nustos ore, uultu, oculis, denique in omnibus gestibus et actionibus decorum ingenuë obseruantes: vt quiuis permoueantur ad honeste de ipsis sentiendum et loquendum. Concluduntur breuiter hoc uerbo omnes uirtutes externae seu foris prominentes, quae in eo qui iam praeficiendus est ecclesiae requiruntur.  \pend\pstart Huiusmodi olim in ecclesia erat ordinatio. Primo loco constituebantur episcopi, summam autoritatem obtinentes. Proximi erant eorum ministri, qui Graecis uocantur διάκουοι. Presbyterorum nomen, vt Ambrosius ait, tam episcopis quam diaconis commune erat. Porro quoniam Apostolus hic bis loquitur de diaconis siue miDuplexordo nistris, colligere licet, ministros siue diaconos seu presbyteros ab episcopis institudiaconorum si tos, in duos ordines distingui solitos. Alter autem erat, de quo Actor. sexto legitur, ue presbyte quod creati sunt septem diaconi spectatae probitatis viri, et pleni spiritu Sancto ac sarorum. pientia, qui ministrarent mensis, siue qui facultates ecclesiae distribuerent. Atque hos creari voluerunt Apostoli, ne ipsi hoc ministerio impedirentur, sed possent deprecationi et ministerio verbi liberius vacare. De his ergo duobus gencribus diaconorum Apostolus sermonem instituit ac primùm de illis qui ecclesias debebant gubernare, qui uerbi et sacramentorum minsterio forent dediti. Huc pertinent et ver ba Ambrosii dicentis, quod in qualibet ciuitate sic ista omnia quae de diaconis hic di cuntur, accipienda sunt vobis, tanquam si ad vos proprie pertinerent, quippe qui in hoc datis operam sacris, vt consecuti doctrinam seu donum prophetiae, postea Αλόνος α a Deo et ab hominibus vocemini ad hoc di aconorum amplissimum munus.  \pend\pstart Non bilingues, uu διmλοyος.) Certe ethnici hos supra modum detestati sunt. Argumentum enim est leuis animi, facileque mouetur suspicio doli aut fraudis, aut mentis male sibi consciae, etiam perfidiae, vbi quis modo ait, modo negat, aliud hodie, aliud dicit cras. Sic Ciceronem perstringit Laberius, dicens, eum, duabus sedere sellis solitum. Et Satyrus vt est in Apologis, hominem contempsit et fugit, quem vidit ex eodem ore calidum spirare et frigidum. Ac dicti sunt duplices maplicesui? viri vulgo et in scripturis, qui lubrica et insyncera fide, verbis suis contraria dicunt et faciunt. Multa ad hanc rem in Adagiis referuntur. Omnium vero minime conueniet, eos promoueri in ministros verbi Dei, qui bilingues, id est, inconstantes, leues, varii, dictis non stantes inuenti fuerint. Quis enim talium doctorum orationi fidem est habiturus, quos nouit Autolici more candida in nigrum vertere, vel, ut Esaias cap.quinto, ait, bonum dicunt malum, et malum bonum, ponentes tenebras lucem, et lucem tenebras. Instar arundinis vento agitatae quidam vel quae stus vel gloriae caussa mutant sententiam, modo Christum et Euang. praedicantes, modon Mosen etlegem. Hodie quoque multi δiuoyes passim ecclessiam labefactant, ad recantandum  \pend
\vspace{0.5cm}\noindent
\vspace{0.2cm}\rule{1cm}{0.2pt}\\ 
\hspace{0.2cm}\textit{mg}
\footnotesize nu qud. 
\normalsize| 
\hspace{0.2cm}\textit{mg}
\footnotesize ri. 
\normalsize| 
\hspace{0.2cm}\textit{mg}
\footnotesize 
\normalsize| 
\section*{AD TIMOTHEVM CAP. III. }
\marginpar{[ p.83 ]}\pstart Pa quae recte docuerunt, quauis de causa parati. Sed tales hic damnant̃ ab Apostolo, et tota scriptura in eos inuehitur. Nam et Esaias dicit, Vae qui dicunt bonum malum, etc. Prouerbus octauo, dicit sapientia, Arrogantiam et superbiam, et viam prauam, et os bilingue detestor. lbidem capite 12. in tales multa.Et Eccl.6.item lacobus 3. grauiter detestatur eos, qui per linguam benedicunt Deo et patri, et per eandem maledicunt. liomini, qui ad similitudinem Dei est conditus. Exeodem ore procedit maledi, ctio et benedictio. Non oportet fratres mei haec ita fieri. Nunquid fons ex eodem foramine emtttit dulcem et amaram aquam? Nunquid potest ficus oleas gignere aut vi se tis ficus? Sic nullus fons salsam ac dulcem potest edere aquam. Tales proinde bilingues et in verbis vani et instabiles, non sunt idonei ad constanter docendam ueritatem, quae sempaer eadem est, et eodem modo tractari pure debet.  \pend\pstart Non multo vino deditos.) Vini quidem moderatus usus non damnatur, sed uino esse deditum et addictum. Certe qui frequenter et multum bibunt, etiam si non inebrientur (quae tamen fere semper solet vini ingurgitationem comitari) tamen minus idonei fiuntrebus gerendis, maximeque obtundunt ingenii sui aciem, quam oportet illaesam seruart iis quibus docendi munus incumbit. Prouerbus 31. prohibentur reges bibere intemperantius, ne forte inebrientur, inquit, et obliuiscantur iudiciorum, et mutent causam filiorum pauperis. En vinum inducit obliuionem recti iudicii sic etiam in ministris ecclesiarum euertit intellectum verbi Dei. Vnde et prophetae saepe contra ebrietatem et crapulam sacerdotum concionantur. Sed Amos 2.Dominus grauissime increpat eos, qui Nazaraeos delectos ad munus propheticum ad vinum bibendum prouocassent. Suscitaui, ait, de filiis vestris prophetas et de, iuuenibus vestris Nazaraeos. An non ita se habet res ? Vos autem fecistis Na-, zaraeos bibere vinum. Recte dixerunt quidam, ebrietatem esse voluntariam insaniam, eamque tandem si frequentior sis, abire in manifestum furorem. Alia multa quotidie sequuntur ebrietatem incommoda. Nam prouerbus 22. Cui vae? Cuius patri Vaec Cui rixae? Cui foueae? Cui sine caussa vulnera? Cui suffossio oculorum * Nônne his, qui commorantur in vino et student calicibus epotandis? Sed nunquam satis multa contra ebrietatem dici possent: maxime vero apud eos, qui uel iam educant instituunturque vt ecclesiis gubernandis reddantur idonei, vel qui iam amodo ad id munus sunt euecti. Verendum autem ne ex his nimium multi passim vitio hoc laborent.  \pend\pstart Non turpiter lucri auidos.) Qui tale de se praebet specimen, vt vel auidus lucri videat tali profecto non erit ecclesiae cura committenda. Est quidem permissa honesta ratic prospiciendi rei familiari, quae sine emptione, vendirione, commutatione, et similibus vlx bene procuratur: Verùm in his actionibus cauendum, ne turpiter lucro inhiare videa tur. Grit autem luerum multo turpissimum, si quid praetextu pletatis a subditis exigaem vel impetrerur. Sane non decet vllo sanctimoniae titulo aliquid praeter stipendium, qu modo sufficit familiae frugaliter et honeste alendae, accipere. Hinc 1. Thess.2.auaritiae etiam suspicionem omnem a se reiicit Apostolus, et maluit laborare, quam exigere ab illis, quod iure poterat. Non abs re2.reg.5. Gechazi puer siue discipulus Helisaes percus sus est lepra, quod a Naaman Syro argentum et vestes mutatorias iniussus accepisset. Ac reuera nostro hoc seculo compertum est, nulla resic irritatos omnium animos ad nocendum sacerdotibus, ad diripiendas templorum et Coenobiorum opes regias, ata que sacerdotum manifesta auaritia.  \pend\pstart Tenentes mysterium fidei cum pura conscientia.) Ab aliis quoque vitiis alienos esse decet hos diaconos, non minus quam episcopos. Sed satis est illa tanque principalia enumerasse, quae etiam a quibusdam vix pro peccatis habebantur. Nunc requirit, quod praecipuum est, et per quod doctrinae puritas conseruatur. Nam vt fortassis naeuus aliquis conspiciatur in diaconi cuiuspiam moribus, est tamen danda opera, vt doctrina semper inco lumis conseruetur. Vult itaque diaconos tenentes mysterium fidei esse cum pura conscientia, id est, qui ipsi in fide sint confirmati, et omnem doctrinam fidei probe teneant ac constant tueant? deinde qui omnia syncera, recta, coram Deo et hominib, agant conscientia. Idem sibs  \pend
\vspace{0.5cm}\noindent
\vspace{0.2cm}\rule{1cm}{0.2pt}\\ 
\hspace{0.2cm}\textit{mg}
\footnotesize Ftnu obest omnibus studiis et actionibus. 
\normalsize| 
\hspace{0.2cm}\textit{mg}
\footnotesize Solon iudica uit principem bruomn mor Le mulctanlumm. 
\normalsize| 
\hspace{0.2cm}\textit{mg}
\footnotesize Etrietas isntauolmtaria 
\normalsize| 
\hspace{0.2cm}\textit{mg}
\footnotesize uj dxsiupdes anat. 
\normalsize| 
\section*{IN I. EPIST. D. PAVLI. }
\marginpar{[ p.84 ]}\pstart volunt haec verba, quod illa in primo capite: Milites bonam militiam, habens fidem et bonam conscientiam, qua repulsa nonnulli circa fidem naufragium fecerunt. In ean dem fere sententiam dictum erat Actor.6.eligendos viros spectatae pietatis, plenos spi ritu Sancto et sapientia. Igitur mysterium fidei, id est, intima religionis nostrae arcana, praecipueque doctrinam illam de iustificatione perfidem, quae vere mysterium dicitur, quippe humani ingenii captum excedens, ii plene et plene cognoscere, tanquam di gitos suos debent quibus ecclesiae cura committetur. Non satis est vulgariter eos t in religione esse doctos, sed oportet excellere omnes, et omnia mysteria eos tenere, vt iuxta illud i. Petri3. semper parati sint ad respondendum cuilibet petenti: vt loquantur de ea, quae in ipsis est, spe, cum mansuetudine et reuerentia conscientiam habentes bonam. Sed iidem habere debent puram conscientiam, id est, vere pie omnoi agere, non fucate, non hypocriticem, non inuite, aut spe tantum lucri: sed serio, ex animo, propter Christum, cum fide et timore, solum Deum remuneratorem ac iudicem ob oculos sibi ponentes. Sic autem dictum est et in primo capite huius Epistolae: Finis praecepti est charitas, ex puro corde, et conscientia bona, et fide non simulata. Porro quoniam non statim admittendi sunt ad munus docendi, qui satis docti apparent, quique purita tem conscientiae ad tempus prae se ferunt:nam et astute quidam norunt pietatem omnem simulare:ldeo addit conditionem summopere necessariam, dicens:  \pend
\phantomsection
\addcontentsline{toc}{subsection}{\textit{Atque hi probentur prius, deinde ministrent, sic vt nemo possit illos criminari. }}
\subsection*{\textit{Atque hi probentur prius, deinde ministrent, sic vt nemo possit illos criminari. }}\pstart Si huius praecepti habita fuisset diligens ratio, non tot lupi passim sub pelle ouina la tentes in ouile Domini irrupissent, a quibus Christus Matth.7. et Apostolus Act. 20. cauendum monuerunt. Sane si nemo pater familias seruo aliquid serii commitet in familia, nisi prius fidem eius satis explorarit: at multo magis necessarium est, eos probari et variis modis tentari, quibus committenda est cura totius sanctae congregationis. Atque hac caussa noluit Apostolus admitti ad Episcopi functionem hominem vεόφυιον seu nouitium: quia talis nondum satis est probatus. Oportet autem et longo tempore probari, et nonnunquam caussis grauibus et periculosis, ( de quo 2. ad Timoth.4.) et quidam alii diu fuerunt fideles discipuli Apostoli Pauli, et haud dubse multum in mysterio fidei sunt consecuti: sed tamen non fuisset vtile tales praeficere ecclesiis, quia nimirum nondum satis fuerant probati, cito enim in aduersis defecerunt, Eam ob caussam in hac Epistola cap.5. Paulus monet Timotheum, ne cui manus cito imponat, neque communicet peccatis alienis. Qui enim eum praeficit ecclesiae, quem non satis probauit et explorauit, videtur, dum talis postea palam delinquit, peccatis huius fauere. Vbi ergo probati et explorati qui fuerint, tum ministrent, inquit, atque ita ministrent, vt nemo possit illos iure criminari. Quin mali eos criminentur vspiam, fieri non potest: at potest fieri, vt nequaquam iure criminentur. Neque vero vlli tam pii, tam sancti, qui nusquam sint culpabiles: possunt tamen ita viuere, vt a manifestis criminibus, et illis omnibus non crassis tantùm, sed quae aliquo modo deformant epi scopi functionem, sibi temperent: ne vllo pacto offendiculum praebeant infirmis, vel etiam malignis ansam obloquendi.  \pend
\phantomsection
\addcontentsline{toc}{subsection}{\textit{Uxores similiter modestas, non calumniosas, sobrias, fidas in omnibus. }}
\subsection*{\textit{Uxores similiter modestas, non calumniosas, sobrias, fidas in omnibus. }}\pstart Non paruam esse rem, episcopum agere, hic satis patet. Non satis est ipsum episcopum probum esse et spectare fidei: sed oportet insuper pari probitate, vxorem, liberos, denique omnes familiares esse praeditos. Quam ob caussam Apostolus non obiter quidem aut παρόργικῶς, sed vt rem valde necessariam adiicit, quales esse deceat tùm episcoporum tùm diaconorum vxores. Fit plerunque , praesertim hoc tempore, cùm vulgus etiam minutissima queque rapit ad calumniam, vt etiam si doctor Euangelicus syna  \pend
\vspace{0.5cm}\noindent
\vspace{0.2cm}\rule{1cm}{0.2pt}\\ 
\hspace{0.2cm}\textit{mg}
\footnotesize a seculis. Si couss. Euang pra dicationem uocat mys rium quod re conditum fa 
\normalsize| 
\hspace{0.2cm}\textit{mg}
\footnotesize t. Petris. 
\normalsize| 
\hspace{0.2cm}\textit{mg}
\footnotesize A minori 
\normalsize| 
\hspace{0.2cm}\textit{mg}
\footnotesize Demas. 
\normalsize| 
\hspace{0.2cm}\textit{mg}
\footnotesize Probation nistrorum ne cessaria. 
\normalsize| 
\hspace{0.2cm}\textit{mg}
\footnotesize Domus epi scopi quali esse debe
\normalsize| 
\section*{AD TIMOTHEVM CAP. III. }
\marginpar{[ p.85 ]}\pstart cere doceat, et synceriter iuxta viuat, nihilominus statim omnes obloquantur et sinistre iudicent, si ei domi sit vxor parum bene morata. Et totam familiam ipsius a vitiis immunem esse decet, qui volet a vitiis purgare ecclesiam. ( Ne quis putet secundum Theophylactum et Primasium, Apostolum hicloqui de dlaconissis, quales creari Cataphrygae haeretici volebant, contra quos Ambrosius.) Qui ne vxorem quidem suam potest ad temperantiam et modestiam adducere, quomodo suadebit am plectendas virtutes infreni multitudini? Sed periculum insuper est, ne vir pius, cui vxor sit immodesta, calumniosa, perfida, ipse quoque ad vitia quaedam per eam alliciatur. Certe subinde viri existimatio nonnihil imminuitur, ob improbitatem mulleris: sicut e confrario probitas viri commendatur, cui et vxor proba contigerit. Socrati vitio versum fuit, quod domi morosam Xantippen aleret. Eccles.25. Beatus vir, qui habitat cum muliere sensata. Et cap.26. Mulieris bonae beatus vir. lbidem, pars bona mulier bo na, in parte bona timentium Deum dabitur viro pro factis bonis. His itaque de caussis merito praescribit Apostolus, quales episcoporum ac diaconorum vxores esse deceat.  \pend\pstart Vxores, inquit, similiter modeslas, σεμνας.) Paulo ante de miniitris loquens, vertit compositos. Oportet igitur castas, pudicas, graues, compositas, bene moratas esse: nam haec omnia ea voce significantur. Vnde Eceles.26.Gratia super gratiam mulier sancta et pudorata.lbidem: Sicut Sol oriens in mundo in altissimis Dei, slc et mulieris bonae species in ornamentum domus eius. Extatnon contemnenda descriptio mulieris probae apud Plautum in Amphitrione. Sic ibi quaedam cum aceibius obiurgaretur, lo quitur: Non ego illam mihi dotem duco esse, quae dos dicitur: sed pudicitiam et pudorem, et sedatam cupidinem, Deum metum, parentum amorem, et cognatum concordiam: atque vt munifica sim bonis, prosim probis. En talem eam volebant ethnici, quae modestae et probae matronae nomen ambiret. Multo igitur magis talem eam esse decet, quae vxor erit episcopi seu diaconi.  \pend\pstart Non calumniosas:) λυδιαθολος: Quae calumniosae sunt, lites, rixas, odia serunt: quare necessum est talis mulieris improbirate, per quam multi alii commoueri queant, perturbari etiam maritum episco pum. Familiare quidem vid etur vitium mulierculis calumniari et cauillari:at nullo pacto ferendum in vxore eius, qui palam omnes docere debet, quae beneuolentiae sunt et Charitatis. Et verendum, ne vxor calumniatrix mari tum quoqu calumniatorem efficiat. Eecles.25. Sicut ascensus arenosus in pedibus veteras ni: sic mulier linguata homini quieto.  \pend\pstart Sobrias, νηφαλιος :) Id prius vertit interpres vigilantes, cùm episcopis loqueretur Sobrietatis igitur nomine non tam cibi potusque temperantia venit intelligenda, quam vigilantia cireumspectioque in omnibus actionibus. Quae et dicit et facit omnia attente, circumspecte, moderate, ea sobria esse dicitur. Antea de episcopo loquens, sobrium uertit, ubi Graece habebatur σώφρινα.  \pend\pstart Fidas inomnibus ) 4. Verbis et factis. Fidas esse decet et in causa Coniugii, ne uidelicet adulterii uel suspicione notentur: et in rerum domesticarum tractatione, ne uidelices dilapident uel prodigant inutiliter rem famillarem. Vel in altero peccare ingens uitium est, acnotae etiam inde nascuntur in marito: qui sentit se mulieris perfidia, et nomine siue existimatione sua, et re etiam tandem spoliari. In summa, omnis probitas, pudicitia, sobrietas, modestia, tam in uxoribus ministrorum uerbi, quam in ipsis ministris requiritur: Vt, sicut hi uiris, ita illae mulieribus sint probitatis norma et exemplum. lmô uero et in omnibus omnium uirorum uxoribus eadem uitae puritas merito re quiritur. Mirum est, quod Theophylactus et Primasius hoc loco adnotarunt, ista de mulieribus dicta, accipienda esse de iis, quae ministrandi officio in ecclesia fungerentur, quales in ecclesia orientali, et in iure pontificio dictae sunt diaconissae: quasi uero, ut Cataphrygae senserunt, ( quod Ambrosius refert et refutat) mulieres ad talia munera in ecclesia sint admittendae: cum tamen Paulus huius Epistolae capite 2. et 1. Corinth. 14. mulieribus docere non permittat, neque autoritatem ullam in ecclesia usurpare. Prior itaque interpretatio nostra simplicior est, et Ambrosio probata.  \pend
\vspace{0.5cm}\noindent
\vspace{0.2cm}\rule{1cm}{0.2pt}\\ 
\hspace{0.2cm}\textit{mg}
\footnotesize vxoredum niosa pestilens aninal. 
\normalsize| 
\hspace{0.2cm}\textit{mg}
\footnotesize Sobrietas iu primis decet mulierculas. 
\normalsize| 
\hspace{0.2cm}\textit{mg}
\footnotesize Fides conuer. 
\normalsize| 
\section*{IN I. EPIST. D. PAVLI. }
\marginpar{[ p.86 ]}
\phantomsection
\addcontentsline{toc}{subsection}{\textit{Diaconi sint vnius vxorus mariti, qui liberis recte praesint et propriis familiis. Nam qui bene ministrauerint, grandum sibi bonum acquirunt et multam libertatem in fide, quae est in Christo lesu. }}
\subsection*{\textit{Diaconi sint vnius vxorus mariti, qui liberis recte praesint et propriis familiis. Nam qui bene ministrauerint, grandum sibi bonum acquirunt et multam libertatem in fide, quae est in Christo lesu. }}\pstart Nunc de eo genere diaconorum agit, qui ministrabant mensis, seu distribuebant ecclesiae facultates. Ambrosius obster hoc loco notat, in unaquaque ciuitate debere septem esse huiusmodi diaconos, presbyteros siue alterius speciei diaconos, duos: episcopum uero unum. In his autem requiri omnes propemodum ulrtutes, quae in episcopis requiruntur, certum est:sed tamen, ut sint episcopis inferiores uita, sicut sunt inferiores functione, conuenit nihilominus eos esse unius uxoris maritos, et, ut uno uer bo dicam, bonos probosque patres familias.  \pend\pstart Sint vnius vxoris mariti. )lgitur non scortatores, non adulteri, non item ludaica intens perantia plures uxores habentes, sed una honesta et pudica contenti.  \pend\pstart Qui liberis recte praesint, et propriis familiis.) Neque enim uerosimile est, eos recte adminutraturos in distributione facultatum ecclesiae, qui proprias facultates non recte dispensarint in sua famslia. Neque multitudinem redigent in ordinem, qui domi non sciunt imperare paucis liberis. Doctrinam in his non admodum requirit, sed praecipue prudentiam oeconomiae usu comparatam: ut colligi hinc possit, tales diaco nos fuisse uiros graues, quique iam diu egerant patres familias. Eccles.4. Noli esse sicut leo in domo tua, euertens domesticos tuos, et opprimens subsectos tibi. Disciplina senper colenda inter liberos est. Prouerbus 13. Qui parcit uirgae, odit filium suum. Et capag. Virga atque correctio tribuit sapientiam, puer autem, qui dimittitur uoluntati suae, con , fundit matrem suam. Deut.21.filius contumax et inobediens iubetur lapidari. Additueluti per occupationem argumentum hortatorium, ut diligentes sint. Nam ubi forte di xissent quidam, munus istud satis abiectum esse, ut nemo id appetere debeat, aut magnopere curare, respondet nequaquam ita esse.  \pend\pstart Nam qui bene minissrauerint, gradum sibi bonum acquirunt.) Quasi diceret: Non  est tam leue, tam humile istud munus, quin qui ex animo et diligenter istud subibunt, * magnum pietatis profectum possint facere. Habet enim suos gradus et incrementa pietas in quocunque officii genere, et spiritus vnicuique dona sua ad mensuram fidei largitur, vt dicitur 1.Corinth.12. Qui diuidit peculiariter unicuique sicut vult: et Matth. 25.Omni habenti dabitur et abundabit: qui vero non habet, etiam id quod habet, auferetur abeo. Illud de auctione donorum intelligitur. Vel bonum gradum sibi acquirunt: quia dum probe et strenue gerunt se in hoc munere, specimen tale de se praebent, vt, si contingat eos illustrari dono prophetiae et δηιδιακτικος fieri, ad sublimiorem gradum, nempe alterius diaconatus siue tandem episcopatus possint promoueri, et sic ecclesiae semper magis ac magis prodesse. Vnde in parabola Matth. 25. Euge serue bone et fidelis, super pauca fuisti fidelis, super multa te constituam, intra in gaudium Domini tui. Quemadmodum in ciuili republica qui aedilem bene egerit, tandem eligitur in praetorem, ex praetore tandem fit tribunus aut consul: lta fieri potest, vt qui diaconum bene egerit in ministerio mensae, tandem vocetur ad ministerium verbi et sacramentorum, denique et ad episcopatum. Nam hi tres ordines duntaxat olim erant in ecclesia, quantum ex scripturis et etiam ex Am brosio colligitur.  \pend\pstart Et multam libertatemin fide, quae est in Christo Iesu.) Graece est πολιιύ παρ́ρησίαυ: interpres reddidit, multam libertatem, vt quodammodo sit grata Antithesis in vocibus, ministrauerit, et libertatem:imon subesse videtur occupatlo: quasi significaretur munus illud, quia ministerium vocatur, nequaquam seruile esse, licet etiam graue sit distribuere pauperibus, inter quos semper sunt querull et morosi, quibus difficillime fatisfacias: sed potius liberale esse, et solis liberis hominibus dignum, ac per id libertatem multam, spiritualem videlicet, comparari. Ambrosius vertit, multam fiduciam in fide, vt sit sensus: certam et indubitatam mercedem iis esse apud Deum repositam, ac propte rea nequaquam ceslandum in eo munere, autid ceu humile contemnendum. Disci\pend
\vspace{0.5cm}\noindent
\vspace{0.2cm}\rule{1cm}{0.2pt}\\ 
\hspace{0.2cm}\textit{mg}
\footnotesize DonoN n crementuncam tingit piis ac fidelibus
\normalsize| 
\section*{AD TIMOTHEVM CAP. III. }
\marginpar{[ p.87 ]}\pstart mus ergo hic, munera illa ecclesiastica summo in honore habenda, non solùm ratioa ne ministerii, quod tamen praecipue debet nos mouere, sed etiam ratione personarum illa subeuntium, et ad praescriptum Apostoli omnem vitam suam instituentium. Atque  admonentur interea omnes Illi, vt ad haec verba Pauli ceu speculum respicientes, vitam moresque componant, vt homines videntes opera eorum bona, glorificent coelestem Patrem, nec non imitari eandem vitae innocentiam conentur.  \pend
\phantomsection
\addcontentsline{toc}{subsection}{\textit{Hec tibi scribo, sperans fare, vt veniam ad te cito: quod si tardius venero, vt noris quomodo oportet in domo Dei versari, quae est ecclesia Dei viuentis, columna et slabilimentum veritatis. }}
\subsection*{\textit{Hec tibi scribo, sperans fare, vt veniam ad te cito: quod si tardius venero, vt noris quomodo oportet in domo Dei versari, quae est ecclesia Dei viuentis, columna et slabilimentum veritatis. }}\pstart Conclusionis loco ista in hacparte de vita episcoporum domestica seu priuata esse possunt, quasi diceret: Haec volui in praesentiarum de vita priuata episeopt, ac de tota oidi natione ecclessastica ad te scribere. Non quidem quod non sim te posthac inuisurus: nam spero per gratiam Dei breui istic adesse: sed si fortasse praeter spem differatur me? aduentus, habeas interim certam hanc formam de ordinatione ecclesiastica quam tuto sequaris. ltaque ne Timotheus discipulus nimis angeretur desiderio videndi magistri, commode haec verba subiecit Apostolus. Etre uera omnia perstrinxit Apostolus abum dem satis his paucis, quaeae ad ecclesiae ordinationem rectam requiruntur: vt his instructus Timotheus ecclesiam rite administrare potuerit. Vnde discimus, parum recte eos facere, qui vel plures, vel, quod deterius est, prorsus alias ordinationes in ecclesias volunt introducere. Vtinam his paucis contenti, saltem paucula haec plene obseruaremus. Si quibus parua res videtur, habere in ciuitate vna unum episcopum, ministros duos, qui verbum et Sacramenta distribuunt: deinde ad summum septem alios ministros, qui ministrant mensis: At certe non est parua res, hos habere tales, quales hic ab Apostolo describuntur. Hic vero omnis fastus et pompa tot inutilium ministrorum et graduum, quos superbi quidam homines excogitarunt prorsus concidit. Rectem autem dicit Apostolus se sperare adire Timotheum, quasi dubitans. Nam semper voluntate Dei agebatur nec also ibat, quam quo spiritus ipsum ducebat seu mittebat. Alias impediuit eum Satanas, ne inuiseret Thessalonicenses: et saepem habuerat in animo inuisere Romanos, sed saepe fuit impeditus. lraque tanquam ambigens, et ex voluntate Dei totus pendens, de aduentu suo nil certi statuit lnsigni tamen charitate permotus et solicitudine eccle siarum, voluit Timotheum de omnibus, quae ad munus episcopi attinent, per literas cer tiorem facere. lacobus 4. Si Dominus voluerit, et si vixerimus, faciamus hoc aut illud. Neque venit Paulus Ephesum, sed Romae captiuus postea secundam scripsit ad Timotheum ex vinculis.  \pend\pstart In domo Dei, quae est ecclesia Dei viuentis.) Per μἐιωνυμίαν domum Dei vocat ecclesiam quemadmodum ipse interpretatione exponit: Tota congregatio credentium quasdomum vnam efficit, dum videlicet singuli credentes sunt velut singuli lapides, e quib parieres ac tectum, sicque tota domus spiritualis, quam Deus dignatur inhabitare, effici tur. Ac Christus huius aedificii lapis est angularis parietes connectens in vnum: fundamentum quoque et apostoli similiter et prophetae fundamenti lapides sunt I.Petri.2. Gustastis quod benignus sit Dominus, ad quem accedentes, qui lapis est viuus, ab ho minibus quidem reprobatus, apud Deum vero electus ac pretiosus: ipsi quoque veluti viui lapides aedificemini, domus spiritualis, sacerdotium sanctum, ad offerendum spirituales hostias, etc. Ephes.2. Vos estis conciues sanctorum ac domestici Dei, superstructi super fundamentum Apostolorum ac Prophetatum, summo angulari lapide ipso lesu Christo, in quo quaecunque structura coagmentatur, ea crescit in templum sanctum in Domino, in quos et vos coaedificamini in habitaculum Dei per spiritum. Etin Psalmo. de hac domo dicitur: domum tuam domine decet sanctitas inlongitudinem dierum. Vocat vero ecclesiam Dei viuentis, simul quia ex viuis lapidibus çonstat:nam et Christus et credentes a Petro dicuntur viui lapides: simul, quia proprie est domus Dei, qui est Deus viuorum, et autor vitae. Deinde per amplificationem eam exornat aptis Meta? phoris, quasi epithetis: cùm dicitiecclesia esse columnam et stabilimentum veritaig Re\pend
\vspace{0.5cm}\noindent
\vspace{0.2cm}\rule{1cm}{0.2pt}\\ 
\hspace{0.2cm}\textit{mg}
\footnotesize 2Thesαα Rom.i. 
\normalsize| 
\hspace{0.2cm}\textit{mg}
\footnotesize Eeufador mus det que modo. 
\normalsize| 
\hspace{0.2cm}\textit{mg}
\footnotesize I. Petriι 
\normalsize| 
\hspace{0.2cm}\textit{mg}
\footnotesize Ephesa 
\normalsize| 
\hspace{0.2cm}\textit{mg}
\footnotesize Psal.95.. Sunt Metaphorae trattae ab aedife cis. 
\normalsize| 
\section*{N I. EPIS T. D. PAVLI }
\marginpar{[ p.88 ]}\pstart cte autem Columna et stabilimentum ueritatis congregatio fidelium appellatur, quia in hac ueritas pure docetur et colitur, et ipsa uicissim per ueritatem confirmatur. Veritas enim Christus est cui ecclesia innititur, et qui ecclesiam semper praedicat fundatam supra firmam petram, quam nullus descendens imber, non venientia flumina, non flantes venti deiicere possunt, sicuts ostenditur in parabola Matth.7. Credentes siquidem nullis errorum turbis, nullis procellis persecutionum, nullis diaboli insultibus commouentur: imon ne pertae inferorum quidem aduersus congregationem credentium praeualebunt. Itaque congregatio illa quamdiu in fide persistit, quamdiu pro veritate inconcussa et infracta stat contra omnem falsitatem, columna et stabilimentum veritatis dicitur. Porro cum in ecclesia primum locum obtinent episcopi, admonent̃ hisce verbis, vt se strenue gerant in hoc spirituali aedificio, et se quoq meminerint quasdam columnas esse et fulcra infirmarum conscientiarum, se quoque pro veritate stabilienda tuen daque in omnibus periculis debere persistere. Sic autem et alibi dixit Apostolus permetapho ram, Petrum et loan. columnas ecclesiae: Galat.2. Cephas et loannes qui videbant colum nae, dexteras dederunt mihi ac Barnabae societatis. Nonnulli hic putant tacitam allusionem subesse ad templum Hierosolymitanum siue ad synagogam, in qua figurae atque vmbrae suum habebant locum:verùm in domo Domini, quae est ecclesia, locum habent res atque ipsa veritas: Ac proinde, quanto res praestant vmbris, tanto ecclesia eminentior et praestantior censeri debet ipsa Synagoga. In domo enim Dei non vmbrae bonorum ostendunt, sed ipsemet Christus, qui est veritas, praedicatur: et extra eam domum sunt merae tene brae et errores.  \pend
\phantomsection
\addcontentsline{toc}{subsection}{\textit{Et citra controuersiam magnum est pietatis mysterium. Deus manifestatus est in carne, iustificatus est in Spiritu, visus est Angelis, praedicatus est gentibus fides illi habita est in mundo, receptus est in gloria. }}
\subsection*{\textit{Et citra controuersiam magnum est pietatis mysterium. Deus manifestatus est in carne, iustificatus est in Spiritu, visus est Angelis, praedicatus est gentibus fides illi habita est in mundo, receptus est in gloria. }}\pstart Quoniam dixit ecclesiam esse columnam et fulcimentum veritatis, ideo per affectum suo more adiicit argumenta sex, ad comprobandum Euangelicam veritatem, quae est in ecclesia, haber̃ et colitur, certam esse et indubitatam (sic enim solet aliquando vbi res excellentes occurrunt, pluribus verbis, et quasi digressione, eas vtcunque illustrare:) addit, inquam, ista, ad significandum, quaenam sit illa veritas, quae in ecclesia manifestatur, retinetur et colitur. Vtitur autem verbis ad asseuerandum valde idoneis. Citra controuersiam, inquit, magnum est pietatis seu Veritatis mysterium, et tam euidens, vt quilibet confiteri debeat, ita rem se habere. Pietatem autem hic pro veritate decet accipi autore Ambrosio. Proinde haec est illa veritas, quam nemo potest negare, et cui tuto innitit ecclesia, eam senperiactans et praedicans: hoc est illud amplum et augustum veritatis mysterium et summarium totius religionis nostrae, quod antea semper occultum fuit, nunc autem nemo credentium non clare intelligit, nempe, quod Deus manifestatus estin carne. Sane his verbis orium vocat, quia antea semper latuit, saltem nunquam sic plene innotuit, vti nunc per Euangelium reuelata est. Sex autem rationibus veritatem hanc confirmat approbat́ vt merito certissimum esse debeat, veritatem de salute nostra in ecclesia retineri et praedicari. Primo,  \pend\pstart Deus, inquit, manifestatus estin carne.) ld quidem verissimum est. Nam, ut dicitur Phi lipa. Qui in forma Dei erat, formam serui accepit, et habitu inuentus ut homo. Ioan.1. Verbum caro factum est. 1. Ioann. 1. Quod erat ab initio, quod audiuimus, quod uidimus oculis nostris, quod perspeximus, et manus nostrae contrectauerunt de sermone uitae, et uita manifestata est, et uidimus, etiam testamur et annunciamus uobis uitam aeternam, quae erat apud patrem, et manifestata est nobis. Itaque palam manffestata estueritas de Christo in ecclesia, cùm uidelicet Christus natus est, et diuinitatis es ius gloria multifarie multisque modis palam est declarata uniuerso mundo, ut per eum lalutem consequeremur.  \pend\pstart Secunda ratio manifestatae ueritatis, quae docetur in ecclesia, Iustificatus estin spiritx.) Id est, per Spiritum sanct. declaratus est uerus Deus, qui sie  \pend
\vspace{0.5cm}\noindent
\vspace{0.2cm}\rule{1cm}{0.2pt}\\ 
\hspace{0.2cm}\textit{mg}
\footnotesize Ioan.14. Matth.16. Ecclesiae si tilltas et d stantia Eccles. qu modo dica tur stabilim tun uerita tis. 
\normalsize| 
\hspace{0.2cm}\textit{mg}
\footnotesize In ecctesgt Christ sunt res et uert tus, umbra chorterint 
\normalsize| 
\hspace{0.2cm}\textit{mg}
\footnotesize Comproba touerttatis Euangelice 
\normalsize| 
\hspace{0.2cm}\textit{mg}
\footnotesize 1.Argum. al siesoscse. 
\normalsize| 
\hspace{0.2cm}\textit{mg}
\footnotesize 2. argatef monio Sp.S. t. Corinth.1. 
\normalsize| 
\hspace{0.2cm}\textit{mg}
\footnotesize ygitas d Christo ma nifestata i ecclesia. 
\normalsize| 
\section*{AD TIMOTHEVM CAP. III. }
\marginpar{[ p.89 ]}\pstart lustus, imon ipsa iustitia et iustificatio nostra, per quem vnum ex peccatoribus iusti red duntur, quotquot in eum vere credunt. Sic alibi quoque iustitia et innocentia Christi nec non potestas eius iustificandi alios per spiritum Sanctum tum scripturis, tum apertis signis comprobata est. Huc pertinent illa Rom.3.Proposuit Deus Christum reconciliatorem per fidem, interueniente ipsius sanguine, ad ostenslonem iustitiae suae propterre, missionem praeteritorum peccatorum quae Deus tolerauit, ad ostendendum iustitiam, suam .. in praesenti tempore, in hoc vt ipse sit iustus, et iustificans eum, qui est ex fide lesu.ltem Spiritus sanctus idem de Christo testatus est apertis signis. Vnde Io.1.Vidi spiritum, inquit Baptista, descendentem specie columbae de coelo, et mansit super eum, et ego non noueiom eum: Sed qui misit me, vt baptizarem aqua, is dixit mihi: super quem videris spiritum descendentem ac manentem super eum, hic est, qui baptizat S.S. Et ego vidi, et testatus sum, hunc esse filium Dei.ltem Rom.1.qui declaratus fuit filius Dei per potentiam et spiritum sanctificationis. Atque haec simplieissima videt interpretatio, maximeque accommoda ad explicandum myiterium veritatis, de quo verba Apostoli decet accipi. Nam verem amplissimum est hoc mysterium, et haec est indubitata veritas, seu pietas totius nostrae religionis: quod videlicet Deus manifestatus est in carne, quando Christus nostra caussa dignatus est homo fieri: et quod idem iustificatus est in spiritu, quando iustitiam suam nostram fieri voluit, ac nos ipse qui sua natura iustissimus erat, iustos red didit, id quod Spiritus sancti testimonio abunde est comprobatum, vt nullus non possit clare intelligere. Multi durius hunclocum interpretantur. Nonnulli de assumpta humanitate intelligi illa volunt: ita, vt, quia Deus et homo vnus est Christus Deus, etiam iustificatus esse dicatur, non quidem in abstracto siue in substantia, (alioqui Deitas ipsa iu stificata diceref, quod absurdum est) sed in natura assumpta, videlicet humana: cui naturae Deus pater addidit omnem perfectionem, ita vt in Christo inhabitet, vt dicitur Coloss.2.omnis plenitudo Deitatis corporaliter, et esset etiam in quantum homo, ple nus gratia et veritate.  \pend\pstart Tertia ratio veritatis, quae in ecclesia retinet et praedicatur, haec est: Visus est Ang lis, quasi dicat: Vera esse quae docentur in ecclesia, de Christo et salute per eum parta id etiam Angelorum testimonio est comprobatum. Nam Angeli quoque illum vt Dominum et salutis autorem sunt venerati, et mysterium salutis nostrae in eo agnouerum Huc pertinet quod Christi natiuitas per Angelum, Gabrielem videlicet, est praenuncit ta. Luc. 1.Christo nato Angelorum multitudo cecinit in coelo gloriam Deo, ad pastores nuncium tulit. Luc. 2. Christo post tentationem diaboli victori Angeli ministrarunt: Matth. 4. oranti in horto antequam caperetur, apparuit Angelus de coelo, corroborans eum: Luc. 22. Postquam resurrexit, Angeli de eo testrmonium ferebant mulieribus ad sepulchrum: Luc.24. Denique ascendente eo in coelum angeli Apostolos al locuti sunt Actor.1.  \pend\pstart Quartum argumentum veritatis, Praedicatus est gentibus: a testimonio Apostolorum Nam et Prophetae antea multa de hac veritate vaticinati fuerunt, quodque ea gentibus, id est, omnibus populis et nationibus foret manifestanda: et nunc demum Apostoli toti mundo eam annunciarunt: ldeo Simeon Luc.2.Christum vocauit Lumem ad reuelationem gentium. Luc24. Christus iubet praedicars in nomine suo poenitentiam acre missionem peccatorum in omnes gentes.  \pend\pstart Tiaes Ini habita estin mundo: Quintum argumentum a fide eorum, qui audierunt; quasi dicat: si non esset veritas quae traditur ab ecclesia, quifieri posset, vt tam multi es assentirent ? Homines enim feruntur ad consentiendum iis, quae sunt vera.lgitur quod uera fuerint, quae praedicatores illi docuerunt, hoc indicio est, quod statim quam pluri mi fidem ei doctrinae adhibuerint. Mirum est quôd homines illa credere potuerint de Christo, ad quae alioquin nullis armis, nullis rationibus humani ingenii induci potuissent: sed Deus fuit efficax in pectoribus eorum, ut sic ueritati assentirent, Marci 16.prae dicaueruntubique Domino cooperante, et sermonem confirmante per sequentia signa; Receptus est in gloria. Sextum argumentum a glorioso exitu: quasi dicat: Hoc quoque ces tum est indiciu, quod Christus sst ueritas et certa salus quia ipse et in coelo et in terr.  \pend
\vspace{0.5cm}\noindent
\vspace{0.2cm}\rule{1cm}{0.2pt}\\ 
\hspace{0.2cm}\textit{mg}
\footnotesize Qomodo Christus iusifirauus sie laspariut. 
\normalsize| 
\hspace{0.2cm}\textit{mg}
\footnotesize 2.Argumentuma testimonio Ange lorum 1. Petri.1. 
\normalsize| 
\hspace{0.2cm}\textit{mg}
\footnotesize LAgmmen tum a testimonio Apostolorum. Esa.25.49. 60.Zacbus 9. 
\normalsize| 
\hspace{0.2cm}\textit{mg}
\footnotesize 2 Arguments̃ a fide corum qui audierunt 
\normalsize| 
\hspace{0.2cm}\textit{mg}
\footnotesize 8. Argume um a foeliei seu glorioso ext tu, 
\normalsize| 
\section*{IN I. EPIST. D. PAVLI. }
\marginpar{[ p.90 ]}\pstart obtinet summam gloriam. lu coelis quidem, nam eo gloriosus ascendit, vbi sedet ad dex » teram Dei patris, patri per omnia par et aequalis. Philip. 2. Deus in summam extulit » illum sublimitatem, ac donauit illi nomen, quod est supra omne nomen, vt in nomine » Iesu omne genu se flectat, coelestium ac terrestium et infernorum, omnisque lingua confiteatur, quod Dominus sit lesus Christus ad gloriam Dei patris. Act.4. Nec aliud nomen est sub coelo datum inter homines, in quo oporteat nos saluos fieri. Hae igitur rationes huc tendunt, vt probetur, veritatem et mysterium pietatis siue salutis in ecclesia contineri, extra eam vero non nisi falsitatem esse et perditionem: Ac proinde admonentur episcopi omni ope in id eniti, vt veritatem hanc tot euidentiss imis rationibus comprobatam in ecclesia retineant ac tueantur. Admonentur quoque omnesesubditi, vt huic veritati sic in ecclesiis traditae indiuulsem adhaereant. Quô autem clarius intelligantur verba Pauli, in arctum syllogismum licet omnia constringere, Summa itaque  argumentorum haec esse videtur. Id pro veritate haberi debet, quod sensu deprehenditur, quod a Spiriru sancto comprobatur, de quo Angeli ferunt testimonium, quod Apostoli vbique terrarum pro certo annunciarunt, cui et fidem plurimi in toto mundo adhibent, denique quod tam in coelo, quam in terris suam habet gloriam. Sed haecomnia ei competunt, de quo agitur in ecclesia. Igitur in ecclesia veritas retinetur ac praedicatur, eique est adhaerendum. Et per consequens recte dicitur ecclesia columna et sta bilimentum veritatis. Agitur autem in ecclesia de Christo, qui est salus nostra, et ueritas ipsa, cui omnia illa Pauli uerba conueniunt.  \pend
\phantomsection
\addcontentsline{toc}{subsection}{\textit{Sequitur quomodo episcopi se gerent in ecclesia erga cuiuscunque ordinis homines. }}
\subsection*{\textit{Sequitur quomodo episcopi se gerent in ecclesia erga cuiuscunque ordinis homines. }}
\endnumbering\beginnumbering\section{CAPVT IIII}\pstart \huge\textbf{S}\normalsize Piritus autem certo loquitur, quod in posterioribus temporibus desciscent quidam a fide, attendentes spiritibus impostoribus ac doctrinis daemoniorum per simulationem falsiloquorum cauterio notatam habentium conscientiam, prohibentium contrahere matrimonium, etc. vsque ad eum locum cap.6.o Timothee. Actenus de vita episcoporum et diaconorum domestica seu priuata: sequi tur iam, quomodon episcopi se gerent publice in ecclesia erga cuiuscunque ordinis homines. Ac primùm quidem hoc 4. cap.docet Apostolus, quomo do se habebunt in genere erga omnes: capitibus deinde 5. et 6. quomodo erga diuersae conditionis homines, ualde quidem ordinate omnia. Porro cum prima cura episcoporum esse debeat, cauere uidelicet, ne in ouile Christi irrumpant falsi doctores, qui gregem decipiant et seducant: et non satis est episcopum docere quae recta sunt, nisi possit etiam refutare et euertere quae falsa sunt: ldeo statim hic praecipit Apostolus, quomodo episcopus suos praemonebit de uitandis falsis doctoribus.  \pend\pstart Spiritus autem cerio loquitur.) Exorditur ab autoritate, more Prophetarum, spiritus S.Sic autem solent prophetae praefari. Haec dicit Dñs, etc. Fidem itaque suis uerbis conciliat ab autoritate spiritus, ne quis ea quae dicturus est, tanquam leuia et incerta uilipendat. Eandem ob causamtρφατικο͂ς addit uocem, certon: quia nimirum indubitato eueniet id, qd' spiritus praenunciat. Vel certo, quia non adumbrate, non sub figuris, uti solent Prophetae facere, ista significantur:sed certo, clare, simpliciter, nude, citra ullum uelamen. Inposterioribus temporibus: desciscent quidam a fide.) Imitatione Prophetica uaticinatur de posterioribus temporibus. Sic et Esa.2.Et erit, inquit, in nouissimis diebus praeparatus mons Domini in uertice montium, et eleuabitur super colles, etc. Comprehenditur autem nomine posteriorum temporum, omne tempus quod inter primum aduentum Christi in carnem est, usque ad secundum aduentum eius ad iudicium. Ideo et Apostoli sua tempora uocauerunt nouissima tempora. Vnde alibi Apostolus dicit: Nos sumus, in quos fines aetatum deuenerunt. Et loannes: Nouissima hora est. Et nostra tempora quibus nunc viuimus, itidem nouissima appellantur, et quae nos sequuntur. In Euangelio Matth.2o. nouissima tempora per horam designantur vndecimam, qua otiosi in vespera mundi in vineam coelestis patrisfamilians conducuntur.  \pend
\vspace{0.5cm}\noindent
\vspace{0.2cm}\rule{1cm}{0.2pt}\\ 
\hspace{0.2cm}\textit{mg}
\footnotesize In ecclesia ueritatis et salutis cont netur myste. rium, extra quam saluus fieri nemo potest. 
\normalsize| 
\hspace{0.2cm}\textit{mg}
\footnotesize Maiorpropositio. Minor. Conclusio. Ecclesia uet ritatis stabi limentum. 
\normalsize| 
\hspace{0.2cm}\textit{mg}
\footnotesize Episcopi u tapublica. 
\normalsize| 
\hspace{0.2cm}\textit{mg}
\footnotesize Aucuio et 
\normalsize| 
\hspace{0.2cm}\textit{mg}
\footnotesize fides. Euangelium caret umbris Tempor posteriora quenam si 
\normalsize| 
\hspace{0.2cm}\textit{mg}
\footnotesize Ab autorita te Sp. Sanct Veritas ora culorum. 
\normalsize| 
\hspace{0.2cm}\textit{mg}
\footnotesize 2Corint.1o t Ioan 
\normalsize| 
\section*{AD TIMOTHEVM. CAP. IIII. }
\marginpar{[ p.91 ]}\pstart Caeterùm hic discimus, insignem charitatem fuisse in Paulo, qui prodendo ista, quae spiritus ei reuelauerat, voluerit omnibus retrô seculis, adeoque nostris temporibus bene consulere. Ac propterea non debemus ista verba Pauli otiose legere, quasi parum ad nos pertinentia, cum ex his uel maxime constet nunc, apud quos vera doctrina retineatur. His ergo posterioribus temporibus, etiam iam nostris praesentibus desciscen quidam a fide: vere summa defectio est, vbi disceditur a doctrina fidei, quae plenilsime in apostolicis literis est tradita:vt reuera iudicari possit a fide deficere, quicuque id uult quasi ad salutem necessarium inducere, quod ' euangelicis siue apostolicis scriptis non est proditum. Sicut imprudentiae iudicari queat neglecto fonte adire cisternas: Itaimpieta tis est deserta doctrina fidei falutem velle quaerere ex doctrinis hominum: ideo discimus, merito nobis orandum, vt in fide seruemur, et fides nobis semper augeatur: sicut orabant apostoli Luc.17. et ne fides Petri deficiat, orat Christus Luc.22. Simon, Simon, ecce Satanas expetiuit uos, ut cribraret sicut triticum: Sed ego rogaui pro te, ne deficiat fides tua, et tu aliquando conuersus confirma fratres tuos. lgitur et nos debemus interdum inuicem confirmare, ne in fide deficiamus.  \pend\pstart Attendentes spiritibus imposloribus, ac doctrinis demoniorum.) Caussa formalis liue et ficiens, cur desicient a fide: quia uidelicet attendent et aurem accommodabunt spir tibus impostoribus. Graece προσέχεντες πνουμασι πλάνοις. προσεχω idem est, quod atten do, aduerto animum, totus adhaereo. Significat igitur totos addictos fore spiritibu impostoribus, fallacibus, seductoribus, erronibus. Nemo porest seruire duobus Dñis: Non audiet aliquis simul Christum et Belial: Qui in tenebris est, odit lucem: Qui non audi spiritum veritatis, qui est spiritus Dei, talis atrendit spiritui impostori, ac propterea in fi de labitur, et assentitur doctrinis demoniorum, id est, quas daemones, aduersarii Christ et veritatis inspirant et instillant impiorum animis. Simili prope modo dicit 2. Thessal. Mittet illis Deus efficaciam illusionis (sicut per praestigias illudunt oculi, vt putent se vi-- dere quae non uident) ut credant mendacio, ut iudicent omnes, qui non crediderunt ueritati, sed tapprobauerunt iniustitiam. Magnum autem est, qu Apostolus uno uerbo et magna autoritateuocat omnem falsam doctrinam generaliter, doctrinam daemoniorum. Sicut em omnis sans et uera doctrina recte dicit doctrina spiritus: nam spiritus est qui docet nos omnemueritatem, ut alibi ait Christus: lta omnis falsa doctrina recte dicit doctrina daemoniorum si ue diabolica: quia diabolus mendax est, et pater mendacii. Huc pertinet, quod lacobus cap.3 Sapientiam quandam uocat daemoniacam, quippe in qua aemulatio est et contentio.  \pend\pstart Per simulationem fassloquorum, cauterio notatam habentium conscientiam.) Modus quo fallent spiritus impostores. Fallent, inquit, per simulationem seu per hypocrisin falsiloquorum: quia nulla ratio ad fallendum simpli ces magis idonea, quam simulata sanctitas. Ideoque data est a Deo angelo Satanae hec potestas, ut in angelum lucis sese transfiguret:2. Cor.1. Hincnul, la opinio proponit ta absurda, am aliena a recto, quae non aliquo fuco et colore ueri po sit tegi et uestiri. Et in omnibus fere erroribus et superstitionibus, quae unquam in eccle siam irrepserunt, hypocrisis dominata est. Quam ob caussam creberrime admonet scriptu ra hypocritas summo studio deuitandos. Caeterùm utut foeliciter agant suum negocium hypocritae coram mundo, et apud simplices habeant pro sanctis et piis: tamen coram Deo habent insynceram et ream conscientiam, habent conscientiam notatam cauterio. Estque elegans metapho raad significandum contaminatam conscientiae maculam, qua cauterium a corpore ad animum trans, fert. καυτήριον graeca uox, instrumentum ferreum quo Chirurgi adurunt putridas cames siue ulcera, unde semper nota aliqua remanet. Etiam maleficis interdum notae in facie, uel alia corporis parte ad contumeliam huiusmodi instrumento inurunt. Recte ergo cauterio notatam habere dicunt conscientiam, qui quamuis foris simulare norint mirabilem sanctimonia, tamem intus conscientiam habent contaminatam et impuram per hypocrisin, sicut nunquam elue re possint impietatis notas, non magis que elui potest stigma cauterio inustum curi. Et per eam notam rei significant coram Deo, et a bonis piisque quibus pura est conscientia, separant. Sicut in piis requirit Apostolus fidem et conscientiam puram seu bonam, ut iam saepe dictum est:lta impiis hic dicit conscientiam esse cauterio notatam, id est, impuram, non bonam. Danda unicuique opera, ne in horum numero reperiatur.  \pend
\vspace{0.5cm}\noindent
\vspace{0.2cm}\rule{1cm}{0.2pt}\\ 
\hspace{0.2cm}\textit{mg}
\footnotesize A fde deftcere est sum ma descetio. 
\normalsize| 
\hspace{0.2cm}\textit{mg}
\footnotesize lermma. 
\normalsize| 
\hspace{0.2cm}\textit{mg}
\footnotesize Maith.8. 2Corin.6. 
\normalsize| 
\hspace{0.2cm}\textit{mg}
\footnotesize 2.1b841. 
\normalsize| 
\hspace{0.2cm}\textit{mg}
\footnotesize Dodtrina denoniort̃ est doctrina fasa Acontrario. 10an.16. 1041.8. Ipocrist et simulata sanctimonia fimplices seducit Satan. Doctrina Sa tanae fucata. 
\normalsize| 
\hspace{0.2cm}\textit{mg}
\footnotesize Arectapbora a corpore ad animum. καυτιριάιω contamino, Cauterio inficio. Pailus hic habet uατηραναι ras. 
\normalsize| 
\section*{IN I. EPIST. D. PAVLI. }
\marginpar{[ p.92 ]}\pstart Prohibentium contrahere matrimonium, iubentium abstinere a cibis.) Veldu ngna quaedam addit a genere doctrinae sumpta, quibus dignosci possint falsi illi doctores hypocritae. Pauca autem signa posuit pro multis, ratus satis esse exempli gratia protulisse. Christus itidem Matth.7.23. et alibi describit quomodo hypocritae dignosci possint, a genereuitae siue a morbus, vtshic a genere doctrinae describuntur: vt nullo pacto excusari possint, qui dicunt, se non posse distinguere veros doctores a falsis. Certe nun quam sermone quis consequatur, quantum Apostolo debeamus, quod haec signa expresserit. Nam uel ex his potissimum clare iudicare possumus apud quos pura doctri na sit conseruata hoc nostro tempore, et e contra qui eam oppresserint, siue qui sint spiritus impostores, falsiloqui, hypocritae, conscientiam habentes cauterio notatam. Quis autem existimaret ista esse contra doctrinam fideic quis cegitaret esse doctrinas daemoniorum, nisi Paulus ipse ita pronunciasset ? Nam vxorem habere, vel non habere, hoc cibo uesci aut non uesci, putaremus inter ἀδιάφερα seu μέσα, id est, indifferentia seu media numeranda. Imo facile persuaderi queat, ista amplissima quaedam et sanctitate plena opera, uidelicet coelibem uitam agere, non ducere uxorem, carnem abstinentia certorum ciborum macerare. Matth. 19. Beatus, qui castrauerit semetipsum propter regnum coelorum: Et coelibatus angelicam puritatem referre uidetur. Seiendum autem et haec et alia multa, si simpliciter considerentur aδιιάφορα seu Lεσα recte appellari, utpote quae neque facta neque neglecta prosunt aut nocent: Sed tum demum in his peccatur contra doctrinam fidei et fiunt doctrinae daemoniorum, quando, ut Paulus hic ait, prohibentur, quasi per se peccata, et eduntur praecepta ad obligandas conscientias. Erroris ergo uis deprehenditur in uerbo prohibitionis. Nam prohibere illa quae in scripturis approbantur, et ita prohibere, quasi qui diuersum faciant, grauissime peccent coram Deo: tum praecepta istiusmodi aequiparare praeceptis Dei: denique asserere, eorum obseruatione Deum coli: hae conditioncs additae efficiunt doctrinam daemoniorum. Nos in libertatem spiritualem uocati sumus per Christum, ac stare debemus in ea libertate, qua Christus nos per fidem donauit, ut habetur Galath.5. Itaque non debemus sieri serui hominum, non debemus attendere spiritibus impostoribus, qui nobis uolunt obtrudere et cogere nos ad istiusmodi humanarum traditionum obseruationem. Quod si in hanc seruitutem coniicimur, et sic ista obseruamus, quasi ad salutem necessaria, aut quasi per ea promereremur gratiam coram Deo: Sane Christus nobis fiet otiosus, et quasi frustra pro nobis mortuus. Caeterum quid scriptura definiat de coniugio uel coelibatu, non est obscurum. Damnarunt nuptias et certos cibos Manichaei, Encratitae seu Tatiani, nuptias in primis, Esseni, Monachi, Apotactitae: quasi dicas extraordinarios: denique pontifices quidam suis prohibuerunt. De Essenis losephus de bello iudaico liber 2.cap.7.De Apotacti cis Augustinus. De Pontificibus, Calixto, Innocentio primo uide Platinam. De Gregorio septimo alias Hildebrando vide Sigibertum Gemblacensem. Scriptura autem nuptias permittit omnibus liberas, coelibatum nullis praecipit. Nam de coelibatu Sapient. 8. dicitur: Scio quia non possum esse continens, nisi Deus det. Matth.19.cum Apostoli dixissent: Si ad istum modum habet causa hominis cum uxore, non expedit contrahere matrimonium, dixit Christus: Non omnes capaces sunt huius dicti, sed ii, quibus datum est. 1.Corinth.7. Bonu est homini uxorem non attingere: attamen propter stupra uitanda suam quisque vxorem habeat, et suum quaeq uirum. Ibidem. Dico inconiugatis et uiduis, bonum est eis, si manserintut et ego. Quod si se non continent, contrahant matrimonium. Nam satius est matrimonium contrahere, quam uri.lbidem. De uirginibus praeceptum Domini non habeo. Quod si duxeris uxorem, non peccasti. Et si nupserit uirgo, non peccauit. Quae uero ibi dicit in commendationem uirginitatis, sic uult accipi, ut non iniiciat ulli laqueum. (sicut non ideo contemnitur argentum, quod maioris precii aestimatur aurum: ita nec matrimonium est improbandum, etiam si ei anteferri uideatur Virginitas:) Sed plurima adhuc passim occurrunt in commendationem coniugii:ut Hebr.13. Matth.19. Genes.2. etc. Quae omnia apertissime damnant istos, qui prohibent matrimonium contrahere.Sed et reru euentus docuit grauissima  \pend
\vspace{0.5cm}\noindent
\vspace{0.2cm}\rule{1cm}{0.2pt}\\ 
\hspace{0.2cm}\textit{mg}
\footnotesize sgu.x9 bas agnosci tur doctrin bypocritae et same. 
\normalsize| 
\hspace{0.2cm}\textit{mg}
\footnotesize ααρςι seu μiσα 
\normalsize| 
\hspace{0.2cm}\textit{mg}
\footnotesize quae sintis de Gellium libus 2 cap.71 et D. Philit pum in loc communibie de traditie nibus bum nis. Origenes ceptus. Auuphor rum discrimeniu proi bibitione. Libertas Christianin adiaph ris. 
\normalsize| 
\hspace{0.2cm}\textit{mg}
\footnotesize ErerNa nichaeorum Encratitaram, Essens runt, Mons chorum, A. potactitarum, Nuptie lib raeuscript ris pmisse. Castum, et coelibem est se non estne stri captus e potestatis. 
\normalsize| 
\hspace{0.2cm}\textit{mg}
\footnotesize A fiuilt 
\normalsize| 
\hspace{0.2cm}\textit{mg}
\footnotesize Commend to conuae 
\normalsize| 
\section*{AD TIMOTHEVM CAP IIII. }
\marginpar{[ p.93 ]}\pstart scandala in ecclesiis eurenisse ex illo coelibatu contra Dei verbum praecepto. Er longu foret enarrare, qualia, quam abominanda scelera sint tum Romae tum alibi per inuoluntarios coelibes perpetrata, et quodammodo impune lata. Quia autem manifestiss mus est error, prohibere matrimonium, ideo Apostolus hoc loco nullum contra eum verbum facit, quasi indignum, qui multis argumentis confutetur: et manifestiorem, quam vt opus sit multum impugnari. Nam et naturae ipsi repugnare a quouis intelli gitur. Tollant vitium in natura positum et concupiscentiam, ac tunc prohibeant nuptias.  \pend\pstart De altero itaque errore est dicendum. ls autem est, Iubentium abstinere a cibis: Et hos quoque falsiloquos doctores senserunt nostra secula. Nam illi ipsi qui nuptias prohibuerunx, etiam coegerunt abstinere a cibis: quasi vero pletas esset sita in certorum ciborum abstinentia, et impietas in eorundem vsu. Vide Platinam in vita Gregorii septimi: etiamsi nonnulli interdictos cibos referant ad Thelesphorum. Vide exemplum Spiridionis episcopi Cyprii in historia tripartita liber ..cap.1o. Sed quae contra hos in scripturis sunt prodita, non opus est per nos adferri: quando Apostolus ipse breuem quandam eius erroris confutationem adiicit.  \pend
\phantomsection
\addcontentsline{toc}{subsection}{\textit{Quos Deus creauit, ad sumendum cum gratiarum actione fidelibus, et his qui cognouerunt veritatem: quod, quicquid creauit Deus, bonum sit: et nihil reiiciendum, si cum gratiarum actione sumatur: sancti ficatur enim per sermonem Dei ac precationem. }}
\subsection*{\textit{Quos Deus creauit, ad sumendum cum gratiarum actione fidelibus, et his qui cognouerunt veritatem: quod, quicquid creauit Deus, bonum sit: et nihil reiiciendum, si cum gratiarum actione sumatur: sancti ficatur enim per sermonem Dei ac precationem. }}\pstart Tria sunt argumenta. Primum a dignitate creatoris et rerum qualitate. Quod bonum est, id nequaquam est reiiciendum. Sed quicquid Deus creauit, id bonum est: (sic enim et Genes.1. dicitur: Et vidit Deus cuncta quae fecerat, et erant valde bona) lgitur cibus quem Deus creauit, quique ideo bonus est, nequaquam est reiiciendus. Huc pertinet quod Ambrosius dicit, cibos per id, quod despiciuntur, quodammodo signi ficars malos esse, et a malo autore inuentos. Sic autem existimarunt Hieracitae et Manichaei.ldeo generaliter dictum esti. Corinth.o. Omne quod in macello venditur, man  ducate, nihil interrogantes propter conscientiam. lbidem. Si quis vos vocat infidelium ad coenam,  et vultis ire, omne quod vobis apponit manducate, nihil interrogantes propter conscientia. .e  \pend\pstart Secundum argumentum a caussa finali sic habet. Cibi sunt condiri, vt sumantur cum grauiarum actione fidelibus. Igitur nihil rerjciendum, si cum gratlarum actione sumatur. Si itaque in hunc vsum sunt conditi, cur illis non vtamur? Huc pertinet, quod Deus omnia condidit ad vsumn hominis, et hominem omnibus praefecit, vt omnibus libere vtatur, Genes19. et Pialm.3. Omusa subiecisti sub pedibus eius, oues et boues etc. Et de gratiarum actione Rom 14. Qui non vescitur vescentem non iudicet. Qui vescitur, Domino vescitur: gratias enim agit Deo. ldeo a Paulo reprehensus Petrus Galath2. apud Antiochiam, quod cum prius sumpsisset cibum cum Gentibus, postea aduenientibus ludaeis sub duxit se ac separauit a Gentibus. Caeterùm non sine caussa addit, Fidelibus, et per in terpretationem, his qui cognouerant veritatem. Hos enim censet vere fideles, qui nat creaturam per se bonam, et creaturae vsum: tum quia sibi vendicat ex abstinentis Illa immoderata iustitiam, ex opere uidelicet suo iustificari uolens coram Deo, et ali, os non ita facientes arroganter aspernatun, et si quas agit gratias, tantum hypocritics id facit: neque etiam habet illum finem abstinentiae, quem deberet, uidel cet maceradignum reddit, et grauissime peccat, qui creaturam Dei bonam aspernatur. Fideles interim omni creatura utuntur sine peccato, non quidem ad ostentationem  \pend
\vspace{0.5cm}\noindent
\vspace{0.2cm}\rule{1cm}{0.2pt}\\ 
\hspace{0.2cm}\textit{mg}
\footnotesize Ex coelibata grauissima enata scanda la. 
\normalsize| 
\hspace{0.2cm}\textit{mg}
\footnotesize Probibere matrimomt̃ 
\normalsize| 
\hspace{0.2cm}\textit{mg}
\footnotesize mfdeui ges rum usus est noxtus, propter ipfiusmet maltiam G dnerias. 
\normalsize| 
\hspace{0.2cm}\textit{mg}
\footnotesize rideles recte sancteque ́z atuntur donis Dei. 
\normalsize| 
\hspace{0.2cm}\textit{mg}
\footnotesize 
\normalsize| 
\hspace{0.2cm}\textit{mg}
\footnotesize 
\normalsize| 
\section*{IN I. EPIST. D. PAVLI. }
\marginpar{[ p.94 ]}
\marginpar{[ p.95 ]}\pstart et Des glorlam. Hinc dicitur Tita. Omnia pura puris: porro pollutis et inndelibus, nihil est purum, sed polluta est illorum mens et consclentia. Coloss.2. Nemo vos iudicet in cibo aut potu, aut parte diei festi, quae sunt vmbra futurorum, corpus autem Christi. Rom.14.1. Corinth.8.ltaque fideles omnia liberem sumunt cum gratiarum actione.  \pend\pstart Sanctificatur enim per sermonem Dei ae precationem.) Tertium argumentum habes ine clusam occupationem. Sanctificatus est cibus: igitur licite sumitur. Quod si quis neget sanctificatum, dicens esse immundum, ac propterea eo non vescendum: id rursus diluit, dicens: eum sanctificari per verbum Dei. Sanctificatur autem cibus, non quod antea non fuerit bonus, quandoquidem quicquid creauit Deus, bonum est: vel, quod fiat sanctior, quam antea: Sed ad vsum sanctificatur, vt fidelis eo vtatur sine peccoto, qui infideli sanctus et purus nequaquam est Sanctificatur autem per sermonem Dei, id est, per Dei institutionem. qui voluit omnia pura esse puris, et sancta esse omnia recte cum fide et gratiarum actione vtentibus: ldeo Actor. 11. Petro dictum est, Quod Deus purificauit, tu ne commune dixeris. Purificata sunt autem omnia, vel obhoc, quod a Deo sunt condita, et bona, et eo sermone, quo Deus illa condidit, eodem sanctificantur. Si autem bona et sanctificata, unde dicentur impura, praesertim ei, dui sis tans quam bonis cum gratiarum actione utitur? Romah.14. dicitur. Noui et persuasum ha" beo per Dominum lesum, nihil esse commune perse, nisi ei qui existimat aliquid esse  commune, illi commune est, qui videlicet infidelis est, et non credit ea verbo Dei esi  se sanctificata: Marci 7. Christus aiebat, nihil esse extra hominem introiens in eum, quod possit eum coinquinare, nimirum quia omnia simul, dum conduntur verbo Dei, etiam per verbum Dei sanctfficantur.1. Coilath.8. Esca vos non commendat Deo, neque, si comedamus, alsquid vobis superest, neque, si non comedamus, quicquam vobis deest. Praeterea adhuc sanctificantur ad vsum nostrum, per precationem, quae fit ex fide. Duplies ratione sanctificantur: primùm generaliter verbo ipsius Dei, quo sunt condita, et instituta omnia vt vtilia et pura sint bene vtentibus: Deinde accedente etiam nostra precatione per fidem. Eam ob caussam Christus exemplo suo frequenter docuit, cum fracturus esset panem, gratlas agere, deinde et cibo sumpto hymnum dicere: ita, vt etiam si quid putaretur impuritatis vlli cibo inesse, interueniente tamen nostra precatione, purus effieeretur. Atque ex hoc loco discimus, doctrinam hanc valde necessariam omnibus Christianis, vt nullum conuiuium auspicentur vel finiant, quin praeeat consequaturque gratiarum actio, quam rationem Tertullianus saepe celebrat: Estque hoc institutum etiam a teneris inculcandum et obseruandum vt sic cum honore et summa reuerentia cibum a Deo creatum et sanctificatum sumamus. Et certe natura impulit etiam ethnicos ad ciborum tractationem honorandam. Hac namque de caus sa potissimum mensae ornantur mundis mappis, non tantum digitorum extergendorum gratia, sed magis veluti ad religionem et puritatem significandam, quae et cibis inest, et in nobis quoque requiritur, quo tempore dona Dei quotidiana accipimus. Caeterùm cum haec ita dicuntur contra praeceptum coelibatum, contra praeceptam abstinentiam a certis cibis: non damnatur quidem virginitas aut coelibatus, non item damnatur carnis castigatio et moderata abstinentia: sed hypooritica illa sanctimonia, et ad has res amplectendas coactio, et conditiones illae additae, quod sint obseruatu necessaria, quod sint cultus Dei: quod gratiam promereantur coram Deo: quod, qui ista obseruant, praestantiores et meliores sint caeteris, qui eadem non obseruant: denique omnia illa scandala et pericula conscientiarum, quae ex huiusmodi doctrinis daemoniorum exoriuntur: veluti ex coacto coelibatu scimus grauissima scandala enata, ex ciborum delectu periclitatas simplicium conscientias: sed quaestum cessisse auaris episcopis.  \pend
\phantomsection
\addcontentsline{toc}{subsection}{\textit{De his si commonefeceris fratres, bonus eris minister Iesu Christi, enutritus in sermonibus fidei bonaeque doctrina, quam usque secutus es. }}
\subsection*{\textit{De his si commonefeceris fratres, bonus eris minister Iesu Christi, enutritus in sermonibus fidei bonaeque doctrina, quam usque secutus es. }}\pstart Infert ipsius esse, de cauendis eiusmodi spiritibus, impostoribus et doctrinis daemoniorum pios admonere, atque vt id strenuë faciant, addit argumenta suasoria. Prius  \pend
\vspace{0.5cm}\noindent
\vspace{0.2cm}\rule{1cm}{0.2pt}\\ 
\hspace{0.2cm}\textit{mg}
\footnotesize 5. Argumentum ab Occu patione. Verbo Dei cibus sancti ficatur. 
\normalsize| 
\hspace{0.2cm}\textit{mg}
\footnotesize Fidelibus e mnia sunt pura ex san ctificatione Dei. 
\normalsize| 
\hspace{0.2cm}\textit{mg}
\footnotesize Duplex fat ctificatioci borum. 1. Sanctific tio manat es uerbo Dei. 2. Sanctific. tio fit per precationem ex fide. Conuiuia Christianorum ichoari et finiri de bent gratia rum actionc. 
\normalsize| 
\hspace{0.2cm}\textit{mg}
\footnotesize Pontificiort conditiones peruersee hpσsuuose. 
\normalsize| 
\section*{AD TIMOTHEVM CAP. IIII. }
\marginpar{[ p.95 ]}\pstart argumentum, Bonus eris minister Iesu Christi, a ratione officii et a pio: quasi dicat, pie feceris, et id quod debes praestiteris, si tuos in tempore commonefacias. Alterum argumentum, Enutritus in sermonibus fidei bonaeque doctrinae, uidetur esse a facili; quasi dicat: facile erit tibi animos illorum praemunire et confirmare contra tales errores ingruentes, quia ab infantia sacras literas nosti. Nota quod sermones fides semper urget quia semper fides illa respicienda est. Omnis doctrina ad fidem et charitatem ceusco; pum debet tendere, uti in principio dixit: Finis praecepti est charitas ex puro corde, et conscientia bona, et fide non simulata. Latet et in illis uerbis argumentum ab honesto siue ab aequo: quam vsque secutus es: quasi dicat: Hactenus semper probe es tutatus sanam doctrinam, turpissimum foret si nunc eius partes desereres. His itaque ra uionibus animat Apostolus suum Timotheum, ut omni diligentia suos commonefa ciat (animaduerte in hoc uerbo modestiam requiri maximam) de vitandis doctrinis daemoniorum. Animaduertant ergo, his argumentis se quoque admoneri, quicumque ecclesiarum ministri, et sua interesse, vt inuigilent sedulo, ne gregem dominicum inuadant lupi rapaces, adferentes doctrinas daemoniorum.  \pend
\phantomsection
\addcontentsline{toc}{subsection}{\textit{Caeterùm profanas et aniles fabulas reiice, quin potius exerce temetipsum ad pietatem. Nam corporalis exercitatio paululum habet vtilitatis. At pietas ad omnia vrilis est, vt que promissiones habeat vitae praesentis ac futura. }}
\subsection*{\textit{Caeterùm profanas et aniles fabulas reiice, quin potius exerce temetipsum ad pietatem. Nam corporalis exercitatio paululum habet vtilitatis. At pietas ad omnia vrilis est, vt que promissiones habeat vitae praesentis ac futura. }}\pstart Generaliter praecipit, quicquid non facit ad ueram pietatem, reiiciendum, neque alio loco habendum, quam quo haberemus profanas, id est, alienas ab omni religione: et andles, id est, vanas, verbis molestas, fabulas. Adeo episcopi interest, non solùm ea resu care, quae manifeste sunt impia et cum uerbo Dei pugnantia, verum illa quoque quae licet ir speciem uideantur bona, tamen ad ueram pietatem haud sunt necessaria. Proponebantur  haud dubie multa et praeclara a pseudoapostolis etiam ex lege Mosaica petita, saltem eilegi consentanea: cuiusmodi erant obseruationes dierum, delectus ciborum, mundicies externa, ieiunia, deinde et alia quaedam quae humanae rarioni probabilia, in admiratione habenda, laude digna et Deo grata fore videbantur: et tamen illa omnia Apostolus rencit ceu abiecta et prorsus indigna, vocans ea profana, et prouerbsali locutione aniles fabulas, quasi illis contemptius nihil posset iudicari. Sicut enim solent prouerbiali figura, ea quae nullius momenti sunt, quaeque pro meris deliriis habentur, ridicula et contem nenda, aniles sabule appellari: quia nimirum anus, natura et aetate loquaces, res absur; das et uanas magno taedio audientium narrare consueuerunt, aut, si oblectant, sola oblectantuanitate. Vnde graecis prouerbium est, γραωυ ἐιλω, id est, anicularum delyramenta Profanum quoque irreligiosum, impurum, et a sacris arcendum significat: Vnde illud, Proeul hinc, procul ite profani. Breuiter itaque et simplicissime, omnia illa quae ad veram pietatem non faciunt, Apostolus nihilo praestantiora duci vult profanis, id est, impuris, sordidis et anilibus fabulis, id est, ridiculis, deliris commentis. Et quo nomine commodius significare posset summum earum traditionum contemptum, quam vocando eas fabulas ? Apparet autem ideo Apostolum audacem hac in partefuisse et magnanimum, in adeo vilisaciendis his, vt nos quoque eo audacius omnem doctrinam, quae non necessaria est ad ueram pietatem, aspernemur, neque patiamur nos aut nostros vlla doctrina, etiamsi non aperte impia, a recto abduci. Sed videamus quid nomine pietatis intelligi conueniat. Vbi interpres reddidit. Pietatem, graece habetur δὐσέθεαυ. Εύσέθεα, si Etymon respicimus, idem est quod bonus siue uerus cultus. Et ita pietatem vertit etiam Hieronymus Epistola29. Pietas, inquit, quam Graeci uel οὐσεθειαν, vel expressius et plenius θεοσέθοαν uocant, aliud non est, quam iustus Dei cultus. Et quomodo proprie Latinis auribus accipiendum sit nomem pietatis, Cicero in oratione pro domo sua ad Pontifices indicauit. Equidem sic accepi, inquit, Pontifices, in religionibus suscipiendis caput esse interpretari quae uoluntas deorum immortallum esse uideatur. Nec est ulla erga deos pietas, nisi honesta de numine sorum ac mente opinto, cum expetinihi ab iis quodsit iniusta et inhonestum  \pend
\vspace{0.5cm}\noindent
\vspace{0.2cm}\rule{1cm}{0.2pt}\\ 
\hspace{0.2cm}\textit{mg}
\footnotesize L Argumnen tum suasorium a ratione of ficii et apio. 2. Argumen tum a facili. Scopus doctrinae fidei et charitati est designatus. 
\normalsize| 
\hspace{0.2cm}\textit{mg}
\footnotesize Ab bonesto et equo. 
\normalsize| 
\hspace{0.2cm}\textit{mg}
\footnotesize Preceptum generale de anilibus et uauis fabulis Iugienaus. 
\normalsize| 
\hspace{0.2cm}\textit{mg}
\footnotesize Anus anitia Fractant. Prouerbium. 9τάσῆθ́κος. 
\normalsize| 
\hspace{0.2cm}\textit{mg}
\footnotesize EiMha
\normalsize| 
\hspace{0.2cm}\textit{mg}
\footnotesize Feusax, 
\normalsize| 
\section*{L. EPIST. D. PAVLI }
\marginpar{[ p.96 ]}\pstart Pietupe arbitrere. Pietas ergo ad Dei honorem respicit, eaque complectitur quae ad verum, bogloriam cul tongrept num, legitimum cultum Dei pertinent. Cum autem Deus nequaquam colatur mandatis hominum, siue cum in traditionibus humanis non situs sit cultus Dei: sequitur ergo solis mandatis Dei eum coli. Mandatorum vero Dei summa haec est, ut credamus in eum quem misit ille: ut dicitur loan.6. Deinde propter eum diligamus nos mutuo: Coutar De loan.1. Quare pietas omnis in fide et charitate consistit, et quicquid ex his oritur, ad us potissimi pietatem recte referetur. Quicquid uero ad has proprie non pertinet, profanum et anifide et Cha libus fabulis simile aestimabitur. Pietas vera in spiritu et uerstate cultum, timorem, amoritate iisque ́ rem, confessionem includit: deinde patriae, parentum, amicorum, inimicorum, superiooperibus qua rum et inferiorum ex aequo dilectionem: In animo praecspue sita est, et fide charigteque  exillis oriun potissimum praestatur: deinde exterius etiam iis operibus, quae suapte natura ad Dent gloriam et proximi utilitatem faciunt. In vera itaque pietate hac debemus non solum aEmphasieid hos, sed nosmetipsos, ut caeteris exemplum praebeamus, reipsa, non tantum uerbis, exerutraque uoce. cere. Exerce, inquit, temetipsum. Exercere significat frequenter, sedulo et cum diligenExercitus tia in aliqua re occupari. Vnde exercitum appellatum apparet, quia in eo copiae miliinde didtus tum assidue occupantur in armis. Et dicuntur exercitia bonarum artium, quia his toto animo et omni tempore intentos esse conuenit. Graece est γυμναζε σεαυτον, a uerbo yυuνέζω, proprie de exercitatione corporum, qualis erat in ludis et gymnasiis athletarum, ubi ιύενοι, id est, nudi se exercebant. Exigitur itaque, ut opera illa pietatis magno nisu et assidue a nobis ipsis expleantur, dum oportet nos in illis exerceri, et totos, ut ita dicam, occupari.  \pend\pstart Nam exercitatio corporalis paululum habet vtilitatis, at pietas ad omnia viilis est.) Sicut antea duobus exemplis in medium adductis ostendit, qualesnam essent doctrinae daemo niorum: lta hic exemplo uno subindicat, qualia ceu aniles fabulas uelit reiici, nempe corporalia illa quae nonnulli splendida tradunt de exercitatione corporall, cuiusmodi sunt, ieiuuana si adfit nia, uigiliae, corporis molestatio, et similia, quae ad refrenandum carnem excogitantur: saperstitio. quae omnia quamlibet mundus admiretur pro difficilibus et Deo placentibus, tamen Comparatio si aliquid habent utilitatis, certe non multum. Confert autem corporalem exercitatiopietatis et nem cum pietate, et hanc ratione utilitatis illi longe praefert: spectandaque est hic comcorporalis munis sententia, quod ex duobus utilibus, id quod magis utile est, eligendum. Sed uiexercitatio de obsecro, qumo Paulus pro exemplo, earum rerum, quae alienae sunt a pietate, quaeque  nihilo praestantiores anilibus fabulis, adfert exercitationem corporalem, quam tamen putarent multi rem magnam et sanctam et praedicandam. Sed quam nugalia diceret et uana, etiam praecepta edita et excogitata de his, deinde persuasiones, quod promererentur gratiam apud Deum? Deinde quod nonnulli alios facerent suorum meritorum ex istiusmodi exercitationibus participes, id quod monachos quosdam fecisse uidemu, qui literas uendebant, quibus fatebantur, se participes facere emptores omnium eiusmodi operum et meritorum: Apostolus omnia esusmodi dicit paululum utilitatis habere, uel fere nihil Reuera autem quid cedit ad Gloriam Dei, quod quis inedia caput suum laedit, aut corpus suum alia ratione afficit iniuria, per superstitionem, ut dicitur Coloss..ac humilitatem animi et laesionem corporis? Quantulum etiam ex istiusmo di iuuatur proximus tuus egens et inopia pressus ? Sane neque Deus his honoratur, nes que fides exercetur, neq insuper exercetur charitas, qua iuuaretur proximus: imone Exempiñ de illi quidem ipsi prosunt, qui se sic vexat et fatigat. Ideo Dominus uolens illud insinuateuane. re per Esaiam cap.58.aiebat: Nunquid tale est ieiunium quod elegi, per diem affligere » hominem animam suam Anut incuruet sicut iuncus caput suum, et saccum et cinerem " sternat? An istud uocabis seiunium et diem acceptabilem lehouae ? Nonne hoc est ie" iunium quod elegi, ut dissoluas colligationes impietatis, soluas fasciculos graues, et di  mittas confractos liberos, atque ut omne iugum disrumpatis ? Nonne ut frangas esurienti panem tuum, et pauperes uagos introducas in domum? quum uideris nudum ope* rias eum, et a carne tua ne abscondas te? etc. En aperte prophetae uerba Apostoli sen» tentiam explicant. Dissolue igitut colligationes impieratis, perinde ac si diceret, exerce te in uera pietate et omitte las exercitationes corporales, stultas et inutiles. Pietate  \pend
\vspace{0.5cm}\noindent
\vspace{0.2cm}\rule{1cm}{0.2pt}\\ 
\hspace{0.2cm}\textit{mg}
\footnotesize cit. Matth.is. Est29. tur. 
\normalsize| 
\hspace{0.2cm}\textit{mg}
\footnotesize Exercitia uis
\normalsize| 
\hspace{0.2cm}\textit{mg}
\footnotesize 
\normalsize| 
\hspace{0.2cm}\textit{mg}
\footnotesize 
\normalsize| 
\section*{AD TIMOTHEVM CAP. IIII. }
\marginpar{[ p.97 ]}\pstart autem ibi aperte constituit in operibus fidei et charitatis, in miericorula, eieeinorynas noxarum condonatione. Hoc visum est exempli caussa de inutili ieiunio ex propheta adducere: Et simili modo iudicandum est de aliis exercitationibus, quae cum ieiunio vix sunt conferendae. ltaque constat exercitationes corporales etiam exacte et diligenter sactas, parum vel nihil habere vtilitatis, pietatem vero, id est, opera ex fide et charitate prodeuntia, quibus honoratur Deus et iuuat proximus, ista multam vrilitatem adferre, nenpe Deo gloriam, proximis commodum, et illi qui ea facit, quietemin conscientia. Ambrosius: Omnis summa disciplinae nostrae in misericordia et pietate est. Prodest autem in haec uerba Pauli frequenter intueri, ut prudenter discernamus, quaenam opera ad veram pjetatem facians et quae non. Videtur autem non temere dixisse Apostolus, exercitatione corporalem paululum habere vtilitatis: quasi aliquid utilitatis illis inesse concederet. Reuera autem, pdest interdum piis, si adhibeant, ad extirpandas, uel saltem opprimendas, quantum fieri licet, re liquias peccati nobis inhaerentes, corpori quasi frena quaedam per exercitationes corporales, sicut Apostolus dicebat 1.Corinth.9. Subiicio corpus meum et in seruitutem redi go:ne quo modo siat, ut cum aliis praedicarim, ipse reprobus efficiar. Et praeparantur per istiusmodi exercitationes corpora nostra, quo minus rebellent spiritui ad bona enitenti. Sed ista non facient pii quasi ad iustitiam sibi parandam, aut quasi legis conscientia ad eiusmodi adigerentur. lmo semper iudicabunt ista, etiam cum integerrime fiunt, cum uera pietate, quae in fide et charitateuersac̃, collata, nulla propemodum esse, imon pene sordida, et nun quam parem utilitatem allatura. Non sunt respuenda illa, quae ad uirtutum incitamentum faciunt uel animum praeparant: Sed ipsis uirtutibus postea prodeuntibus nullo pacto aequanda Sicut praeparata quae eriguntur cum domus aedificanda est, sunt quasi necessaria, nor tamen postea suntutilia, perinde ut domus ipsa: ita de omnibus illis actionibus, qua ad continentiam, ad sobrietatem, praeparant, est iudicandum, qu ipsa quidem sunt ueluti necel saria effrenis animis, sed in illis pietatis utilitas nequaquam est collocanda, sed de his satis.  \pend\pstart Vt quae promisseiones habcat vitae praesentis et futura.) Probatio quod lit utilis pietas. Et subintelligitur argumentum hortatorium huiusmodi. Illi rei potius intendendum esti quae habet adiunctas promissiones, quam quae non habet: Sed pietas haber adiunctas promissiones: exercitatio autem corporalis non habet. lgitur pietati potius intendendum, quam exercitationi corporali. Caeterum quod exercitatio corporalis non habeat ad iunctas promissiones, satis liquet, quia nusquam legitur quempiam ratione eius iustificatum, uel alias consecutum aliquid praemii a Deo: nisi forte dicamus illud praemium, qu hypocrita recipiunt propter difficilia opera sua, gloriam inanem apud homines: sed hanc accipiendo amittunt gloriam apud Deum: et quam imperiti in illis honorant pro pietate, id apud Deum damnatur pro summa impietate. At uero pietas habet adiunctas promissiones. Nam cum pietas constet fide et charitate: sequitur omnes promissiones factas credentibus et charita tem exercentibus ad pieratem debere referri. Quomodo autem fide et charitate sancti patres obtinuerint promissiones, plena est scriptura. Vnde Hebr.11. dicitur generaliter, sancti assecuti sunt perfidem promissiones. Roman.4. Promissio contigit Abrahae non per legem, sed perfidem. Longum esset exempla enumerare. Dicit autem promissiones esse vitae praesentis et fuiurae: Estque Amplificatio utilitatis, quod etiam aeterna sit, non momentanes tantùm. Constat enim, qu pietatem exercentibus non solum praemia sint proposita 1r futura uita, de qua re plurima extant testimonia, sed etiam in praesenti uita proposita sunt bona temporalia. Huiusmodi erant promissiones illae ad Abraham Genes.12. faciam te ir gentem magnam et benedicam tibi, et magnificabo nomen tuum, erisque benedictus: Benedicam benedicentibus tibi, et maledicam maledicentibus tibi. Et Genesas.ad lacob: Terram in qua dormis, tibi dabo et semini tuo, eritque semem tuum quasi puluis terrae. Dilataberis ad orientem et occidentem, et septentrionem et meridiem. Et Deuteras. praemia diuersa et temporalia promittunturiis qui obseruabunt legem Dei:sicute contrario cap.27. proponuntur horrenda mala transgressoribus. Et in nouo testamento Matth.6. Nouit pater uester coelestis, qu opus habeatis his omnibus: quin potius quaerite primum regnum Dei et iustitiam eius, et haec omnia adiicientur uobis. Atque idcirco proposita sunt praemia etiam in hacuita, et externa seu camnalia, simulut rudiores his ad bene agendu mouerentur,  \pend
\vspace{0.5cm}\noindent
\vspace{0.2cm}\rule{1cm}{0.2pt}\\ 
\hspace{0.2cm}\textit{mg}
\footnotesize 5ριs εαι ue Quues. 
\normalsize| 
\hspace{0.2cm}\textit{mg}
\footnotesize Exercitatio corporalis ministra est. Similitudo. 
\normalsize| 
\hspace{0.2cm}\textit{mg}
\footnotesize Pietatisuuilitas. Argumentuin ab utili. 
\normalsize| 
\hspace{0.2cm}\textit{mg}
\footnotesize Corporalis exercitatio non habet promissiones iustitiae aut prenii diuini. Matth.6. 
\normalsize| 
\hspace{0.2cm}\textit{mg}
\footnotesize Amplificatio a tempotre. Pietatis prae mia utramque  quta mancit. 
\normalsize| 
\section*{LIN I. EPIST. D. PAVLI }
\marginpar{[ p.98 ]}\pstart qui parum mouentur spiritualibus quae non percipiunt: simul ut omnes intelligerent, se quicquid habent boni, id a Dei bonitate accipere, ac propterea semper illi debere gratias agere, et in solo Deo spem ac fidem omnem collocare. Ex his itaque manifestum est, cur potius danda opera pietati quam corporali exercitationi, et cur, quae non sunt coniuncta pietati, reiicienda sint tanquam profanae et aniles fabulae.  \pend
\phantomsection
\addcontentsline{toc}{subsection}{\textit{Indubitatus sermo, dignusque qui omnibus modis approbetur. Nam in hoc et laboramus et probris afficimur, quôd spem fixam habemus in Deo viuente, qui est seruator omnium hominum, maxime fidelium. Praecipe hac et doce. }}
\subsection*{\textit{Indubitatus sermo, dignusque qui omnibus modis approbetur. Nam in hoc et laboramus et probris afficimur, quôd spem fixam habemus in Deo viuente, qui est seruator omnium hominum, maxime fidelium. Praecipe hac et doce. }}\pstart Confirmatio quod ita res se habeat de uera pietate, ut dictum est. Alibi quoque saepe utitur hac confirmationis forma Sunt enim uerba ad confirmandum valde idonea. Sic enim et initio tertii capitis, Indubitatus sermo, etc.  \pend\pstart Namin hoc et laboramus et probris afficimur.) Ratio confirmationis a signo. Quasi di cat: Quod etiam uera sint, quae asserimus, id vobis indicio certo esse potest, qu hac de caussa laboramus et pericula ingentia adimus. Ideo autem constantes sumus in tot aduersis, qu spem fixam habemus in Deo uiuente, (qui est seruator omnium hominum, quamlibet abiectorum et grauiter periclitantium in hoc mundo, dummodo tamen sua perfidia se ipsos a salute non excludant:) et certi sumus de doctrina nostra, ut nullis prorsus aduersis ab hac ueritate, quae in fide et charitate consistit, commendanda, deterreri queamus aut abduci. Caeterùm animat Timotheum his uerbis, ut strenue tueatur partes suae doctrinae contra falsos doctores, nullisque periculis a confessione ueritatis moueatur. Atque  efficit uera pietas cognita, ut pericula non reformidemus. Vbi enim quis certus est de genere doctrinae suae facile contemnit omnes aduersariorum machinationes, et retinet spem fixam in Deo uiuente. Hinc illa occurrunt, quae animos in certa spe confirmant: » Matth.5 Beati qui persecutionem patiuntur propter iustitiam. Et, Beati cum probra iecerint in vos homines, et insectati fuerint, et dixerint omne malum verbum aduersus uos,  mentientes propter me.1. Petri 4. Si Christianus affligitur aliquis, ne pudefiat, imô glo  rificet Deum in parte hac. Rom.5.Gloriamur super afflictionibus, scientes qu afflictio  patientiam pariat, patientia vero probationem, probatio autem spem, spes non pudefacit Vocat Deum viuentem, quia autor est uitae, tunc etiam vitam nobis dans, cùm in mundo videmur mori. Quoniam vero Deus viuens est, et ipse viuificans, debet spes nostra magis erigi, et fixior semper fieri.  \pend\pstart Qui est seruator omnium hominum, maxime fidelium) Deus seruat omnes homines, qui propter veritatem laborant vel probris afficiuntur, quique in fide permanent. Si qui salutem per eum non assequantur, id accidit, quia seipsos sua malitia et peccatis a salute excludunt. Porro in solo Deo omnem spem salutis suae defigant, eundem per ueram pie tatem colentes: caetera autem omnia quae adferuntur ceu ad salutem parandam conuenientia aut necessaria, si ad veram pietatem non pertinent, ceu aniles fabulae et profanae reiiciant.  \pend\pstart Praecipe haec et doce.) Breuis conclusio ad praecedentia, de vitandis doctrinis daemoniorum et amplectenda sana doctrina de uera pietate. Sunt uero haec omnia sic accipienda ab ecclesiarum ministris, quasi ad ipsos proprie dicta, omnibusque ; incumbit hoc pacto inuigilare gregi dominico, ne doctrinae daemoniorum siue aniles fabulae locum in pectoribus simplicium habeant.  \pend
\phantomsection
\addcontentsline{toc}{subsection}{\textit{Nemo tuam iuuentutem despiciat, sed esto forma fidelium in sermone, in dilectione, in spiritu, in fide, in puritate. }}
\subsection*{\textit{Nemo tuam iuuentutem despiciat, sed esto forma fidelium in sermone, in dilectione, in spiritu, in fide, in puritate. }}\pstart Sequunt uariae et breues paraeneses, inprimis episcopis et omnibus quoque ecclesiarum ministris mire salubres. Cum autem summopere cauendum sit episcopo, ne quid in ipso deformet muneris sui dignitatem: primo loco admonet Apostolus Timoth, ut sic se gerat in ecclesia, ne iuuentus ipsius despici queat. Mali homines in ministris ecclesiarum quiduis norunt reprehendere, et uitio plerumque uertunt, qu reuera uitio nequaquam esset uertendum Et in magnis uiris, iisq qui in publico uersant, uel minutissima errata statim no tant, ut Plutarchus in politicis docuit. Vnde Scipioni obiecta est Somnolentia: Alcibi adi optimo principi asscripserunt luxum: Cneo Pompeio, quunico digito caputscalpe\pend
\vspace{0.5cm}\noindent
\vspace{0.2cm}\rule{1cm}{0.2pt}\\ 
\hspace{0.2cm}\textit{mg}
\footnotesize Ratio confirmationis a figno. 
\normalsize| 
\hspace{0.2cm}\textit{mg}
\footnotesize Doctrinae certitudo spem et fide in aduersis corroboral 
\normalsize| 
\hspace{0.2cm}\textit{mg}
\footnotesize A Deo pen det salus noestra. 
\normalsize| 
\hspace{0.2cm}\textit{mg}
\footnotesize Nostro ipse rum uitio fit ut indigni fiamus salut te et gratia csse 
\normalsize| 
\section*{AD TIMOTHEVM CAP. II }
\marginpar{[ p.99 ]}\pstart Iet: cum aliud non posset clarissimis viris obiici. Sic vero in Timotheo si quid reprcheur potuisset, quem tamen non dubiu vixisse inculpatissime, id fortassis erat iuuentus. ltaque Apostolus cupiens prorsus nihil in eo posse culpari: et qa fortassis nomnulli ipsum quoque  Paulum notassent quasi parum prouidum, qui tam iuuenem virum tantae ecelesiae, Ephesiorum nem pe praefecisset: monet esit ita vitam omnem instituere, vt aetas ipsa dignitati elus nihil detraheret. Videri possit et occupatio huic loco subesse. Verosimile est Timotheum quasi inuitum subiuisse tantum munus, sicut tunc omnes potius refugiebant, ab eo munere que ambie bant: maxime autem hac ratione voluisse se excusare, qu esset iuuenis. Sed Paulus reiicit hanc excusationem, dicens: lta versate, vt nemo occasionc̃ habeat despiciendi tuam aetatem quoel quispiam caussari vellet diminutum esse inaetate, id compensato probis moribus et senili grauitate: Esto forma fidelium in sermone, in conuersatione, etc. atque hoc pacto fiet, vt nemo iuste aetatem tuam sit contempturus. Caeterùm iuuentutis nomine, vt Varoni placet, intelligi conuenst vitae spacium ab anno tricesimo, ad quadragesimum quintum: dictaque ; iuuentus, q ea aetate rempublicam militari opera possent iuuare. Hippocrates iuuentu. tem numerauit ab anno duodetricesimo ad tricesimum quintum. Itaque apparet limotheum florida aetate vixdum quadragesimum attigisse annum, quando ecclesiae Ephesinae praefuit. Proinde ea aetas nequaquam est contemnenda. Tametsi enim senectus ob multarum rerum experientiam utilis esse iudicetur ad gubernandum vel respublicas vel ecclesias: tamen, quia eam plerunque circumueniunt multa incommoda, iuuentus euam est ad magnarum rerum tractationem admittenda, maxime si studium, proba indoles, vitae integritas eam commendent. Extant autem exempla abunde multa, quae testan tur de hominum ingeniis aliter, atque de vlno, quod aetas solet potissimum commendare, iudicandum. Leguntur iuuenes quidam praeclaras et vtiles res gessisse, aetate ipsa nihil eos impediente. Sicute contra senes quidam imprudenter et improuide multa gesserunt, annorum numero nihilo facti meliores. Dona Dei non concluduntur annorum vel dierum spatiis. Deus enim linguas infantium facit disertas. Daniel.1.de tribus illis pueris, Anania, Azaria, Mizaele, dicitur, quod omne verbum sapientiae et intellectus, quod sciscitatus est ab eis rex Nabuchodonosor, inuenit in eis decuplum su per cunctos ariolos et magos. Et ipse Daniel tenera aetate Spiritu sancto afflatus, seniores impios et impuros mira prudentia conuicit et damnauit i.Samuelis3. Samueli natu inferiori apparet Dominus, et arguere iubet summum sacerdotem Heli nonagenarium. Hierem. I. Hieremias audit a Dño: Noli dicere, quia puer ego sum: ad omnia quae mittam te ibis. Tobus 1. De Tobia legitur, cum esset iunior omnibus in tribu Naphthalim, nihil tamen puerile gessit in opere. loannes Baptista admodum adhuc iuuenis do cebat et praedicabat poenitentiam et baptizabat. Christus tricesimo aetatis anno docendi munus est aggressus. Sed et in reipublicae administratione Solomon aetatis anno 12. mulierum pro paruulis quaestionem dissolust. losias quoque octo annorum aras idolorum subuertit, lucos ac templa excidit, sacerdotes idolorum iussit interfici. Et sa plent. 9. scriptum est, senectus venerabilis est, non diuturna neque annorum numero computata: Cani sunt autem sensus hominis, et aetas senectutis vsta immaculata. Sed  si vulgato quoque prouerbio dicitur, quod saepe etiam est olitor valde oportuna locutus, quid ni vir suuenis bene instiltutus a teneris, cum probis et bonis versatus possit mul ta vtiliter et cum laude perficere ? Haec placuit ideo diligentius commemorare, non quod vllo pacto senilis aetas sit minoris pendenda, quae etiam gentibus semper magno habita est in precio, vt natura ipsa docuit: vnde Ouidius: Magna fuit quondam capiti reuerentia cano: Sed vt adolescentiam quoque ad bene de se sperandum animarem, adderemque  uelut calcar ad gnauiter decurrendum in puscherrimo studiorum et uirtutum stadio, neque  diffidant, se ubi profecerint, etiam ad rerum magnarum administrationem prouehendos. lam uero quibus rationibus iuuenilis aetas reddi queat omnibus uenerabilis, apta enumeratione congerie ostenditur.  \pend\pstart Sed esto forma fidelium.) En qua ratione possit efficere, ut iuuentus non contemneret, quasi dicat: Hoc pacto consequeris, ut nemo iuuentutem tuam sit contempturus: si uelut exemplar quoddam temetipsum in omnibus pponas.Oportet namo epscopumuelut  \pend
\vspace{0.5cm}\noindent
\vspace{0.2cm}\rule{1cm}{0.2pt}\\ 
\hspace{0.2cm}\textit{mg}
\footnotesize vaecag rinĩ de die natai.. 
\normalsize| 
\hspace{0.2cm}\textit{mg}
\footnotesize Afnut 
\normalsize| 
\hspace{0.2cm}\textit{mg}
\footnotesize Phcue 
\normalsize| 
\hspace{0.2cm}\textit{mg}
\footnotesize EPaa I4 
\normalsize| 
\hspace{0.2cm}\textit{mg}
\footnotesize Quomodo iuuentus fiat honorabiuis. 
\normalsize| 
\section*{IN I. EPIST. D. PAVLI. }
\marginpar{[ p.100 ]}\pstart speculum esse, in quod omnes intueant, ibique perspiciant, quae ipsos deceant uel dedeceant Plerunque ad superiorum uitam solent inferiores suam quoque formare et fingere. Ques pastores suos sequuntur. 1.Petri5.Pascite, qui in uobis est, gregem Christi, curam illius agen tes non coacte, sed uolentes, non turpiter affectantes lucrum, sed propenso animo, neq ceu dominium exercentes aduersus cleros, sed sic, ut sitis exemplaria gregis. Debens autem episcopus forma fieri fidelium, et non tantum fidelium, sed etiam ut insideles eius exemplo ad pietatem amplectendam commoueantur.  \pend\pstart In sermone.) id est in sana doctrina, ut purem et syncere sermonem siue uerbum Dei administret. Vel potius, in sermone, id est, in omnibus uerbis, in colloquiis, ne fucate cum ullis agat, ne mentiatur, ne rustice aut nimis dure homines alloquatur. lacobus 3. Si quis in sermone non labitur, hic perfectus estuir. Coloss.4. Sermo uester semper cum gra, tia sit sale conditus, ut sciatis, quomodo oporteat uos unicuique respondere.  \pend\pstart In conuersatione.) id est, uita et’externis moribus, que decet esse sanctam, sobriam, modestam, quietam: in summa, omni ex parte irreprehensibilem, tam domi quam foris. Et certe doctrina eorum siue uerba non magni fient, quord uita deprehenditur improba. Christus Matt.s.uo cabat apostolos lucem mundi, sed hoc addidit: Sic luceat lux uestra coram hominibus,  , ut uideant opera uestra bona, et glorificent patrem uestrum qui est in coelis.  \pend\pstart In dilectione.) neminem aspernetur, omnes ex aequo in Christo, et propter Christum diligat, eaque praestet ofneia charstatis erga pauperes, erga quoscunque , ut omnes intelligere possint, eius mentem spiritu charitatis abunde insfructam. Christum in diligendo referat, qui paratus fuit etiam animam pro grege ponere.  \pend\pstart Inspiritu) declaret se spiritu Dei agi, non humanis affectibus: reluceat in eo spiritus prudentiae, timoris, intellectus, fortitudinis, etc.  \pend\pstart infide.) Constantem em̃ in hac esse decet, neque vnquam deficere, etiansi grauissima pericula immineant. Debent quidem omnes fide pollere: Sed episcopi debent omnibus quasi praeire. In puritaie.) Vita debet esse pura, casta, illibata, candida, simplex: Graece est, μάγνέα, quae epitasis quaedam est omnis temperantiae et continentiae: ideoque apte postremo loco ceu concludere uolens eam posuit. Ex his colligitur, quod, qui autoritatem suam integram con stare uolunt apud plebem, suumque munus aut etiam personam suam nequaquam uulgo contenni, neque uel potestatem, vel aliam conditionem in ipsis damnari: eos istud consecuturos, si ita se gerant, uti Apostolus ĥe praescribit: si forma sint fidelium in sermone: Haud dubiem omnes eos summo in honore sunt habituri. His moribus fuit Ambrosius, qui Theodosium imperatorem arguendo inducere potuit ad publicam poenitentiam. Sin aurem secus se gerant, sibi potius imputent, quam populi malignitati, quod paruipendantur.  \pend
\phantomsection
\addcontentsline{toc}{subsection}{\textit{Donec venero, aitende lectioni, exhortationi, doctrinae. }}
\subsection*{\textit{Donec venero, aitende lectioni, exhortationi, doctrinae. }}\pstart Non dubium quin Timotheus homo adhuc iuuenis licet probe exercitatus, tamen aliquando in perplexis causis quae ad eum per credentes de religione perferebantur, uti solet sieri, angeretur ac uehementer optaret Paulum praeceptorem suum videre praesentem, vt ab eo plenius de omnibus doceretur. Et Paulus eam ob caussam habuit in votis Ephesum aliquando redire, sicut in fine tertii capitis huius epistolae indicauit. Nunc vero iubet Timotheum interea temporis totum intentum esse literarum sacrarum lectioni, e quibus perinde discere queat omnia scitu necessaria, imo et plenius et certius, atque a quocunque praeceptore., Attende inquit, lectioni: Graece προσέχεπ ητ ἀυαγνώσει. προσέyς, id est, totus sis addictus, toto animo adhaereas: id em significat προσέγσ. Et ἀνάγνωσις non simpliciter lectionem significat, sed diligentem, cum recognitione et examinatione ac iudicio. Vult ergo Timotheum non oscitanter legere aut otiose obiterque ́, aut tantum succisiuis horis, qd' nonnulli putant satis: sed reuera desiderat omni tempore intentum, et quidem accurate legentem. Porro si ita Apostolus praecipit Timotheo, quem tamen iam paulo ante dixit enutritum fuisse in recta doctrina, et a teneris institutum in sacris literis: quid nobis faciendum? quid qbusdam aliis, qui dum Timothei esse volunt muneris dignitate, tamen eruditione vix vmbram Timothei referunt? Sed ideo commendat Apostolus tanto viro lectionem scripturarum, qu ea non solùm auertit animum a multis malis, uerumetiam confert homini quicquid ad omnem pietatem est scitu necessarium. Si hominem consulas vel angelque , possunt ii falls et fallere:  \pend
\vspace{0.5cm}\noindent
\vspace{0.2cm}\rule{1cm}{0.2pt}\\ 
\hspace{0.2cm}\textit{mg}
\footnotesize Epifcopu est speculum et exemplo gregis Chri stian. 
\normalsize| 
\hspace{0.2cm}\textit{mg}
\footnotesize Conterfauïe Episcopi quais esse debeut. 
\normalsize| 
\hspace{0.2cm}\textit{mg}
\footnotesize Dilectio. 
\normalsize| 
\hspace{0.2cm}\textit{mg}
\footnotesize Iohanu 
\normalsize| 
\hspace{0.2cm}\textit{mg}
\footnotesize Spiritus Dei affectibus carms reluctatur. tas ex fide e spiritu pro feiscitur. 
\normalsize| 
\hspace{0.2cm}\textit{mg}
\footnotesize Puet constat tasonper ad salut.m necessgria. vitae puri
\normalsize| 
\hspace{0.2cm}\textit{mg}
\footnotesize Saioura fe cra fidus ma gister et do ctor impert torum. πρησέαο qud. αναγωος est exquisit. sectio et recognitio. Iohan.5. Scrutamint scripturas. A maiore. vsus lectionis feriptirum mudtiplex
\normalsize| 
\hspace{0.2cm}\textit{mg}
\footnotesize 
\normalsize| 
\section*{AD TIMOTHEVM CAP. IIII. }
\marginpar{[ p.101 ]}\pstart si adis simplex uerbum Dei, id clare et citra dolum aperit tibi omne ueritate. Vnde eit hic egregia et laus Scripturarum, et exhortatio ad earum tractationem et obiurgatio nostrae segnitiei proposita. Addit:  \pend\pstart Exhortationi et doctrinae.) Adiuncta haec lectioni esse debent, uel potius est causa finalis siue usus lectionis. Exigua res est, si episcopus multa legat intra domesticos parietes, si ipse solus legat et intelligat multa: sed debet in hunc usum legere, ut lectione adiuuet̃ ad suum munus dexterius obeundum, ut etiam caeteri fructum aliquem ex lectione eius capiant. Capient autem fructum, si diligenter exhortet et doceat. Hec sunt duo precipua officia episcopi, ad haec duo debent studia eius omnia tendere, nisi haec itidem faciat, grauis ei poeha a Deo imminet infligenda. Exhortatione opus est tum iis, qui deiectiore sunt animo, intuentes in peccatorum suorum magnitudinem, neq ad fidendum in misericordia Dei satis sunt erecti:tum iis, qui tardius ad bene agendum progrediuntur, uel etiam a recte exorbitant: itaut exhortationis nomine intelligi conueniat, et consolationem, et ad monitionem, et increpationem, et obsecrationem, et omnia illa in quibus suadere uel dissua dere conuenit. Doctrina autem tunc adhibenda, cum, qui de ueritate dubitant, uel non satis eintelligunt, sunt plenius instruendi. Qui igit exhortando ac docendo cessator est, grauiter hic reprehenditur. Qui non possunt ista efficere, frustra ornantur inani titulo ministro, rum ecclesiarum: qui uero possunt, et non faciunt, omnium peccant grauissime. Neque uerosim le est eos fore suo gregiformam siue exemplar bene uiuendo, qui formam proponere nor dignantur saltem exhortando uel docendo. Animaduertant denique hoc loco omnes, quibus eccle siarum cura est demandata, quod praecipuë ab iis requiraf, nempe iugis lectio sacrarum literaru sine qua nec exhortari, nec docere poterunt in summa, nihil recte in officio suo facere  \pend
\phantomsection
\addcontentsline{toc}{subsection}{\textit{Ne neglexeris quod in te est donum, quod datum est tibi per prophetiam, cum impositione manuum, autoritate sacerdotii. }}
\subsection*{\textit{Ne neglexeris quod in te est donum, quod datum est tibi per prophetiam, cum impositione manuum, autoritate sacerdotii. }}\pstart Argumentis aliquot Timotheum hortatur ad diligentiam, vt omni conatu studeant suo muneri satisfacere. Primum argumentum est a ratione officii.  \pend\pstart Ne neglexeris quod in te est donum.) Donum nimirum vocat publicum illud et sublime munus episcopi. Debet autem vnumquemque monere ad diligentiam muneris sui ratio. Miles es: strenuus igitur sis in militia. Senator es, consule igit, quantum in te est, rebus politicis. Sic episcopus es, prospice igitur ecclesiae. Et saepe occurrunt huiusmodi argumenta a vocatione, a dignitate ministerii. Caeterùm quoniam ad episcopatum non prouehitur, neque rite obibit id munus, nisi qui illustratus fuerit spiritu prophetiae, quae est facultas enarrandi scripturas: deinde nisi imponantur ei manus, id est, publice et legitime uocetur ordineturque ́: ideo adiecit.  \pend\pstart Quod datum est tibi perprophetiam cum impositione manuum.) Ita vt hic duo itidem subsint argumenta: alterum a dono prophetiae, quo ostenditur eum esse idoneum ad id praestandum:alterum a publica ordinatione, quo ostenditur, ei legitime commissum id munus, atque ideo iusta autoritate posseac debere id explere. De propheria semel dictum est. Quod ad impositionem manuum attinet, sciendum impositionem manuum ex vetustissimo Hebraeorum instituto recepta fuisse obseruatamque Apostolorum temporibus, ad testandum exterius, qu talis, cui impe nebat manus, iudicaret aptus ad publicam functionem, sicque Deo ipsi commendabatur, et tacita interdum uota pro eo fiebant. Sic legimus Actor. 6.multitudinem congregatam elegisse u ros sapientia et Spiritu S.plenos ad diaconatum, eisq imposuisse manus. Et Genes.48 Iacob imponit manus duobus filiis Ioseph, et bene illis precaf. Et Leuit.1. manus ponitur super caput hostiae immolandae. Christus pari instituto manus imponit paruulis Idem de credentibus dicebat, et super aegros manus imponent et bene habebunt. Fiebal autem haec impositio manuum, in electione episcopi uel diaconi, ad approbandam person?̃ multitudini, per seniores, penes quos erat haec autoritas. Vnde hicadiectum est, cum imi positione manuum autoritate sacerdotii: vel, ut graece significantius dicitur et dilucidius, uexe ἐπιδεσεως τίῦ́γορῶν τπορισθυτερίς. πρεσθυτερᾶον significat ipsam seniorum congregationem. Et conueniebat plurium consensu accedente, eum, qui eligebatur, idoneum hoc externo signo comendarineo faclle cuiqua hoe pacto imponebat manus, nis cujus fide syncema,  \pend
\vspace{0.5cm}\noindent
\vspace{0.2cm}\rule{1cm}{0.2pt}\\ 
\hspace{0.2cm}\textit{mg}
\footnotesize 4 e8pgaiis verbum Dei caret fuce, uno estueritas ipsaIobana7. Psal.119. caussa finalis et usus le ctionis sacra rum scripturaeum. OfficumEpiscopi in exhortatione et doctri na uerjatur. Exbortatione quando, utendum. Doctrina quando adhibenda. Aminori. 
\normalsize| 
\hspace{0.2cm}\textit{mg}
\footnotesize 1.Argumentum bortato, num, aratione officit. Douum hoc loco munus significat etpiscepi. 
\normalsize| 
\hspace{0.2cm}\textit{mg}
\footnotesize 2. Argumentum a dono prophetiae. 3. Argumen tum a publicaordinatione siue le gitima uoca Rionc. ampositio. manuum. Marciro. Luc.18. Maitbus 19. Axcite. 
\normalsize| 
\section*{IN I. EPIST. D. PAVLI. }
\marginpar{[ p.102 ]}\pstart temque haberent, pbe exploratam:ideoque Paulus cap. sequenti: Manus ne cito cui imponas  \pend\pstart Caeterùm si vllas huiusmodi ceremonias conuenit in electione ministrorum in ecclesia obseruare: quanto satius esset saltem hac simplicissima forma et ceremonia esse nos contentos, quam tot superuacanea et propemodum ridicula, neque intellecta a populo accersere? Proinde quicunque in publico ecclesiae ministerio sunt constituti, cogitent se quoque  per haec argumenta ad diligentiam admoneri, videlicet per officii ipsius rationem, per de num prophetiae et per publicam suam ordinationem.  \pend
\phantomsection
\addcontentsline{toc}{subsection}{\textit{Haec exerce, in his esto, vt tuus profectus manifestus sit in omnibus Attende tibi ipsi et doctrinae persiste in his. Nam id si feceris, teipsum seruabis, et eos qui te audierint. }}
\subsection*{\textit{Haec exerce, in his esto, vt tuus profectus manifestus sit in omnibus Attende tibi ipsi et doctrinae persiste in his. Nam id si feceris, teipsum seruabis, et eos qui te audierint. }}\pstart Conclusio huius capitis, quo demonstratu est, quomodo se geret episcopus in ges nere erga omnes.  \pend\pstart Hec exerce, in his esto. Appositissimis verbis eum vrget, vt in his quae iam dixit, praesertim de gerendo se pro exemplari probitatis, et de legendo, exhortando, docendoque totus sit occupatus. Prius diximus quid sit proprie exercere se in aliqua re, nimirum se tontum alicui rei dedere, et assiduo ei intendere. Graece est, rαῦταμελέτα, a verbo μελετάω, exerceo, excogito, meditor, vt plane deceat episcopum toto animo his rebus esse iugiter intentum, Vt tuus profectus manifeslus sit in omnibus.) Extimulat eum argumento a causa finali, quae oritur ex charitate, per quam debemus omnia facere ad aedificationem proximorum: quod maxime pertinet ad cos, qui, vt Christus aiebat ad discipulos suos, sunt lux mundi, et. lucere debet eorum lux, vt homines videant opera eorum bona, prouocenturque ad vitae emendationem. Ideo Matth.5. Candela iubetur poni super candelabrum, ut luceat omnibus qui sunt in domo. Ac vide quomodo vocet profectum in virtutibus, et profectum ipslus, Vt tuus, inquit, profectus. Vberrimus autem est profectus qui in spiritualibus bo nis acquiritur: lta namque et capite sexto huius epistolae pietas vocat magnus quaestus ac lucrum, quouis lucro rerum terrenarum incomparabiliter praestantior: similiter ua φατικῶς addidit, in omnibus.  \pend\pstart Attende tibi ipsi et doctrinae.) id est, da operam, vt et vitae probitatem et doctrinam senpoffet iuseri per retineas, neque patiaris harum duarum rerum quicquam in te desiderari. Annectit rationem ab vtili.  \pend\pstart Nam id si feceris. I. pariter recte viuis et recte doces, teipsum seruabis, et eos qui te auannulo. dierint.) Conduplicatum bonum sequitur, quando et personae ipsius, et personae aliorum, diligentia illa erit salutaris. Neque enim tam suae propriae saluti inuigilare debent episcopi, quam saluti omnium pariter subditorum. Recte viuendo ac docendo, conscientiae suae in primis satisfaciunt, eamque probant Deo: deinde alios ad doctrinam suam amplectendam ac mores insequendos erigunt. Hebr.13. de episcopis dicit, quod inuigilant pro animabus vestris tanquam rationem reddituri. lacobus 5. Qui conuerti fecerit peccatorem ab errore viae suae, saluam faciet animam a morte. Daniel.12. Qui ad iustitiam erudiunt multos, quasi stellae fulgebunt in seculum et in aeternum.  \pend
\endnumbering\beginnumbering\section{CAPVT V}
\phantomsection
\addcontentsline{toc}{subsection}{\textit{\huge\textbf{S}\normalsize Eniorem ne sauius obiurges, sed adhortare vt patrem, iuniores vt fratres, mulieres natu grandiores vt matres, iuniores vt sorores, cum omni castitate. }}
\subsection*{\textit{\huge\textbf{S}\normalsize Eniorem ne sauius obiurges, sed adhortare vt patrem, iuniores vt fratres, mulieres natu grandiores vt matres, iuniores vt sorores, cum omni castitate. }}\pstart \huge\textbf{N}\normalsize VNC quomodo episcopus aget cum cuiuscunque sortis hominibus. Ac primum qua ratione cuiuscunque aetatis vel sexus homines adhortabitur. Certum autem est non Timothei tantùm caussa hanc scriptam epistolam, qui haud dubie satis nouit et vidit in ipso Paulo, quomodo humaniter et ciuiliter, sed praeterea quomodo excellenti quadam modestia Christiana cum omnibus foret transigendum: Sed multo magis propter alios, qui pari munere in ecclesiis essent functuri. Neque enim omnes sunt tam politici eaque ingenii dexteritate et candore praediti, vt videre possint, quid ipsis decorum sit vel indecorum. Itaque Apostolus, quoniam iamiam dixit, quod episcopo incumbit exhortari, ac docereincipit praeseribere, quomodo cum decoro id faciet.  \pend
\vspace{0.5cm}\noindent
\vspace{0.2cm}\rule{1cm}{0.2pt}\\ 
\hspace{0.2cm}\textit{mg}
\footnotesize Conclusio hu tus partis, qua ostensum est, quomo do segeret episcopus ii genere erg omnes. Causa finalis ex charitate. Matth.5. 
\normalsize| 
\hspace{0.2cm}\textit{mg}
\footnotesize ἱππεσασις posset inseri bi tanquam symbolum in annulo. 
\normalsize| 
\hspace{0.2cm}\textit{mg}
\footnotesize Quomodo se gerat epscopus erga omnis condi stionis homines 
\normalsize| 
\section*{AD TIMOTHEVM CAP. }
\marginpar{[ p.103 ]}\pstart Seniorem, inquit, ne sauius obiurges, sed adhortare vt patrem.) In consideratione aetatis es sexus primo loco posuit seniorem virum. Ac recte noluit huiusmodi hominem saeuius obiurgarenon solùm quia parum decuisset, quod Timotheus homo iuuenis asperius increpasset aliquem natu grandem: verùm quia in genere conuenit episcopum inreprehendendis seu corrigendis vitiis summa vti modestia et mansuetudine. Praeterea quia id quoque exigit hominum, qui errarint, ingenium, quorum maior pars duci potius vuit quain cogi.lgitur seniorem, inquit, ne laeuius obiurges: Graece est πραυθυτεσω χη πιλπλητης. πρευθυτορet hic proprie de aetate intelligi conuenit, non de officio vel publicq ministerio, id quod liquet, tùm ex eo, quod mox adiicit per Antithesin de iunio, ribus: tùm ex eo, quod, postquam de viduis egit, subiicit de presbyteris ecclesiarum ministris. Ετεπλκσιω significat castigo, obiurgo quodammodo et verbere addito: ut Metaphora subesse uideatur: quomodo etiam uerbera linguae interdum metaphoricas usurpant pro saeua maledicentia. Nullo pacto igitur in eos, qui aetate sunt maiores, utendum asperiore reprehensione: sed admonendi suntut patres: id est, affectus exprimendus est uerae minimeque fucarae charitatis:atque eodem pacto agendum est co illis, ut filius quispiam acturus est cum suo patre. Natura autem iubet filios cum parentibus agere modestissime semperque honorem quendem deferre, etiamsi contingat fihos ad amplissimos honores euectos, parentibus obscuris permanentibus Itaque uolens Apostolus demonstrare, quomodo episcopus integerrim o affectu reprehendet siue adhortabitur cuiuscunque ordinis homines, proponit similitudines ab affectibus naturalibus: Id quod etiam uidere est in sequentibus uerbis: Iuniores vt fratres, mulieres nata grandiores vt matres, iuniores vt sorores: Ne quis putet cul adulatoria arte per λατέχρησιν huiusmodiurendum blandis appellationibus. Caterum quia beneuolentia iuuenis uiri praesertim episcopi, erga uirgines uel iuniores mulierculas facile posset degenerare, aut reuera in turpem et impudicam beneuolentiam, aut saltem m impudicitiae suspicionem: ldeo prudenter attexit Apostolus circumstantiam cum omni casti tate: ut nimirum intelligat episcopus, cui cum iunioribus delinquentibus mulierculis perinde negocium esle debet, atque cum senioribus, sibi maximopere cauendum esse ne quid cum his committat, unde uel sinistram suspicionem incurret. Rom.12. Parare debemus honesta in conspectu omnium hominum. Etiam tuta plerunque sunt suspecta. Tametsi nonnulli, quoniam Graece habetur, πάση αγνέα, atque αγνεα sonat nobis purita tem, candorem, simplicitatem, summamque ; sobrietatem ac temperantiam, posteriora ista uerba ad omnia praecedentia referri malunt: quandoquidem et cum senibus, et cum iuuenibus, tam uiris quam foeminis, pari modestia agi par est, neque apud ullos licet saeuitiam, aut odium, aut blanditias ulla nota ostendere. Proinde his omnibus id demum sibi uult Apostolus, quôd ministri ecclesiarum in exhortando siue corrigendo, debent ob omni uel specie saeuitiae, inuidiae odii auaritiae, arrogantiae incontinentiae, abesse quam longissime: Sed potius summam charitatem, eamque minime fucatam, et prudentiam spi rirus ubique ostendent. Amari magis uelint a suis, quam timeri: et, cum medici animorum sint, curare et sanare debent, non exulcerare aut inficere. Potestas enim datur a Do mino non in destructionem, sed in aedificationem, ut dicitur 2. Corinth. 10. Et seruun Dominiz. ad Timoth.2. non oportet pugnare, sed placidum esse erga oës, propensun ad docendum, tolerantem malos cum mansuetudine. Quanquam dum ita agendum est cum iis, qui per ignorantiam imprudentiamue lapsi sunt, et in quibus signa apparentuitae me lioris seu resipiscentiae: tamen qui relabunt́, maxime si accedat malitia, et tanquam, (ut 2.Pet 2.dicif.) canes ad uomitum, et sues reuersae ad uolutabrum, iram Dñi in se prouocant, et uident aegre blandis monitis reducendi in uiam rectiorem, tales, inque , conuenit asperic re obiurgatione urgeri. Vulneribus eorum post frustra adhibitum oleum infundendum estunum, addendus Sal, qui punget, imo tandem ferrum cadens, ne uicinae carnes putrefiant, te tumque corpus paulatim inficiat. Certe enim qui piis monitionibus prouocati contemnunt redire ad frugem, ii uel ad aliorum tùm timorem, tum consolationɇ̃ autoritate diuini ue bi, quod est gladius ecclesiae, sunt coërcendi, et si expediat, ab ecclesiae corpore resecandi, ne peruicatta sua et pteruia uniuer sum corp' morbidum reddant Sicuero Christus ipse  \pend
\vspace{0.5cm}\noindent
\vspace{0.2cm}\rule{1cm}{0.2pt}\\ 
\hspace{0.2cm}\textit{mg}
\footnotesize Actais ba bendaest ra de
\normalsize| 
\hspace{0.2cm}\textit{mg}
\footnotesize Valerius de honoranda senectute li. 2.cap.1. Ab affau naturdli. 
\normalsize| 
\hspace{0.2cm}\textit{mg}
\footnotesize CRuxcae scopi. 
\normalsize| 
\hspace{0.2cm}\textit{mg}
\footnotesize Quomodo episcopus se geret erga reprobos et malitiose Deum prouocantes. Metaphora a medicis de riuata Rebus desperatis utendum est amputatione membro rum putridorum. 
\normalsize| 
\section*{IN I. EPIST. D. PAVLI. }
\marginpar{[ p.104 ]}\pstart magno ardore obiurgauit praefractos et tumidos Scribas et Pharisaeos, uo cans eos pro geniem uiperarum, filios diaboli, etc. Alibi Petro dicitur: Vade post me Satana. Nam placide leniterque agere cum inferioribus, ubi opus est acrimonia, non doctoris sed seductoris et corruptoris est. Sic Apostolus quoque ad Galat, cap.3.uocat αὐοήτος, stula tos, amentes: in qua Epistola multa magno affectu et aspere pronunciata occurrunt. Et i. Corinth.4. Quid uultis, inquit, cum uirga ueniam ad uos, an cum Charitate et spi ritulenitatis? Sed interim circumspectum utique decet et prudenter discernere, ubi, et quando, et apud quos quid deceat.  \pend
\phantomsection
\addcontentsline{toc}{subsection}{\textit{Uiduas honora que vere viduae sunt. Quod si qua vidua liberos aut nepotes habet, discant primùm propriam domum pie tractare, et vicem rependere maioribus. Hoc enim est honestum et accoptum coram Deo. }}
\subsection*{\textit{Uiduas honora que vere viduae sunt. Quod si qua vidua liberos aut nepotes habet, discant primùm propriam domum pie tractare, et vicem rependere maioribus. Hoc enim est honestum et accoptum coram Deo. }}\pstart Praecipit, quomodo episcopus aget cum uiduis mulierculis. Inde uero praeceptio haec nata, quod in primitiua ecclesia, pauperes, uiduae praesertim, quae et aetate fractae, et fortassis Euangelii caussa uiris et facultatibus spoliatae, ex ecclesiarum facultatibus quas pii ditiores conferebant, frugaliter alebantur, cuius instituti rationem uero simile est ludaeis, quorum lex uetuit mendicitatem, etiam antea obseruatam. Quoniam uero nonnihil molestiarum habebat ea facultatum ecclesiasticarum dispem satio, prudenterque uidendum erat per episcopum, qualibus eleemosyna distribueretur: neque ita dandum erat citra delectum quibusuis, ut tenues ecclesiae facultates tandem non sufficerent alendis uere egentibus: ideo paulo exactius de uiduis hic disserit, docens: quales nam uiduae forent ad eleemosynarum partem admittendae, et quales rursus excludendae. Sic autem primo ait in genere.  \pend\pstart Viduas honora, quae vere viduae sunt.) Vides rectissime adiectam particulam, quae uere uiduae sunt, ad discrimen illarum, quae quidem nomine et honore uiduae uidentur, at sic interim agunt in luxu et uoluptatibus, ut nulla coniugata magis. Has siquidem non decet inter viduas numerare. Eas autem, quae vere viduae sunt (quae quomodo agnosci possit, mox describetur) iubet non statim admittendas ad ele emosyna rum participationem: sed generaliter dicit honore afficiendas, id est, omni beneuolentiae affectu prosequendas, et, si opus fuerit, consilio, intercessione, etiam auxilio omniadiuuandas. Sic namque vocabulum honoris saepe accipitur, vt in praecepto illo: Honora patrem et matrem: sic namque Christus ipse Matth.15. vsum vocis illius explicat. Porro ut constaret, quales nam viduae etiam ex his, quarum vita esset probitate spectabilis, forent ecclesiae facultatibus adiuuandae, disertis verbis, ex ratione circumstantiarum caussam hanc expendit.  \pend\pstart Quod si qua vidua liberos aut nepotes habet, discant primùm propriam domumpie tractare, et vicem rependere maioribus.) Ratione numeri variati, (nam et in Graecis ita quoque sonant verba) diuerse exponunt hunc locum: Μαυθανετωσαν, id est, discant, quidam ad viduam retulerunt, alii ad liberos seunepotes. Ambrosius itamaluit reddere. Si qua autem vidua filios vel nepotes habet, discat primum domum suam pie tractare. Ex qua versione haec sequitur sententia: Si qua est vidua, cui sunt liberi vel etiam nepotes, ea nequaquam ob id putet sibi otiandum, ac deinceps viuendum eleemosyna ecclesiae, quasi quae aetate exacta sese satis honeste gesserit: imo adhuc ea se occupabit in administratione domus, eo more quo ante egit, eamque ad honorem Dei pie instituet: et vicem rependet maioribus, eandem curam gerens promouendis suis li beris et nepotibus, quam maiores ipsius impenderunt in ipsa promouenda. Reuera enim quantum ei benefecerunt maiores, qui ipsam benigne et honeste semper instituerunt: tantundem debet ipsa et suis, eadem praestando piae institutionis officia: ac nisi istud faciat, ingrata quodammodo erga suos maiores existit, neque agnoscit, quam multa bona a suis maioribus acceperit, dum suis liberis paria non rependit. Atque hac ratione (si ad viduas refertur) intelligitur, quomodo ςοργαι φυσηκσὴ semper debe\pend
\vspace{0.5cm}\noindent
\vspace{0.2cm}\rule{1cm}{0.2pt}\\ 
\hspace{0.2cm}\textit{mg}
\footnotesize Mattbus 3. Ioann.8. Math. 10 
\normalsize| 
\hspace{0.2cm}\textit{mg}
\footnotesize Pradentiae piscopi. 
\normalsize| 
\hspace{0.2cm}\textit{mg}
\footnotesize Manota uerari uiduarum el Fse-- 
\normalsize| 
\hspace{0.2cm}\textit{mg}
\footnotesize Exod E 
\normalsize| 
\section*{AD TIMOTHEVM CAP. V. }
\marginpar{[ p.105 ]}\pstart ant per familias propagars, neque unquam possint affectus parentum erga filios sopiri aut intermitti. lam hic subinrelligitur, has viduas, quae sic se gerunt erga liberos aus nepotes, merito apud hos victum fabituras, proindeque ecclessae eleemosynis nequaquam alendas. Alii retulerunt verbum μαυlανέτοσαν, id est, discant, ad liberos seu nepotes, id quod facit Theophylactus: vt sit sententia:liberis illis seu nepotibus viduae dan dam operam, ut pie sancteque familiam suam instituant, id quod ex ipsa uidua discere possint: interea autem merito viduam ipsam alant, quae in docendo administrationem familiae ipsis prodest, et cui tenentur vices rependere. Ac subest argumentum a lure naturae, Rependant vicem maioribus. In natura enim positum est, et leges etiam latae sunt a sapientibus, vt liberi alant senes parentes. Si enim ciconiae tale pra bent nobis exemplum (vnde illud αὐτιπελαργο͂ν, quod alias quoque dicitur yuοΒοτκοιο turpissimum fuerit, si nos eandem pietatem homines hominibus non praeste. mus. Et lex erat, vt liberi parentes alerent aut vincirentur: qui in calamitate parentes descrerent, insepulti abiicerentur. Homerus de quodam lliad.4. pronunciauit, eun ideo diuinitus e vita sublatum esse ante tempus, quod nutricationem denegasset. Ideoque in scripturis senes vocant filios baculos suae senectutis, vt habetur Tobiae 5. propterea quod sperent se in grauida aetate per filios tanquam baculos sustentandos. A quomodocunque accipiantur verba Pauli, et quamcunque interpretationem velis sequi, adiecta est duplex ratio Suasoria, cur ita facere seu viduas, seu liberos conueniar.  \pend\pstart Hoc enim est honestum et acreptam coram Deo.) Ab honesto et a voluntate Dei. Honestum est id coram mundo, praeterea gratum ipsi Deorac proinde nullo modo debent viquae officia illa negligere in liberos seu nepores, neque libri in ipsas viduas. Proinde ex his Pauli verbis perspicuum fit, huiusmodi viduas, quae liberis suis seu nepotibus vsui esse possunt in administranda familia, vel quarum liberi seu nepotes iam ipsi familias instituerint, merito ab his alendas, neque ecclesiam earum cura grauandam. leaq iure de lectum habendum eorum, quibus impendetur publica eleemosyna, nec temere quibusuis dandum: tùm episcoporum interesse, prouidere, vt diaconi, qui praesunt distribuendis eleemosyais seu ministerio mensarum, in officio non sint cessatores. Praeterea pium et Deo valde acceptum opus, propriam domum pie tractare: vbi Graece habetur δὐσεθωρ, quod pprie significat pse, religiose, adeoque in flde et timore Dei instituere. Neque enim satis est, si curetur sedulo, ne quid decedat in re familiari, sed multo ma gis opus est, imo istud primo loco requiritur, vt doceantur omnes familiares reuereri Deum, sancte agere priusquam quaestuose. Postremo pium esse et aequum, vt iuniores vicem rependaut maioribus, omni honoris officio etiam victui necessar ria suppeditando.  \pend
\phantomsection
\addcontentsline{toc}{subsection}{\textit{Porrô que vere vidua est ac desolata, sperat in Deo, et perseuerat in obsecra tionibus ac precationibus noctu dieque . Porro quae in deliciis versatur, ea viuens mortua est Et haec praecipe, vt irreprehensibiles sint. Quod si que suis et maxime familiaribus non prouidet, sidem abnegauit, et est infideli deterior. }}
\subsection*{\textit{Porrô que vere vidua est ac desolata, sperat in Deo, et perseuerat in obsecra tionibus ac precationibus noctu dieque . Porro quae in deliciis versatur, ea viuens mortua est Et haec praecipe, vt irreprehensibiles sint. Quod si que suis et maxime familiaribus non prouidet, sidem abnegauit, et est infideli deterior. }}\pstart Nondum explicat, qualibus viduis ecclesla subuenire debeat: Sed quoniam dixit eas quae verem viduae essent, honorandas: ideo hic instituit descriptionem eius viduae quae pro vera vidua haberimeretur. Describit autem eam quasi a Signis: arque vt ma gis elucescat eius conditio, addit et per αυτιθεσιμ descriptionem eius, quae falso vidua reputatur:sicitaque ait:  \pend\pstart Porron quae vere vidua est ac desolata, sperat in Deo, et perseuerat.) quasi dicat: Haec sunt signa quibus discernere possis vere viduam, qualem nos diximus honorandam, ab ea, quae falsd viduam se gloriatur. Si qua est desolata: id est, mariti, liberorum, neporum, consanguineorum, denique totius mundi auxilio destituta. Acsperat in Deo: id est, ex solo Deo vult pende re, nulli placere studet, nisi soli Deo, mundum et omnia, quae mundi sunt, contemi nens t atque id quo testatum faciat, Perseuerat in elseerationibus ac precationibus no\pend
\vspace{0.5cm}\noindent
\vspace{0.2cm}\rule{1cm}{0.2pt}\\ 
\hspace{0.2cm}\textit{mg}
\footnotesize luoepreugs lionis uarielus. 
\normalsize| 
\hspace{0.2cm}\textit{mg}
\footnotesize  NtuxeA minori, siue a brutis. αππασnae 
\normalsize| 
\hspace{0.2cm}\textit{mg}
\footnotesize Τηροβοσuσῖο, Lexde alen as parenubus  Liberorum officia erga parentes cost que senes
\normalsize| 
\hspace{0.2cm}\textit{mg}
\footnotesize At boxesia et uoluntaleDet
\normalsize| 
\hspace{0.2cm}\textit{mg}
\footnotesize Deledu dlendaerum uiduarum. 
\normalsize| 
\hspace{0.2cm}\textit{mg}
\footnotesize pe domuas ministrare Deo grais. num. 
\normalsize| 
\hspace{0.2cm}\textit{mg}
\footnotesize Vxe niue desempuo a signis. 
\normalsize| 
\hspace{0.2cm}\textit{mg}
\footnotesize Difermem et signauerae ac Else uiduae, 
\normalsize| 
\hspace{0.2cm}\textit{mg}
\footnotesize A distributione tenpo rs, 
\normalsize| 
\hspace{0.2cm}\textit{mg}
\footnotesize 
\normalsize| 
\section*{IN I. EPIST. D. PAVLI. }
\marginpar{[ p.106 ]}\pstart ftuδιεαuς; id est: omni genere pietatis, quod ei per aetatem et facultatem licel praestare, omni tempore intenta est. Est ergo egregia descriptio uiduae Christianae, per quam intelligi potest, eam, quae hoc pacto id uitae genus traducit, meriton omnibus modis honorandam. Sunt autem in scripturis proposita aliquot huiusmodi viduarum exempla et viuae imagines, ex quibus cognosci potest, quantum decus sit earum, quae vere viduae sunt. Iudith.8. Erat ludith eleganti aspectu nimis, cus vir suus reliquerat diuitias multas et familiam copiosam, ac possessiones armentis boum et gregsbus ouium plenas: sed nihilominus praeterquam quod illa eleemosynas largiter conferebat, tamen de ea scribitur, quod in superioribus domus suae fecit sis bi secretum cubiculum, in quo cum puellis suis clausa morabatur: et habens super lum bos suos cilicium, ieiunabat omnibus diebus vitae suae, praeter Sabbatha et Neomenias, et festa domus lsrael.ln tota historia eius mirabslis in Deum pietas, et in patriam mirus amor ostenditur. ludic. 4.celebratur vidua relicta Lapldóth Deborah, spiritu pphetiae clara, cuius prudentia lsrael victoriam contra Sisaram ducem labin regis Chanaan obtinuit. Ruth. 1. Noemi et uiri et filiorum solatio destituta, quae nurum suam Ruth Moabitidem, vt patrios Deos abiuraret, illexit, et natum ex Ruth nepotulum Obed instar nutricis ac gerulae pie sancteque educaust.1.Reg.17. Vidua Sarepthana hospitalitatem exercet, Heliam Prophetam in summa omnium rerum penuria pascere non dubitauit. Lucae 2.Hamna vidua filla Phanuelis annorum fere 84. ieiuniis ac deprecationibus die ac nocte in templo operam dabat, ac tandem Dominum vidit et confessa est, locutaque est de illo omnibus, qui expectabant redemptionem Hierosolymis. Fuerunt haud dubie et senctae viduae inter illas, de quibus Macth.27. dicitur, quod ministrabant Domino et discipulis eius de facultatibus suis. Lucae 21.commendat Dominus pauperculam viduam, quae minuta duo aera misit in gazophylacium, cuius affectus pluris factus est a Deo, quam dona omnium diuscum multo amplissima. In summa, sumul ex hisce smaginibus, simul ex descriptione Pauls, cum qua haec exempla relata per omnia conueniunt, manifeste licet colligere, quaenam viduae ceu vere viduae sint honorandae: et quam sit illa vitae conditio laude digna, vbi Christianae pietatis simulachra adiuncta habent, et talia praestant, qualia sanctas illas viduas praestitisse scriptura commemorat. Acnullius tam exigua facultas, nullius tam inidonea ad benem agendum aetas, quin vel orando, vel eleemosynas, nempe aera minuta, aut calicem aquae frigidae (quem Marcus capite 9. magni facit) conferendo, vel mores iuuentutis formando uel alla quaecunque charitatis officia queat exercere. Atque his ratsonibus se vere viduas potuerunt declarare. Sed adhuc clarius euadit, qualem esse vere viduam deceat, si sequentem fucarae viduae descriptionem intueamur.  \pend\pstart Porro que in deliciis versatur, ea viuens mortua est.) Vidua, cui adsunt non solum viltae neceslaria, verum ea quibus abuntatur ad delscias, talis licet marito sit destituta, ne quaquam tamen in illarum viduarum numerum est recensenda, quae vere uiduae haberi debent. Viduae verae agunt desolate: viduae fucatae agunt in deliciis, neminem reuerentur, nulli parent, rem domesticam tractant negligenter, ad libertatem ac licentiam sunt intentae, otio deditae curam quarumcunque rerum a se abiiciunt. Porro quae huiusmodi est certe nequaquam venit honoranda.  \pend\pstart Ea viuens mortua est.) Per αὐτιμειαθολιι seu commutationem uitam elus lignificatur nihilo praestantiorem morte. Neque enim uere uiuit, quae carni duntaxat uiuit, non Deo. Saepe etiam mors in scripturis non solum carnis obitum, sed animae interitum spiritualem damnationemque significat. Vnde mors secunda dicta est mors spiritualis siue mors animae in Apocalypsi. Et mortui dicuntur, qui mundo quidem uiuunt, caeterum quia non habent cognitionem Dei, et Deum non curant, mortuis nihilo sunt coram Deo praestantiores. Vnde Matth.octauo dicebat Christus adolescenti, sine ut mortui sepeliant mortuos suos. Et usitato more loquendi etiam uulgo dicuntur uiua cadauera et uiua sepulchra, qui sie uiuunt, ut nihil agant, quo se uiuere testentur, uel a mortuis parum dis  \pend
\vspace{0.5cm}\noindent
\vspace{0.2cm}\rule{1cm}{0.2pt}\\ 
\hspace{0.2cm}\textit{mg}
\footnotesize Enumeratid exemplorum, 2. Iudith. 
\normalsize| 
\hspace{0.2cm}\textit{mg}
\footnotesize Epetori 2. Noemi
\normalsize| 
\hspace{0.2cm}\textit{mg}
\footnotesize 4. Viduas rephibana 5.Hamna. 
\normalsize| 
\hspace{0.2cm}\textit{mg}
\footnotesize 6. Maliere Chrislo mint strantes. 7. Vidua pd percula. 
\normalsize| 
\hspace{0.2cm}\textit{mg}
\footnotesize Pdfe ueix descriptio. 
\normalsize| 
\hspace{0.2cm}\textit{mg}
\footnotesize vnuru. 
\normalsize| 
\hspace{0.2cm}\textit{mg}
\footnotesize amurtas 123. 
\normalsize| 
\hspace{0.2cm}\textit{mg}
\footnotesize Hors gininteritumsg ritualem saepe significat Apoc. 2e 
\normalsize| 
\section*{AD TIMOTHEVM CAP. V. }
\marginpar{[ p.107 ]}\pstart stant. Itaque Apostolus viduam viuentem in deliciis comparat mortuae, nihiloque meliorem ducit, quam si mortua esset: quippe inutilia, et ad mortem animae ten dentia sunt omnia, quae ipsa facit: ac proinde nequaquam instar eius, quae vere vidua est, honoranda.  \pend\pstart Et hec praecipe, vt irreprehensibiles sint.) Proplus volens ostendere, quale viduam esse deceat, exigit aciubet eas prorsus irreprehensibiles, adeoq omni ex parte inculpabiles esse, ne vel aduersarii earum vitam possint calumniari. Atque ut ita se gerant, vult Apostolus per episcopum id praecipi ac iuberi, ταῦτα παράγγεπε.  \pend\pstart Quod si quae suis et maxime familiaribus non prouidet, fidem abnegauit, et estinfideli deterior. Ad at vnam hanc rationem quae iustissime ab omnbus viduis exigitur, quamque nu lo pacto a se reiicere possunt:nempe vt curam gerant suorum et maxime familiarium, Id autem natura ab omnibus exigit, et nemo non intelligit iustissimum esse vt fiat. Quae autem id nolit praestare, ad quod natura ipsam compellit, certe talis merito viuens mortua dicitur, quippe quae Christum abnegat, et infideli quouis deteriorem sese exhibet: cum fides non modo exigat ea, quae sunt naturae legibus praecepta: verùn ettam requirat altiora quaedam ex Euangelica charitate prodeuntia. Et in infideli nor extinguuntur ςοργαὶ φυσιααι, quae et praeceptae sunt a Deo:quare quae illas abiicit, re cte infideli deterior iudicatur. Nomine autem fidei, dum dicitur abnegare fidem, pe Synecdochen intelligitur quicquid ad fidem pertinet, nec non virtutes illae, quae ex si de oriuntur, per quas solas obtinere possunt, vt vere viduae sint et dicantur. Proinde hinc colligitur, grauissimum esse peccatum, negligere curam suae familiae, etiamsi prae texantur honestae et specietenus piae caussae:cuiusmodi sunt, quod quidam aiunt, se religionis ergo Romam adire aut Hierosolymam, vel monachos uelle fieri. Et si peccatum est suos negligere, quid dicemus de iis, qui suis etiam iniuriam inferunt? Prouidenda autem familiaribus sunt, non tantum ea, quae ulctui sunt necessaria, quaeque externae uitae conducunt, sed multo magis, quae ad ueram attinent pietatem, qua ratione ante dictum est, uiduam debere domum ὁυσεBαμ pie tractare, id est,in iis, quae pertinent ad religionem et ueram pietatem, debere instituere. Et si haec cura hoc loco exigitur etiam a uiduis quantopere incumbit eadem cura ipsis patribus familiansPostremo euidens argumentum hic habetur, inter primas dotes probae uiduae habendum, si ipsa strenuë prouideat suis siue consanguineis, siue aliis, maxime fa milsaribus et eius opera iure adiuuandis. Atque haec quidem de ea, quae inter probas uiduas sit habenda: Nunc qualis digna iudicetur, quae ecclesiae facultatibus alatur:  \pend
\phantomsection
\addcontentsline{toc}{subsection}{\textit{Uidua allegatur non minor annis sexaginta, quae fuerit vnius viri vxor, in operibus bonis hominum testimonio comprobata, si filios educauit, si fuit hospitalu, si sanctorum pedes lauit, si afflictus subministrauit, si in omni opere bono fuit assidua. }}
\subsection*{\textit{Uidua allegatur non minor annis sexaginta, quae fuerit vnius viri vxor, in operibus bonis hominum testimonio comprobata, si filios educauit, si fuit hospitalu, si sanctorum pedes lauit, si afflictus subministrauit, si in omni opere bono fuit assidua. }}\pstart De iis, quae uel ipsae adhucpraesunt familiis, uel habent liberos aut nepotes familiam instituentes, dictum est, satisque indicatum, nequaquam opus esse, ut tales subleuentur ecclesiae facultatibus. In ea uerô, quae recensebitur in numero alendarum ab ecclesia. circumstantiae diuersae sunt considerandae. Primo debet esse uere uidua, uti paulo ante descripta est: nempe desolata, id est, destituta omni auxilio, sperans in Deo, pie perseuerans in obsecrationibus ac precationibus nocte dieque. Sed praeterea, ut paucis omnia concludantur, duo adhuc sunt animaduertenda, nempe aetas et ui ta anteacta.  \pend\pstart Vidua, inquit, allegatur.) Graece est καιαλεγεθα.ααιαλέγω significat deligo, conscribo, centurio, sicut milites solent conscribi et in album referri.  \pend\pstart Ron minor annis sexaginta.) En aetas uere commiseratione et ope digna:nam iam effoeta emenita, non potens ample obire aliquos labores parando uictui tùm de qua nequaque   \pend
\vspace{0.5cm}\noindent
\vspace{0.2cm}\rule{1cm}{0.2pt}\\ 
\hspace{0.2cm}\textit{mg}
\footnotesize Qudis ugra uidua esse debeat. 
\normalsize| 
\hspace{0.2cm}\textit{mg}
\footnotesize Signum quo cognoscitur uiuens mortua 
\normalsize| 
\hspace{0.2cm}\textit{mg}
\footnotesize Infidelis non exuit eneσοργίαν. Iista aĩtean exuit. Igitur est infideli detcrior. Panüie curam depone re fidei et pietai est contrarium. Aminori ad maius. Dosuidue probae est pia et diligens prouidentiain re fimiliari. 
\normalsize| 
\hspace{0.2cm}\textit{mg}
\footnotesize Confideratio circĩstan tiarum. 
\normalsize| 
\hspace{0.2cm}\textit{mg}
\footnotesize Aetauiduuerae, quam ali ecclesiae fa cultatibus con uenit. 
\normalsize| 
\section*{IN I. EPIST. D. PAVLI. }
\marginpar{[ p.108 ]}\pstart metuendum, quod vel impudicitie, vel cuiuscunque improbitatis notam incurrat. Quo argumento et declaratur mirabilem esse humanam fragilitatem, cui omnis propemo dum aetas suspecta esse debet: et omni studio cauendum, ne vsquam impudicitiae notam incurramus. De vita anteacta plurima sunt inquirenda: ldque ob has caussas. Primô enim si statim eae quoque admitterentur, quae antea vixissent improbe, id scandalo esset piis et benemeritis, qui existimarent virtuti non suum esse honorem, sed aeque bene fie riindignis atque promeritis. Quinetiam redderentur aliquo modo securi, qui vitiis assueuerint, ac sic quasi occasionem haberent pergendi in veteri sua praua consuetudine, dum spes esset se tandem, vbi egerent, ecclesiae facultatibus benigne alendos. Tertio, esset id ingents offendiculo aduersariis religionis, si viderent qui fuerint per omnem vitam apud nos infames, vt scurras, mimos, lenones, decoctores, aleatores, lenas, ambubaias, ali: vt non dicam, quod pii diuites animum a dando et conferendo in gazophylacium ecclesiae remitterent, iudicarentque parum recte collocatas suas eleemosynas, si id genus colluuiem quandam ea eleemosyna pasci cernerent. Et certe quae semel lenas siue ambubaias egerint, facilime recidunt, quamuis vetulae, ad designandum quod viduarum professionem dedecet. Nemo credat quid mali possint anus comminisci, quae antea malis assueuerint. Postremo si quaeuis aniculae, cuius cuius conditionis antea fuerint, suscipientur alendae ex aerario ecclesiae, metuendum ne istud non sit suffecturum. Nequaquam ergo mirum videri debet, quond Paulus, vt erat mirem prudens, tam scrupulose hanc causam viduarum tractet. Haec porro sunt, quae de vita anteacta vult inquiriac considerari.  \pend\pstart Que fuerit vnius viri vxor. Id est, quae non visa sit veluti carnis cupiditate aestuans mortuo priore marito, nouas nuptias appetere, sed aequanimiter suam tulit sortem, non gestiens placere mundo, fallacibus et ineptis curis, non solicita vt viris vllis magis pla ceret, quam vni Christo. Non solent incontinentiae notam effugere mulieres, quae nuptias cito repetunt. Vnde et in iure infamiae labe asperguntur, quae ante annum nubunt: Et ludith.i5.in ipsa laudatur, quod post virum suum alterum nescierit. Tametsi nemo hinc putet ita commendari monogamiam, vt cum Montanis damnentur secundae nuptiae.  \pend\pstart In operihus bonis hominum testimonio comprobata.) Requirit de eius pietate in vita anteacta testimonium aliorum, et commune bonorum de ea iudicium, nam consideran dum etiam quales sint qui testimonium ferunt. Recenset et diuersas species bonorum operum, quae in vita anteacta viduarum solent animaduerti.  \pend\pstart Si filios educauit.) Non tantum nutriendo corpora, sed etiam pascendo animas, sana videlicet doctrina, quemadmodum decet familiam δυσεssᾶν.  \pend\pstart Si fuit hospitalis.) Prius quidem bene mereri debent de suis, nempe filiis et domesticis, postea etiam de aliis. Insignis hospitalitas resplenduit in vidua Sarepthana.  \pend\pstart Piliorum eduSt sanctorum pedes lauit.) Sancti sunt quicumque pii vere credentes et timentes Deut. Pe dumlotione per σμυεκδιογιων intelligunt̃ et minima et maxima quaeque pietatis humanita tisque officia in sanctos, praesertim peregre venientes, et de via fessos exhiberi solita. Pe des autem lauare aduentanti hospiti, insignis est humanitatis pariterque valde officiosi pectoris, vt hoc exemplo significetur uidua, quae plane nihil non fuerit prompta in gratiam sanctorum facere. Erat autem ueteribus mos hospitum aduentantium pedes  statim lauare. Vnde Genes.18. Abraham Angelis illis, quos hospitio excepit, afferam, inquit, pauxillum aquae et lauentur pedes uestri. Et cap.19. Loth. ad Angelos similirer: obsecro domini, declinate in domum pueri uestri, et manete ibi: Lauate pedes uestros, et mane proficiscemini. Et Christus Lucae7. Pharisaeo ait, aquam pedibus meis non de disti, haec autem lachrymis rigauit pedes meos, et capillis capitis sui extersit.  \pend\pstart Si afflictis subministrauit. )Id est, si aegrotis, uiduis, captiuis, aduersaque fortuna ad inopiam et miseriam redactis opem tulerit:ltem si afflictos, moestos, solata fuerit, consilii inopes consiliis sanis adiuuerit, denique aliis modis, uerbis, factis, affectu, charitatis argumenta praebuerit.  \pend\pstart Si in omni opere bono fuit affidua.) Conclusio paucis uerbis complectens omnia. Si,  \pend
\vspace{0.5cm}\noindent
\vspace{0.2cm}\rule{1cm}{0.2pt}\\ 
\hspace{0.2cm}\textit{mg}
\footnotesize 1.Causa Apersona piorum pa perum. 2.causa Apersonmalorue pa perum. 3.causa A persona aducrsariorum. 4.cause A persona diuitum. 
\normalsize| 
\hspace{0.2cm}\textit{mg}
\footnotesize 5.causa Aergis ehaustio. 
\normalsize| 
\hspace{0.2cm}\textit{mg}
\footnotesize Libido et is continentia uiduas dede cet. 
\normalsize| 
\hspace{0.2cm}\textit{mg}
\footnotesize Tepus seck darte nupte rum iuri conuenieus. 
\normalsize| 
\hspace{0.2cm}\textit{mg}
\footnotesize Montanus se cundas dam nat nuptias. Testimonuum pietatis et Iudicium be norum. 
\normalsize| 
\hspace{0.2cm}\textit{mg}
\footnotesize Opera bona uiduarum. 
\normalsize| 
\hspace{0.2cm}\textit{mg}
\footnotesize catio. Hospitalita Lotio pedi 
\normalsize| 
\hspace{0.2cm}\textit{mg}
\footnotesize Mosuetus 
\normalsize| 
\hspace{0.2cm}\textit{mg}
\footnotesize Affictort ie creatio. 
\normalsize| 
\hspace{0.2cm}\textit{mg}
\footnotesize Exgctαe in omni opl rebono. 
\normalsize| 
\section*{AD TIMOTHEVM CAP. V. }
\marginpar{[ p.109 ]}\pstart Inquit, fuit assidua, id est, semper et magno animo occupata, in omni opere bono (Em phatice) ita vt nullum prorsus opus sit verae pietatis, in quo illa se non gnauiter exertuerit. Plurima itaque et vere magna raraque in ea requirit vidua, quae digna iudicabitur, vt ecclesiae facultatibus alatur: Caussas iam ante recensuit: fortassis ideo quoque Apostolus ita delectum haberi voluit, quod perquam tenue tùm temporis esset ecclesiarum aerarium, nec sufficiens alendis multis. Certe nostris temporibus, quibus raro et vix vsquam tam integras liceat inuenire viduas, quando etiam auctiores aliquam to sunt multis in locis ecclesiarum facultates, liceret ad quaedam conniuere, nempe etiam sl aetas illa ( quam paucae nunc attingunt) non adsit: etiamsi aliquando secundum maritum experta, dummodo tamen vere desolata sit et bonorum testimonio vi ta eius anteacta probata fuerit: videtur non excludenda a communibus ecclesiae beneficiis. Sed quid ista dico? Vtinam facultates amplissimae ecclesiarum non prodigerentur, per inutiles, otiosos, scortatores, adulteros, praebendarios, qui opes ecclesiarum effundunt alendis scortis, leonibus, venatoribus, aucupibus, mulionibus, etc.  \pend
\phantomsection
\addcontentsline{toc}{subsection}{\textit{Porro iuniores viduas reiice: cum enim lasciuire coeperint aduersus Christum, nubere volunt, habentes condemnationem, quod primam fidem reiecerint: simul autem et otiose d'scunt circumire domos, imo non solum otiosae, verum etiam garrule et curiosae, loquentes quae non oportet. }}
\subsection*{\textit{Porro iuniores viduas reiice: cum enim lasciuire coeperint aduersus Christum, nubere volunt, habentes condemnationem, quod primam fidem reiecerint: simul autem et otiose d'scunt circumire domos, imo non solum otiosae, verum etiam garrule et curiosae, loquentes quae non oportet. }}\pstart Per Antithesin ostendit iunsores seu minores sexaginta annis non recipiendas in carum numerum, quae ecclesiae facultaribus sunt alendae. Addit et caussas varias, a vitiis sumptas, quibus minor aetas plerunque est obnoxia.  \pend\pstart Cum enim lasciuire coeperint aduersus Christum, nubere volunt.) Videtur ab experientia argumentari, quasi iam ita compertum fuisset in iunioribus viduis, quod vbi semel foi tassis quaedam participes fuerunt publicae eleemosynae, et factae securiores, coeperunt ecclesiae bonitate abuti, et cum antea simulassent se deinceps nunquam nupturas, breui terpres reddidit, Cum lasciuire coeperunt, Graece est καταςρlυηάσωσ, a verbo ςριυνάω, intemperanter et immodeste ago, infrenis sum, quast soluto vel rupto freno ago, vt metaphora subesse videatur a iumentis, quae, cùm bene pasta commode aluntur, ferocire incipiunt, et lora omnia frangunt effugiuntque ́. Huiusmodi autem viduarum intemperantiam vocat lasciuiam aduersus Christum, quia nimirum cum antea visae sunt se totas voluisse dedere vni Christo, quasi de maritis postea nunquam cogitaturae, atque  eo praetextu receperunt ecclesiae eleemosynas:recte, quando iam iterum captant viros, et vel fornicantur, vel nuptias nouas meditantur, dicuntur aduersus Christum lasciuire. Et malihomines tùm interpretarentur, ab ecclesia quodammodo alimenta praebita lasciuiae: vnde nomen et secta Christianorum protinus male audiret: neque sine nominie. Christiani probro tales vidue conspiciunt lasciuire. Nubere quidem per se caret culpa, sed id non caret culpa, quond, quae antea pollicitae sunt se velle deinceps per omnem vitam coelibes agere, atque eo praetextu obtinuerunt publicam eleemosynam, iam fallunt palam vniuersam ecclesiam, et religionem ea leuitate non parum infamant. Vnde sequstur,  \pend\pstart Habentes condemnationem, quod primam fidem reiecerint.) Turpi illa mutatione prioris propositi praebent occasionem, vtab omnibus damnentur reprobenturque ceu leues, vane, perfidae, et libidinosissimae, quod fidem siue promissionem illam, qua prius pollicebantur, se vitam uelle traducere in coelibatu, irritam fecerint. Nam eae duntaxat ad mittebantur ad participationem communis eleemosynae in ecclesia, quae sese affirmabant ac promittebant caste in viduitate permansuras. Nam duntaxat in quantum uiduae erant et uiduae permanerent, iudicabantur subleuande ecclesiae facultatibus. Quae m' trimonium praeeligebant, libere quoque id faciebant, acrelinquebantur alendae a marl tis. Illa itaque promissio de casta uidutate nequaque talis uoti speciem habebat, qualit  \pend
\vspace{0.5cm}\noindent
\vspace{0.2cm}\rule{1cm}{0.2pt}\\ 
\hspace{0.2cm}\textit{mg}
\footnotesize Ab expguena. 
\normalsize| 
\hspace{0.2cm}\textit{mg}
\footnotesize Ετρlυαν Metaphora 1iumentis. 
\normalsize| 
\hspace{0.2cm}\textit{mg}
\footnotesize OfErdE gριαnο͂ sipu ficat dam  nationem anine sed dan n. tionem per Iniquum boni num iuatcium si ue damnatio nen ac notam infimiae Promissa et conditiones ut duarum receptarum ad elemosynae communis partricipationem. 
\normalsize| 
\section*{IN I. EPIST. D. PAVLI. }
\marginpar{[ p.110 ]}\pstart vota perpetuon obligantia conscientiam quidam confinxerunt: quippe quaedam post illa promissionem siue votum, si ita velint, nupserunt, vt hic Apostolus commemorat. De inde si maxime vellent affirmare vota per virgines omittenda:at videant, qualem velit esse Apostolus, quae castitatem uoueat, nempe sexagenariam. Quanto prudentius ca uit Apostolus 1.Corinth.7.ubi de uirginitate ac coelibatu disserens sic uoluit semper accipi sua uerba in commendationem uirginitatis dicta, ut non iniiceretur laqueus conscientiis. Sed aliud uitium accusat Apostolus in uiduis iunioribus, ob quod minus alendae essent ecclessae facultatibus.  \pend\pstart Simul autem et otiosa discunt circumire domos.) Ideo illis non dandum, quia ex liberalitate ecelesiae captant occasionem turpis otii, quod mater est omnium uitiorum. Ecclaesiast. 33. Multam malitiam docuit otiositas. Prouerbus 12. Qui sectatur otium, stultissimums est. Porro habentes otium discunt circumire domos. Cum enim speciosus esset et ueneram dus uiduitatis titulus, hoc nomine factae arrogantes et ambitiosae se in diuitum domos inuerecunde ingerunt, et, quod fere solet huic rei adiunctum esse, impudenter a quouis petunt, et a quouis accipiunt. Hoc enim significabat Christus, cùm obiiceret Pharisaeis Matth.23.qu comederent domos uiduarum, qui sine dubio circumibant domos illarum, et sem per petebant: qualem et monachum quendam describit Hieronymus in quadam Epistola ad Eustochium de uirginitate seruanda. Prouerbus 7. describitur meretrix circumcursitans, nunc foris, nunc in plateis, nunc iuxta angulos insidians. lerem.14. Dilexit mouere pedes suos et non quieuit, et Domino non placuit. Sed quoniam eiusmodi uiduae oben rantes per diuitum aedes solent nihil omittere, quo sperant se diuitibus placituras: hinc quoqu quaeuis meditantur, confingunt, commentantur, quae grata illis fore arbitrantur. Vnde sequitur:  \pend\pstart Imo non solùm otiosae, verum etiam garrulae et curiosae, loquentes quae non oportet.) Exagge ratio est perCorrectionem, qua ostendit, accedere turpi otio et cursitationibus illis garrulitatem et curiositatem, denique et obscoenitatem. Prouerbus 7. Garrula, uaga, quietis impatiens. Loquunt, inquit, quae non oportet, id est, quae non solum non decent uiduas, sed nec quascunque  mulieres: imô loquunt quae nocent, dum uaria deferendo et puarietate ingeniorum in aedibus  diuitum quaeuis contingendo, odia, iniqua iudicia, obloqa serunt, quae nonnunque in apertas rixas et lites abeunt. Quia itaque hoc pacto sese gerunt, certe nemo cordatus et uere uiduas honorare uolens, tales iudicet ex ecclesiae facultatibus alendas. Ac grauissime his uerbis feriuntur, quaecunque hoc nostro seculo praetextu religionis, et castae uiduitatis cir cumcursitant, et quaeuis a diuitibus extorquent, interea otiosae, garrulae, curiosae. In Grae co est allusio et αὐταυέκλασις in uerbis illis, ουμένου ἀργαὶ, αλλα κααι πόρίσργοι. De his uide commentarios nostros in 2. Epistola ad Thessalonicenses, caput tertium.  \pend
\phantomsection
\addcontentsline{toc}{subsection}{\textit{Uolo igitur iuniores nubere, liberos gignere domum administrare, nullam occasionem da re aduersario, vt habeat maledicendi caussa. la n.nomnullae deflexerunt secutae sathanam. }}
\subsection*{\textit{Uolo igitur iuniores nubere, liberos gignere domum administrare, nullam occasionem da re aduersario, vt habeat maledicendi caussa. la n.nomnullae deflexerunt secutae sathanam. }}\pstart Per occupationem adiecit, quid iunioribus faciundum statuat. Nam ubi quis dixisset forte, Non uis, ô Paule, iuniores uiduas sustentari ecclesiae facultatibus, propter cer ta uitia quibus plerunque laborant, quid igitur restat illis faciundum? Respondet, nubant et mariti eas alant, et in eo uitae genere se pie exerceant, liberos alendo, familiam admi nistrando, non discurrendo otiose: in quibus rebus honestas et probitas sita est non poeVolo, inquit, si viduae q̃ se continere non possunt, et cauere a scandalis, denuô nubant, po tius quam si promittant se, spe accipiendi uictui neccessaria ab ecclesia, uicturas coelibes. Atque haec est honestissima ratio iudicanda, qua iuniores uiduae possint sibi prospicere, si uidelicet denuo nubant, potius quam vel extrema laborent inopia, uel alias se parum honeste gerant. sequitur ratio, a uitando scandalo.  \pend\pstart Nullam occasionem darc aduersario, ut habeat maledicendi caussam.) Si uidua uult perma  onem dare maleuolis hominibus obloquendi. Quod si nubit, ipsa omnem occasionem adimit. Nam que nubit, id nemo iure potest reprehendere. Porro uariis modis uidua iuue  \pend
\vspace{0.5cm}\noindent
\vspace{0.2cm}\rule{1cm}{0.2pt}\\ 
\hspace{0.2cm}\textit{mg}
\footnotesize Oumtuu nis mater. 
\normalsize| 
\hspace{0.2cm}\textit{mg}
\footnotesize Circuire a nos quid. 
\normalsize| 
\hspace{0.2cm}\textit{mg}
\footnotesize Exaggeratio per oce pationem. 
\normalsize| 
\hspace{0.2cm}\textit{mg}
\footnotesize Mores fall rumuidus 
\normalsize| 
\hspace{0.2cm}\textit{mg}
\footnotesize Amanea cis. 
\normalsize| 
\hspace{0.2cm}\textit{mg}
\footnotesize Nuptiae permisse et pre positɇ umic nbusuduis. 
\normalsize| 
\hspace{0.2cm}\textit{mg}
\footnotesize Occasio cα uonntɇ suxda. 
\normalsize| 
\section*{AD TIMOTHEVM CAP. V. }
\marginpar{[ p.111 ]}\pstart nis praebet occasionem maledicendi aduersariis: si se comit paulo elegantius, si cum ullis habet consuetudinem, si a quibusdam salutetur resalutetque ;: in summa, sexcentae circumstantiae sunt, quae statim eam ceu impundicam suspectam reddunt: et dicta et facta, et gestus et omnia ipsius a malis hominibus obseruantur. Omnia ergo haec his verbis interdicuntur et prohibentur ei, quae pro vera uidua haberi velit.  \pend\pstart Iam enim nonnulle deflexerunt secutae sathanam.) Comprobatio rationis praecedentis ab exemplo siue ab experientia: quasi dicat: Nemo miretur, me sic iubere, vt viduae nu-4 bendo sibi consulant, et aduersario occasionem obloquendi praeripiant. Nam iam amodo compertum est, quasdam quae coelibes se fore promiserant, et ea conditione ecclesiae liberalitate enutritae sunt magnorum scandalorum caussam extitisse, dum cursitanses, garrulae, curiosae, tandem et in fornicationem prolapsae, multorum animos conmouerunt ad sinistre iudicandum de vniuersa nostra religione, et simul labefactarunt pulcherrimum hoc alendarum vere pauperum viduarum institutum.  \pend\pstart Secutae sunt sathanam.) Id est, opera sathanae, per Metonymiam. Sunt autem opera sa thanae, otiari, vagari, deferre, obloqui, calumniari, lasciuire, et scandalorum caussam prae bere. Ergo iunioribus viduis satius est vt nubant, quam vt in viduitate committentes opera sathanae occasionem praebeant malis sinsstre iudicandi. Expendant istud iudicium Pauli, qui immodice efferunt coelibatum, quique damnant secundas nuptias.  \pend
\phantomsection
\addcontentsline{toc}{subsection}{\textit{Quod si quis fidelis aut si qua fidelis habet viduas, suppeditet illu, et non oneretur ecclesia, vt his, quae vere viduae sunt, suppetat. }}
\subsection*{\textit{Quod si quis fidelis aut si qua fidelis habet viduas, suppeditet illu, et non oneretur ecclesia, vt his, quae vere viduae sunt, suppetat. }}\pstart Conclusio huius causae de alendis viduis per occupatione Nam si qui forte iterum o biecissent ad verba praecedentia, viduas quasdam tam miseras esse, vel aliqua ratione tam infoelices, vt non videantur inuenturae maritu. Respondet, ac definit in summa, eas quae amicos habent et consanguineos, nequaquam ecclesiae obtrudendas. Perspicua sunt verba Pauli: Dicit autem: Si quis fidelis aut si qua fidelis. Nam credentibus vult istud persuasum esse debere, ac misericordia mouert hos vult, vt viduam sibi sangumne vel amicitia iunctam foueant Incredulis seu gentilibus praescribere non putat ad se pertinere. Et vere declarat se fidelem, quicunque huius rei curam ad se benigne recipit. Neq hic omisit Apostolus certam rationem, quando per ἀτιολογίαν ait,  \pend\pstart Et non oneretur ecclesia, vt his, quae vere viduae sunt suppetat.) A difficili argumentatur. et ab inconuenienti quod sequeretur: quasi dicat: si quaeuis reciperentur, certe ecclesla oneraretur, atque exhaustis eius facultatibus eae, quae vere viduae sunt, quae neque ae-e tate, neque amicis, neque vlla ratione semetipsas possunt iuuare, destituerentur, quod non n recte esset factum. Aequius est ergo, his semper prospici, quam illis, quae alio pacto ii vel denuo nubendo, vel per amicos iuuart possunt. Vtinam autem haecratio alendarum o viduarum in primitiua ecclesia obseruata, nunc quoque in vsum reuocaretur: imo uti-e nam omnium pauperum ratio haberetur. Sane ad instituendam subuentionem paupe rum, quaedam argum enta et ordinationes commodae, ex his uerbis Pauli excerpi pos sent, nempe qualibus potissimùm conferenda eleemosyna. Studeant ministri ecclesiarum hanc et huiusmodi alias bonas ordinationes suis petsuadere, et opera misericordiae diuitibus iugiter commendare. O nunquam satis laudatos meos Hyperios, qui i non solum uiduis, sed omnibus pauperibus cuiuscunque sortis uel aetatis, nec domesti-, cis suis tantum, uerum etiam peregrinis aduentantibus consultum uoluerunt.  \pend
\phantomsection
\addcontentsline{toc}{subsection}{\textit{Qui bene praesunt presbyteri, duplici honore digni habeantur, maxime ii, qui lab orant in sermone et doctrina. Dicit enim sertptura Boui trituranti non obligabis os. Et, dignus est operarius mercede sua. }}
\subsection*{\textit{Qui bene praesunt presbyteri, duplici honore digni habeantur, maxime ii, qui lab orant in sermone et doctrina. Dicit enim sertptura Boui trituranti non obligabis os. Et, dignus est operarius mercede sua. }}\pstart Sequitur, quomo do episcopus se geret erga presbyteros. Curari autemuult per eu ut his honeste prospiciatur, deque eare diligenter doceri populum. Rectissime autem postquam egit de alendis uiduis ex ecclesiae facultatibus, agit causam et ministrorum ecclesiae, quos oportet et aequissimum est, simili modo, quoniam illis non uacat propter publicum ministerium rebus terrenis et negotiationibus ad quaestum facientibus dare operam, et ecclesiae facultatbus sustentart.  \pend
\vspace{0.5cm}\noindent
\vspace{0.2cm}\rule{1cm}{0.2pt}\\ 
\hspace{0.2cm}\textit{mg}
\footnotesize Ab exemplo 
\normalsize| 
\hspace{0.2cm}\textit{mg}
\footnotesize Metonymic 
\normalsize| 
\hspace{0.2cm}\textit{mg}
\footnotesize Conclusio 
\normalsize| 
\hspace{0.2cm}\textit{mg}
\footnotesize Occupatio. 
\normalsize| 
\hspace{0.2cm}\textit{mg}
\footnotesize Fideles curam gerere debent pauperum uiduarum. 
\normalsize| 
\hspace{0.2cm}\textit{mg}
\footnotesize A difficilem et inconuenientiEcclesia non oneranda aut exhaurienda. 
\normalsize| 
\hspace{0.2cm}\textit{mg}
\footnotesize Laus Hyperiorum. 
\normalsize| 
\section*{IN I. EPIST. D. PAVLI. }
\marginpar{[ p.112 ]}\pstart Qui bene praesunt presbyteri.) Hoc loco presbyteri vocabulum officii, non ae tatis iim cludit rationem. Vocantur autem presbyteri generali et communi nomine omnes, qui ministrant in ecclesia:id quod antea quoq admonuimus. Nam et episc. ipse presbyrer quoque dicitur, testante Ambros at non omnis presbyter episc. De his ergo omnibus hic ei sermo est, qui episcopo inferiores sunt, quos Ambrosius enumerat in qualibet ciuitate, duos diaconos, qui administrant verbam et Sacramenta, et septem diaconos, qui ministrant mensis, distribuentes ecclesiae facultates. Animaduertendum vero, Paulum, cum agere institueret causam ministrorum, preposuisse id primo loco, Qui bene praesunt: Nam mul ti quidem ambiunt appellari presbyteri, et hocnomine facultatib, ecclesiarum ali: Sed vel paucissimi sunt, qui bene praesunt. Apostolus autem eos commendat qui bene praesunt. Sa ne otiosi et inutiles, qui Epicuri de grege porci magis videntur, quam pastores animarum, non videntur his presbyteris, quos Paulus commendat, annumerandi. His itaque verbis admonentur omnes presbyteri, vt bene praesint, hoc est, ita praesint, vt prosint, secundum voluntatem spiritus Des omnia dispensantes. Bene aut praeerunt, si officium suum obeant vri Apostolus hac Epistola praescripsit, in qua quales eos esse, tùm domi tùm foris quomodo item et quae docere debeant, quomodo ministrare, plenissime demonstratur. Proinde istiusmodi presbyteri duplici honore digni habeantur. Nomine honoris per Metonymiam intelligi conuenit, quicquid eis ratione officii inferiores debent. Quod vt clarius intelligeret, vocauit honorem duplicem. Duplicem autem honorem interpretantur: Alterum videlicet honorem obedientiae, quo honoratur in illis verbi Dei autoritas, cuius administratio semper salutaris et necessaria, tanto pluris debet fieri prae qua cunque functione, quanto sublimior est anima quam corpus. Alter honor est subsidii et iusti itipendii, quod conuenit honeste alendae familiae. Etiamsi Theophylactus duplicem honorem exponat, abundans, amplum, et duplicatum stipendium, quod sufficiat et familiis sustentandis ministrorum, et aliis bonis operibus in pauperes exercendis, si cut olim Leuitis decimae conferebantur. Porro de priore illo honore non admodum laborat Apostolus, neque res magna est, si impendatur, neque ceu grauem possunt vlli detrectare. Nam et magistratus ciuiles honorandos esse ratione officii publici, nemo non concedit. Et medicum honorandum praeterque qu scriptura precepit, quae ait: Honora medicum, propter necessitatem enim creauit eum Dominus: tùm Homerus quoque dixit: labos γαρ αυηρ πολαν αὐτάβιω ἀλλων. Iam si honorandi ratione officii, qui vtiles sunt in rebus ciuilibus et vita externa: si honorandi, qui curant salutem corporum: aequissimum est honorari ratione officii eos, qui vtiles et necessarii sunt in rebus ecclesiasticis, et vita interna seu spirituali, quique curant salutem animarum. Sed de altero honore potissimum Sermo est. Hunc uero omnibus in genere ministris debitum esse affirmat, idque argumentis efficacibus mox comprobat. At interim vt diuersa sunt in ecclesiis munia neque cuiuis ministro quaeuis conueniuel ta aliis plus, aliis minus honoris huius videt deberi: Vnde ait:  \pend\pstart Maxime ii qui laborant in sermone et doctrina:) Sunt illi quoque presbyteri siue diaconi alendi ab ecclesia, qui ministrant mensis: sed ii praecipue qui fungunt munere docendi quique administrant Sacramenta. Discimus hinc multo honorificentissimum esse munus ac iustissime eos liberaliter ali in ecclesiis, qui populum instruunt verbo et sana doctrina. Alia munia suum locum habere debent, nec sunt contemnenda: sed istud multo excellem tius. Sunt autem coniuncta sermo et doctrina: quia ille sermo maximam habet vtilitatem, non qui tumide ostentat insignem peritiam oratoriam, vel multam subtilitatem Philosophicam: sed qui docet veram pietatem et vitam format. Nam huc potissimùm omnes conditiones ministrorum debent tendere. Scientia inflat, charitas aedificat.1.Cor.8. Praeterea discunt hic ministri illi ecclesiarum, quod se sedulo debeant in sacraium literarum studiis exercere, cum hac parte, videlicet docendo, iudicentur praecellere, et ob hanc potissimum caussam iudicant digni, qui ab ecclesia sustententur. Ideoque ; Apostolus vsus est verbo Iaborandi, simul quo ostenderet, ministros esse debere admodum diligentes in tractando verbo, nullosque studiorum labores recusare:simul etiam vt obstruerent eorum ora, ꝗ putant sine multo labore recte doceri posse, neque laborem interpretant, nisi ꝗ manibus et corporis nisu exercet. Na ita est. Valde laborat studioss ministri ecclesiart, etiasi vulgo semp  \pend
\vspace{0.5cm}\noindent
\vspace{0.2cm}\rule{1cm}{0.2pt}\\ 
\hspace{0.2cm}\textit{mg}
\footnotesize Ambrofiut Duo diaco Septem dis cont. 
\normalsize| 
\hspace{0.2cm}\textit{mg}
\footnotesize Emphasiti Bene. 
\normalsize| 
\hspace{0.2cm}\textit{mg}
\footnotesize Presbytero ram est ben preesse. 
\normalsize| 
\hspace{0.2cm}\textit{mg}
\footnotesize Honor det tupreibyt nis. 
\normalsize| 
\hspace{0.2cm}\textit{mg}
\footnotesize A pari. 
\normalsize| 
\hspace{0.2cm}\textit{mg}
\footnotesize Primgruam presbytero rum munus in uerbo et de ctrinacollo catur. 
\normalsize| 
\hspace{0.2cm}\textit{mg}
\footnotesize Presbyteri ecclesiae se stentatione digmi. 
\normalsize| 
\hspace{0.2cm}\textit{mg}
\footnotesize Studiamini strorum eccte siae non contcuenda. 
\normalsize| 
\section*{AD TIMOTHEVM CAP. V. }
\marginpar{[ p.113 ]}\pstart feuarl videantur. Habent literarum lectio et veri sensus indagatio maximum laborem inclusum, habent item curae et solicitudines pro bona ordinatione in Ecclesia mo lestias grauiores, quam quiuis labor manuarius. Ecclesiarum ministri laborant, sicut et laborant magistratus ciuiles, vtriq autem animo magis que corpore: et tanto prestantius debet eorum opus iudicari, quanto pauciores inueniuntur qui illud opus queunt rite efficere. Interim otiosi et inutiles in ecclesiis, qui preter nomem nihil habent, quo se pro mini. stris venditant, quales heu nimium multi ecclesiis praeficiunt, viderint ipsi num ita laborent, vt suo victu digni iudicari possint. Certe Dominus per Ezech. cap.34-iis, qui se tantum pascunt et gregem negligunt verbo Dei pascere, grauissime comminat. Etin tales, ꝗ nec ver bo, neg aliis actionibus laborant, competit haud dubie illa lex, QyI NON LABORAT IS NEC EDAT. Praeterea videre hic licet quid maxime in his ministris ecclesiarum requiratur, nempe vt laborent in sermone ac doctrina. Sicubi aliquid desiderari videatur in vita, et fortasse minus videntur docere ipsis actionibus: at dent operam, vt, quod praecipuum est, laborent strenue in sermone et sana doctrina. Videant ne huius partes deserant. Sequuntur argumenta, cur ministri ecclesiarum sint alendi.  \pend
\phantomsection
\addcontentsline{toc}{subsection}{\textit{Dicit enim scriptura Boui trituranti non obligabis os: Et, dignus est operarius mercede sua. }}
\subsection*{\textit{Dicit enim scriptura Boui trituranti non obligabis os: Et, dignus est operarius mercede sua. }}\pstart Vtrunque argumentum ab autoritate sumitur. Prius a Mose, Deut.25. Posterius a Chri sto, Matth.1o. Illud quidem sic habet a minore. Boui trituranti non est obligandum os: Ergo nec ministro docenti neganda est alimonia.1. Corinth.9. ostenditur allegoricam es se locutionem. Et frequenter a re rustica similitudines et allegoriae in Euangeliis adi ducuntur: vt Matth.9. Messis quidem multa, operarii autem pauci. Confertur auten minister docens et enarrans scripturas Boui trituranti:sicut enim bos triturando excu tit grana ex aristis, facitque vt a paleis separentur grana: lta qui enarrat scripturam, sen sum et mysteria educit e suis folliculis facitque vt sana doctrina a noxia, et virtutes a vi tiis discerni queant. Alterum argumentum Christi est a communi sententia etiam vul, go in ore habita. Eadem argumenta et plura 1.Corinth.9. Galat.6.1.Thess... Nec erat o, pus plura adferre apud Timotheum qui ea bene nouerat, et potuit ipse populum plenius de eadem causa docere. Non debent quidem ministri affectare amplum adeo stipendium, ne videantur se magis velle pascere quam gregem. Et Matth.1o.quae gratis acceperunt a Deo, iubentur gratis dare. Et 1. Petri,. lubentur non affectare lucrum, sed propenso animo pascere gregem Christi. Et vt Act.2o. commemorat Paulus dicens, nihil se ab Ephesus accepisse: vtrum magis est iuxta Christi verbum dare quam accipere? At interim conferendum est in ministros omnes quidem sed praesertim eos qui docent, quantum satis est familiae honestae pascendae Tanta merces, inquit Ambrosius, debet esse Euangelizantis regnum Det, quaneque contristet̃, neque extollat. Pariter.n. et inopia illum premi et abundare periculosum est. Diligentior fit magisque assiduus esse potest in stu diis sacraru literarum, si non grauet inopia. Et crescit, inquit idem Ambrosius, in illo autoritas, cùm videt se etiam in praesenti laboris sui fructum percipere, non vt abundet, sed vt ?.  non desit: satrapas opibus aequare non debet, at probum ciuem et in honesta existi mationae? habitum merito aequabit, quantum ad sumptus attinet in sua familia expendendos. Pauper. tas etiam siduciam loquendi interdum adimit, et minuit autoritatem, praesertim apud rudes. I. taque ne in contemptum veniant ministri, ne etiam penuria pressi cogantur verum dissimu lare, vel adulari, decet eis prospici, vnde honeste viuant: ac tali mediocritate, vt neque popes videant fieri audaciores, neque per inopiam timidores, sicque nulla spe vel auiditate quaestus, nullo quoque metu potentiorum, intrepide et synceriter doceant. Propter has autem causas, nempe et propter viduas alendas caeterosque pauperes, et propter mins. stros ecclesiarum, maiores nostri opes non paucas ecclesiis contulerunt, imitantes pios illos diuites in primitiua ecclesia, qui ad pedes Apostolorum sua attulerunt pauperibus ex aequo distribuenda. Distribuebantur autem illa ecclesiis publice collata stipendia in qua tuor partiones, ut ex quadam Epistola Gregorii Magni ad Augustinum episcopum Can mariensem colligiur. Vna pars cpiscopon cedebac et̃ familiae, propter hospitalitazea  \pend
\vspace{0.5cm}\noindent
\vspace{0.2cm}\rule{1cm}{0.2pt}\\ 
\hspace{0.2cm}\textit{mg}
\footnotesize ludacu bores animi corporis labonibus sunt praestabiliores. 
\normalsize| 
\hspace{0.2cm}\textit{mg}
\footnotesize Lex contra pigros et otiosos. 
\normalsize| 
\hspace{0.2cm}\textit{mg}
\footnotesize Duo argumenta ab autoritate. 1. a mose. 2. a Christo. similitudo a re rustica. Bos triturans. A communi sententia. 
\normalsize| 
\hspace{0.2cm}\textit{mg}
\footnotesize Ministri uerbi debent esse lieni ab auaritia et cupiditate. Ministri uerbi honeste alendi. 1. Caussa. 
\normalsize| 
\hspace{0.2cm}\textit{mg}
\footnotesize Mediocritas optima. 4. 
\normalsize| 
\hspace{0.2cm}\textit{mg}
\footnotesize Exemplum liberalitatis maiorum nostrorum Quatuorpartiones eleemosynarum. 
\normalsize| 
\section*{IN I. EPIST. D. PAVLI. }
\marginpar{[ p.114 ]}\pstart atque susceptionem, alia clero (.1.pauperibus studentibus sacris) tertia pauperibus, quar ta ecclesus reparandis. Ac bene quidem tunc adhuc facultates ecclesiae administrabantur: sed breui post hypocrisi et superstitione in ecclesia praeualentibus, et quiduis adhuc ab ignaro diuitum genere impetrantibus, cum iam opes mirum in modum auctae essent, vera coepit siert illa Bernhardi querela, quôd religio quidem peperisset diuitias, sed filia deuorauit matrem. Diuitiae extinxerunt in ecclesia omnem religionem, quando per hypocritas et superstitionum magistros, ad omnem fastum et superbiam sunt accommodatae, nullaque amplius habita est pauperum alendorum ratio: Sed omnia qui dam ad suum priuatum vsum et pompam conuerterunt, ad haec nostra vsque tempora. Orandus Deus, vt abusus omnis extirpetur, et verus vsus facultatum ecelesiis restitua tur. Nam ecclesiasticis istis satrapis et episcopis aulicis miseri et famelici doctoies Euangelici frustra canunt.  \pend
\phantomsection
\addcontentsline{toc}{subsection}{\textit{Aduersus presbyterum accusationem ne admiseris nisi sub duobus aut tribus testibus. Eos, qui peccant coram omnibus argue, vt et caeteri timorem habeant. }}
\subsection*{\textit{Aduersus presbyterum accusationem ne admiseris nisi sub duobus aut tribus testibus. Eos, qui peccant coram omnibus argue, vt et caeteri timorem habeant. }}\pstart Apte, postquam egit de honorandis presbyteris, qui bene praesunt, adiungit ista de ar guendis, qui deprehensi fuerint culpandi. Indicat autem primo non facilem admittendam accusationem aduersus presbyterum. Vtque ita statuat, variae caussae eum mouere pote rant. Certum est, maligna esse multorum ingenia, et promptissima omni tempore ad lae dendum!eos, qui praesunt, id quod secula nostra nos plus satis docuerunt. Sed multo magis, quoniam veritas odium parit, laborant mali homines ad nocendum ministris veritatis, per quos iplorum vitia reprehendunt et produnt. Aegre enim ferunt se res prehendi, ac propterea oderunt censores presbyteros. Sicut puori pleruntque oderunt paeda gogum suum, et vulnerati aegrotique auersantur medicos, et fures ac scelerum conscii suspectos habent magistratus, pauperes inuident diuitibus Hinc ergo non statim sunt audiendi quiuis accusatores, cum plerunque soleant plebei presbyreris inuidere. Deinde si quiuis, et ob quamuis leuiculam rem audiatur, non parum imminuetur presby terorum autoritas, contra quos liberum est cuiuis et perditissimo dicere. Et quid non possunt confingere mali et colore ueri persuadere ubi male uolunt? Notum est subinde etiam subornatos falsos testes plures simul conuenientes, quorum astu ac dolis boni uiri et immerentes oceiderentur. Nam ad uotum lezabelis 1.Reg.21. duobus testibus filiis Belial uictus, ac ueluti blasphemus in Deum, iniuste lapidatus est Naboth. Et Christus Matth.26.apud Catapham grauissime accusatur blasphemiae a duobus falsis testibus. Stephanus item Actor.7. per falsos testes conuictus lapidatus est. Quae cum ira se habeant, non abs re Apostolus difficulter uult audiri accusationem aduersus prebyte rum, quem uix uerosimile est publici criminis reum posse peragi, cum non sit admissus ad munus illud, nisi qui per omnem uitam uixerit inculpatissime. Deinde ex eius accusatione macula quaedam uidetur inurenda toti ecclesiae, et uniuersus ministrotum ordo non parum dehonestari. Si itaque contingat aliquid contra presbyterum quempiam deferri, non admittatur accusatio, nisi hac conditione: Nisi sub duobus aut tribus testibus: Nam haec forma accusandi in foro praescripta olim Deuter.17. et 19. etiam in accusatione apud episcopum siue apud ecclesiam seruanda est: siquidem et Matth.18. iubentur prius duo aut tres adhiberi, quam ad ecclesiam crimen alicuius est deferendum. Tametsi propterea, quod etiam falsi testes possunt inter se conuenire plures, uti prius ostendimus, magna prudentia hicopus est.  \pend\pstart Caeterùm quoniam presbyteri etiam homines sunt, ac proinde et errare possunt et resipiscere, ideo ubi plures adfuerint accusatores, et certo comprobatum fuerit, presbyterum offendisse: hanc legem seruandam statuit: Eos qui peccant coram omnibus arque: Adeo ministerii functio illos non excusat, non adfert impunitatem: quin ubi publice peccauerint, publice quoque arguantur, et resipiscentiae exemplo omnibus, quos offenderunt errando, satisfaciant poenitendo et ueniam deprecando. Ac certe nist ita fiat, grauia scandala orientur, ac uidebitur episcopus peccatis illorum fauere, caeterique mox iudicabunt sibi eadem impune licêre: quin ipsa quoque persona, ad cusus peccata conniuetur, fortassis hinc in deteriora prola  \pend
\vspace{0.5cm}\noindent
\vspace{0.2cm}\rule{1cm}{0.2pt}\\ 
\hspace{0.2cm}\textit{mg}
\footnotesize Bernbgd Diuitias f perit relig at filia de rauit matt 
\normalsize| 
\hspace{0.2cm}\textit{mg}
\footnotesize Abusus faci taum eccles. estcarum. 
\normalsize| 
\hspace{0.2cm}\textit{mg}
\footnotesize Ingeniorum malitia Ns gque. cusandase presbyter. Feausa 
\normalsize| 
\hspace{0.2cm}\textit{mg}
\footnotesize 2. Veritaso dium parit. Matth. 14. oscae 4 1. Reg. 22 A Simili.  2. Inuidiaus gi. 
\normalsize| 
\hspace{0.2cm}\textit{mg}
\footnotesize 4. Autorita. tis uiolatio. 
\normalsize| 
\hspace{0.2cm}\textit{mg}
\footnotesize 5.Mendacia et calumniae: 6. Fassi testes. 
\normalsize| 
\hspace{0.2cm}\textit{mg}
\footnotesize 7. Auitae pro bitate et innocentia. 
\normalsize| 
\hspace{0.2cm}\textit{mg}
\footnotesize s.Matmate clesiae inuri tur ex accui fatione pres byteri. 
\normalsize| 
\hspace{0.2cm}\textit{mg}
\footnotesize Non conni. uendum pecc. tis publice admissis per presbyteros 
\normalsize| 
\hspace{0.2cm}\textit{mg}
\footnotesize Forma accus fandi presb, terum. 
\normalsize| 
\hspace{0.2cm}\textit{mg}
\footnotesize ScHdat per culam. 
\normalsize| 
\section*{AD TIMOTHEVM CAP. V. }
\marginpar{[ p.115 ]}\pstart betur, existimans crimen suum uix pro crimine reputandum. Apostolo satis fuit caussam finalem ab utilitate publica adiicere, dum ait: Vt et caeteri timorem habeant: Incutitur omnibus salutaris quidam timor ne quid tale audeant committere, ob quod vi dent publice reprehendi ipsos etiam presbyteros. Et quia plus periculi imminet a malc exemplo superioris, sicut magis dolent ob laesionem capitis quam alicuius alterius membri:ldeo etiam aequum est propter alios publicem redargui presbyteros, qui publice dell querint. A fonte infecto inficitur totus riuus. ltaque discimus hic nequaquam excipien das personas in ecclesiasticis censuris, sed simpliciter, qui publice deliquerint, public arguendos. Quod autem de publice admissis ista intelligenda veniant, id liquet exratione testium: nam quod plures viderunt, id quasi publicum aestimat. Discimus item cauendum mini. stris omnibus, ne vlla ratione offendiculum praebeant caeteris, qui ipsorum exemplo sun regendi, coercendi, in timore et sancta disciplina retinendi. Nomnulli hanc legem de publice arguendis qui publice peccarint, non contra presbyteros tantum latam intelligi uolunt, sed contra quoscunque peccantes: ld quidem verum est, ita eam exercendam: tametsi proprie hic contra presbyteros lata est, et sunt omnia haec uerba praecedentibus comnectenda. Neque  quae hic dicuntur, pugnant cum illis quae Matth.18. de correptione fraterna priuatim agenda dicuntur. Nam hic publica agitur, quaeque ; iam delata est ad ecclesiam. Quibus autem de rebus presbyteri possint apud episcopum accusari atque argui, longum foret disputare.  \pend
\phantomsection
\addcontentsline{toc}{subsection}{\textit{Obrestor in conspectu Dei et Domini lesu Christi, et electorum, vt haec serues, sine praecipitatione iudicii, nihil faciens iuxta propensionem animi. }}
\subsection*{\textit{Obrestor in conspectu Dei et Domini lesu Christi, et electorum, vt haec serues, sine praecipitatione iudicii, nihil faciens iuxta propensionem animi. }}\pstart Concludit obtestatione, ac uidetur alludere ad morem quo solent Sacramento astrin gi qui in republica iudicis uel alterius magistratus personam suscipiunt. Porro hac obtestatione indicat se grauissimam legem de forma accusationis et redargutionis presby terorum tulisse, eamque summa diligentia sequendam et seruandam.  \pend\pstart Obtestor, inquit, vi haecserues.) Graece διιαμαρτύρομαι quod significat protestor, praesentibus testibus depono Deos homines que contestor: et διιαμαρτυρο͂ν lurisconsultorum uerbum est: contestari in lite aliqua de scripto, et ut vulgo dicunt, coram omnibus protestari. Hic proprie significat Deum ipsum contestor, quod verba sequentia indicant: vtitur et huiusmodi obtestatione 2. Timoth.4.  \pend\pstart In conspecTu Dei et domini Iesu Christi, et electorum angelorum.) Vide qualem et quam vehementem obtestationem faciat. Per Deum patrem, et Christum dominum, et electos angelos, qui omnia intuentur et spectant, quos oportet semper imaginari, sicuti et vere sunt, omnibus actionibus nostris, praesertim iudiciis, praesentes esse et intentos: obtestor te, ô Timotheô, vt haec serues. Electos angelos vocat ad discrimen eorft qui ceciderunt a sua dignirate per peccatum, qui electi non sunt, sed reprobi et dannandi. Merito his uerbis eius pectus commoueatur, cui iudicium aliquod praesertim in ecclesia est tractandum. Ac dignissima essent haec uerba, quae non solum vbi episcopi audiunt ecclesiasticas causas, uerum etiam ad omnia fori tribunalia describerentur, ur semper in oculis iudicum uersarentur. Addit aliam conditionem in iudicio retinendam Sine praecipitatione iudicii, γω, ἰι προκριματet, sine praeiudicio, quemadmodum si dicere mus iudicandum, sed ita, vt nihil fiat in praeiudicium aliorum. Praeiudicium est, cùnex opinione, quam primum concepimus temere, pronunciamus de aliquare, causa ipsa nondum plene cognita, qua cognita aliter foret pronunciandum. Quare et recte interpres reddidit, Sine precipitatione iudicii. Prae cipitatio enim et festinatio facit, ul priusquam totam causam satis perspexerimus, parum recte pronunciemus. Eam erge merito Apostolus a iudiciorum forma abesse uult.  \pend\pstart Mmil Jacientes iuxta propensionem animi.) καταπρόσκλισιν, per inclinationem, ueluti in unam partem magis quam alteram inclinatus: Nos uno uerbo dicere possemus uulgari more, nihil faciundum in iudicio ex affectu. Excludendus ergo est a recto iudicio affectus, iudicandum sine ulla personarum acceptione, ex aequo et bono, non uni parti magis fauendum quam alteri, locum non habeat autoritas cuiusquam uel gratia Conueniunt haec uerba cu ilis Deuter.7. ludicent popuia iusto iudicio, nec in altera  \pend
\vspace{0.5cm}\noindent
\vspace{0.2cm}\rule{1cm}{0.2pt}\\ 
\hspace{0.2cm}\textit{mg}
\footnotesize Offendiculum cauendum mimstres. 
\normalsize| 
\hspace{0.2cm}\textit{mg}
\footnotesize Cassafira lis ab utitita Personae non debentrespici. te publica. Similitudo. 
\normalsize| 
\hspace{0.2cm}\textit{mg}
\footnotesize Ouesii
\normalsize| 
\hspace{0.2cm}\textit{mg}
\footnotesize Angeli electi. 
\normalsize| 
\hspace{0.2cm}\textit{mg}
\footnotesize Praeiudicium fugiendum. Praecipitatio iudicii stitiae. obest cognitioni et iu
\normalsize| 
\hspace{0.2cm}\textit{mg}
\footnotesize affectus in iudicio non multo minus huic obtenperandumi. Ratio cur consulendus: affectus non sit audiendus u iudiciis. 
\normalsize| 
\hspace{0.2cm}\textit{mg}
\footnotesize Angeli reprobi. 
\normalsize| 
\section*{IN I. EPIST. D. PAVLI. }
\marginpar{[ p.116 ]}\pstart partem declinent. Non accipias personam necmunera, quia munera excaecant oculos sapientum, et mutant causas iustas: luste quod iustum est persequeris. Tria ergo insignia documenta ex his verbis Paulinis colligere licet. Primo iudicandum vt Deo et angelis inspectoribus. Secundo, sine praecipitatione, id est, matura inquisitione et deliberatione. Tertio, sine affectu seu personarum acceptione. Quae tria non solum in ecclesiasticis actionibus, verum et in ciuilibus et forensibus merito sunt obseruanda.  \pend
\phantomsection
\addcontentsline{toc}{subsection}{\textit{Manus cito ne cui imponas, neq communices peccatis alienis. Temetipsum purum serua. }}
\subsection*{\textit{Manus cito ne cui imponas, neq communices peccatis alienis. Temetipsum purum serua. }}\pstart Commode, postquam dixit non cito iudicandum, adiicit similiter, non temere aut cito quenquam praeficiendum ecclesiae. Nam manus imponere, quemadmodum semel a nobis dictum est capite 4.nihil aliud est, quam externo eiusmodi indicio ac cere noniola indicare, testari, idoneum esse, cui publica functio in ecclesia committatur. traque per Metonymiam perinde est, atque aliquem ecclesiae praeficere: sicut sceptrum ulicui in manus dando solent creare regem, pileum dando seruis, solent liberos reddere. Satis ergo declarat, non paruum esse munus, praefici ecclesiae, ac valde accurate ab episcopis animaduertendum, qualis nam ad tantam functionem sit admittendus. Opus est prius bene explorari, id quod arguit voxilla ἐμφατικῶς legenda cito. Ideo. Ambrosius praecedentem obtestationem, non ad causam illam de accusatione et redargutione presbyterorum, sed ad hanc de non cito ordinandis retulit. Sed sequitur hic quoque ratio a periculoso:  \pend\pstart Neque communices peccatis alienis.) quasi dicat: si indignum et inidoneum praeficis, qui uitiis suis illam tunctionem deformat, iam vidêris illius peccatis communicare et fauere, fitque ́, vt eius vitia tibi imputentur: Istud, inquit, ne facias, Temetipsum purum serua: id est, vide ne tibi ipsi nota inuratur, propter improbam vitam eorum quos ordinaueris: Vide tuam conscientiam hac in parte liberes, eamque probatam, puram, Deo et hominibus reddas. Constat autem hinc perinde peccare eos, imo grauius, qui inidoneos ordinant, atque illos, qui cum inidonei sint, ordinantur, et post ordinationem suam in manifesta scandala prolabuntur. Communicant tales episcopi alienis peccatis, in eos peccata aliorum recidunt, atque ipsis imputantur: quia videntur approbando personam, simul approbasse vitia eius et peccata, et sic suam ipsorum conscien tiam commaculant. Huc pertinet illud Roman I. Digni sunt morte non solum qui ea faciunt, verumetiam qui assentiuntur iis qui faciunt.Item deligendos esse idoneos constat ex illo Numer.11. vbi dominus ad Mosen, congrega mihi septuaginta viros de senioribus Israel, quos tu nosti quôd senes sint populi ac magistri. Matth.5. Quomodo caecus caeco dux erit? Et, Si sal infatuatus fuerit, etc. Item omnia quae dicta sunt de diaconorum doctrina et vita.  \pend
\phantomsection
\addcontentsline{toc}{subsection}{\textit{Ne posthac bibas aquas, sed vino paululo vtere propter stomachum tuum, et crebras tuas infirmitates: }}
\subsection*{\textit{Ne posthac bibas aquas, sed vino paululo vtere propter stomachum tuum, et crebras tuas infirmitates: }}\pstart Quasi πάριργον obiter inspersum de tuenda valetudine, tametsi tacita potius lateat reprehensio superstitiosae abstinentiae, tanquam si respiciens ad illa verba:temetipsum purum serua: subindicaret, vsu vini hominem nequaquam impurum reddi, sed potius promouendo eos qui sunt inidonei. Atque ideo magis cauendum, ne quis polluat se bromouendo inidoneos, quam ne polluat se vtendo vino. Hac vero ratione comnectas ista cum praecedentibus, alians non. Fortassis autem vel a parentibus institutus, vel proprio ludicio adductus Timotheus agebat abstemium. At veron si ita abstinuisset quadam superstitione, tanquam si usu uini putaslet se polluendum, et aequa abstinentia uini fore mundiorem, merito iudicium id reprehendendum fuit. Quod si etiam adeo sibi in terdixisset usum uini uel etiam ciborum, ut ualetudini aliquid inde decederet, erat et huius intemperatae temperantiae admonendus. Certum est Paulum non sine iustissima caussa admonitionem istam apposuisse.  \pend\pstart Nepeshac bibut aquam, sed vine paululo vtere.) in calidis regionibus receptum est,  \pend
\vspace{0.5cm}\noindent
\vspace{0.2cm}\rule{1cm}{0.2pt}\\ 
\hspace{0.2cm}\textit{mg}
\footnotesize rauast semper sun expendenda 1. Se coras Deo et in tospectu An gelorum iadicare. Psal. 82. 2Sine preci pitatione, et maturo çonsilio examinationeque di ligentires agenda. 
\normalsize| 
\hspace{0.2cm}\textit{mg}
\footnotesize 3. Sine affectu et per sonarum ac ceptatione iudic is prae sidendum. Manus imponere quie hoc loco si gnificet. 
\normalsize| 
\hspace{0.2cm}\textit{mg}
\footnotesize Metonymia Ratio a periculoso. 
\normalsize| 
\hspace{0.2cm}\textit{mg}
\footnotesize Consćienti, puritas insit episcopo in ordinandis ministris. 
\normalsize| 
\hspace{0.2cm}\textit{mg}
\footnotesize Quomodo episcopi pe£ catis alienis communicem in praeficie. dis malis e incptis. A commination. A praecepte diuino. Vfisuinime 
\normalsize| 
\hspace{0.2cm}\textit{mg}
\footnotesize deratus expedit ualeti dini, nec pol luit honuix nisi acceda abusus, aus superftitio Superstitio sa abstinen tia auino. 
\normalsize| 
\section*{AD TIMOTHEVM CAP. V. }
\marginpar{[ p.117 ]}\pstart luuentutem, quae et ipsa calore naturali abundat, aqua eaque interdum cocta, vti ad sedandam sitim. Et paucis admodum in vsu est vinum non dilutum aqua. Olim Romanis matronis vini vsus in totum interdictus, ideoque institutum, oculo excipere aduentantem maritum, vt sentiretur num quae vinum olêret. Et nunc imnuptae in Gallia pletunq abstinent vino. Multa apud veteres in commendationem huius temperantiae sunt prodita. Plato de legibus pueris interdicit vino ad annum usque decimumoctauum, quod non sit necesse ignem igni addere: inde ad quadragesimum permittit modicum: post eam aetatem plusculum ad vitae taedium leuandum. ln summa, usus uini non dannatur, modo temperate sumatur. Augustinus modico vino lubenter vtebatur, ratus eo acui ingenium. Atque habet sine dubio vinum mirabilem medicandi uim, magnusque eius ad tueidum ualetudinem usus est: qua de re Aegineta in praeceptis salubribus cap.95. Qui sa nitati studet, quanta sit uini potentia, ignorare non debet. Plinius hbus 23.cap.1. idem docet: Idem liber 14. cap22 ostendit, quantum incommodet, si immodice sumatur: et Paulus Aegineta eodem capite. Paulus uero hic addit rationes ab utili et necessario petitas ere medica. Propter siomachum tuum et crebras infirmitates.) Aegineta stomacho gratum esse iudicat uinum qd'Paulo est acerbius. Fuluum uinum maxime ad id commendat. Et cum uinum, vt ibidem dicit, calorem excitet, sanguinem reddat utiliorem, et concoctionem a diuuet, colligi potest per contrarium, ista impediri, si quis uino in totum abstineat, ac proinde in aegritudines facile incidere. Eccles.31. Exultatio animae et cordis uinum moderate sumptum. lam uero ut ad institutum regrediamur, si quis ita potu aut etiam cibo abstinet, ut incidat in aegritudinem, hic iniurius est suo corpori, ac proinde culpandus non minus, imo magis, quam si iniuriam inferret corpori alieno, uulnerando aut alias laedendo, ut hac ratione reddat se homo ante tempus inutilem et inidoneum ad praestanda ea, quibus uel ecclesiam uel rempublicam iuuare porerat. Vna cum uiribus corp oris debilitatis debilitantur et uires animae. Constat ergo grauissime eos peccare, qui uel immoderata abstinentia, ieiuniis, uigiliis, etiam studiorum laboribus et similibus sese eneruant et excerebrant. Deo ista nequaquam placent, qui uoluntaria quidem, sed intra prudentiae limites posita requirit officia et obsequia: sicut periculosum est ita indulgere carni, ut extinguatur spiritus, ita periculosum est, sic indulgere spiritui, ut extinguatur caro. loannes abstemius nequaquam Christo sanctior, qui uino usus est. Sed nunc ita comparatum est, ut pauci admodum ministri ecclesiarum, uel alii, his praeceptis Paulinis egeant: Potius in diuersum praecipiendum est, ne uino sic, ut faciunt, deinceps utantur, idque propter stomachum quidem, seilicet quem grauant, et uinum erodit, uel infirmitates, quas abusu uini conciliant.  \pend
\phantomsection
\addcontentsline{toc}{subsection}{\textit{Quorundam hominum peccata ante manifesta sunt praecedentia ad iudicium, quosdam vero et subsequuntur. Consimiliter et bona ante manifesta sunt, et quae secus habent, occultari non possunt. }}
\subsection*{\textit{Quorundam hominum peccata ante manifesta sunt praecedentia ad iudicium, quosdam vero et subsequuntur. Consimiliter et bona ante manifesta sunt, et quae secus habent, occultari non possunt. }}\pstart Epanaphora est seu Regressio. Regreditur enim ad id, quod dicere coeperat, cauen dum uidelicet, ne communicaret peccatis alienis, inidoneos eccleslis praeficiendo. Ac uidetur latêre occupatio. Nam ubi Timorheus dixisset forte, qui possum cauere, quan do ignoro, quales homines sint, et mihi non constat de alienis peccatis? Respondet Pau lus: ideo, inquit, dixi, ne cui cito manus imponas, neue communices peccatis alienis. Nam uaria et fallacia sunt hominum ingenia. Quorundam peccata et enormis uita manifesta sunt, praecedentia ad iudicium, et sam ante nota, ut possis in tempore iudicare, tales nequaquam praeficiendos. Quosdam uero subsequuntur, quia uidelicet sic malitiam suam no, runt tegere et dissimulare ad tempus, ut non nisi postquam promoti sunt, sese prodan horrenda peccata designantes. Proinde si prioris generis hominem praeficias ecclesiae; plane in te culpa reiicitur qui nosse potueris antea illum prorsus fuisse inidoneum, et sic peccatis illius communicas. Quod si posterioris generis hominem contingat per te promotum, est quod excuseris, nec potest tibi imputari. Nam praeuidêre qualis quis futurus sit, id solius est Dei. Vt enim dicitur 1.Reg.16.Homo uidet ea quae comparent, Dominus autem intuetur cor. Hiobus 12. Dominus nouit et decipientem, et eum qui decipitur. Hinc finxit antiquitas Momum id in Deo creatore hominis reprehendisses  \pend
\vspace{0.5cm}\noindent
\vspace{0.2cm}\rule{1cm}{0.2pt}\\ 
\hspace{0.2cm}\textit{mg}
\footnotesize Vdoια la 1.cap.1. 
\normalsize| 
\hspace{0.2cm}\textit{mg}
\footnotesize Raio ab na tili et neces saio. 
\normalsize| 
\hspace{0.2cm}\textit{mg}
\footnotesize Apart 
\normalsize| 
\hspace{0.2cm}\textit{mg}
\footnotesize Icontraris, 
\normalsize| 
\hspace{0.2cm}\textit{mg}
\footnotesize Epanaphora. 
\normalsize| 
\hspace{0.2cm}\textit{mg}
\footnotesize Quina non fint admittendi ad ecclesiasticas fuctiones siue ad ministerium uerbi Dei. 
\normalsize| 
\section*{IN I. EPIST. D. PAVLI. }
\marginpar{[ p.118 ]}\pstart quod non addidisset fenestras humano pectori, per quas licuisset videre, quae’nam ho mo in pectore gereret, et num integrum et candidum foueret agnum, an potius vulpem seu lupum.1. Corinth. 12. Nemo nouit, quae sunt hominis, nisi spiritus hominis qui in eo est. Sententia ergo est, eos, de quorum peruersis moribus manifesto constat et potest iudicari, nequaquam promouendos ad ecclesiasticam functionem: Quod si promouentur, episcopos grauiter peccare et communicare illorum peccatis. Quod si vero per episcopum promoueantur, qui antea quidem semper se bene gesserint, sed postquam euecti sunt ad munus ecclesiasticum, tunc primum incipiunt hypocritae peccata commit. tere, episcopus non peccauit, neque communicat peccatis eorum: subsequuntur em̃ peccata illorum tempus illud, quo erat de ipsis iudicandum, bonine essent an mali. Itaque vocabulum iudicii non refertur ad iudicium Dei in fine seculi futurum, sed ad iudicium humanum quo examinantur homines in mundo bonine sint an mali. Ex quibus discimus, de iis, quorum peccata lam manifesta sunt mundo, libere debere iudicari, quod sint inidonei et indigni. Matth.7. Ex fructibus eorum cognoscetis eos. Neque potest arbor mala bonos fructus facere, neque colligunt de spinis vuam, neque de tribulis ficus. Rursus discimus eos qui mali sunt et hypocritae, tolerandos etiam in ministerio et ueros ministros esse, neque prius deiiciendos, nisi ubi certa malitiae suae signa aediderint. Sic Christus ludam tulit, per eundem signa et miracula fieri uoluit, donec ille saepe frustra admonitus et nolens resipiscere, qualis nam esset reipsa monstrauit. Tunc primum Christus aiebat illi, loan.13. Quod facis, fac citius: quibus uerbis non praecipiebat Christus illi peccatum, sed tantum indicauit ei Christus reiectionem ab electorum consortio. Postremo ob has caussas, discimus eos diu et uariis modis explorandos, qui ad ecclesiasticam functionem sunt euehendi.  \pend
\phantomsection
\addcontentsline{toc}{subsection}{\textit{Consimiliter et bona opera ante manifesta sunt: et ea quae secus habent, occultari non possunt. }}
\subsection*{\textit{Consimiliter et bona opera ante manifesta sunt: et ea quae secus habent, occultari non possunt. }}\pstart Per ἀντίθεσιν eadem ratione iudicari uult de bonis operibus. Quidam, inquit, comprobati sunt palam et omnium testimonio boni, et si per hypocrisin tantum bene age rent, ea non posset diu celari, sed tandem manifestaretur aliqua ex parte egrum malitia. Proinde qui hoc pacto semper comprobati sunt boni, neque unquam compertum estin eis aliquid diuersum, tales sunt omni commendatione dignissimi, ac merito promouendi Sententia satis est perspicua,  \pend
\endnumbering\beginnumbering\section{CAPVT VI}
\phantomsection
\addcontentsline{toc}{subsection}{\textit{Q Vicunque sub iugo sunt serui, suos dominos omni honore dignos ducant, ne nomen Dei et doctrina male audiat. Qui vero fideles habent dominos, ne contemnant, quod fratres sint: sed magis inseruiant, quod fideles sint ac dilecti, qui beneficentiae participes sunt Hec doce et exhortare. }}
\subsection*{\textit{Q Vicunque sub iugo sunt serui, suos dominos omni honore dignos ducant, ne nomen Dei et doctrina male audiat. Qui vero fideles habent dominos, ne contemnant, quod fratres sint: sed magis inseruiant, quod fideles sint ac dilecti, qui beneficentiae participes sunt Hec doce et exhortare. }}\pstart \huge\textbf{E}\normalsize PISCOPVS quomodo se geret erga seruos, siue quae docebit seruos factos Christianos: Ac de officio seruorum erga dominos alibi quoque diligenter egit Apostolus, Ephes.6. Coloss.4. et 1.Corinth.7. Sermo autem est deiis, qui uel natiuitate, uel captiuitate, uel alio casu ciuiliter seruiebant. Apparet autem quosdam seruos Christianos, auditis multis de libertate euangelica quae per Christum parta spiritualis est, libertas uidelicet a peccato, a morte, a diabolo, uoluisse doctrinam istam interpretari de libertate carnali, ciuili, et corporali, sicque se a suorum dominorum potestate subducere: quae res facile traduxisset et magnam inuidiam conflasset Christiano nomini: eamque ob caussam Apostolus fuit diligentissimus, ut libertatem euangelicam commode interpretaretur, ostenderetque ́ eam nequaquam tollere seruitutem ciuslem. Vndei. Corinth. 7. Vnusquisquein ea  \pend\pstart \huge\textbf{}\normalsize  \pend
\vspace{0.5cm}\noindent
\vspace{0.2cm}\rule{1cm}{0.2pt}\\ 
\hspace{0.2cm}\textit{mg}
\footnotesize Behaas humano in hac uita lo quitur Apl stolus, non de iudicioinfine seculi. 
\normalsize| 
\hspace{0.2cm}\textit{mg}
\footnotesize mridun. 
\normalsize| 
\hspace{0.2cm}\textit{mg}
\footnotesize Officium seruorum fidetium ergado minos. 
\normalsize| 
\hspace{0.2cm}\textit{mg}
\footnotesize Falsa libertatis interpretatio. 
\normalsize| 
\section*{AD TIMOTHEVM CAP. VI. }
\marginpar{[ p.119 ]}\pstart vocatione, in qua vocatus fuit maneat: seruus vocatus es, ne sit tibi curae. Quin etiam si potes liber fieri, potius vtere. Etenim qui in Domino vocatus est seruus, li bertus Domini est. Idem reipsa ostendere uoluit, quando Onesimum, quem Chusto genuit Romae in vinculis, domino suo Philemoni remisit. Euangelium siqui, dem non abolet eiusmodi politicas constitutiones. Caeterùm videntur quidam callidi serui huiusmodi rationibus voluisse a se iugum excutere, dicentes: Si dominus nobis est impius, merito ab eo discedimus, quia non est aequum, vt Christianus, Christi sanguine redemptus et factus liber seruiat impio. Sin pius est domi nus, uel ideo iustius est discedendum, quia iniquum est agere contra quenquam im periose, quem nouit sibi per Christum factum parem et fratrem. At Paulus hic tan quam per αντιςρι φωτα has rationes uoluit euertere, hoc fere pacto argumentans Domini quibus seruitis, aut sunt fideles, aut sunt impii. Si impii sunt, at ideo illi sun honorandi, ne videlicet nomem Dei et doctrina male audiat. Si fideles sunt, at eo iustius sunt honorandi, quia fideles sunt ac dilecti et beneficentiae compensatores. Haecbreuiter est Pauli sententia. Nunc singula uerba excutiamus.  \pend\pstart Quicunque sub iugo sunt serui.) Plenius significare uolens seruos suis dominis maxime deunctos et subiectos esse debere, sicque in obedientia cupiens ipsos retinere, dicit, eos sub iugo agere, idque metaphora sumpta a bobus, ut apparet, quibus proprie iugu imponitur, quando ad labores educuntur, atque ut boues imposito iugo facile ad quaeuis ab aratore flectuntur: lta homo factus seruus ad omnia facienda ab hero impellit: Tales itaque facti serui, suos dominos omni honore dignos ducant : τος ἰδινε δεσπότας, id est, proprios, priuatos illos dominos suos: nimirum generoso animouult esse seruos Christianos in infideles dominos, ut uidelicet dominos suos sic aestiment, tanquam quibus ipsimet gene rositate animi imperare potius debeant, saltem quibus in uirtutum omni genere uelini praeferri et praeire. Quo animo legimus de Aesopo, quod cum uenundareturut seruus, ro gatus quid sciret facere, respondit se scire imperare liberis. Quidam aiunt, quod iusserit praeconem edicere, Si quis uelit emere dominum. Aesopus emptus a Xeniade dixit huic, uide ut iussis meis pareas: cumque hic diceret, id praeposterum fore: Si medicum, inquit, emisses aeger, non illi obtemperares? Hoc itaque pacto Christiani seruisuos dominos infideles, generoso animo suos et proprios iudicabunt. Caeterùm multum exigit Apostolus a Christianis seruis, quando ait, dominos suos Omnihonore dignos ducant: non dicit simpliciter, honorent eos, sed praeterea dignos ducant honore: deinde etiam, Omni honore: existiment reuera deberi illis honorem ceu dignis. Per honorem autem hunc omnia officia humanitatis, et ciuilis subiectionis intelligi conuenit: neque uult Apostolus ullum huiusmodi officium a seruo Christiano relinqui.  \pend\pstart Ne nomen Dei et doctrina mase audiat.) Ratio ab inconuenienti, quod ex negle cto officio sequeretur. Sane si post fusceptam religionem Christianam serui ull redderentur insolentiores praefractioresque ́, ubi antea fuerant modesti et morigeri, hanc mutationem in deterius protinus imputabunt nomini Christiano et doctrinae euangelicae, statim omnes uno ore pessime de ea iudicabunt, ac dicent, perniciosam esse doctrinam, quae sic etiam seruorum peruertit animos. Imo si quid an plius auderent serui, iudicium ullum proferendo de paranda sibilibertate, iam pla ne seditionis reum peragetur Euangelium, et in perniciem omnium piorum im pii coniurabunt. Quae mala ne eueniant, aequum est, ut serui Christiani dominos suos etiam infideles honore dignos existiment. Nec leue malum est, id committere, propter quod male audit statim uniuersa religio, et quam plurimi ab ea alienantur, quos magis decebat nostra modestia et officiositate ad amplectendam religionem nostram allicere. Vnde Roman. 2. Apostolus exprobrat ludaeis, quod nomen Dei propter eos male audiret inter gentes. Contra magnis laudibus dignissimum est, et ad gloriam Dei propagandam conuenientissimum, si quis ita uitam moresque instituat, ut etiam hostes mirari, laudare, approbare religionem cogantur. Vnde si modeste et benigne et liberaliter seruiant serui suis heris infidelibus, feri  \pend
\vspace{0.5cm}\noindent
\vspace{0.2cm}\rule{1cm}{0.2pt}\\ 
\hspace{0.2cm}\textit{mg}
\footnotesize ubertate Christiana non tollitur seruiuus cus lis. 
\normalsize| 
\hspace{0.2cm}\textit{mg}
\footnotesize Dilemma. Anabaptista rum argumentum. 
\normalsize| 
\hspace{0.2cm}\textit{mg}
\footnotesize Argumentem ab muersione. 
\normalsize| 
\hspace{0.2cm}\textit{mg}
\footnotesize Metaphora a bobus, 
\normalsize| 
\hspace{0.2cm}\textit{mg}
\footnotesize Cbristani infidelibus etiam dominis seruire tenentur. 
\normalsize| 
\hspace{0.2cm}\textit{mg}
\footnotesize Emphasis onn. 
\normalsize| 
\hspace{0.2cm}\textit{mg}
\footnotesize Ab inconuenienti. 
\normalsize| 
\hspace{0.2cm}\textit{mg}
\footnotesize Ruangetis male audit propter seruorum iuobedientiam et improbitatem. Scandalum. 
\normalsize| 
\section*{IN I. EPIST. D. PAVLI. }
\marginpar{[ p.120 ]}\pstart potest vt heros suos Christo lucrifaciant: quae res foret ad nominis Christiani et de: ctrinae gloriam.  \pend\pstart Qui vero fideles habent dominos, ne contemnant, quod fratres sint, sed magis inseruiant.) Egit de seruis, qui nacti fuerant dommos infideles, nunc de iis qui nacti fideles. Hi ergo non debent contemnere suos dominos propterea, quod ratione fidei fratres et pa res coram Deo facti sunt, eundem patrem coelestem inuocantes in eodem spiritu, eiusdem filii lesu Christi sratres facti et cohaeredes. Ista quidem recte sicse habent coram Deo et in spiritu. Galat.3. Quandoquidem in Christo non est mas neque foemina, neque seruus neque liber, neque Graecus neque ludaeus. Ac coram mundo interim manet conditio externa et ciuilis, neque magis ex seruo fiet per religionem liber, quam ex ma re fiet foemina. ltaque ius, quod per politiam antea consecutus est herus in seruum, nequaquam tollitur per religionem Christianam. Ac proinde non habet Christianus seruus, quod contemnat suum herum eo nomine, quodis quoque sit Christianus et fidelis. Imo hoc nomine, magis, promp tius, diligentiusque ei inseruiet. Vnde sequitur ratio per inuersionem: quasi diceret: Contemnere vultis dominos quia fideles? Imo, ideo, inquam, pluris eos facietis, quod fideles sint. Fides communis debet animorum affectus copulare: Si seruus nollet se contemni ab hero, quod ipse esset fidelis: iniquum est, quod seruus contemnet herum, quia similiter est fidelis. Et si quae mutatio in melius facta est in seruo, per religionem, par quoque mutatio contigit in hero perreligionem: quare sicut, quod ad internam uitam et spiritualem attinet, omnia videntur inter seruum et herum facta paria: ita, quod ad externam et spiritualem uitam attinet, omnia priora manent paria: tantùm mitius, modestius, hilarius, candidius, quam antea, erga inuicem sese gerere herus et seruus poterunt. Huiusmodi mutatio ratione religionis exterius potest sieri, alia non. Et si iudicare non potes aliter, quam aequum esse, vt seruias domino tuo, si foret infidelis, quanto aequius est, eidem seruias sam fideli? Sed alia adhuc mouere animos seruorum possunt: quod fideles sint, ac dilecti. Vbi nos habemus dilecti, Graece legitur ἀγαπηrοι, a verbo αγαπάω, quod significat diligo, bont consulo, contentus sum, acquiesco alicui rei. Et Hieronymus in epistola ad Philemonem, dicit αγαποτον significare dignum qui diligatur. Proinde inspecta significatione vocabuli, cum heri fideles facti sint ἀγαποτι id est, dilecti et vere digni qui diligantur, cum etiam facile contenti sint minusque morosi, et boni consulant ea quae fiunt, nec non lubenter acquiescantiis, quae recta sunt et honesta: plane aequissimum, talibus ex animo morem geri, ac nullo modo veniunt contemnendi. Atque indicat istud verbum, quod seruus iam non amplius seruili metu debeat Christianum herum suum reuereri et obseruare, sed magis beneuolo et amico animo omnia facere, sicut ipse herus quoque non amplius herili saeuitia est vsurus: tantum nimirum in vtroque efficiente religione et fide, per charitatem non fictam.  \pend\pstart Qui beneficentiae participes sunt.) Tertia ratio adiecta, ob quam decet illis per omnia obedire. Eorundem beneficiorum Dei participes sunt, videlicet gratiae et uitae aeternae propter Christum promissae et exhibitae. Si itaque fidelibus eadem beneficentia Dei com munis sit, non est quod quisquam fidelium alterum contemnat, multo minus, quod seruus quispiam fidelis vilipendat suum herum fidelem. Caeterùm quae hic dicuntur de officio seruorum Christianorum, erga heros suos siue fideles siue infideles, accipi debent in genere de officio quorumcumque Christianorum erga suos magistratus, siue fideles siue infideles.  \pend\pstart Hac doce et exhortare.) Breuis et grauis Conclusio, qua praecipit ea omnia, de quibus hactenus disseruit, sedulô et magna cum diligentia inculcari. Nihil relinqui debet intentatum ab episcopo, vt suorum animos in sana doctrina confirmet. Docere debet cum summa autoritate eos, qui sunt ignari: et exhortari debet cum summa modestia eos, qui iam cognitionem aliquam sunt consecuti.  \pend
\phantomsection
\addcontentsline{toc}{subsection}{\textit{Si quis aduersam sequitur doctrinam, et non accedit sanis sermonibus Domini }}
\subsection*{\textit{Si quis aduersam sequitur doctrinam, et non accedit sanis sermonibus Domini }}\pstart  \pend
\vspace{0.5cm}\noindent
\vspace{0.2cm}\rule{1cm}{0.2pt}\\ 
\hspace{0.2cm}\textit{mg}
\footnotesize Aminon. maiut. 
\normalsize| 
\hspace{0.2cm}\textit{mg}
\footnotesize A condul ne immut. bili. 
\normalsize| 
\hspace{0.2cm}\textit{mg}
\footnotesize Fides animos bomil consociat. Apari. 
\normalsize| 
\hspace{0.2cm}\textit{mg}
\footnotesize A minori. 
\normalsize| 
\hspace{0.2cm}\textit{mg}
\footnotesize Quanta of cia heri de beant seruis et serui hi ris, si fides utrosque per chartatis spiritum obuget. 
\normalsize| 
\section*{AD TIMOTHEVM CAP. VI. }
\marginpar{[ p.121 ]}
\phantomsection
\addcontentsline{toc}{subsection}{\textit{nostri lesu Christi, et ei, quae secundum pietatem est, doctrinae, is inflatus est, nihil sciens, sed insaniens circa quaestiones ac disputationum pugnas. }}
\subsection*{\textit{nostri lesu Christi, et ei, quae secundum pietatem est, doctrinae, is inflatus est, nihil sciens, sed insaniens circa quaestiones ac disputationum pugnas. }}\pstart Ad animum Timothei confirmandum, magisque excitandum, instituit orationem con tra hypocritas et pseudoapostolos, tanquam si dicat: Bono sis animo Timothee: Auda cter perge docere,uti et alias ex me audisti, et hac epistola iam perscripsi. Quod si qu aliud doctrinae genus exigant, uel nonnulli aliud proferant, et hanc doctrinam, quam tradidi, miris artibus impugnent atque opprimant: tu interim nihil moueare, sed potius contemnito illos ut inflatos, turgidos, uanos, insanos sophistas, seiungere ab illis.  \pend\pstart Si quis diuersam sequitur doctrinam ἐτις ἐτεροδιιδασκαλα.) Est autem ετεροδιιδασκαλο͂υ, uel dillersam doctrinam amplecti, ut et ad praeceptores et ad discipulos referri queat. Ac sine dubio iam amodo multi huiusmodi falsi doctores et discipuli irrepserant in ec clesiam Ephesinam: sicut et Act.2o. Apostolus futurum praedixerat: ac propterea merito contra tales Apostolus hic grauissime disserit, eosque suis depingit coloribus, acuitandos docet.  \pend\pstart Et non accedit sanis sermonibus Domini nostri Iesu Christi.) Doctrinam quam exposui hac epistola, uocat sermones Christi, quia nimirum tota a Christo est tradita, neque usquam uel latum unguem discedens a Christo, ideoque merito ei adhaerendum. Vocar autem Christi sermones sanos ύγιαίνος, id est, sanos et falubres, ab utilitate eos commer dans, ut tanto auidius eos amplectamur. Hinc dicit Sapientia prouerbus 8.lusti sunt sermones mei, non est in eis prauum quid neque peruersum.  \pend\pstart Et ei quae secundum pietatem est doctrinae.) Per interpretationem, cuiusmodi sint sermones Christi exponit: uel obliquem perstringit doctrinam de exercitatione corporali:nem huic paulo ante pietatem opposuit. Saltem intelligit per doctrinam secundum pietatem, eam, quae est de fide et charitate, in quibus commendandis non solum haec epistola, uerum omnia scripta Pauli uersantur. In summa igitur, qui ueram pietatem nequaquam docet, uel accedit alteri, is inflatus est, nihil sciens: En quibus titulis eos ornat. Inflatus est, quamlibet doctus appareat, et sibi et imperitis suis sectatoribus: quamlibet splendeant omnia coram mundo, quae dicit et facit: sit licet sapientia carnali seu mundana optime instructus: at si quod uerum est licet dicere, ac simpliciter respicimus ad veram ac diuinam philosophiam euangelii, talis nihil scit. Scientia sine chars. tate inutilis: quia 1.Corinth.8. Scientia inflat, charitas aedificat. Verascientia pacifica et modesta est, arrogantiam nescit: contra inscitia siue opinata scientia, qua quis praesumit et falso credit se scire, quae reuera ignorat, talis inflat. Inflati non possunt habere veram scientiam, quia, vt 1.Petris. dicitur, Deus superbis resistit, humilibus dat gratiam. Sedinsaniens circa quaestiones et disputationumpugnas:) νοσῶν significat aegrotantem, averboνοσέω, ac refertur ad animum, ideoque interpres nobis reddidit insanientem Mente sani non sunt, dum sese sicinuoluunt quaestionibus et disputationum pugnis ut nunquam possint se extricare. Ac verem languent, vbi in cetta ac solida doctrina fi dei nequeunt consistere, neque habent, quo se sustentent, quamdiu vt dubii et animo deficientes in scabie quaestionum volutantur. Quibus fidles deest, illis necessario inest mira animi aegritudo. Et si tales sunt ipsi doctores ac medici, id est, aegroti, quomodo hialios curabunt ? Curent semetipsos prius ab aegritudine animi sui, ac tum conentur mederi aliis. λογομαγίαε sunt et verborum et disputationum pugnae. Huiusmodi itaque animo languentes sunt omni studio deuitandi, quia plurimum obsunt veritati, aliud docent praeter pietatem, et quaestiones inutiles, imo noxias et periculosas obtrudunt. Et omnibus seculis aliqui sunt inuenti, multi quidem ex ludaeis, (Galat. 1. Si quis uobis praedicauerit Euangelium praeter id, quod accepistis, anathema sit.) et philosophis tempore Apostolorum: deinde haeretici diuersi, qui quaestiones et λογομαyιας excogitarunt innumeras, vix seipsos intelligentes. Monstrosae mentis monstrosa est ratio: et non possunt ex insana mente alii quam insani sermones proficisci, verborum pugnas excitantes. Stoici accusantur inter philosophos, quod verba inusitata consinxerint, ut προνγμυνα et απoxpoνγμυνe, allaque eiusmodi. Valentinus  \pend
\vspace{0.5cm}\noindent
\vspace{0.2cm}\rule{1cm}{0.2pt}\\ 
\hspace{0.2cm}\textit{mg}
\footnotesize Oralo con tra bypocritas et pseudospostolos. 
\normalsize| 
\hspace{0.2cm}\textit{mg}
\footnotesize Mαα- πκαλdς qd. 
\normalsize| 
\hspace{0.2cm}\textit{mg}
\footnotesize 1.vfus. Doctrina Christi est sermosanus, ad antmae sa diens. uio ab vtilitate. 2. vfus doctrinaesanae a causa formali. lutem expeCommenda 
\normalsize| 
\hspace{0.2cm}\textit{mg}
\footnotesize Doctrina se cundum pic tatem est do Arina fidei et charitatis. 
\normalsize| 
\hspace{0.2cm}\textit{mg}
\footnotesize Mctapbora a corpore ad arimum. 
\normalsize| 
\hspace{0.2cm}\textit{mg}
\footnotesize curateipsum. Fidei destitu tio animi pa rit egritudinem. A prouerbio. Medice    
\normalsize| 
\section*{IN I. EPIST. D. PAVLI. IN I. EPIST. D. PAVLI. }
\marginpar{[ p.122 ]}\pstart excogitaust quadrigas αἐωνων: Alii alia absurdissima sunt commenti. Et nostro seculo, qualia ipsi aliquando audiuimus, in theologorum scholis ueluti Scotistarum realitates obiectiuas et subiectiuas, vniuersalia, formalitates: tum ridiculas quaestiones, quas longum foret et inutile recitare: Licet ea uidere in quorundam scriptis, quae adhuc alibireperiuntur. Omnia autem a uera pietate aliena, et ab Apostolo hic damnata: Et merito: Nam sequitur:  \pend
\phantomsection
\addcontentsline{toc}{subsection}{\textit{Ex quibus nascitur inuidia, contentio, maledicentia, suspiciones male, superuacaneae conflictationes hominum mente corruptorum, et quibus adempta est veritas, qui existimant quaestum esse pietatem. Seiungere ab his, qui eiusmodi sunt. }}
\subsection*{\textit{Ex quibus nascitur inuidia, contentio, maledicentia, suspiciones male, superuacaneae conflictationes hominum mente corruptorum, et quibus adempta est veritas, qui existimant quaestum esse pietatem. Seiungere ab his, qui eiusmodi sunt. }}\pstart Deteltatur logomachias, argumentis sumptis a periculoso et damnoso: oster ditque ́ per congeriem malorum inde enascentium genera. Sic autem mala haec proueniunt. Scientia opinata adducit fastum, reddit inflatum: qui inflatus est, contemnit mox alios prae se: contemptus iste parit inuidiam: Inuidia excitat contentiones et rixas:E rixis erumpunt conuicia seu maledicentia: Ab his nascuntur suspiciones malae et machinationes: Sic denique immortale odium et nunquam remittendae conflictationes fouentur ac retinentur.  \pend\pstart Suspiciones: ὑπόνοιαε πσοοηραε.) ύπόνοια significat suspicionem, praeterea et superbiam, confidentiam, Item allegoriam siue siguram, qua aliud dicitur, quam audiens expectabat. Nonnulli interpretantur suspiciones malas, impias suspiciones ac dubitationes de rebus diuinis, quae tunc oriuntur in animis, dum singula, quae ad pietatem pertinent, vocamus in dubium, ita vt modo hoc, modo illud suspicemur, ac tandem quo diutius cogitauerimus, eo minus certi habeamus: Id quod Simonidi accidit, qui rogatus ab Hierone Tyranno, vt quid Deus esset definiret, impetrato aliquot dierum ad cogitandum spacio, tandem respondit, se minus deprehendisse, quam antea. Vt Cicero liber de natura Deorum 1.  \pend\pstart Superuacaneae conflictationes.) παραδιιατριθαι. διατριθὰ sunt congressus et disputationes philosophorum: addita praepositio facit significari inutiles et superuacaneas. Merito vero superuacaneae et inutiles, quia frustra quaeritur cognitio eius rei, per spinosas quaestiones et verborum ambages, quae compendiose et simplici oratione veritatis satis cognosci potest: si tamen talis sst, quam cognoscere sit necessarium. Nam si non sit cognitu necessaria, superuacaneum quoque est de ea inquirere. Itaque multis nominibus ceu multorum malorum causae improbantur quaestiones et λογομαχίαε. Vbi animus promptus est ad cognoscendum veritatem, nihsl opus est argutiis, sed sufficit simplex veritatis oratio. Argutiae laedunt et obscurant potius, quam explanant vel illustrant. Illius tantum est quaerere et quaestiones mouere, qui dubitat. At vera fides nihil dubitat: quare nec agnoscit quaestiones. Et praecipua religionis mysteria praestat nude, vt proponuntur, credere, quam scrupulosem vocare in quaestionem. Non damnabitur quispiam quia nescit in totum qualis sit natura Dei: quomodo Deus pater genuerit filium, etc. Sed damnabitur, si non humiliter credat, Deum esse simpliciter optimum potentissimumque ́: si non credat, filium ex patre genitum, Spiritum Sanctum exutroque procedentem: etc. Hoc ipso nomine quod quaedam inquiris, ostendis te dubitare, et fidem tuam ceu imperfectam prodis, et accusas, et damnandam prostituis. Quid quod quaedam vt maxima diligentia inquirantur, eô sunt intellectui nostro minus peruestigabilia ? Ac Deus certo consilio, ut nimirum fragilitatem nostram semper agnosceremus, quaedam latere nos uoluit. Prouerbus 25. Qui scrutator est maiestatis, opprimetur a gloria. Ergoin suiaumenius, quaererencel, quod cogmtu eit meecharmam,mmuthtlitaiciaHetitquod inuentae facit credere, creditam facit diligenter obseruare. Quicquid ultra hos limites quaeritur, curiositatis est, non necessitatis, et maiore fructu ignoratur, quam queritur:ut non dicam, quod intelligi fortassis non potest:aut, si putes te intelligere, intelligis tamen inutiliter, et multa mala ex curiosis quaestionibus tuis enascuntur. Sed sequitur adhuc altera improbatio quaestionum et logomachiarum, a qualitate personarum.  \pend
\vspace{0.5cm}\noindent
\vspace{0.2cm}\rule{1cm}{0.2pt}\\ 
\hspace{0.2cm}\textit{mg}
\footnotesize Argumen a periculo so. Enumeratit malorum enascentium ex opinata et fasa se entia 
\normalsize| 
\hspace{0.2cm}\textit{mg}
\footnotesize Simonides Deum defiture nescit. 
\normalsize| 
\hspace{0.2cm}\textit{mg}
\footnotesize Derebus c gnitu necessaeiis disputationes in stituendae sunt Christi anis. 
\normalsize| 
\hspace{0.2cm}\textit{mg}
\footnotesize Graio uer tatissimples excludit argutias. Quenan de Det natura sint abscon dita homin intelligentiae Quaena co gnitu ad sali tem necessa ria. Mysteria recondita. 
\normalsize| 
\section*{AD TIMOTHEVM CAP. VI. }
\marginpar{[ p.123 ]}\pstart Hominem mente corruptorum, et quibus adempta est veritas.) Saepem ex conditione doEtorum solet eorum doctrina aestimari. Sed neque verosimile est, eos sanum aliquid posse tradere, qui ipsi mente sunt corrupti. Si corrupta est eorum mens et conscien tia: Igitur et corruptum est siue ad corruptionem tendens, quicquid ipsi in medium proferunt. Et cum ne in philosophia humana quidem quid rectem statuere possint, nisi qui integra et sana sunt mente, quanto minus id hi facient in philosophia diuina? Sec quid faciant quibus adempta est veritas ? Nimirum altercando amittitur veritas, inquit ille. Iraque hi dum multitudine quaestionum omnia videntur velle intrestigare, tandem hoc consequuntur, vt nihil certi retineant, imo vt amittant veritatem, si quam prius habuerint.  \pend\pstart Qui existimant quaestum esse pietatem.) Videtur tribus de causis improbare falsos doctores: a malignitare ingenii: ab ignorantia veri: et tertio ab auaritia. Ac reuera comprehenduntur his tribus omnia, quibus solent homines ad falsam doctrinam serendam moueri: Sed latissime regnat auaritia, a qua Apostolus ostendit se semperalienissimum, prout atque vnumquemque cuangelicum doctorem decet. Nam vna haec fere nota est certissima, qua verum doctorem a falsis discernere queas. Falsi doctores existimant quaestum esse pietatem: quicquid ad ipsorum lucrum facit, nouerunt id interpretari ceu pertinens ad cultum Dei: Sophistico lenocinio persuadent hominibus, quicquid sentiunt sibi suoque ventri vtile: de quibus dici potest illud Sapient. 15. Existimauerunt lusum esse vitam nostram, et conuersationem vitae compositam ad lucrum, et oportere vndecunque etiam ex malo acquirere. Aetas nostra habuit harum imposturarum certiora et plura argumenta, quam vt opus sit copiosius hac de reagere. Contra auaritiam autem doctorum, tum prophetae, tum Christus, tum Apostolus, commemorans suos labores, creberrime sunt concionati.  \pend\pstart Sciungere ab his qui etusmodi sunt.) Non opus est callidas eorum quaestiones dilsoluere, id quod philosophi quoque nonnulli satis animaduerterunt, irridentes interdum potius absurda parandoxa et dilemmata captiosa, quam dignantes respondêre Neque etiam contra mente corruptos, et a ueritate alienos auarosque ; sophistas verbis est contendendum, et clauus clauo fundendus. Sed hoc faciet pius minister verbi, suosque vtidem faciant, admonebit, postquam aliquoties viderit se frustra monitorem factum! neque spem superesse, ad meliorem frugem reuocandos hypocritas, ipse eos vitabit, nevel ipse uel sui ab eis tanquam contagiosis pecudibus inficiantur Grex totus in agris, Vnius scabie cadit et prurigine porci. Quibus nocere nequa quam volumus, sicuti nec aequum foret, id salrem vel nostri caussa est faciendum, nos ab eis seiungemur. Tit.3. Sectarum autorem hominem post vnam et alteram admo, nitionem fuge, sciens quod euersus sit qui eiusmodi est, et peccet per se damnatus. Pro uerbus 2o Honor est homini, qui se separat a contentionibus.  \pend\pstart Est autem quaestus magnus pielas cum animo sua sorte tontento. Notans hypocritas et falsos doctores auaritiae, dixit, eos existimare, quaestum esse pietatem. Proinde parans nunc se ad disserendum contra auaritiam ministris verbi maximopere uitandam, statim per inuersionem siue ἀντίςρέφοy, dicit, ipsam pietatem esse quaestum, maxime cum adiacet animo, qui sua sorte contentus esse norit, qualem animum et ipsa pietas ef, tilitas. ficit. Et uere ita est. Nam ex pletate utilitas prouenit et ad ipsos, qui eam habent, et ad alios qui ab his per Charstatem adiuuantur. Quaestus autem quantum confert uni auaro, tantum aliis inique detrahit, si tamen ipsi auaro aliquid confert. Et pietas habet promissiones uitae praesentis et futurae, quod semel dictum est. Pietas itaque nouit uni Deo fidere, statuereque sibi nihil defuturum, neque in hac uita, neque in altera: ac proinde nequaquam sibi inhiandum terrenis et fluxis istis opibus. Atque ita facile reddit ipsa animum sua sorte contentum, vtneque praesentia bona, quae obueniunt, fastidiat, neque absentia ualde anxie desideret. Sic autem ipse Apostolus dese Philipp.4. Ego didici in quibus sum, iis contentus esse: noui et humilis esse, noui et excellere, ubique et in omnibus institutus sum et saturari et abundare et penuriam patl, omnia possum per Christum, quime corroborat-Prouerbm. Qui  \pend
\vspace{0.5cm}\noindent
\vspace{0.2cm}\rule{1cm}{0.2pt}\\ 
\hspace{0.2cm}\textit{mg}
\footnotesize u **4t4 4 Raict personarti. Qui mente sunt corrupti. ii nibil. sanum trade Le ualent. A minore 
\normalsize| 
\hspace{0.2cm}\textit{mg}
\footnotesize mmooë causis falsi; doctores tiprobantio. L Ingenii ma lignitas. 2. Ignoraiio ltert. I. Auaritia. Sora falsotum doctorum. 
\normalsize| 
\hspace{0.2cm}\textit{mg}
\footnotesize uauritia religioni nihil est perniciosius. 
\normalsize| 
\hspace{0.2cm}\textit{mg}
\footnotesize Clauum clauo tundere. 
\normalsize| 
\hspace{0.2cm}\textit{mg}
\footnotesize Praeceptum, quod tenebit minister uerbi contra hypocritas. Prouerbium agrege. 
\normalsize| 
\hspace{0.2cm}\textit{mg}
\footnotesize Pietatis utilitas. 
\normalsize| 
\section*{1 N I. EPIST. D. PAVL I }
\marginpar{[ p.124 ]}\pstart confidit in diuitiis suis corruet, iusti autem quasi virens folium germinabunt. Neque quenquam destituit Deus, nisi qui sibi ipsi deesse velit. Psal.36. Non vidi iustum derelictum, nec semen eius quaerens panem.1. Petri5. Illi cura est de vobis. Matth.6. Videte volatilia coeli, quia non serunt, neque metunt, neque conuehunt in horrea, et pater uester » coelestis alit illa. Cognoscite lilia agri, quomodo crescant, non laborant neque nent: » Attamen dico vobis, ne Solomonem quidem in vniuersa gloria sua sic amsctum fu» isse, vt vnum ex his. Ibidem. Nouit pater vester coelestis quod opus habeatis his omnibus. Quin potius quaerite primùm regnum Dei et iustitiam eius, et haec omnia adiicientur vobis. Praeterea animaduertit pietas, diuitias quodammodo obesse potius, quam prodesse. Nam Christus alibi eas interpretatur spinas Luc.8. Et Marci 1o. ait, facilius Camelum per foramen acus ingredi, quam diuitem in regnum Dei. Et Eccles. 31. Qui aurum diligit, non iustificabitur. Prouerbus 28. Melior est pauper ambulans in simplicitate sua, quam diues in prauis itineribus. Psal.36.Melius est modicum iusto super diuitias peccatorum multas. In summa, tametsi Abraham, Dauid, et quidam alii diuitias cum pietate obtinuerint, tamen paucis admodum id datum est: et pietas senper nouit sua sorte contenta esse, id quod opes non norunt. Grauissime itaque peccant illi ministri verbi, quibus persuaderi nequaquam porest, pietatem esse quaestum, id est, abunde sufficere omnia, neque quicquam ei deesse posse, cui adsit fides et charitas. Pudeat, pudeat illos nominetenus Christianos, quibus, quod ad hanc rem attinet, quidam ethnsci rectius sapuerunt. Praeclare definierunt Stoici, Solum sapientem esse diuitem, videlicet quod nihil concupiscat. Et apud Ausonium Bias: Quae nam summa boni? Mens quae sibi conscia recti. Quis diues ? Quinil cupiat. Quis pauper? auarus. Horatius etiam eleganter disputat sermonum liber 1. Saryra prima. Qui fit Mecoenas, vt nemo quam sibi sortem, seu ratio dederit, seu sors obiecerit, illa contentus viuat, etc. Seneca item ad Lucilium epistola no, et alias saepe. Sed idem Apostolus diserte adiicit.  \pend
\phantomsection
\addcontentsline{toc}{subsection}{\textit{Nihil enim intulimus in mundum, videlicet nec efferre quiequam possumus: sed habentes alimenta, et quibus tegamur, his contenti erimus. }}
\subsection*{\textit{Nihil enim intulimus in mundum, videlicet nec efferre quiequam possumus: sed habentes alimenta, et quibus tegamur, his contenti erimus. }}\pstart Sequuntur argumenta a quaestu abducentia et suadentia, ut sua quisque sorte contenti simus: hic vero duplex ratio proponitur:altera a miseria humana in nariuitate et in morte: altera est a naturae sufficientia, quae paucis contenta est. Primum vt conditi sumus, miseri sumus, nudi, omnium egeni : pari conditione hinc emigramus: quid igitur medio tempore vitae, quod exiguum est et incertum, tem anxie incertis et fallacibus diuitiis corrogandis inuigilabimus? lobus 1. Nudus egressus sum de utero matris meae, et nudus reuertar illuc. Iobus 27. Diues cum dormierit, nihil secum auferet, aperiet oculos suos et nihil inueniet. Eccles.5. Noli anxius esse in diuitiis iniustis, non enim proderunt tibi in die obductionis et uindictae. Sola pietas nos a morte comitatur: Ei ergo potius intendendum, quam fluxis opibus, quas hic debemus velimus nolimus relinquere, neque post hanc vitam iuuare queunt. Certe si habeamus communi vitae necesfaria, id satis pio animo esse debet. Natura paucis contenta est: quae obtingunt supra necessitatem, superuacanea sunt, subinde noxia, abusus et luxus caussa. Idcirco apostolus ait,  \pend\pstart Sed habentes alimenta et quibus tegamur, his contenti erimus.) Partitur in duo, omnia, quae sunt nobis in hac vita necessaria, quibus contenta est natura: videlicet in alimenta. et vestimenta. His duobus potest vita, sanitas, caeteraque conseruari: et qui haec habet, habet vnde viuat. Eccles. 29. Initium vitae hominis aqua et panis, et vestimentum et domus protegens turpitudinem. Proinde, vt paucis concludamus, commendatur hic nobis studium pietatis, et cum pietate colenda frugalitas, quae cum omnibus, tum Christianis in primis debet merito esse commendatissima. Quanta credis mala inuehi per abusum opum, per tantum denique in vestitu et victu non necessarium appara ta Vt non immerto ettam ethnicileges vestiaras, sumptuariasque tulerint: quas vereor  \pend
\vspace{0.5cm}\noindent
\vspace{0.2cm}\rule{1cm}{0.2pt}\\ 
\hspace{0.2cm}\textit{mg}
\footnotesize Pxraexe Solus Sapiens diues. 
\normalsize| 
\hspace{0.2cm}\textit{mg}
\footnotesize quibus de2. Ratio a miseria con ditionis. 
\normalsize| 
\hspace{0.2cm}\textit{mg}
\footnotesize Agunenta 
\normalsize| 
\hspace{0.2cm}\textit{mg}
\footnotesize bortatur pi os ab auari tia. 
\normalsize| 
\hspace{0.2cm}\textit{mg}
\footnotesize 2 Ratio a Naturae suffi centit. 
\normalsize| 
\section*{AD TIMOTHEVM CAP. VI. }
\marginpar{[ p.125 ]}\pstart nimium etiamm necessarias multis Christianis. Vide apud Suetonium. Et quam multi ne ipsi quidem uti uel abuti opibus suis audent, solis ingratis et dilapidatoribus haeredibus diuites, sibi vero miserrimi et iniurii?  \pend
\phantomsection
\addcontentsline{toc}{subsection}{\textit{Caeterùm qui volunt ditescere, incidunt in tentationem et laqueum, et cupiditates multas, stultas ac noxias, quae demergunt in exitium et interitum. }}
\subsection*{\textit{Caeterùm qui volunt ditescere, incidunt in tentationem et laqueum, et cupiditates multas, stultas ac noxias, quae demergunt in exitium et interitum. }}\pstart Argumenta a periculoso per Enumerationem. Enumerat enim diuersa et grauia pericula, quae quaestui deditis eueniunt.  \pend\pstart Qui volunt ditescere, inquit, incidunt in tentationem et laqueum.) En primum pericu, lum nequaquam leue aut contemnendum.  \pend\pstart Oue volunt, inquit, ditescere.) Non simpliciter, qui sunt diuites.. Nam diustem esse id quidem per se malum non est, sed res media, utilis bene utentibus, ut patet de sanctis, Abrahamo multisque aliis, qui praeterea quod non excludantur a salute, id argumento est, quod post pauca subiungit Apostolus, quid Timotheum uelit praecipere di uitibus ut salui fiant: sed diuitem uelle fieri cum non sis. Ita ut uelle significet hic, ualde appetere, per fas et nefas conari, omne studium in hoc ponere, ut ex pauperibus diuites fiant:sicut apud Satyricum: Nam diues qui fieri uult, Et cito uult fieri: sed asscribam uersus: sic autem de auaritia loquitur:  \pend
\phantomsection
\addcontentsline{toc}{subsection}{\textit{Inde fere scelerum causse, necplura venena Miscuit, aut ferro grassatur sepius vllum Humanae mentis vitium, quam seua cupido Indomiti census. Nam diues qui fieri vult, Et cito vult fieri, sed quae reuerentia legum, Quis metus aut pudor est vnquam properantis auari? }}
\subsection*{\textit{Inde fere scelerum causse, necplura venena Miscuit, aut ferro grassatur sepius vllum Humanae mentis vitium, quam seua cupido Indomiti census. Nam diues qui fieri vult, Et cito vult fieri, sed quae reuerentia legum, Quis metus aut pudor est vnquam properantis auari? }}\pstart Habere opes, honestis artibus eas seruare, paulatimque augere, non est malum, seq cum habeas, eas immodice amare, uel, cum non habeas, quascunque rationes etiam iniquas inire, ut illas pares. Nam Eccles.1o. dicitur: Nihil est iniquius, quam amare pecuniam Ecclesiastes cap.5. Qui amant diustias, fructum non capient ex eis. Etuixfier potest, ut qui eis parandis ampliandisque ; intendit animum, pariter seruare queat uera pietatem. Matth. 6.Nemo potest seruire Deo et Mammonae. Et Matth.16. Quid prodest homini, si mundum uniuersum lucretur, animae uero suae detrimentum patiatur? Itaque cum sine detrimento animae uix parentur:rectem dixit Apostolus: Eos, qui ditesce reuolunt, incidere in tentationem et laqueum. Solet autem scriptura per tentationem et laqueum intelligere grauissima pericula, quae a Diabolo et mundo imminent: Siquidem diabolus est qui tentat excitatque ad male agendum. Vnde Matth.4. Diabolus antonomatice uocatur ille qui tentat. Et tentat iple plerunque per mundum et ea quae in mundo sunt. Tentat itaque diabolus eos qui ditescere uolunt, suggerens eis, ut parandarum opum gratia perfide agant, famudent, peicrent, furentur, iniurias inferani quibusuis, quacunque ratione: in summa, nullum flagitii genus praetermittant. Quid non mortalia pectora cogis, Auri sacra fames ? Adhaec perpetranda qui plus quam mille machinis impelluntur, sanem grauiter tentantur, et in laqueum damnationis incidunt, unde se nunquam possint extricare. Et cum diabolus habeat innumera tendicula ad decipiendum, non facile potest quisquam laqueos eius euadere, et tentationis bus resistere, praesertim ubi quis cor habet auaritia intectum, quod facilem velit omnibus suggestionibus eius consentire. Exemplum in luda, cuius auaritia pro origine proditionis et omnium scelerum in Euangelio notatur. Vnde addit, quod incidunt in tupiditates multas, stultas ac noxias. Vbi diabolus eos nunquam desinit tentare, tandem quaeuis cupiunt perficere, quae ille suggerit. Atque hae cupiditates stultae sunt, nec hoc solum, uerum etiam noxiae. Vnde et diues ille stultus appellatur Luc.12.cum diceret: Anima, habes multa bona reposita in annos multos, requiesce, comede, bibe, gaude:  dixit illi Deus: Stulte, hac nocte animam tuam repetunt abs te. Quae uero parasti, cui ec cedent ? Sic qui recondit sibi, nec est erga Deum diues. Et prodigalitas diuitum, e  \pend
\vspace{0.5cm}\noindent
\vspace{0.2cm}\rule{1cm}{0.2pt}\\ 
\hspace{0.2cm}\textit{mg}
\footnotesize Arσte4x diuitum in uundo. Decupatio. 
\normalsize| 
\hspace{0.2cm}\textit{mg}
\footnotesize in duiujs parandis et feruandis quid reprebendendum. Difficile est questui et pietatiuna studere. Laqueus pe ricula Satauae significat. 
\normalsize| 
\hspace{0.2cm}\textit{mg}
\footnotesize Diaboli uis in filiis buius seculi. rels saane
\normalsize| 
\hspace{0.2cm}\textit{mg}
\footnotesize Ioban.u2. 
\normalsize| 
\hspace{0.2cm}\textit{mg}
\footnotesize Capiduaes stultae. 
\normalsize| 
\section*{IN I. EPIST. D. PAVLI. }
\marginpar{[ p.126 ]}\pstart incredibilis labor opes terra marique noctes ac dies venantium nonne stultitiae est argumentum certissimum ? Noxias autem esse diuitias multi testatum fecerunt. Ecclel.5. Duritiae conseruatae in malum dominisui. Et fere solet eis, qui iniuste opes parant, calamitosus exitus euenire, siue in hac vita siue in altera, vt patet de diuite epulone Luc. 16. Et sicut, cum eas pararent, omnibus insidiabantur, ita ipsis postquam pararint eas, omnes vicissim insidiantur. Vnde recte dixit philosophus ille rogatus, cur aurum palleret, id fieri, quia omnes ei insidiantur. Saepe ingratus et dilapidator suca cedit haeres. Araneae labori in texenda tela quae mox frangitur, similis est labor diuitum in congerendis opibus. Sicut hydropico, quo plus sunt potae, plus sitiuntur aquae. Et sicut in scabioso frictio excitat frictionem: ita in auaro  \pend
\phantomsection
\addcontentsline{toc}{subsection}{\textit{Interea pleno cum turget sacculus aere }}
\subsection*{\textit{Interea pleno cum turget sacculus aere }}
\phantomsection
\addcontentsline{toc}{subsection}{\textit{Crescit amor nummi quantum ipsa pecunia crescit. }}
\subsection*{\textit{Crescit amor nummi quantum ipsa pecunia crescit. }}\pstart Quid opus est ad breuem vitam, et ne vnsus quidem diei certam, longis et multis et fluxis, similiter vnoque momento furto interdum perituris opibus? Vt autem cu piditates illas supra modum noxias ostendat, addit ueluti per incrementum et ualdεἐαφατικῶς, quod demergunt hominem in exitium et interitum. Βυδίζοσι: uerbum Bubiτα, demergo, in profundum eo: Estitaque metaphora sumpta ab iis, qui aquis obruuntur: qua significatur, auaros plane perire, fine ulla salutis spe. Sicut iis, qui in mediis aquis sunt, imperiti natandi, nulla spes affulget auxilii. Certissimum ergo imminet exitium omnibus, non solum ministris uerbi, sed quibuscunque aliis, qui quaestui dediti fuerint. Sequitur amplificatio periculorum, habens adiunctum exem plum siue argumentum ab experientia.  \pend
\phantomsection
\addcontentsline{toc}{subsection}{\textit{Siquidem radix omnium malorum est studium pecuniae, quam quidam dum appetunt, aberrauerunt a fide et seipsos implicuerunt doloribus multis. }}
\subsection*{\textit{Siquidem radix omnium malorum est studium pecuniae, quam quidam dum appetunt, aberrauerunt a fide et seipsos implicuerunt doloribus multis. }}\pstart Quasi dicat: In summa, quoniam difficile esset enumerare omnia mala et pericula, quae ab auaritia proueniunt, ipsa est omnium malorum radix et origo. Sicut enim ex radice prouenit truncus, rami, fructus: lta ex auaritia omne malorum genus enascis tur. Permetaphoramitaque, quam Cicero uocat breuem Similitudinem, auaritia dicitur malorum radix. Simili figura dixit quidam φιλαργυριαν μντροπολιν πάσης κα- κίας. Rectissimem autem φιλαργυριαν reddidit interpres studium pecuniae, significans exprimensque affectum illum amoremque erga pecuniam, quem Graecauox euidentius significat, quam Latina uox auaritia. Adiicit eius rei probationem quandam ab experientia?  \pend\pstart Quam quidem dum appetunt, aberrauerunt a fide, et seipsos implicuerunt doloribus multis.) Paueis significat, compertum esse, quod qui pecuniis auide inhiarunt, aberrarint a fide, mala sibi accersiuerint etiam spiritualia et coram Deo: deinde etiam dolores et pericula externa coram mundo. Sicut pieras habet promissiones uitae futurae et praesentis: Ita e contrario auaritia habet adiuncta pericula uitae futurae et praesentis. Exempla sunt de Pseudoapostolis et hypocritis, qui quaestum putauerunt pietatem. Ac reuera fieri non potest, quin qui auide intendunt lucro, ueritatem peruertant, dissimulando multa in gratiam diuitum, a quibus speratur praeda: multa item proferendo quae credunt illis profutura. Quod qui faciunt, necessum est eos a fide aberrare, Christum negare, ut mundo placeant, ueritatem imminuere ut suam commendent fal sitatem. Vbi autem a fide aberrant, quae est salutis ancora, quid aliud possunt, quam in omnibus, quae sunt religionis, fluctuare tandemque et mergi? Ideo alibi dicitur auaritia idololatria. Colos.3. Rursus autem ubi a fide aberrant, Deum deserentes, et a Deo deserti, quomodo sperent suocessus in hoc mundo? ltaque implicant se multis doloribus, curis, molestiis, cursitationibus, laboribus se confieiunt, et dum in conscientia quietem non possunt habere, multo minus obtinent eam foris. Ita enim paratum est, opes magnis laboribus acquiruntur, ingentibus curis retinentur, summis doloribus amittuntur. Huc pertinet ilud lobus 20. Cum fatiatus fuerit diuitiis, arctabitur, aestuabit, et omnis  \pend
\vspace{0.5cm}\noindent
\vspace{0.2cm}\rule{1cm}{0.2pt}\\ 
\hspace{0.2cm}\textit{mg}
\footnotesize Cpsutxe noxiae. 
\normalsize| 
\hspace{0.2cm}\textit{mg}
\footnotesize Aurum cus palleat. Possessiont incertae. A simili. 
\normalsize| 
\hspace{0.2cm}\textit{mg}
\footnotesize Vitabreuis et inceria multis opibus non mdiget. 
\normalsize| 
\hspace{0.2cm}\textit{mg}
\footnotesize Danum cre ant diuttiae cupidis. 
\normalsize| 
\hspace{0.2cm}\textit{mg}
\footnotesize Metaphor. a pericluamn bus in aquis. lone - 
\normalsize| 
\hspace{0.2cm}\textit{mg}
\footnotesize Amplificatio periculo rum ab experientia. 
\normalsize| 
\hspace{0.2cm}\textit{mg}
\footnotesize Metaphora ab arboribus  
\normalsize| 
\hspace{0.2cm}\textit{mg}
\footnotesize sumpta siue a radicibus plantarum. Auaritia m torum radix. onαρvρία Ab experientia. Acontrario, siuc a dssim li. Auarorum o pseudoapostolorum de etrina peruersa. Poeiculaa uarorum. 
\normalsize| 
\section*{AD TIMOTHEVM CAP. VI. }
\marginpar{[ p.127 ]}\pstart doloriruet in eum. Volo id a diuitibus ipsis interrogart, num opes ipsorum quietem, hilaritatem, securitatemque omnem adimant. Diuitiae ipsis non inseruiunt, sed ipsimet sunt infelicium opum miserrima mancipia: et non habent ipsi opes, sed ab opibus habentur.  \pend
\phantomsection
\addcontentsline{toc}{subsection}{\textit{Tu vero homo Dei ista fuge, sectare autem iustitiam, pietatem, fidem, charitatem, patientiam, mansuetudinem. Certa bonum certamen fidei, apprehende vitam aternam, ad quam et vocatus es, et professus es bonam professionem coram multis testibus. }}
\subsection*{\textit{Tu vero homo Dei ista fuge, sectare autem iustitiam, pietatem, fidem, charitatem, patientiam, mansuetudinem. Certa bonum certamen fidei, apprehende vitam aternam, ad quam et vocatus es, et professus es bonam professionem coram multis testibus. }}\pstart Sequitur adhortatio, continens argumenta suasoria, ut quaestum fuglat, pietatem verô amplectatur. Vocat autem Timotheum hominem Dei, phrasi Hebraica:nos pos semus teddere, hominem diuinum quemadmodum in ueteri testamento Helias.3. Reg17.uir Dei dicitur. 1.Reg 9. Samuel uir Dei et homo Dei appellatur.3. Reg.12.uir De mittitur ad leroboamum sacrificantem in Bethel.lud 13.uxor Manoae uocat Angelum, qui praenunciarat natiuitatem Samsonis, uirum Dei. Per contrarium uocantur uiri Be lial impii. Latet ergo in hac appellatione ceu stimulus quidam, quo Timotheus ad officium suum excitetur, uelitque reipsa testari, se Deo magis addictum, quam Mammonae: se potius esse uirum Dei, quam uirum Belial.  \pend\pstart Tu vero homo Dei, inquit, ista fuge.) Videlicet quaestum et omnia, quae ad quaestum pertinent, omnia illa pericula grauissima, quae pecuniae studium sequuntur: tom etiam quaestiones illas et logomachias, de quibus ante diximus  \pend\pstart Sectare autem iuslitiam, pietatem.) Has per congeriem opponit uirtutes illis uitiis qua ut dictum est, ex quaestionibus et mrpis lucri studio oriuntur. Nomine iustitie accipi enda est generaliter uirtus, qua quod aequum est fit, siue unicuique suum tribuitur ut di scepationum, rixarum, nulla sit occasio: tùm qua ostenditer animus ab omnibus auar tiae sordibus alienus. Quaestus enim semper fere habet adiun stam iniustitiam. Qui au tem iustitiam uolet colere, nihil sibi iniuste, per fraudem aut dolum parabit.  \pend\pstart Pietatem.) Hanc iam ante opposuit questui, tùm falsae doctrinae. Est autem propriq pietas, ut alias a nobis dictum est, cultus Dei, et religionis pura doctrina obseruatioque . Fidem.) Haec quia non dubitat, non utique mouet quaestiones: sed neque inhiat tempo ralibus bonis et cito perituris, sed spiritualibus tantùm et aeternis.  \pend\pstart Charitatem.) Neque haec fert contentiones uel aemulationes, iurgia, lites, etc. multo minus aliquid ab aliis conatur fallaciter obtinêre ad quaestum augendum: lmo pacem componit inter litigantes, sua impendit egentibus.  \pend\pstart Patientiam.) Qui hanc quoque habuerit, non quaerit inanes logomachias, facile etiam erit sua sorte contentus: imô perferet etiam, si quid ei iniuriae inferatur. Luc.21.In patsentia uestra possidebitis animas uestras. Luc.8. Semem uerbi affert fructum multum in patientia. Mansuetudinem.) Quantun falsorum doctorum ferocia arrogantiaque obest, tantum vbique prodest pii doctoris modestia et mansuetudo: et ad agnitionem amoremque veritatis pertrahas quospiam sola mansuetudine, quos nunque moueas vllis argutiis sophisticis, quamlibet Chrysippeis. In his igitur virtutibus sese exerceat episcopus, imô et quicunque minister ecclesiae, potius quam vel quaestionibus et logomachiis, aut quaestui sese dedat. Imo his virtutibus putet se omnes falsos doctores quaestionum et errorum magistros, auaros, superbos, mente corruptos, superaturum. Hae sunt artes, quibus superabit episcope falsos doctores: hae sunt artes, e quibus q̃stum coram Deo faciet, animas Christo seruatori lueri fi ciens: Hae sunt note, per quas verus doctor a falso dignoscet: in his vera scientia, verae opes sitae sunt. Quidam ita ordinant, vt referant iustitiam, pietatem, fidem ad Deum, tres reliquas exercendas erga proximum. Sequuntur argumenta hortatoria.  \pend\pstart Certa bonum certamen fidei, apprehende vitam eternam.) Quemadmodum in acie suni quidam vocati optiones, qui animum militibus addunt: et in certaminibus non desunt, qui accla mant, Euge, macte virtute et eiusmodi:ita Apostolus subinde optionem hoc pacto agit, piosque ; ad virtutum opera diligenter animat. Ait itaque : Certa bonum certamen fidei. Alludit ad institutum publicorum certaminum olympiacorum, quibus apud olympiam et alibi, cursu,  \pend
\vspace{0.5cm}\noindent
\vspace{0.2cm}\rule{1cm}{0.2pt}\\ 
\hspace{0.2cm}\textit{mg}
\footnotesize De fugienda auanua et sectanda pie tate. 
\normalsize| 
\hspace{0.2cm}\textit{mg}
\footnotesize Viri Dei pis. viri Belial impii. 
\normalsize| 
\hspace{0.2cm}\textit{mg}
\footnotesize Virtutum et uitiorum αὐτίθισις. 
\normalsize| 
\hspace{0.2cm}\textit{mg}
\footnotesize Iustitiarepu gnat aurttie. 
\normalsize| 
\hspace{0.2cm}\textit{mg}
\footnotesize Pictg. 
\normalsize| 
\hspace{0.2cm}\textit{mg}
\footnotesize Fide 
\normalsize| 
\hspace{0.2cm}\textit{mg}
\footnotesize Charitas 
\normalsize| 
\hspace{0.2cm}\textit{mg}
\footnotesize Patientia 
\normalsize| 
\hspace{0.2cm}\textit{mg}
\footnotesize Molicuudo 
\normalsize| 
\hspace{0.2cm}\textit{mg}
\footnotesize 2 hortatoria 
\normalsize| 
\hspace{0.2cm}\textit{mg}
\footnotesize   
\normalsize| 
\section*{IN I. EPIST. D. PAVLI. }
\marginpar{[ p.128 ]}\pstart saltu, lucta similibusque certabatur. Et quemadmodum diceremus alicui suscipiendum certamem cursus aut huiusmodi:ita ait, Certa bonum certamenfidei: id est, contende, omnes fide, siue tuendo doctrinam fidei superare. Et apprehende vitam aeternam. Pro brauio videlicet et praemio quod consequeris, si in certamine fidei te strenue gesseris. Falsi doctores sectantur ea, quae demergunt ipsos in exitium et interitum: pii debent illa secta ri quae adferunt vitam aeternam. Significat autem hac allusione, non facile esse, neque cu iuis datum, illis virtutibus enumeratis praestare, sed multo labore in id enitendum, nec minus conandum, que vincere conantibus olympia. Ac reuera perpetuum certamem est nobis ac lucta quamdiu in hoc mundo versamur: caro depugnat aduersus spiritum, spiritus ad uersus carnem vicissim. Et falsa doctrina contra ueram: et multa alia sunt certamina spi ritualia, quibus ministrum ecclesiae decet plus quam athletico labore exerceri. Dandatz opera in hoc ministris ecclesig̨,ne fides siue fidei doctrina vsque labefactetur. Hac enim incolumi facile omnis falsa doctrina euertitur... loan.s. Haec est victoria, quae vincit mum dum, fides nostra. Huc pertinet illud 2. Timoth.4. quod Apostolus de seipso dicit: Certamen bonum decertaui, cursum consummaui, fidem seruaui: quod superest, reposita est P mihi iustitiae corona, quam reddet mihi Dominus in illo die, qui est iustus iudex.1.Co ' rinth.o. omnis qui certat, per omnia temperans estillli quidem, vt perituram coronam accipiant, nos autem vt aeternam. Ecclesiast. I.vsque ad mortem certa pro iustitia. Sequitur,  \pend\pstart Ad quam et vocatus es.) A vocatione. Si vocatus es ad aeternam vitam apprehendendam, aequum est, te ita geras in hoc uirtutum certamine, vt eam foeliciter apprehendas neque fine vocationis tuae frusteris. Vocatus es non ad subtiliter disputandù, non ad fallendum populum, non ad quaestum exercendum: sed ad apprehendendam vitam, per certamen fidei.ld igitur praesta, ad quod vocatus es. Saepe autem vtitur Apost.his argumentis a Vo catione ductis, cum vult exhortari: acreuera maximum pondus apud omnes habere deberent.  \pend\pstart Et professus es bonam professionem coram multis tesiibis.) Haec ratio est a vitae instituto et publica professione. Quemadmo dum. n.vnusquisque quod genus vitae suscipit, debet cu rare, vt illi respondeat: qui se militem, pfitetur, militet: ꝗ medicum, medeat: qui architectum, aedificet: ita aequissimum est, qui ministrum verbi se profitet, suam quoque functionem gnauiter exequaf. Timotheum ergo, ꝗ suscepit hocuitae institutum, ut ecclesias Dei gubernaret, pe mittens se diligentem fore, idque pfessus est coram multis testibus, partim qui ordinationi eius, cum imponerent ei manus, interfuerant: partim, qui eum iam saepem audierant doctrinam Dei strenue asserentem et defendentem: hortat argumento sumpto a uitae instituto et professione. Ac non paruam uim habet circumstantia addita, coram testibus, et quidem multis. Ergo admonet unusquis qu sui instituti et professionis habere ratione, dareque operam, ut uocationi diuinae quae praeit, et professioni hominis, quae subsequit, quasi uocationem di uinam approbans, eiq se astringens gnauit respondeat. Nonnulli per professionem intelli gunt professionem seu confessionem illam fidei, quae praeit baptismum, in qua renunciantes sathanae, mundo, et pompis eius, iuramus in uerba Christi, et iam amodô nos fideles fo re promittimus et profitemur. Et certe uel huius professionis commemoratione debent hominum animi meriton ad diligentiam commoueri. Vtinam Deus illud animi robur nobis inderet, ut saepe possemus in memoriam reuocare, utque in tentationibus potissimum nobis occureret, quod in baptismo sumus professi et pollicitimnec dubium, quin plurima peccata eius professionis consideratione impedirentur.  \pend
\phantomsection
\addcontentsline{toc}{subsection}{\textit{Praecipio tibi in conspectu Dei, qui viuificat omnia, et Iesu Christi, qui te statam fecit sub Pontio Pilato bonam professionem, vt serues praeteptum, immaculatus, irreprehensibilis, vsque ad apparitionem Domini nostri lesu Christi. }}
\subsection*{\textit{Praecipio tibi in conspectu Dei, qui viuificat omnia, et Iesu Christi, qui te statam fecit sub Pontio Pilato bonam professionem, vt serues praeteptum, immaculatus, irreprehensibilis, vsque ad apparitionem Domini nostri lesu Christi. }}\pstart Adnortationi addit praeceptum, quo Timotheum magis urgeat. Magna est autem prae ceptorum uis, neque temere sunt ea contemnenda, maxime si ex uerbo Dei trahant originem, sicut certum est Apostolum nihil uoluisse praecipere, qu non et a Christo praeceptum fuerat. Alias si quid precipiant episcopi, quod alient sit uel pugnans cu ueibo, certe iudiciu Dei ipsi non euadunt. Apostol' igit Timoth. etiam precepto urget, ut seruet ea q̃ hacten, de tuem da doctrina fidei contra falsos doctores prescripta sunt: sed et q uires addat precepto dicit: In conspectu Dei qui viuificat omnia, et Iesu christi.) Quasi dicat, autoritate ipsis Deiego  \pend
\vspace{0.5cm}\noindent
\vspace{0.2cm}\rule{1cm}{0.2pt}\\ 
\hspace{0.2cm}\textit{mg}
\footnotesize Argumenti a magnitudi ne praemii. 
\normalsize| 
\hspace{0.2cm}\textit{mg}
\footnotesize Certamem pia rum in terris. 
\normalsize| 
\hspace{0.2cm}\textit{mg}
\footnotesize In fideuicte ria Christianorum cerni tur. 
\normalsize| 
\hspace{0.2cm}\textit{mg}
\footnotesize Argonens Auocatione 
\normalsize| 
\hspace{0.2cm}\textit{mg}
\footnotesize A uitae insituto et publica profes fione. 
\normalsize| 
\hspace{0.2cm}\textit{mg}
\footnotesize A circunse tia. 
\normalsize| 
\hspace{0.2cm}\textit{mg}
\footnotesize Professioud tations. 
\normalsize| 
\hspace{0.2cm}\textit{mg}
\footnotesize Preceptan 
\normalsize| 
\hspace{0.2cm}\textit{mg}
\footnotesize Episcoporin praecepta quatenus re cipienda. Ab autorita te Dei. 
\normalsize| 
\section*{AD TIMOTHEVM CAP. VI. }
\marginpar{[ p.129 ]}\pstart ista praecipio, non mea tantùm. Et verem ita dicere potuit, quando praeceptum istud ad glo riam solius Dei promouendam pertinebat. Vocat interim Deum patrem, cuncta viui Fcantem:nam et illi propria virtute vita adest, etab eo omnia vitam accipiunt. Quod erge non solùm proponitur faciendum et suadetur, verùm etiam praecipitur, et quidem au toritate Dei patris cuncta viuificantis, id nullo pacto licet negligere.  \pend\pstart Qui teslatam fecit sub Pontio Pilato bonamprofessionem) Subest argumentum ab exem plo: proponit enim Timotheo Christum in tuenda veritate imitandum: quasi dicat? Christus constanter tutatus est veritatem, et quidem vbi apud impium iudicem inno center delatus perielitabatur. lgitur debes et in hac parte Christum insequi, animoque ́ impererrito ad mortem vsque pro veritate stare. Constat autem, quomodo Christus liberrime sit confessus apud Pilatum, se veritatem tantum docuisse, et in hoc fuisse natum et in mundum venisse, vt testimonium ferret veritati, et omnes qui essent ex veritate, audire vo cc̃ eius. Tum de regno suo coelesti, quod nequaquam ex hoc mundo, simpliciter disseruit. ltaque ad extremam vsque mortem doctrinam veritatis Christus defendit, maluitque mori,que eius patrocinium deserere, aut a sua professione, ob que in mundum missus fuerat, discedere. Hanc autem constantiam Christi in tuenda veritate decet omnes, presertim ministros ecclesiarum, imitari. Facit etiam illa commemoratio passionis dominicae nomnihil, vt Ambros.ait ad terrorem incutiendum, cuius passionis reus necesse est fiat, ꝗ hoc spernit, propt quod passus est ChristVtserues praeceptum, immaculatus, irreprebensibilis.) id est, pura conscientia, ex animo et summa cum diligentia, ne quid in te possit desiderari:ne patiaris te corrumpi per falsos doctores, neque hypocrisi quadam seruabis preceptum istud, sed immaculatus et irreprehensibilis. I sa ad apparitionem Domini nostri Iesu Christi.) Argumentum a Seueritate et certitu. dine diuini ludicii: quasi diceret: ita hac in parte officium facito, vti vis Deo respondere in die iudicii. Ac valde mouet haecratio: Sciunt. n. omnes sibi tunc reddendam rationem factorum ac dictorum omnium, autem grauissime puniendos, qui in officiis suis fuerint negligentes, quique veritatem pro uirili non defenderint. Significatur autem extremq̃ illud iudicium, per apparitionem Christi. Nam Christo datum est omne iudicium a patre: et tunc ve niet cum maiestate magna iudicaturus viuos et mortuos: quo tempore, qui pro verita te stererint in vita sua, accipient aeternam gloriam qui vero cam impugnarint, aeternis destinabuntur suppliciis. Discimus itaque hic, sic omnia nobis agenda, vt parati simus occurrere Domino venienti ad iudicium, sicut parabola docet de quinque prudentibus virginibus, et possimus actionum nostrarum rationem reddere. Nonnulli interpretantur, Vst ad apparitionem Domini, vsque ad tempus mortis, per Metalepsin. Mortem enim consequitur visio Christi, siue vbi mortui fuerimus, tunc apparebit nobis Christus, spes et gloria nostra. Vunusquisque igitur adillum diem fatalem sese paret, et sic agat partes veritatis, tanquam si quotidie eum expectet, sequitur de illa apparitione.  \pend
\phantomsection
\addcontentsline{toc}{subsection}{\textit{Quam temporibus suis ostensurus est ille beatus et solus princeps et rex regnantium, et dominus dominantium, qui solus habet immortalitatem, lucem habitans inaccessam, quem vidit nemo hominum, neque videre potest, cui honor et imperium sempiternum. Amen. }}
\subsection*{\textit{Quam temporibus suis ostensurus est ille beatus et solus princeps et rex regnantium, et dominus dominantium, qui solus habet immortalitatem, lucem habitans inaccessam, quem vidit nemo hominum, neque videre potest, cui honor et imperium sempiternum. Amen. }}\pstart Faciunthaec, vt de apparitione Christi in extremo iudicio uerba precedentia intelligenda uideantur. Ait autem qu apparitionem illam Christi, beatus ille Deus, qui solus per se beatus est, et ipse omnia beans imo ipsamet beatitudo, ostensurus est temporibus: id est:state et definito tempore quod ei placebit. Nam de die et hora illa neque angeli neque aliquis mortalium nouit, sed solus Pater. Dies ille Domini ut fur in nocte ueniet.1. Thessal.5.1. Petri3. Tantum nostra interest uigilare, quia qua hora non putamus, filius hominis ue niet. Porro ut conslderatione eius diei quilibet magis moueretur ad agendum quod sui est officii, pleniusque significaret, quam sit omnibus transgressioribus dies ille metuendus, placuit Apostolo selectis uerbis diuini numinis maiestatem, gloriam, potentiame que describere. Possunt autem, cum diunarum personarum Patris et filii par sit facultas et potentia, haec omnia tam patri quam filio tribui, ut siue de Christo siue de paueaccpuas omnia,nhureserat De parte quidamalunt Acgomma quide propre  \pend
\vspace{0.5cm}\noindent
\vspace{0.2cm}\rule{1cm}{0.2pt}\\ 
\hspace{0.2cm}\textit{mg}
\footnotesize Ab aonte Dei. Iohan... Adt.i7. 
\normalsize| 
\hspace{0.2cm}\textit{mg}
\footnotesize Ab exenplo christi. Comparatio Christi et Ti notheiin do ctrina fidei 1oum.s. Hattbus 27. Iuc-n. 
\normalsize| 
\hspace{0.2cm}\textit{mg}
\footnotesize Audey εturi certitudine et seue Pitate. lons 2Cors. Mattbus is. 1.Cor.1s. Matth.25. t.Thess4. 
\normalsize| 
\hspace{0.2cm}\textit{mg}
\footnotesize Matth.25. 2 Petris. 
\normalsize| 
\hspace{0.2cm}\textit{mg}
\footnotesize Aurciu. 
\normalsize| 
\hspace{0.2cm}\textit{mg}
\footnotesize Luc. 2117. 
\normalsize| 
\section*{IN I. EPIST. D. PAVLI. }
\marginpar{[ p.130 ]}\pstart competunt patri, sed habet eadem a patre nihilominus et filius. loann. 1o. Ego et pater vnum sumus.  \pend\pstart Ille beatus.) Per Antonomasiam vocatur Deus beatus, a seipso immensam habens beatitudinem. Et habet gratiam atque emphasin pronomen, Ille beatus, inquit.  \pend\pstart Et solus princeps, rex regnantiu et dominus dominantiü.) Vocabulis his a potestatibus hur mundi sumptis, summa potentia, dignitas, maiestas Dei significatur: et proprie conue niunt illa soli Deo, sed participatione et ratione similitudinis cuiusdam accommodantur inferioribus potestatibus quae ipsae ab illo solo instituuntur, reguntur, gubernantur, defenduntur, elus nutui parêre debent, tanquam humiles eius ministri. Apocal.19.  \pend\pstart Qui solus habet immortalitatem.) Sicut autem pater vitam habet in semetipso: ita dedit et filio vitam habere in semetipso. Iterum Christus: Ergo sum via, vita, veritas Ioan.6.1.14.  \pend\pstart Lucem habitans inaccessam.) Lucem, inquam, humanis sensibus inpenetrabilem:li oa culus humanus non potest defixe intueri claritatem Solis: quanto minus intellectus humanus intuebitur contemplabiturque illam lucem diuinae maiestatis ? Idem propemodum significant et illa:  \pend\pstart Quem vidit nemo hominum, neque videre potest.) Nemo quidem vidit, vti est, secundum suam naturam. Leguntur quidam dixisse, sibi visum Dominum: vt Genes.32.lacob ait: Vidi Dominum facie ad faciem vbi postea significatur Angelum fuisse. Et Esa.6. vidi Dominum sedentem super solium excelsum! Sed intelligi conuenst, visam ipsis Dei maiestatem, quantum datur percipere humanae imbecillitati, hoc autem perspeciem et similitudinem quandam Substantia siue natura in hac vita eatenus est homini com prehensibilis, quatenus sensui obiectis rebus et imagsnibus quibusdam adumbratur. Alias locum habetillud: Qui scrutator est maiestatis, opprimetur a gloria. Et, Nunc quidem quasi per speculum et per aenigma videmus eum, tunc autem, videlicet in altera vita, videbimus eum sicuti est. In Christo quidem iam carne induto multa de diuina naturam didicimus, et qui vidit eum, vidit et patrem, vt ipse ad Philippum dicit loan. ti4. et ia dem Coloss.2. Hebr.1.dicitur expressa imago substantiae eius: sed tamen nos non ma gis omnem naturae diuinae conditionem ex Christo potuimus cognoscere, quam ex ima gine quis cognoscit totam naturam eius, cuius imaginem pictam intuetur.  \pend\pstart Cui honor et imperium sempiternum.) Comprecatio est Apostolo non insolens. Solet enim ita saepe claudere, vbi quid ad laudem Dei magnifice recensuerit. Sic enim et capite primo huius postquam de gratia egerat: Regi seculorum immortali, inuisibili, soli  \pend\pstart Amen.) Vsitata vox Hebraeis, cùm bene precantur: estque approbantis, acclamantis, beneque optantis Septuaginta reddiderunt γούοιτο, Fiat. Ambrosius hoc loco interpre tarur verum. Caeterùm omnibus ecclesiarum ministris incumbit, hanc Pauli praeclaram exhortationem, simulque praeceptum cum argumentis insertis, alta mente reponere, nec non dare operam, vt iuxta praescriptum eius sanam doctrinam fidei constanter tueantur, neque patiantur gregem dominuscum a falsis doctoribus seduci. Sane vbi episcopi neque hac adhortatione Pauli, neque praecepto commouentur, non est mirum, si in illis ecclesiis falsa doctrina regnet et superstitiones, haereses, scandala multa conspiciantur.  \pend
\phantomsection
\addcontentsline{toc}{subsection}{\textit{His, qui diuites sunt in praesenti seculo, praecipe, ne elato sint animo, neque spem ponant in diuitiis incertis, sed in Deo viuente. }}
\subsection*{\textit{His, qui diuites sunt in praesenti seculo, praecipe, ne elato sint animo, neque spem ponant in diuitiis incertis, sed in Deo viuente. }}\pstart Satis superque praescripsit, quomodo episcopus se erga cuiuscumque sortis homines gereret, siue quae eos doceret:istud autem de instituendis admonendisque diuitibus nondum exposuerat. Quia autem iam multa dixit aduersus cupiditatem diuitiarum grauissimaque pericula a diuitiis imminere: ideo placuit ei hoc loco, quomodo diuites agerent, formam quandam adiicere ne desperationis laqueus diuitibus iniiceretur, quasi a regno Dei prorsus essent exclusi, et ne diuitem esse per se malum damnationique obnoxium putaretur. Porro quemadmodum qui agellum aliquem volunt colere, prius extirpantmalas et noxias herbas, postea vero iaciunt bona semina: lta Apostoius primo prae seribit, quae sint diviuis.  \pend
\vspace{0.5cm}\noindent
\vspace{0.2cm}\rule{1cm}{0.2pt}\\ 
\hspace{0.2cm}\textit{mg}
\footnotesize Deus Beatus dicitur per Antonoma fian. 
\normalsize| 
\hspace{0.2cm}\textit{mg}
\footnotesize Quomodo Deus uisus dicatur patribus. 
\normalsize| 
\hspace{0.2cm}\textit{mg}
\footnotesize Prouer. 25 
\normalsize| 
\hspace{0.2cm}\textit{mg}
\footnotesize 2.Corinth.13 Exod.33. Mosi recusa tur uisio glo riae Dei. 
\normalsize| 
\hspace{0.2cm}\textit{mg}
\footnotesize Comprecatio σlσέη* 
\normalsize| 
\hspace{0.2cm}\textit{mg}
\footnotesize siütad. 
\normalsize| 
\section*{AD TIMOTHEVM CAP. VI. }
\marginpar{[ p.131 ]}\pstart bus, quantum ad ipsas diuitias attinet, mala in illis vitanda: deinde quae bona ex illis facienda: Sic itaque ait,  \pend\pstart His qui diuites sunt in prasenti seculo, praecipe.) Diuites in hoc seculo, quia aliae sunt diui tiae in altero seculo, et hic quidam diuites sunt opibus huius seculi caducis, sed iidem pau peres suntopibus spiritualibus et aeternis, non sunt diuites bonis operibus. Tametsi multi etiam opes in hoc seculo possidere queunt, qui tamen non sunt censendi mali. His autem praecepta dari Apostolus et tradi vult perministros ecclesiarum, quomodo sese gerere possint, vt cum huius seculi diuitiis possideant etiam diuitias alterius seculi.  \pend\pstart Primum autem praeceptum de malis vitandis est, ne elato sint animo. Solet fere diuitil as comitari tumor animi siue fastus. Omnia enim vulgo tribuunt diuitiis, et quicquic diuites faciunt, multi quasi recte factum protinus approbant. Hinc veron animo efferuni Vnde Solonem dixisse ferunt: Satietatem nasci ex opulentia, ex satietate ferociam. Proinde qui elato est animo, sequitur eum contemnere proximum. Qui vero contemnit proi ximum, omnes humanitatis et beneficentiae affectus exuit. Itaque de nemine bene meret̃, omnes prae se contemnit, non commiseretur, nemini quicquam opis impertit. Deinde contemptum proximi sequitur etiam contemptus Dei et summa caecitas. Haec vitia videre est Luc.16.in diuite epulone, qui fastu magno splendide agens, interim contem psit proximum egenum, ac propterea damnatus, qui tamennon iniuste parasse opes, non uim cuiquam intulisse, non aliquid huiusmodi commisisse legitur. Et Ezechiel. 16. Sodomapeccata enumerant, superbia, saturitas panis et abundantia, et otium ipsius et filiarum eius et quod manum egeno et pauperi non porrigebant. Et multis aliis scripturae locis fastus et superbia deprimuntur. Fuerunt tamen quidam patres sancti diuites in hoc seculo, sed iidem nequaquam elato animo, verum pauperes spiritu, mites, nihil sibi fastuose ar rogantes. Tales ergo diuites vult Apostolus.  \pend\pstart Alterum praeceptum. Neg spem ponant in diuitiis incertis: Cùm vocat incertas diuitias eo nomine satis ostendit, in illis spem non ponendam. Neque enim spes ponenda sn rebus incertis et perituris. Profecto autem incertae merito dicuntur. Quid enim aliud fa ciunt, quam quotidie cadunt et fluunt, propterea fluxae et caducae appellatae? Quic aliud quam quotidie mutant dominum? Et e manu ad manum demigrant, certam sedem non habentes? Sicut aquae in mare recidunt, et a mari rursus ad fluuios et terram per pluuj am:ita pecuniae a diustibus principibus ad populum, rursus ab hoc ad principes et diuites. Et vbi habentur, momento vno omnes abripiuntur, incendio, aquis, furto, alioue casu omnes auferuntur. Denique vbi illae tibi non adimuntur, tu de improuiso adimeris illis: sicut stultus ille Luc.12.audiuit a Deo. Quomodo ergo spes in iis ponetur? Si vlla in re sperandum est speremus in Deo viuente. Deus est qui vitam dat, et hanc et futuram, in eo ergo sperandum, non in diuitiis, quae cum vitam nulli dent, saepem caussa sunt, cur vira quibusdam adimatur. Et vbi opes desunt, aut si adsint, non tamen iuuant: tunc Deus speranti in eum adest et adiuuat. Maledictus itaque qui confidit in multitudine diuitiarn suarum: contra vt dicitur Hierem.7. beatus vir qui confidit in Domino, et erit Dominus fiducia elus. Qui vero spem ponunt in diuitiis, ii idololatriam manifestariam com mittunt, dum ex diuitiis salutem expectant.  \pend\pstart Ous prabet nobis omnia assatim ad fruendum.) Commodissime subiicit, quod Deus om. nia bona hicnobis suppeditet. Si Deus nobis omnia praebet:non est igit, qu quisque diues elato sit animo, quasi ipse sibi sua solertia bona parasset. Rursus si Deus omnia prae bet, est igitur in Deo potius sperandum, quam in diuitiis. Itaque his verbis ratio non obscura continet, cur elato animo esse non deceat, curque ́p in diuitiis non sit fidendum. Ac reue ra Deus plurimis quotidie argumentis nos vrget ad confitendum, quod omnia ipsius bonitati sint accepta ferenda, non nostris laboribus vel prudentiae. Quam multi.n. etiam anxie laborant, et tamen opes non consequunt ? nimirum quia putant se sua prudentia posse parare et non a solo Deo expectandas. Rursus vbi aliqui opes pararint, quam cu tôrursus illis pereunt? quia non diutius apud eundem commorantur, que Deo visum, Si Des est dare, curnon idem cùm uelit auferat? lobus 1. Dominus dedit, Dominus abstulit sit nomen Dominibenedictum, sicut domino placuit, ita factum est Nechominihus  \pend
\vspace{0.5cm}\noindent
\vspace{0.2cm}\rule{1cm}{0.2pt}\\ 
\hspace{0.2cm}\textit{mg}
\footnotesize Praecepta di uitibus propo sita. 
\normalsize| 
\hspace{0.2cm}\textit{mg}
\footnotesize 1. Praeceptum diuitum, ne surt fastuosi Psd.73.37. contemptus proximi, ex animo clato oritur. 
\normalsize| 
\hspace{0.2cm}\textit{mg}
\footnotesize Contemptus Dei sequitur contemptum proximi. Matth.5. 2 praeceptum, Non ponendaspes in diuitiis incertis. 
\normalsize| 
\hspace{0.2cm}\textit{mg}
\footnotesize Ratio, cur non sit spes fixa collocam da in diaieiis in incertituanc trperr H 
\normalsize| 
\hspace{0.2cm}\textit{mg}
\footnotesize Pfalzs A fimili. Eccles. 41. in Deo fperandum. Ratio prima: qua est Deus uiuës, etc. Ephess. 2-Curin deo sperandum qa omnia bonlargitur. Omnia Deo acceptaferem da. A uanitate et successu hominum in certo. Ab interitu et incertitu dne epum. 
\normalsize| 
\section*{N IN I. EPIST. D. PAVLI. EPIST. D. PAVLI }
\marginpar{[ p.132 ]}\pstart lolum ita prospicit Deus, sed omnibus animantibus, omnibus creaturis, qua pascitvo lucres coeli, pullos coruorum. Sed hoc quoque notandum, quod dicit Deum dare omnia, affatim, abundanter, largiter. Ea est diuina liberalitas, non solum necessaria nobis adfert, sed facit vt etiam abundemus. Tanto igitur plus nos illi debemus, ac minus solicitos esse oportet pro terrenis opibus congerendis.  \pend\pstart Ad fruendum.) Videlicet, ἀς ἀπολαυσιν, a verbo ἀπολαύω, quod significat fruor, delicior, capio fructum vel commodum. ltaque significatur ad satietatem, largiter, et ad magnum commodum, et vt hilariter illis vtamur, omnia nobis a Deo dari. Discimus proinde omnia quae hic nobis obtingunt bona, vni Deo accepta ferenda, et pro illis Deo agendas gratias: neque ea posse nostris viribus parari, vel parta retineri, eoque non in illis sed in Deo tantùm sperandum. Sequitur nunc, quis sit verus et laudatus diuitias um vsus, quomodo bona ex diuitiis facienda, siue quare Deus affatim omnia praebeat.  \pend
\phantomsection
\addcontentsline{toc}{subsection}{\textit{Ut bene faciant, vt diuites sint operibus bonis, vt faciles sint ad impartiendum, libenter communicantes. }}
\subsection*{\textit{Ut bene faciant, vt diuites sint operibus bonis, vt faciles sint ad impartiendum, libenter communicantes. }}\pstart Qui diuites sunt in hoc seculo, vbi cauerint, ne per diuitias ruant in peccata, videndum eis, vt ex diuitiis arripiant occasionem bene agendi. Non in hoc dantur opes a Deo, vt quispiam non vtatur, multo minus vt vnus solus vtatur. Sed vtalii quoque ex eisdem bene habeant. Non possessores diuitiarum proprie sunt dicendi diuites, sed tantum dispensatores. Non minor est virtus bene vti beneque dispensare opes, quam bene acquirere. Non satis est opibus esse diuites, sed oportet etiam esse diuites bonis operibus. Si diuitiae bona dicendae sunt, nemini nocere debent, sed prodesse quam plu rimis. Et qui bona possident debent per bona illa boni fieri. Repugnat enim quond ali qui bona dicuntur possidere, et ipsi boni non efficiantur. Aut fortassis efficiunt hi, vt quae bene vtentibus bona vocantur, ipsis male vtentibus dicantur et sint mala. Et credi par est, diuinitus factum, quod vulgo bona appellantur, vt nimirum notatione nominis admonerentur homines ad bene vtendum, occasionemque benefaciendi inde ar riperent. Sentiant ergo quam plurimi ea vere talia esse:alias non sunt bona appellanda et contra ius omne bona dicentur, si nulli prosunt. Addit etiam modum benefaciendi.  \pend\pstart Vt saciles sint ad impartiendum.) δυμεταδότος ἐναι. Εὐμετάδιοτωsignificat liberalem, non moleste dantem. ltaque poscit hoc rectus vsus opum, vt diuites non inuiti dent, et duntaxat importunis precibus rogati:sed ex animo, statim, hilariter, sine coactione. Dandum est, vt ait Apostolus 2. Corinth.9. secundum propositum cordis, non ex molestia, aut necessitate. Nam hilarem datorem diligit Deus: lbique addit Similitudinem: qui sementem facit parcem, is parce messurus est. Et qui sementem facit libenter ac beni gne largiendo, copiose messurus est. Prouerbus 3.Ne dicas amico tuo, vade et reuertere et cras dabo tibi, cum statim possis dare. Ac propterea foelix ille lob, qui dicere potuit (cap.31.) Si negaui, quod cupiebant pauperes, et oculos viduae expectare passus sum: si comedi buccellam meam solus, et non comedit ex ea pupillus. Et ibidem. Foris non mansit peregrinus, et ostium meum viatori patuit. Hic sane recte vtebatur diuitiis, fecitque ut bona illa et sibi et aliis bona fierent, ac propterea Deus etiam illius opes sem per multiplicabat, et semel amissas omnes conduplicabat.  \pend\pstart Lubenter communicantes.) Κοινωνικὸς. Theophylactus istud ad mores et ingeniumre fert, ut diuites sint non morosi, sed humani, mites, comes alios haud grauatim admittentes. Nam morositas facit, ut minus adeant, minus que audeant suam inopiam miseri eis exponere, sicque illi morositate a se depellunt pauperes. Aut, si tamen dant, sic dant, quasi simul exprobrent, uel dando etiam simul molestent iniuriis pauperem. Ceterum quoniam diuites his rationibus putarent sibi parum consultum, iudicarentque, siomnibus ita expositi forent, dantes tùm petentibus tùm non petentibus egenis, breui exhauriendos loculos, scrinia et omnia: ideo adiicit potissimum rationes, quibus et hanc obiectionem euertat, et eos ad praestandum opera charitatis inuitet. Prior est ab utili.  \pend
\vspace{0.5cm}\noindent
\vspace{0.2cm}\rule{1cm}{0.2pt}\\ 
\hspace{0.2cm}\textit{mg}
\footnotesize Disperfaio opum. 
\normalsize| 
\hspace{0.2cm}\textit{mg}
\footnotesize Acougat 
\normalsize| 
\hspace{0.2cm}\textit{mg}
\footnotesize Quomodo partis diuiuis utendum, 
\normalsize| 
\hspace{0.2cm}\textit{mg}
\footnotesize τυηδαδοrou eivat. 
\normalsize| 
\hspace{0.2cm}\textit{mg}
\footnotesize Duairde ctorum solicitudo uana di exhaustione opan. 
\normalsize| 
\section*{A. D TIMOTHEVM CAP. VI. }
\marginpar{[ p.133 ]}\pstart Recondentes sibi ipsis fundamentum bonum in posterum.) Quasi dicat: Ne putent perii quod in alios conserunt: nam reuera non tam aliis, quam sibi ipsis prosunt, recondentes fundamentum bonum in posterum, in altera videlicet vita, vbi larga praemia recipient. V. detur autem alludere ad communia diuitum studia, quibus hoc pouissimuin voto est, vt ca mode habitent ac magnificas aedes sibi extruant. Itaque dum facultaubus suis iuuant paus peres, dicit, eos sibi recondere fundamentum et principia amplae domus in coelis. Br, si fi de et charitate prosequant se, posse perficere. Ita bonum fundamentum sibi parauit per clecmosynas Cornelius Centurio, Actor. 1o. Emphasis est in voce sibi ipsis. Ne putent sibi perire quicquam̃, aut aliis duntaxat prodesse eleemosynas. Hac igitur ratione ab vtili permoueri vult eorum animos ad lubenter praestandum ea, quae de vero vsu diuitiarum sunt praescripta Hucvero pertinet illud Matth.6.Recondite vobis thesauros in coelo, ybi neque erugo neque tinea corrumpit, et vbi fures non perfodiunt neque furantur. Nanubi fuerit thefaurus vester, illic erit et cor uestrum Et Luca6.facite uobis amicos ex manmona iniusto, ut cùm defeceritis, recipiant uos in aeterna tabernacula.  \pend\pstart Vt apprehendant aternam vitam.) Hic altera ratio latet a magnitudine praemii. Non debet esse graue, impartiri temporalia bona, ut pro illis accipias aeterna. Foelix et iucunda est commutatio, pro carnalibus accipere spiritualia, pro fluxis et momentaneis stabi£ lia et aeterna. Graece est, ο ἀωνίς ζωῆς, id est, aeternae uitae. Ambrosius reddidit ueram uitam. Nam haec praesens uita est, inquit, duntaxat imago uitae. Et quia imago non om nino inanitas est, in hac uita illa acquiritur, etut ad illam festinaretur, ex imagine, quam praecellat ueritas, addiscitur, quae neque superbia, neque spe diuitiarum attingitur, sed humilicate et spe Dei, qui hanc promissst. Ita Ambrosius. Sunt ergo per ministros uerbi diuites his rationibus de uero usu diuitiarum crebro et diligenter instituendi: et non his duntaxat argumentis, uerùm aliis, qualia quotidie ex reipsa et euentu enascuntur, excitandi ad eleemosynas in pauperes largiter conferendas. Et nisisstud faciant, soiant se iudicium Dei sine miserscordia non euasuros.  \pend
\phantomsection
\addcontentsline{toc}{subsection}{\textit{O Timothee depositum serua, deuitans profanas vocum vanitates, et oppositiones falso nominatae scientiae, quam nonnulli profitentes circa fidem aberrauerunt. }}
\subsection*{\textit{O Timothee depositum serua, deuitans profanas vocum vanitates, et oppositiones falso nominatae scientiae, quam nonnulli profitentes circa fidem aberrauerunt. }}\pstart Breuis rotius Epistolae Conclusio, in qua ad perstandum in sana doctrina fidei, de uitandumque inutiles quaestiones et logomachias adhortatur. Ac perapte respondent haec tùm illis, quae in exordio habebantur de fugiendis illis, quae quaestiones magis prae bent, q aedificationem Dei, quae est per fidem: tùm illis, quae alibs deinceps eadem de re occurrerunt: vt hine de statu seu de propositione principali Epistolae promptum sit iudicare. Magno itaque et vere paterno affectu ait!  \pend\pstart O Timothee depositumserua) Pαρακαlαθήκιων φυλαβον. Metaphora vocauit sanom doctrinam siue dotes a Deo collatas, quibus ea conseruari possit, depositum. Quemad. modum autem is, apud quem aliquid depositum est, cauere debet, ne id vel mutetur, vel deterius fiat: ita episcopi est, id quod apud ipsum depositum est, quodque Deus ipse et concredidit, videlicet sanam doctrinam, omni studio integram incorruptamque conser uare, atque ad eam rem dotibus a Deo sibi concessis beneuti. Vel per depositum intelligi potest ipsum episcopi munus, quod a Timotheo doctrina et vita integrum eral conseruandum. Nam vna hac voce Apostolus haud dubiem plurima vult intelligi. Huc pertinet parabola Matth.25. de eo, qui peregre proficiscens dedit e seruis, alteri quinque  talenta, alii duo, alii veron vnum, vnicuique secundum propriam facultatem. et laudati il, li, qui aliquid lucrati sunt atque adiecerunt: ille vero qui defodit, grauiter reprehensus est. Significat ea parabola, quomodo Dei donis vtendum ad augmentum gloriae ipsius. ltem illud i. Corinth.15. Gratia eius quae perfecta est in me, non fuit inanis, sed copio sius quam illi omnes laboraui. Rursus 2. Timoth. I.Egregium depositum seruato pes spiritum sanctum, qui inhabitat in nobis. vbi etiam est Graeca vox παρακαιαθμ́κν. Sec addit quodamodo, qua ratione Apostolus posset depositum seruare:  \pend\pstart Deuitans, inquit, profanas vocum inanitates et oppositiones falso nominatae scientia.) BeΒίλος κποφωνιας. Βiιθιλος profanas, impuras, pollutas: Nam Bihuna, vti admonuimus  \pend
\vspace{0.5cm}\noindent
\vspace{0.2cm}\rule{1cm}{0.2pt}\\ 
\hspace{0.2cm}\textit{mg}
\footnotesize . Raio a utui. 
\normalsize| 
\hspace{0.2cm}\textit{mg}
\footnotesize Et qu da, plus accedit Commodi, quam ei, daccinis. Fundamentis 1aeitur uitae futurae per elecmossnas 
\normalsize| 
\hspace{0.2cm}\textit{mg}
\footnotesize Domus fir: ma et maiens paranda uu coelis. 2. Am4onit. duie praeni, 
\normalsize| 
\hspace{0.2cm}\textit{mg}
\footnotesize Conclusiototie Epistola. 
\normalsize| 
\hspace{0.2cm}\textit{mg}
\footnotesize  stans exhortatione. 
\normalsize| 
\hspace{0.2cm}\textit{mg}
\footnotesize Metapbord́: Depositum quomodo con seruandum. Sanadoctrina, dotes, et munus episcopi, perde posium iutelligi possant. 
\normalsize| 
\hspace{0.2cm}\textit{mg}
\footnotesize Modus seruam di depositti. θερα sacta ad quae iupe n etion dl Imssi. 
\normalsize| 
\section*{IN I. EPIST. D. PAVLI. }
\marginpar{[ p.134 ]}\pstart alias, erant Graecis dicta quaedam sacra, ad quae etiam impuri admittebantut. Theophylactus suspicatur his verbis ab Apostolo notatos tum Nicolaitas, tum gnosticos siue Euchaetas tunc temporis exortos et valde impuros. Nam de his apud Augustinu libro de potestate daemonum legitur, quod 1o. Calend. Apriles, quo die Dominus noster crucifixus, cum puellis etiam affinibus et quibuscunque conueniebant, peractisque ́ neris cum illis coibant. Nono deinde mense ibidem loci natos infantes circumcidebant, cruorem excipientes: dein corpora cremabant, cruorem et cineres permiscebant quae in poculis inuicem propinarent:sicque rati sunt, diuinum characterem sibi imprimi, ac daemonia recipi, per quae vaticinarentur. Multa alia turpia et infanda perpetrabant, vocabulis etiam miris vtebantur in sua doctrina. etγοφωνια vaniloquium: ἀσέθεα, verba aliena a religione, pugnantia cum pietate.  \pend\pstart Et oppositiones, ἀὐτιθεσως.) Huiusmodi itaque omnia, videlicet et κενοφωνίαι, et oppo sitiones falsô nominatae scientiae, id est, monstrosa docendi forma, vocabula et phrases insidiosae, alienae a scripturis, tùm oppositiones ac disputationes sophisticae, licet in speciem videantur ad pietatis commendationem facturae:haec omnia generaliter sunt vitanda:alias si non vitentur, non bene seruatur depositum. Scientia proprie est veros rum et certorum: falsorum autem est duntaxat opinio, non scientia. Ratio subiicitur, cur uitandae, a periculoso: Quam nonnulli profitentes circa fidem aberrarunt. Sic et in primo capite dicit, quod quidam, quoniam a fide non simulata et conscientia bona aber rassent, deflexerunt ad vaniloquium, uolentes esse legis doctores, non intelligentes quae loquuntur, neque de quibus asseuerant. Rursus bona conscientia et sana doctri na repulsa nomnulli circa fidem naufragium fecerunt, quorum de numero Hymenaeus et Alexander. Denique hoc postremo capite uocauit eos, qui non accedant sanis sermo nibus Christi, et ei, quae secundum pietatem est doctrinae, inflatos, nihil scientes, insanientes circa quaestiones et λογομαχίας. Quod autem decipiantur multi a sua ipsorum sapientia et scientia, liquet ex eo lerem.1o. Stultus factus est omnis homo a scientia sua et Esa. 47. Sapientia tua et scientia tua haec decepit te. ltaque intelligitur, Apostolum id maxime exegisse a Timotheo, uoluissepitidem significare omnibus ministris ecclesia rum, assidue in id incumbendum, ne qua doctrina, praeter eam, quae est secundum pie tatem, id est, in qua docentur fides et charitas, in ecclesiam admittatur. Quicquid alio tendit et quibuscunque rationibus, quacunque arte uel scientia id proponatur, est reiicien dum. Audacter dicendum est cum Paulo Galat.1.Etiamsi nos aut angelus e Coelo praedicauerit uobis Euangelium praeter id, quod Euangelistarum et Apostolorum scriptis est ad nos perlatum, anathema sit. Vtinam autem, utinam haec conclusio praesentis Epi stolae omnibus episcoporum domibus, omnibus templorum suggestibus aureis literis inscripta ministrorum ecclesiae animis perpetuo obuersaretur: nec dubium, quin haereses, falsae opiniones, et doctrinae daemoniorum multae propellerentur.  \pend\pstart Gratia sit tecum, Amen ) Subscriptio. Vtipse ait 2. Thess.3. solitus erat ad omnem Epistolam sua manu ita subscribere, ut ex ea suscriptione certum argumentum haberentur, Epistolam ab ipso profectam, non ab alio quopiam pseudoapostolo, qui legem Mofaicam uel inanem philosophiam magis quam gratiam Christi praedicaret. Apostolus totus erat gratia Christi celebranda, quam et spirant et sonant omnes eius epistolae. Graeca exemplaria habent missam ex Laodicea metropoli Phrygiae Pacatinae: nonnulli addunt, per Tychicum diaconum. Caeterum cum ex Epistola ad Collossenses constet Laodiceam illam Phrygiae Pacatianae, in quae erant et Colossae, Paulo non suisse uisam : uerosimile est ex altera quapiam Laodicea. Nam adhuc alteram memorant in Syria, quae sit uinetis celebris et Carica dicta : alteram uero in Galtia ponunt: Ex illa autem hanc missam autumant.  \pend
\vspace{0.5cm}\noindent
\vspace{0.2cm}\rule{1cm}{0.2pt}\\ 
\hspace{0.2cm}\textit{mg}
\footnotesize Suofticort sacra nefanda. 
\normalsize| 
\hspace{0.2cm}\textit{mg}
\footnotesize Scientis et opinio dif ferunt mate ria siue subiecto. 
\normalsize| 
\hspace{0.2cm}\textit{mg}
\footnotesize Ratio a per culoso, cur uitandae 
\normalsize| 
\hspace{0.2cm}\textit{mg}
\footnotesize lισνορωνίαι et αvribiσας Inflati ope mone scien ue deaipium tur a sua i psorum sa pientia. 
\normalsize| 
\hspace{0.2cm}\textit{mg}
\footnotesize Doctrina, quae est citra pietatem in ecclesia no debet recipi. 
\normalsize| 
\hspace{0.2cm}\textit{mg}
\footnotesize subscriptio. 
\normalsize| 
\hspace{0.2cm}\textit{mg}
\footnotesize Vnde missa 
\normalsize| 
\endnumbering
\end{pages}
\end{document}
        